%source: https://learning.edx.org/course/course-v1:HarvardX+TinyML4+1T2022/home
% Applications of TinyML
% google log in

% Bilder TinyML/MLOps

% Reihenfolge
%1. TinyML/Fundamentals
%2. TinyML/Applications
%3. TinyML/deployment
%4. TinyML/MLOps for Scaling TinyML

\chapter{MLOps for Scaling TinyML}

\section{Objectives}

This course introduces learners to Machine Learning Operations (MLOps), which is about following best practices to deploy and maintain ML models in production. The course specializes in MLOps for TinyML.

\bigskip

Topics

\bigskip

\begin{itemize}
  \item What is MLOps
  \item Why \& When is it Necessary
  \item Automating MLOps
  \item MLOps Platform Features
  \item Examples of MLOps Platforms
\end{itemize}

\section{Prerequisites}

This course is designed to be accessible to all learns. As such, there are no formal prerequisites for this course.

\bigskip

That said, a basic understanding of coding in a high level-language like python may be helpful (and will be needed for a few select optional coding examples). Similarly, familiarity with machine learning and/or embedded systems at a high level may also be helpful.

% 1.1: Welcome

\section{What to Expect in This Cours}
    While many organizations recognize that Machine Learning (ML) can drive significant value, successful deployments and effective operations are the main bottleneck for gaining value from AI. To address this, organizations need to build the necessary ML engineering culture and capability by developing processes that unify ML system development (ML) with ML system operations (Ops) – in short, they need to develop an MLOps pipeline.
    
    To support the engineering community to realize MLOps in practice, Google made available to the world the practitioners guide to MLOps to help developers implement MLOps and smart practices during the ML workflow. We will go through this applied pipeline and we will discuss the unique challenges that emerge as a result of deploying TinyML to many embedded devices. Since MLOps requires the collaboration of business leads, subject matter experts, data scientists, data engineers, ML architects, software developers and maintainers, this course is designed to be accessible to all of these different backgrounds.
    
    Prerequisites
    The core of the course is designed as a broad introduction to the many different stages of the MLOps pipeline. We expect learners to enter this course either without any background in TinyML, or after taking Courses 2 and 3 en-route the advanced TinyML certificate, or after taking Courses 1, 2, and 3 en-route to completing both the standard and advanced TinyML certificates.
    
    What This Course is NOT About
    This is not a hands-on lab-based course. TinyML Courses 1, 2, and 3 are fairly hands-on. This course, by design, is meant to paint and give you a broad overview of MLOps. 
    
    Instead, this course is comprised of a broad overview of MLOps with optional coding exercises, real-world case studies, and links to resources to enable learners to build additional depth in the areas in which they are most interested. We hope that the broad approach enables the course to be accessible to all learners and introduce the important topic of MLOps for TinyML!

\GRAPHICSC{1.0}{1.0}{TinyML/4 MLOps/5.1.png}

\section{Welcome Message}
%1.2: Getting Started
VIJAY JANAPA REDDI: Welcome to Machine Learning Operations for Scaling TinyML
course.
Did you know that about 87\% of companies wish to incorporate AI
into their pipeline never succeed.
Now that's a staggering number.
But don't worry, with many of the concepts you learn in this course,
I promise that you'll be able to avoid the common pitfalls that
lead to that statistic.
In this course, I'll help you learn three main things.
Why you need machining operations to be successful in deploying AI
into the real world.
Second, what are some of the unique challenges that
emerge as you try and employ machine learning operations in order
to scale out TinyML deployments?
And third, I'll talk about some of the potential solutions
we have in the marketplace that will help
us address some of these challenges.
With that said, allow me to introduce myself.
My name is Vijay Janapa Reddi, and I'm a Professor at Harvard University.
And I'm very excited and passionate about sharing
the knowledge that I have, which is why I'm here with you today.
Let's get started with the basic definition of what is TinyML.
TinyML is arguably one of the fastest growing fields
in machine learning today.
It's this beautiful intersection between traditional machine
learning algorithms, embedded systems, and embedded software.
It's all about doing on device sensor analytics,
using machine learning right at the endpoint
where the sensors are producing a huge volume of data.
And the intent is to do this with ultra-low power consumption,
so that we can ultimately deploy machine learning continuously.
Always on machine learning, so that we can convert this raw sensor
data into intelligence, which can then be used to produce
a whole new different set of insights.
And our hope is that we can do this efficiently,
so that we can build systems that operate
on batteries, so they can last forever.
One of the key and most interesting things here is about data.
Machine learning is all about harnessing big data.
Now when it comes to TinyML, sensors generate an insane volume of data.
For example, five quintillion bytes of data
are produced every day by IoT devices.
And the truth is that less than 1\% of this unstructured data
is actually analyzed and used at all in production.
And already, we take machine learning for granted.
So imagine the possibilities that we can unlock
if we can actually harness all this data that sensors are continuously
producing.
And if ML is already pervasive today, it's
going to transform the world once we tap into TinyML.
So imagine the possibilities.
Hence, it's unsurprising that we're seeing a wide range of endpoint devices
that are becoming increasingly intelligent.
Here, I'm showing just a handful of consumer-facing devices,
but TinyML cases also span industrial settings.
Predictive maintenance, manufacturing, there
are big opportunities in that space in order
to really generate true value from TinyML capabilities.
In fact, I myself advised several startups
that have innovative ideas in embedding intelligence onto tiny devices
so that they can deploy them in the field.
And as I go through this course, I promise you
that I will interleave the concepts with some of these real-world case studies,
so that you have a tangible experience with how these concepts convert
into the real world.
Now while TinyML is extremely exciting, I
want to talk to you about a crucial challenge that
is facing many entrepreneurs who are trying to build real world products
and solutions based on tiny machine learning capabilities.
You must realize that it's one thing to deploy a machine learning model
in a small little device, but it's a whole different thing when you actually
want to scale it out into the real world,
into production where you have thousands of these different devices.
And that's where MLOps really comes in.
MLOps is something that is extremely important for anyone and everyone who's
serious about doing anything with ML in the real world.
You must know MLOps, period.
No excuses.
Now what is exactly MLOps?
What MLOps, or machine learning operations,
is a collection of different processes that
are aimed at reliably and efficiently designing, testing, deploying,
and maintaining machine learning models in production.
Now consider that you have this amazing black box project, your secret sauce.
And you're really excited about it and you
want to try and take this project to life in the real world.
Now what do you think it's going to take?
Do you think it can be a one-person project?
The reality is it's going to take a whole team in order
to get your project to life.
Now this is going to range from having a data engineer who
can build the right data pipeline, to having the right data
scientists who can actually train the appropriate model.
Then of course, you'll need a product manager,
who can actually establish what the key KPIs are,
so that you can actually make sure you're meeting those key performance
indicators.
You need a DevOps team that actually knows
how to push the model into production.
So simply put, there are many people you need
to take into account in order to get your black box project into production.
And what we're going to do is try and understand who these people are,
why they are important, what their roles are,
what are the challenges they're going to face,
and how are they going to help you realize that project.
Broadly, these folks are all going to fall
into one of these four main buckets.
The data engineering team, the machine learning
engineering team, the development operations team, and the operations
management team.
All of them play a critical role in various stages
or at various times of the project.
Starting from identifying requirements, to building the models,
to setting up the right staging environments,
to pushing things into production.
Now don't worry if any of these terms seem new.
I know it's a lot of mouthful coming at you.
But over the spread of this course, we will go deep into each of these things.
And at the same time, I promise you, as we go deep you will not get lost.
I've spent a lot of time thinking about how to get these concepts across,
so I hope you will be able to keep up with me.
Now, while you're going to be learning a lot of exciting things,
the first and foremost thing I want you to know
is that you should be excited that you're actually in this course.
I can hands-down say that your spending time in this course
is going to be worthwhile, because there's
an incredible amount of interest in the industry around MLOps.
What I'm showing here is a chart for the term MLOps from Google Trends.
It clearly shows that there's a huge volume of interest,
especially recently.
And this is because industries are wanting
to apply machine learning into the production pipes,
but they don't have the necessary chops, or the understanding
of what are all the bits and pieces that are needed.
And that is why I'm excited that you're here today,
because you're going to be one of those people who
understands how to actually build these systems and push them into production.
And that's going to differentiate you.
Now before we get into the details, I want
to leave you with this very important statement by Dean Scully from Google.
"Developing and deploying machine learning models
is relatively fast and cheap, but maintaining them over time
is difficult and expensive."
Now take this to heart, because this is a statement that
actually comes from Google, who has had a lot of experience
in building these things.
Now, this is true be it Google, Microsoft, Alibaba, Amazon.
Pick your favorite enterprise company that's trying to use AI today.
They all face this exact same problem.
So I'm glad you're here.
Stick with me through this course.
And in the first section, I'm going to introduce
the broad general concepts of MLOps and how the course is structured.
And then, over the remaining sections, we'll
go deeper into each of those stages.
I'll see you in the next video.


\section{Who Should Take This Course?}

\GRAPHICSC{1.0}{1.0}{TinyML/4 MLOps/pasted_image_0.png}

Machine learning can be a game-changer for a business, but without systemization, it can soon devolve into a scientific experiment. The major issue isn't building an ML model. It's putting together an integrated ML system and keeping it operational -- that’s where MLOps comes in.

MLOps is a method for ML professionals to collaborate and communicate in order to manage the production ML lifecycle. It is a culture and practice in machine learning engineering that tries to bring together the creation and operation of machine learning systems (Ops).

While developing a model to meet business objectives (item classification or predicting a continuous variable) and deploying it to production is relatively simple, the reality is that only a small portion of a real-world ML system is made up of ML code, as seen in the picture above, and operating that model in production presents a slew of challenges. The surrounding elements are wide and complex, necessitating the effective and systematic collaboration of a team of Data Scientists, developers, and operations professionals. Furthermore, while engineers can improve their machine learning skills by taking one of the thousands of courses available online, there are few resources available for business leads, subject matter experts, data scientists, data engineers, ML architects, software developers and maintainers, etc. to understand MLOps.

Our aim in this course is to provide you with the resources needed to design and deploy integrated ML systems and keep them operational. As such, we first introduce learners to the fundamental concepts of machine learning operations (MLOps). Second, we want to prepare you for the future by introducing you to Embedded MLOps. Embedded machine learning is one of the fastest-growing fields within applied ML and it comes with its own unique set of challenges and opportunities. If you are interested in deploying and scaling Edge AI technologies, then this course is for you!

\section{The Past, Present and Future of ML}

VIJAY JANAPA REDDI: In this course, we're focusing on MLOps with TinyML,
so I want to make sure I give you a little bit of context about TinyML
and why this is an important and emerging area.
Applications of machine learning are pervasive.
More often than not, we're already using ML without ever knowing about it.
For example, if we're dealing with images on social media,
it's highly likely there were already using email.
You can imagine applications of object detection
for applications such as shopping.
If you're using Google Lens, for instance,
nowadays you can point and click, and you can automatically
look up information for the particular item in a picture.
And that relies on object detection.
And segmentation capabilities, for instance,
in order to enhance images that we're viewing on our phones--
or of course, there are applications of this in autonomous driving as well.
If you're using things like recommendation systems,
you're bound to be interacting with services like Netflix or YouTube.
All of them rely on recommendation engines
at the heart and soul of these systems.
Without doubt, if you ever click the Translate button on a web page,
you are using machine learning to translate the web
page from one language to another.
So my point is this we're completely immersed in machine learning today,
and we take it for granted.
And a lot of this really relies on big machine learning models.
And generally, these models are getting bigger and bigger over time.
For example, here I'm showing you the kind of state-of-the-art models that we
have in just 2019 and 2020.
The Megatron model from NVIDIA, which was leading in 2019,
was completely dismissed by the transformer-based generative language
model that Microsoft introduced in 2020.
The 2020 model is nearly two times as large as the state-of-the-art model
in 2019.
It's incredible how fast and how big these things are growing.
So the general point I'm making is that models are getting bigger.
This is quite true, but this is only one half of the story.
What we're also witnessing is that more and more devices are
starting to spring up around us, and a lot of these devices--
be it our phones, for instance, which have evolved from being feature phones,
or what I call is dumb phones--
have transitioned into smartphones.
And some of that exact technology's now moving into our homes
and becoming increasingly ubiquitous.
Now, the reason that these devices are increasingly becoming intelligent
is because it offers some serious perks.
Bandwidth was one of the first great things.
Machine learning algorithms that are the endpoint
can extract meaningful information from data
that would otherwise be inaccessible due to bandwidth constraints.
A lot of the small devices can produce a huge volume of data
that we simply can't beam up to the cloud.
Latency is the next thing.
On-device machine learning models can respond
in real-time to inputs enabling applications,
such as autonomous vehicles with safety mechanisms and so forth,
which would otherwise not be viable if we're
dependent on network latency or guaranteed connection, expectations
that we must have.
Economics-- but processing data on device,
embedded machine learning or tiny machine learning,
avoids the cost of transmitting data over a network
and processing in the cloud.
Processing data in the cloud is actually quite costly.
Reliability-- systems that are controlled by on-device models
are inherently more reliable than which those depend on a connection
to the cloud.
Of course, right?
For example, imagine you have a car that's
driving through a tunnel you can't be sure that you're going to have a GPS
connection right down Privacy--
when data's processed on the embedded system,
and it's never transmitted to the cloud user privacy is protected.
And this is less chance of an abuse there.
This is especially true in the context of medical applications, for instance.
Now, imagine the wide range of different embedded devices
that we already have access to today.
Now, a lot of these different devices are
capable of running different kinds of machine learning
algorithms, those different use cases.
They're largely driven by sensors, be it a vision sensor, an audio sensor,
or some other type of sensor.
For example, if you take a Jetson Nano, which
is a fairly beefy embedded system, at the highest level
you can do something quite sophisticated,
like video classification, where you're continuously
doing inference on a video stream.
However, you might not be able to do that on a Raspberry Pi,
because it does not have the kind of same level of computational capability.
Nonetheless, you might still be able to do some cool things,
like object detection.
And if you shrink the device even smaller,
and you go more towards power efficient and low power capabilities,
like the Arduino pro for instance, which is designed specifically
for industrial applications-- you can do things
like image classification, which is actually very, very
helpful in many applications that are in the industry, especially in agriculture
and so forth.
Now, you get the general idea of.
Low power devices can do very complex things, like video classification,
but they can still do awesome intelligent tasks--
for example, like keyword spotting.
OK, Google, Alexa, and so forth-- those are keyword applications.
Last, but not least, a big shout-out to the Cortex-M0,
which is one of the smallest microcontrollers that is out there.
And this thing has only about 20 kilobytes or so of effective space
in order to run small models, which means the question becomes,
can we turn that into useful?
Yes, we can actually do things like anomaly detection, which
is a huge market in manufacturing for applications
like predictive maintenance.
So the point is this--
across these wide array of different devices,
what we see is that, at every different level, no matter how big or small
it is in the embedded systems context, you
can actually generate true value out of it
by running machine learning intelligence at the endpoint.
Now, as time goes by, these devices will undoubtedly get more powerful.
The Cortex-M0 is going to be more powerful.
So not only is it going to be able to do anomaly detection-- it's
going to be able to do what the Cortex-M4 was able to do,
which was keyword spotting, for instance.
So you get the general idea that, as time goes by,
more and more machine learning capabilities can
be embedded into smaller devices as these devices improve.
And that is going to lead to a huge new future.
And this is why I like to say that the future of machine learning
is tiny and bright.
We're going to see these new devices that are coming out,
and we're going to see lots of new capabilities
that emerge that are all capable of running machine learning
models on them.
So with that said, keep it in mind, and we
will be able to embark on a whole new journey.

\section{Why the Future of ML is Tiny and Bright}

\subsection{A New Paradigm}

Over the past decade, we have witnessed the rise of machine learning and its gradual incorporation into the industry. More recently, we have begun to witness the bifurcation of machine learning into two separate paradigms based on scale. The first of these is cloud-based machine learning, working at thea large scale and utilizing vast swathes of data stored in data lakes alongside scalable frameworks such as Dask or PySpark. Cloud-based machine learning has been the major focus of industry and is largely compute-centric; centralizing data in a single location before performing computing operations.

In contrast, microcontroller-based machine learning, often referred to as TinyML, instead focuses aton the small scale and is inherently data-centric. Embedding machine learning methods within microcontrollers provides new opportunities that have the potential to offer more efficient computation, reduced communication overhead, cost savings, as well as localized computation. 

Here, we will outline some of the key points that make TinyML an important player in the future machine learning ecosystem.

\subsection{Machine Learning at the Edge}

Computing data locally, rather than streaming data to the cloud or an edge server, is a cheaper and more power-efficient solution for intelligently processing data available on endpoint devices from attached peripherals (e.g., sensors). This not only conserves energy by limiting communication overhead, but also unlocks new applications and use cases. Consequently, there is a growing interest in providing machine learning functionality for endpoint devices. In recent years, it has become possible to take noisy signals like images or audio data and extract meaning from them using neural networks, with TinyML allowing us to run these networks on our microcontroller devices. Since sensors themselves use little power, TinyML can be used to perform real-time machine learning computation on our devices without the need for external communication. We can see applications of this in our daily lives; both Apple and Google run always-on machine learning (ML) based neural networks for voice recognition on these kinds of power-efficient chips (think “OK Goole” or “Hey Siri”).

\subsection{Tiny Computers Are Ubiquitous}


The shift towards microcontroller-based machine learning is largely a result of the increasing number of internet-connected endpoint devices. Whereas ten years ago people had at most a phone and laptop, nowadays, a single person can have a dozen or even more personal devices, ranging from smart thermostats and TVs to smart watches and e-readers. Figure 1 (data from IC Insights) suggests that there will be over 30 billion microcontrollers sold this year, with that number increasing year by year -- to put that into perspective, there are only about 10 million servers in use across the planet and only 4 billion smartphone users. Microcontrollers typically do not get much attention because they are often used to replace functionality that older electro-mechanical systems could do, in cars, washing machines, or remote controls, but they are ubiquitous.

\GRAPHICSC{1.0}{1.0}{TinyML/3 Development/5.4.png}


Microcontroller (MCU) Demand Forecast showing actual sales (in millions of units) for 2016-2018 and forecast for 2019-2023. Showing a linear expected increase from 20,000 million units in 2016 to almost 40,000 million units in 2023. Source IC insights.

FIGURE 1

\subsection{Microcontrollers are Cheap}


Microcontrollers are much cheaper than conventional hardware such as phones and laptops, because they aren’t designed to be general-purpose computational workhorses that run complex workloads. Instead, they are often designed to perform specific tasks, slowly, albeit steadily. The average price of a microcontroller unit (MCU) is already less than USD \$1. Figure 2 shows the average sales price (ASP) for a microcontroller is hovering close to 55 cents. As demand continues to rise, it is expected that this will soon fall below 50 cents. On a global scale, the microcontroller market (\$M) is poised to grow by USD \$6.74 Billion between 2020-2024.

\GRAPHICSC{1.0}{1.0}{TinyML/3 Development/tiny_bright.png}



Microcontroller (MCU) Pricing Forecast showing actual average selling price (in US dollars) for 2016-2018 and forecast for 2019-2023. Showing a decrease from \$0.75 in 2016 to ~\$0.63 in 2018 and a linear expected decrease to almost \$0.55 in 2023. Source IC insights.

FIGURE 2

\subsection{MCUs are Resource-Constrained, Ultra-low Power Systems}

The holy grail for almost any embedded device is for it to be deployable anywhere and require no maintenance like docking or battery replacement. Any device that requires tethered electricity faces a lot of deployment barriers. It can be restricted to only places with electrical wiring. Even where electrical wiring capability is available, it may be challenging for practical reasons to plug something new in, for example, on a factory floor or in an operating theatre. Putting something high up in the corner of a room means running a cord or figuring out alternatives like power-over-ethernet. The electronics required to step down or convert the main-line high voltage to low voltage are expensive and waste energy.

But perhaps the most significant barrier to achieving ubiquitous deployment of computers everywhere is how much energy an electronic system uses. Here are some rough numbers for standard components based on figures from Smartphone Energy Consumption, and so it is no wonder that smartphones need to be tethered to the wall and recharged every night:

\begin{itemize}
    \item Active cell radio might use 800 milliwatts.
    \item A display might use 400 milliwatts.
    \item GPS is 176 milliwatts.
    \item Gyroscope is 130 milliwatts.
    \item Bluetooth might use 100 milliwatts.
    \item Accelerometer is 21 milliwatts.
\end{itemize}

A key opportunity with using MCUs is they require a very minimal amount of energy. A microcontroller itself might only use a milliwatt or even less. Still, you can see that peripherals (accelerometer, gyroscope, GPU, etc.) require much more energy to stay on.

A coin battery might have 2,500 Joules of energy to offer, so even something drawing only one milliwatt will have severe consequences as that means it can run continuously for only about 30 days. Of course, most current products use “duty cycling” and power naps to avoid being always on, but we see how tight the budget is even then. The overall thing to take away from these figures is that while processors and sensors can scale their power usage down to microwatt ranges, displays and especially radios are constrained to much higher consumption, with even low-power wifi and Bluetooth using tens of milliwatts when active. The physics of moving data around is generally well-known to require a lot of energy.

The general wisdom is that the energy an operation takes is proportional to how far you have to send the bits. CPUs and sensors send bits a few millimeters and are cheap. Radio sends them meters or more and is expensive. We don’t see this relationship fundamentally changing, even as technology improves overall. We expect the relative gap between computing and radio costs to get more expansive because we see more opportunities to reduce computing power usage.

\subsection{Tiny Machine Learning (TinyML) on MCUs}

Can you imagine a future where every one of the 250B tiny MCUs runs intelligent machine learning algorithms that can sense their surroundings, predict events in real-time, and even make nudges or recommendations based on the sensors’ activity? We can transform the world.

For instance, smart consumer devices like smart toothbrushes will have MCUs that are tightly coupled with sensors that dynamically adjust to your brushing intensity based on pressure. Medical devices in the future may use microcontrollers with biomedical sensors to control drug delivery as and when needed. Microcontrollers in automotive applications can aid functional safety on the road, detecting engine conditions, etc. to ensure our safety. Finally, MCUs will likely have widespread applications in industrial devices for continuous process monitoring and anomaly detection for applications such as predictive maintenance that can save millions of dollars in productivity loss and downtime by forewarning maintenance engineers.

TinyML is a relatively new machine learning paradigm, yet it is producing remarkably astounding results. Several recent examples include audio, visual, and sensor fusion applications, such as voice and facial recognition, voice commands, and natural language processing.

Looking further into the future, we imagine a world where we have a tiny battery-powered image sensor that I could program to look out for things like particular crop pests or weeds and send an alert when one was spotted. These could be scattered around fields and guide interventions like weeding or pesticides in a much more environmentally friendly way.

\subsection{The Future of TinyML is Bright}

We can conjure up a thousand other products. It feels a lot like being a kid in the Eighties when the first home computers emerged. None of us had any idea what they would become, and most people at the time used them for games or storing address books, but there were so many possibilities. A new world emerged out of those burgeoning systems. So we challenge you to invent the future. Come on and embrace the future with us by learning all we have to teach you about tinyML!

\section{Machine Learning Lifecycle}

VIJAY JANAPA REDDI: Before we dive into MLOps,
we need to understand the basic machine learning lifecycle,
as that is where everything begins.
Generally speaking, any AI project has four main stages.
The first is data engineering.
It's all about data collection, making sure you have the right data sets
and data sources ready for you.
Then, with that data you go in to model engineering,
where you actually want to train your network.
Then you go on to actually deploying the model
onto a device or a whole bunch of devices
and then making sure that the model is actually
performing the way you expect it to perform
when it's actually in a real product.
You can further break each of these four components into smaller subtasks.
For example, you can take data engineering,
I can break it down into data collection and data preprocessing.
Then you can take model engineering and break that down into effectively
designing the model or the neural network architecture,
and then training that corresponding neural network architecture
so it does the right thing.
In this process, you're going to be playing around
with various kinds of optimizers, learning rates,
and a whole bunch of different kinds of machine
learning parameters that are extremely important in order to get the model
to achieve high accuracy.
But, of course, you need to actually deploy that trained model
onto a device.
Now, for embedded devices you're going to have to convert the model so that it
can actually run efficiently on device.
Or in some cases, even make sure that the big neural network
is shrunk down in size to quantization, or quantization of our pruning
or sparsity optimizations and so forth in order
to actually be able to get it onto the physical device.
Because the amount of memory sometimes you
have is only on the order of tens of kilobytes.
Once you get it onto the device, then you actually
want to make sure that the model is performing as you expect.
Because sometimes, the data in the real world
might actually be different from your training data.
So you have to have a continuous monitoring pipe.
Now, this particular course that you're in right now
is part of a larger series of courses, which I am not
sure if you're familiar with or not.
So I want to give you a little bit of context, primarily
so that you have the right resources as you need them based on your interest.
Course One of TinyML focuses very heavily on the fundamentals of ML,
where you learn about the language of machine learning.
You also learn how to program basic models, things like MLPs or CNNs.
And then you get some exposure to the concepts
of embedded systems, the resource constraints,
what are the challenges and so forth, just in terms of the devices
by themselves.
Then Course Two really focus on the applications of TinyML,
where we expose learners to various kinds of TinyML applications.
For example, keyword spotting or predictive maintenance
and manufacturing and so forth.
And the whole idea there is to actually try and program these models
in order to do those tasks.
In Course Three, the learners go on to actually physically deploying
the models onto an embedded device.
Such as using TinyML kit that we created for the course.
So through courses One through Three, folks learn the concepts,
how to program the models and physically interface those models
with sensors on a real device.
Now, Course Four, this one is quite different.
It's all about how you take this machine in the life cycle
and you scale it to more than one device.
We learn the operational procedures.
In other words, we learn how to iterate through the machine learning
lifecycle over and over and over.
Typically, when you go through Course One through Three,
you're just kind of understanding the mechanics of actually putting
a model onto a device, which is important and complicated by itself.
But in this course, it's not so much about the details
as much as it is about the bigger picture.
When you scale TinyML to more than one device, as in a real setting,
you will run into all kinds of problems that we will discuss next.
So in some way, the key takeaway of this video
is that we will be focusing a lot on this machine learning lifecycle,
and we'll learn how to operationalize it.
In other words, how to have a push button solution where you can actually
push a button and you can steadily go through this process over and over
and over in a very controlled fashion, using MLOps.
I'll see you in the next video.

\section{Review of Course 1, 2 \& 3}

Throughout this course, we will continue to use the terminology developed in prior courses. For returning students, this will be a recap of the basic language we have been using. For new students (welcome!), this will be an introduction to some of the fundamental parlance used within the machine learning community, as well as an overview of topics studied in courses 1-3.

At a high level, during the previous courses, we have covered the basics of both ML and TinyML, including their lifecycle, learning procedure, and examples of end-to-end deployment through case studies utilizing image-based data, audio data, and sensor data. These examples incorporate all stages of the developmental workflow, including data collection, model training, and model deployment onto a device.

\subsection{Course 1 Recap (Fundamentals of TinyML)}


We began this journey by starting from a blank slate. First, we illustrated the potential for TinyML as a game-changing technology, before then highlighting the challenges associated with its development and incorporation into commercial and consumer applications. Following this, we took a deep dive into machine learning. We learned that machine learning is inherently data-driven, and thus requires high-quality data (and usually lots of it) in order to learn associations. We learned that in supervised learning, we can use feature information (our X variable) in order to predict our output (our Y variable), and can use machine learning to “learn” the mapping between these two sets of variables, using a training set. Once we have developed this mapping, we can assess how well our mapping performs by comparing the performance of our predicted output on an unseen set of data, which we call the test set. Using this procedure, we can perform both regression and classification tasks.

Following this, we took an even deeper dive into the subfields of machine learning focused on neural networks, and, more specifically, deep learning. We learned that neural networks (NNs) are universal function approximators that allow us to approximate any non-linear function by learning from training data. Through a procedure known as backpropagation, we are able to update the weights in our NN to values that more closely match the function we are trying to approximate. This procedure is what is known as “learning” for a NN, and the rate of learning is controlled by a quantity known as a learning rate. NNs can handle any kind of input data, including image, audio, and sensory data. We then delved into convolutional neural networks (CNNs), which work much the same as traditional NNs but are specially designed to accommodate spatial correlations within data. This capability is especially important for images and provides superior performance compared to traditional NNs.

In the final parts of course 1, we saw examples of CNNs for classifying images, discussed TF Data as a method of streaming data, delved into augmentation and regularization techniques to combat overfitting in NNs, and had our first introduction to responsible AI development.

\subsection{Course 2 Recap (Applications of TinyML)}

While course 1 covered the fundamentals of machine learning, course 2 is where our journey into the world of TinyML really began. We started to look at the lightweight version of Tensorflow that Google has designed for microcontrollers, TF Lite, and how to convert our previously developed models into TF Lite versions that could be deployed on microcontrollers. We learned about quantization techniques that can be used for post-training quantization, different optimizations to consider during the conversion process, as well as more advanced techniques such as quantization-aware training. We will recap some of those concepts here as well, but only as they pertain to introducing you to new ideas that are specifically relevant to MLOps.

\GRAPHICSC{1.0}{1.0}{TinyML/3 Development/001.png}


A high level flowchart showing all of the steps in possible keyword spotting applications and how the keyword spotting task is a small part of a much large speech recognition application architecture. We see how KWS can set off a speech recording leading to a speech to text and then natural language processing model which will determine the request from a server which responds with speech via a text to speech model.

This culminated in several end-to-end implementations, the first involving keyword spotting (KWS). This application has become increasingly important in commercial applications, and is used in many smart devices for wake word detection, such as Apple’s “Hey, Siri” and Google’s “OK Google”. In our implementation, we taught a microcontroller to be able to respond to a small subset of words, such as “Yes” and “No”, by training it on a set of audio data. This use case introduced us to many of the difficulties associated with TinyML, such as the importance of data quality, the difficulty of debugging, and the importance of non-functional requirements such as inference time or model size. We were also introduced to the prospect of generating our own dataset, which is fraught with its own difficulties such as producing sufficient variation in the data to prevent overfitting and ensure good performance of the algorithm during inference.

Our next use case introduced visual wake word (VWW) detection: looking at an image to see if a given object is present. This is an important task with many real-world applications, such as person detection, mask detection (particularly prescient since the COVID-19 pandemic), animal detection, and so on. This use case provided us with a nice segue into image data, which is generally more complex to process and more computationally intensive than other forms of data, making it particularly challenging for resource-constrained devices like microcontrollers. To help build CNNs capable of working on microcontrollers, we were introduced to lightweight architectures suchas MobileNet, which utilize depthwise separable convolutions to minimize the number of computations per image. Transfer learning was also discussed, wherein high-performance networks such as VGG-16 and Inception can be adapted to the task at hand in circumstances where the original task is similar to the current task. Anomaly detection was also discussed as a separate use case that touches on the paradigm of unsupervised learning.

These use cases helped us to gain a more visceral feel of the ML workflow, going through the stages of data collection, data preparation, model design, model training, model optimization, and model conversion. We left the final stages, model deployment and model inference, for course 3, since this requires a microcontroller to deploy our models on!

\subsection{Course 3 Recap}

Course 3 is where all the fun began! Not only did we now have our models at the ready, but in this course we learned how to deploy them on a Nano 33 BLE Sense from our TinyML Kit and tick off the last two sections of the ML workflow: model deployment, and model inference.

\GRAPHICSC{1.0}{1.0}{TinyML/3 Development/Review_of_course_1_2__3..png}


The tinyML course kit from Arudino showing the Arduino Nano 33 BLE sense, the OV7675 camera, the tinyML shield, the usb cable, and the box that contains all of it.

This course is where we really delved into TF Lite, and its even lighter framework, TF Lite Micro. We also introduced to hardware-specific terminology, such as Arduino “cores”, embedded frameworks, mbedOS, and bare metal machines, and how each of these technologies are relevant to our TinyML applications. We spent time discussing the memory hierarchy and the trade-offs between memory size, cost, and read time, and how this can influence the performance of our system. Following this, we discussed low-level aspects of the TF Lite Micro, such as the FlatBuffer files used for model storage, the pre-allocated Tensor Arena used for performing intermediate calculations, and OpsResolver to import only the operations that are needed by our given model. After covering these essentials, we moved onto some end-to-end implementations of the projects covered previously, beginning with keyword spotting, then visual wake word detection, and finally the development of a gesturing magic wand that relies on accelerometer data embedded within our Nano 33 BLE Sense.

\subsection{What’s Next?}

Finally throughout Courses 1, 2, and 3 we explored Tiny Machine Learning through the lens of the typical end-to-end workflow from Data Collection to Training to Deployment to better understand all of the steps needed to deploy a successful TinyML application. Through MLOps we will expand beyond the workflow discussed in the previous courses and consider all of the things one would need to do to scale their models and optimize their performance. MLOps is an important aspect of machine learning that is often overlooked, but is one of the most essential skills for developing robust applications for commercial and consumer applications. Hopefully, you are familiar with all of the review material we have discussed. If not, we suggest you go back and look at some of the previous courses to help jog your memory.

\GRAPHICSC{1.0}{1.0}{TinyML/3 Development/002.png}

\section{Scaling TinyML}

VIJAY JANAPA REDDI: The goal of this video
is to help you understand some of the challenges we'll
face when we try to scale TinyML deployments beyond one device.
Recall this life cycle figure that I showed you in the previous video?
Well, this life cycle, which goes from data collection to actually deployed
and making inferences stage, is all tethered to just one device.
Now, that's fairly complicated, in terms of actually getting
the data to training a model to actually deploying it.
All of that is complicated.
I'm not dismissing that fact, but no matter what application you're building
and no matter whether you got it to work on the device or not,
scaling it to more than one device just steps the whole game up
to a whole new level.
Consider this example application of keyword spotting, where
the machine is basically responding to the keywords yes or no.
For yes, the green LED turns on.
For no, the red LED turns on.
Pretty cool, right, given that all of the processing that's needed
is actually literally happening on that tiny little embedded microcontroller,
which is smaller than my finger.
To do this keyword processing, there's an incredible amount of onboard device
processing taking place.
The leftmost figure shows us how an analog signal--
a continuous analog signal has come in from the microphone,
and it is then converted into a digital representation, which is represented
as a spectrogram in the middle.
And then the spectrogram, which is effectively a set of frequencies--
and the intensity of the frequencies is captured by the brightness
of the image--
is then fed into a neural network, which then just thinks
of it as a traditional cat or dog image, and then makes
an inference on whether the predicted word needs to be a yes or no,
and returns the result.
Now, to break this down more systematically,
let's take a look at the workflow that is actually involved.
First, we set up the device to get the data from the microphone.
Then we go through the process of building the machine learning model,
coding it in something like TensorFlow, or PyTorch,
or pick your favorite framework.
Then we work real hard to try and train the model to convergence--
in other words, to try and get it to a point
where it actually has high accuracy.
And then be evaluated for practical deployment
by pushing it onto the physical [INAUDIBLE] device.
At this point, if you've completed the full loop,
you've gone through the process of taking an analog input,
converting into a digital representation,
and then making a high-level framework-- like TensorFlow, for instance--
decide whether the word is a yes or no.
Or maybe it's perhaps just silence.
This in by itself is a fairly sophisticated and complicated process--
no doubt.
So if you've just done this, you should be proud of yourself
and give yourself a pat on the back.
Well, but there's so much more to it.
The question that you should be asking yourself after you've put it
onto the device is, what happens?
For starters, you're going to start making predictions on the device,
so undoubtedly, you're going to have to monitor how the device is actually
doing its predictions.
For instance, you'd want to look at what the false positive rates are
or what the false negative rates are.
In other words, you want to know how frequently is a machine actually
responding correctly or incorrectly to the end users,
because that is really going to have a dramatic impact on the users'
experience for this product.
Hence, it becomes extremely critical for us to manage the models
and keep them updated over time based on the feedback that we're getting.
Now, this is a lot simpler said than done.
Things like distributions and profiles of the users will change over time.
We'll get fairly deep into understanding the drifts,
and how to detect them and cope with them over the course of the remaining
sections, But for starters, let's run ahead,
assuming that we know that the model is performing on device.
Typically, you can decouple the life cycle in two major components.
The first part is what happens on the host side
when we're actually training the models, and the second part
is what happens on the device.
Often, when you're training models, you have the luxury
of using big and beefy workstations and high performance systems
to train the models robustly.
In other words, it's highly unlikely that your computer system is suddenly
going to stop working for some crazy apparent reason.
When was the last time, for instance, this even happened?
Probably when you dropped your notebook or something like that--
well, embedded systems suffer from these types of issues.
Remember, they're commodity parts, and these commodity parts
are what allowed them to be really cheap.
And this is going to affect the life cycle of the device, your ML device.
So in the context of scaling TinyML, the problems
really start on device, once you're done doing the benchmarking.
Now, the first part where I'm going to kick off
is with, what are the issues around just raw deployment?
Deploying models into the embedded ecosystem
is extremely challenging when you're thinking of scaling things up.
This is because there are wide and rich array
of different heterogeneous devices out there.
For example, just take a look at this little sampling
that I'm showing you here.
On the leftmost side, you have the Himax port.
In the middle, you have the Arduino and the SparkFun board.
And on the right side, you have the eSPI.
These devices are vastly different in their offerings and capabilities,
despite the fact that they're widely used
for educational and prototyping purposes.
Let's get a little bit more specific.
For instance, here I'm showing you the major capabilities
for these four different devices.
As you can see, their clock speeds are vastly different, ranging.
Anywhere from 64 megahertz all the way up to 400 megahertz.
That's a pretty big difference.
And that's important, because depending on how fast your processor can run,
it's going to affect how big or small your network is going to be.
The bigger the network is going to be, the more latency you're going to have,
which means, on one device, for instance,
it might actually run really well, because it has a higher clock.
But that exact model deployed on another device
is going to have a totally different impact,
because if the clock frequency is only 64 megahertz,
it's going to run substantially slower, which is going
to impact your end user experience.
So the same model's going to have different experience once put
into production in the field.
Similarly, the memory capabilities, both in terms of flash storage and RAM,
are also vastly different.
Flash is typically in the order of 1 or 2 megabytes.
RAM is in the order of hundreds of kilobytes.
But that said, these systems are still quite different.
If you look at their flash capability and their RAM capability,
there's still an order of magnitude difference.
Why am I bringing this up?
Because if you're able to push a certain binary that has your machine
learning model to one device, then you want
to be able to push that exact model, which you're very proud of, because it
retains a really good accuracy-- you might not
be able to push that exact model to another device that
has much smaller resources.
Typically, these are problems you don't ever
encounter in a high performance [? system, ?]
because they tend to have a huge amount of memory, really fast computing
systems by default.
So in the embedded systems world, these kinds of things
start becoming a major problem for production operationalization.
Now, coming back to the big picture, what I'm trying to really emphasize
is that you have to worry about whether you can achieve portability
across devices while dealing with the heterogeneous aspects when
it comes to deployment.
Now, next is the issue of effectively being
able to monitor the device when it's out in the field.
Remember, I said that you can push the model out onto the field,
but then you really have to know how it's performing in the field,
because that's what ultimately ends up impacting the user experience.
Now, in a traditional computing environment,
you will take the network connectivity for granted.
In other words, being able to have a network connection, internet connection
so you can send data back to the server is
something you don't even sweat about.
It's really quite common that you're going to have that connectivity.
But that is most definitely not the case when
it comes to the embedded ecosystem.
Why is that?
When you look at these embedded devices, network connectivity
is somewhat intermittent.
It really depends on the device.
Cost is a big issue.
So often these embedded devices try and not even support
any kind of network connectivity.
And without radio connectivity, then what ends up happening
is you won't be able to monitor things in the field.
And let's say, if you have a device that may be a little bit more costly,
but has the communication capability, then continuously connecting
to that device in order to monitor how the machine learning model is
performing there is also dangerous, because it
will consume way too much energy.
And remember, our models are supposed to be running on battery.
Let me get a little bit deeper into this.
It's a well-known fact that, when you compare
the amount of energy required for computation which is communication,
computation is really cheap.
Communication or moving data is extremely costly,
in terms of energy consumption.
The plot on the right here is trying to show you
how much energy is consumed between computation and communication.
The blue bars corresponds to various forms
of computation inside a microprocessor.
The red bars correspond to memory accesses or date of movement
within the microprocessor.
And you can clearly see that there's a large difference--
an order of magnitude easily-- between computation versus communication
even within a single system.
Now, when you look at the cost of moving data across the network,
be it through Bluetooth, or LoRaWAN, or some other kind of communication
technology, it is several orders of magnitude costlier, which
means that the more you communicate with this endpoint device,
the more energy you'll be using just to get data from it--
which means you are likely going to be able to run
the device for a very short amount of time, given a finite amount of energy.
And that is an extremely challenging thing you have to address.
I know of production use cases where over 50\% of the energy
is just dedicated for communication alone,
because it is so important to understand how
the model is performing in the field.
So my takeaway message here is and monitoring embedded devices is not
something you can take for granted.
And this is one of those things that's going to have a big impact on the MLOps
pipe, where traditionally, people take communication
with the model for granted.
Last, but not least, let's talk about security and robustness.
It's extremely critical when you're deploying machine learning
models at the endpoint to worry about the security of the model.
Typically, these ML devices are going to be used for important applications,
like monitoring patients' health.
In such a case, where you have access to a lot of that data that
is supposed to be private, you need to make sure that device is secure.
And it also needs to be robust, because you're dealing with embedded devices.
They tend to have physical sensors.
You're physically dealing with an embedded device that can actually
break That becomes a problem.
So as you can see, this device is locked and loaded
the whole different set of sensors.
All we are trying to do is try and convert all the sensor information
into raw intelligence that we can actually harness.
So to this end, if you look at all these different devices,
they have many different kinds of sensors.
Now, there are some problems here as well.
For example, a microphone on the Arduino device
can be quite different from the microphone on the eSPI device.
What this means is that you cannot assume that,
just because you have a microphone, your machine learning model is going
to perform identically on the Arduino versus the eSPI.
Why?
Because the signal processing that is going
to happen in converting the audio signal in from the analog input
into a digital representation is going to have slightly different nuances,
and that's going to affect your model's behavior again.
So yet again, you get into this issue where suddenly this machine learning
model is interacting with these physical attributes that actually
start introducing new challenges.
Another thing is that, given that TinyML devices are
likely to be deployed, say, in an industrial manufacturing plant
or in wildlife conservation area, these sensors shouldn't break--
I mean physically break.
So this raises another question of, how can you robustly deploy TinyML metal
devices into production?
You can't assume you just have one microphone.
You might have to have a backup microphone.
These are all parts and things of MLOps that you
have to do that specifically come up only in the context of TinyML.
And that's why I'm bringing them up here.
Again, the point I'm trying to make is that it's
one thing to benchmark TinyML deployments on your desk,
but it's a whole different thing to take the device out of the production.
Where you're exposed to extreme heterogeneity
in terms of the device's ecosystem, then you
face challenges in terms of monitoring because of the communication
limitations, and then ensuring that these devices can actually
be safe and secure, and that you can actually rely on them robustly.
And this, my friends, is why it's so important
to learn embedded machine learning operations.
TinyML is still a very nascent field.
There are loads of challenges that we have to overcome,
so unless you understand these challenges
and you have a deep appreciation for how to operationalize things,
you will be doomed to be the statistic--
1 out of 87 companies that fail out of hundreds of them.
So let me try and help you understand MLOps so
that you can address these challenges.
Specifically, we look at embedded machine
learning operations in this course.
I'll see you in the next video.


\section{MLOps Overview}

VIJAY JANAPA REDDI: In this video, I want
to give you a little bit of an overview of what MLOps is all about.
Now, recall that MLOps, are machine learning operations,
is a collection of processes that are aimed
at reliably and efficiently designing, deploying, and maintaining machine
learning models in production.
Let's try and understand what this means a little bit
further through the lens of this ML life cycle that I've showed you earlier on.
Now, in order to execute this workflow, you typically
are going to have to write a lot of code.
Doesn't matter who you are or what awesome project
you're working on-- you're going to have to go
through writing this code, period.
And I mean you're going to be writing a lot of code.
That sounds alarming.
If you're anything like me, I'm an extreme programmer, always
eager to prototype, get something out of the door real quick
by programming quick and fast.
Now, that's great for a proof of concept demonstration,
but it's a whole different thing when it has
to be something of production value, which is something
scalable, and robust, and so forth.
And that my friends really is where MLOps comes in.
It's all about automating this end-to-end flow
so that everything is smooth, repeatable, and robust.
Simply put, MLOps is taking that workflow and coupling
with automated flows, flows that are robust
and flows that are repeatable, with the push of a single button.
Broadly, you can bucket MLOps into five focus areas.
Running end to end means going through the entire ML life cycle process
I showed you in the previous two videos in a single swoop.
The managing complexity means being able to work
through all the code you're going to have in that as a collective team.
Evaluating results means repeatability, which is key in machine learning.
ML, due to the nature of data, is not necessarily repeatable.
The data can drift, and that can lead to problems.
Improvement means having flows that can be systematically monitored,
and that monitoring has got to be such that it can actually
trigger data collection or retraining of the machine learning model in order
to perform better at the new trends that we're actually witnessing.
Tracking means having versioning in place
so that you can almost always be accountable for how your model behaved
on some crazy wild day.
In typical flow, you might have a data engineer,
for instance, who's responsible for gathering
some label data that's actually going to be used to train models.
Then you'll likely have a data scientist who's using that data,
writing up a bunch of code, and tuning a whole bunch of parameters
in order to get a good machine learning model to perform well in the field.
Then, ultimately, you're going to have a DevOps team or an ML engineering
team that will push that trained code into production,
so that's serving on the end device.
Now, in reality, of course, there are many more components and lots
of different people involved in this.
But the reality is that a lot of this interaction between multiple people
means that you need to streamline the process.
Effectively, you're going to have this problem where
you have too many cooks in the kitchen.
What that means is that, because you have too many cooks in the kitchen,
you might not actually end up with the end goal that you actually shot for,
and you might end up somewhere else.
That's the nature of the job, if you are not systematic about the processes.
Now, this figure is one of my favorite in this particular course.
The reason is that it helps us understand,
why is it that, compared to traditional systems which we have been building,
machine learning systems are so different that we need to have
an entire MLOps concept around it?
Now, in a traditional system, you have some code that runs on a system.
This code can easily be tested.
You can write a whole bunch of unit tests and integration tests,
and you can monitor that code and determine how it's running on device.
Now, this is true whether it's an embedded device, or a workstation,
or something else.
But on the right, we see an ML-based system.
First thing you notice is it's got a lot more components and pieces in it.
First, you're going to have to train the code, which
results in your neural network that needs
to be trained on the right set of data.
Then you're going to put it into production, where you don't have
to subjected to a whole bunch of data.
But this is where things get really tricky,
because in addition to just having the code for the neural network,
you're going to have data which is actually
going to cause you network to behave differently based
on the characteristics of the data.
In the past, if you had some data and you fed it into some code
that you've written in a traditional system,
you're going to get the exact same output.
But with the neural networks, because the data can actually shift quite a bit
and the network is going to behave based on past trends it's seen,
it can actually behave differently during production from what
you think it might actually do.
So my point is that, in machine learning systems,
you not only have code that you have to worry about
as you did in the traditional sense, but you also have to worry about the data.
This is a new element in the system stack.
OK?
So code and data together, versus just dealing
with code-- that's a whole different level of additional complexity.
And this is why tooling is becoming an extremely important portion of MLOps.
You can piece together MLOps flows through a bunch of Python scripts,
or you can actually build on top of existing
frameworks, such as TensorFlow.
Or more recently, especially in the context of embedded machine learning
systems, we're seeing the emergence of startup companies,
like Edge Impulse and many others, that are
trying to build end-to-end platforms that are dedicated for embedded MLOps.
Let's take a look at some of the more traditional systems.
TensorFlow eXtended is Google's solution to deal with the issue
that they were seeing, where people were taking TensorFlow and using
it to train models, and then trying to put them into production.
And they realize that just TensorFlow for training alone was not enough.
They needed to have an ecosystem that actually
supported the effective deployment of TensorFlow tools.
And therefore, they ended up building TensorFlow eXtended.
Similarly, SageMaker is from Amazon.
It's also an example of an end-to-end framework.
It focuses, much like TensorFlow does, on cloud ML.
Uber developed an eXtended framework called Michelangelo quite a while back,
and they did this because they needed to power their machine learning services.
Even Airbnb, Netflix, Lyft have all developed their own in-house MLOps
frameworks.
So we should be asking ourselves this question--
why is it that these companies are all going and building MLOps tools?
Well, first and foremost, they all realized one thing,
which was that there was no consistency between the existing eXtended workflows
that were there.
People were trying to stick together different boxes and all together,
and then none of which was really gelling together.
So they caused a lot of confusion when different teams were
trying to interact with one another.
Especially when you have a new idea and you're creating a new team,
you don't want them to invest all this effort
into creating a systematic MLOps pipe.
You want to have one nice clean design, where new teams can basically
take that MLOps tool and actually use it, and just get
jump started on what they really want to produce in terms of value,
not recreate a MLOps pipe.
Now, there are also a wide array of different ML use
cases across the company.
So you have to realize that no two ML tool
chains that are developed within a group are going to be identical.
So you have to somehow deal with this, because that heterogeneity, again,
causes problems in terms of communication between teams.
And also, the existing ML workflows are often too slow,
and they're fragmented and brittle, obviously,
because they're just kind of these makeshift solutions every team was
building.
This is really why we ended up seeing that there
was a real need for MLOps tools that were really robust and solid.
They wanted to build it once and forget about it,
and make sure that it just works.
These are the real reasons.
So first, you want to be able to have a way to manage and process data.
You want to be able to train the machine learning models.
You want to version the models, much like we do with GitHub for code.
We want to be able to evaluate and compare different versions of models
quickly and easily so we know which one to deploy.
And when we do deploy, we want to deploy them very robustly into the system.
The devil's, of course, always in the details,
and that's why you end up with all these different flavors of types of machine
learning frameworks that are supporting operations.
So while the existing frameworks that I gave you
examples of are really focused on cloud MLOps,
one thing I do want you to keep in mind is
that those cloud MLOps frameworks don't directly
apply for the embedded ML ecosystem.
So you can't just take something like TensorFlow eXtended or you can't just
take something like SageMaker and assume that it'll just magically
work for embedded MLOps.
And that's what this course is all about.
We're going to be talking about the MLOps piece,
and then we'll try and understand what are
the nuances that make it extremely challenging for the TinyML context.
Now, over the course, we will go through many different sections.
Some of the concepts that I've listed here are what we will be touching on.
Now, the way to kind of go through this could
be that I could just talk to you about this based on my experience
with helping startups.
And so forth.
But what I've chosen to do, which I found works well,
is to actually work in collaboration with industry
in order to help you understand these concepts so
that you can connect the concepts to some
of the real-world practices in industry.
So to this end, I will be having Lara Suzuki, who's actually
one of the leads at Google working on MLOps, and responsible AI,
and so forth, actually lay out how Google actually
approaches MLOps in their lens.
And I will then go into the details of explaining
how it fits or does not fit the issues that surround embedded machine
learning.
So throughout the course, we will be following this waterfall chart.
This waterfall chart has got multiple stages of MLOps laid out,
starting from ML development all the way down
to actually doing continuous monitoring and doing management as a whole.
The whole course is structured around this beautiful picture.
We'll be stepping through each one of these blue boxes.
Every section in the course is laid out as one of these blue boxes.
So the way you should approach this is that you're
going to focus on one blue box, move on to the next one.
Having said that, one thing that's extremely crucial
is that there are lots of strong interdependencies
between these blue boxes.
So it's not that it's a clean waterfall chart.
It's actually a very tightly coupled design.
But for the purpose of giving you some structure,
we're going to follow this waterfall chart.
With that said, I want you to remember that you're always going to MLOps.
Whether you do it poorly or whether you do it really well doesn't matter.
Point is you're doing MLOps.
So if you're always doing MLOps, you might as well
make sure you're doing it well.
With that said, let's try and understand how the course is structured
in terms of readings and so forth, and then
get started with the actual material.
I'll see you in the next couple of videos.

\section{Overview of MLOps}

\GRAPHICSC{1.0}{1.0}{TinyML/3 Development/overview.png}

While many organizations recognize that Machine Learning (ML) can drive significant value, successful deployments and effective operations are the main bottleneck for gaining value from AI. As of today, more than 87\% of data science projects never make it into production.

To support organizations in coming up to speed faster in this important domain, it is important to understand ML Operations (MLOps).

\subsection{A Brief Introduction to MLOps}

Teams across all industries and domains find themselves spending an unnecessary amount of time on the development of ML models, without the necessary people, processes, or technology to efficiently deploy and operationalize them. In the current landscape, we find a high degree of manual, one-off work as many engineers are constantly crafting custom solutions, inflexibility of the resulting deployments as components are not reusable nor reproducible, and a high number of errors as these custom solutions are often poorly documented leading to issues during handoffs between Data Science and IT. As such, ML solutions often fall short of their goals and cost their organizations time and money. In summary, deploying ML models at scale is very challenging, especially due to:

\begin{itemize}
    \item Organization Barriers - Managing different ecosystems like programming languages
    \item Compute Constraints - Continuous and reliable compute resources (dedicated servers or cloud-only options)
    \item Portability Issues - Huge dependencies on legacy systems
    \item Seasonality - ML workloads works in patches; this need auto-scaling capabilities
\end{itemize}


Businesses are looking to reduce the time from change (via data or code) to deployment, improve deployment frequency, speed up time to value with new ML use cases, lower the failure rate of new releases, shorten the lead time between fixes, and facilitate better collaboration (reusing features and sharing models).

To address this, organizations need to build the necessary ML engineering culture and capability MLOps aims to unify ML system development (ML) with ML system operations (Ops). MLOps strongly advocates automation and monitoring at all steps of ML system construction, from integration, testing, and releasing to deployment and infrastructure management. It takes both its name as well as some of the core principles and tooling from DevOps. This makes sense as the goals of MLOps and DevOps are practically the same: to shorten the systems development life cycle and ensure high-quality software is continuously developed, delivered, and maintained in production. Its Machine Learning’s unique challenges and needs – managing the lifecycle of Data, Models, and Code – has led MLOps to quickly evolve as a domain of its own.

To support the engineering community to realize MLOps in practice, Google made available to the world the practitioners guide to MLOps to help developers implement MLOps and smart practices during the ML workflow. We will go through this applied pipeline and we will discuss the unique challenges that emerge as a result of deploying TinyML to many embedded devices.

Amongst many processes and operations, MLOps requires:

\begin{itemize}
    \item Continuous Integration: Is no longer only about testing and validating code and components, but also testing and validating data, data schemas, and models.
    \item Continuous Delivery: Is no longer about a single software package or a service, but a system (ML training pipeline) that should automatically deploy another service (model prediction service)
    \item Continuous Training: This is a new property, specific to ML systems, concerning automatically retraining and serving the models.
\end{itemize}

The maturity of the ML process is defined by the level of integration and automation of these phases shown in the figure below, which reflects the velocity of training new models, quality of the assets generated, and reliability of the overall system.

\section{Course Structure}

VIJAY JANAPA REDDI: In this video, I'm going
to talk to you a little bit about what the structure of the course is.
And the reason I want to do this is that I
want to give you a little bit of a framework for how
to think about ingesting the information that's going to be coming to you.
As I mentioned earlier, everything in this course
is structured around this waterfall diagram.
We will be using this as our platform for actually understanding the MLOps
pipe.
We will be going from ML development to ML training to continuous training
to converting models to model deployment to prediction serving to continuous
monitoring-- a lot of buzzwords at this point, but every single one of them
will be defined systematically.
For instance, the first module is going to be looking at ML development.
We'll open this up, and the way we'll open this up
is by first having Dr. Laura Suzuki from Google,
as I mentioned earlier, come and give you
the big picture of what happens inside the ML development stage.
In every single one of them, we're going to open up the box,
and it's going to be this fairly complicated system inside.
Then we'll close the box up, and then we'll proceed.
But for every single one of them, after Laura opens up
that box, what I'm going to do is then dive deep into MLOps in a way
that is specific to what we see with TinyML.
And I'll give you a principles first sort
of a perspective about the things that actually work
and the things that don't work.
In other words, for every module, we'll start with the big picture overview.
Then we'll talk about how that big picture overview-- which
is something that would be applicable in the context of cloud-based machine
learning-- might not apply when you're talking about TinyML.
And then I'll talk about some potential solutions
for how we can scale TinyML through the lens of MLOps.
And through this sort of a staggered approach, what I'm hoping to do
is, for every single module, you get the big picture of that particular module,
then you go deep into what are the challenges,
and then you understand some potential solutions.
Now, look, one thing I want you to know down here
is that, for some of these problems that I'm going to expose to you,
there are no solutions, my friend.
That's potentially an opportunity for you to create the next startup.
What I'm trying to do down here is that, as this nascent field is starting
to bubble up, I'm trying to identify what are the common pitfalls that most
people will run into, and try and lay that out
so that you can, ahead of time, think of alternative approaches
you might be able to take.
So keep that in mind as you go through this course.
I have a reading that comes up after this that lays out all the course
material that we will be providing, which is in terms of readings.
We'll have case studies, we'll have forum discussions-- a lot
of engaging activity, where you can take some of the concepts
that I'm talking to you about and then try and practice
them, and interact with the students that are on the platform
along with you.
With that said, I'll see you in the next video, where
one of the key things I want to talk about
is the mindset you've got to have as you go through this course
we'll see you there.

\section{Course Activities}

\GRAPHICSC{1.0}{1.0}{TinyML/3 Development/topics_course.png}

MLOps is a broad topic. As such, there is no one size fits all approach for helping you learn. For instance, loading you up with programming exercises does not help if you are a business leader just trying to make sense of the big picture (the operations needed). At the same time, if you are a data scientist you want to understand what bag of tricks are available for you to put your toolbox together so that you can put the concepts into practice. Given the various MLOps topics, we decided to have a sampling of everything. Ultimately, the goal is to help everyone, from business leaders making decisions to software engineers developing the ML maintenance pipeline, to have an appreciation for each other’s role in deploying ML at scale into production. 

To this end, the course is based on Videos, Google Colabs, Readings, Case Studies, Discussion Forums, Tool Explorations, Quizzes and Tests (for taking the course as part of a certificate program).

Videos are the main means by which the instructors communicate directly with you. Each video is generally between 5 to 10 minutes as we find that that’s what is most effective for retaining learners’ concentration. Each video focuses on specific topics so that you can come back to it.

Readings supplement videos. Sometimes they will be used to add more detail to a topic covered in a video. Sometimes they will be used to introduce or summarize additional topics.

Case Studies in the form of readings provide a glimpse into some real-world deployments that we have gathered from starts-ups and other major leading organizations that have put MLOps into practice. These case studies talk about the challenges of deploying ML into production for a given use case. More often than not, as you will find out, MLOps is less about coding and more about understanding what are all the components that must click together. 

Discussion Forums exist so that you can connect with your fellow learners and learn from one another. These forum activities are generally preceded by some open-ended questions that encourage discussion amongst learners. The goal here is to solicit a broad range of opinions. 

Quizzes: will act as knowledge checks and will be short sets of multiple-choice questions (MCQs) sprinkled throughout the course to make sure you are learning the key themes. 

Tests: Summative quizzes (aka tests) will only be available to those pursuing the certificate and will act as formal knowledge tests. They will be a longer test of MCQs that will ask conceptual questions about the course content covered in the videos, readings, and collabs, and will also ask questions about the coding assignments.

Optional: The following are optional, which we provide in the additional resources section for each module.

\begin{itemize}
    \item Some sections contain optional labs. Google Colabs provide hands-on learning with code in the cloud using Google’s Collaboratory environment and host programming assignments reducing compute resources on your local machines. Don’t be alarmed if you aren’t a programming guru, you will be able to follow through by reading the detailed instructions and comments we have provided. But if you are a programmer, you will be able to exercise your Python programming skills and explore further. 
    \item Tool Explorations will give you the opportunity to explore a variety of external tools that can help you with MLOps at your organization. Again, as this course is targeting a variety of learners, we expect some learners will simply want to know what the tools are, while others will dive in and test them out. In most cases, we will try to not only link the tools but also to tutorials that we think do a good job of explaining how to use the tools.
\end{itemize}

\section{Your Mindset: T-Shaped Skills Needed for ML Engineer}


As I mentioned earlier, in this course, we'll
be covering a lot of different MLOps concepts, ranging from ML
development through the continuous monitoring and everything in between.
Now, that's a lot of ground to cover.
And sometimes you might be wondering, do I really need to know all of this?
When in reality, you might be a data engineer, or data scientist,
or product manager, or a software engineer.
Well, clearly, you can't be it all.
But that's OK.
The goal here is not for you to master everything.
It is more for you to have a broad awareness
of how all these different pieces fit together so that you can be better
at the job that you do.
To that end, I want to introduce it to the T-shaped skills mindset
that you should have for this course.
The T-shape is generally used to describe a person
with deep vertical skills in a specialized area, such as TinyML,
with broad but not necessarily very deep skills in other relevant areas.
Now, this is a concept that Tim Brown actually came up with,
where he was a CEO for a California design firm
where they actually actively think about these kinds of processes
in order to help companies be successful.
And companies love T-shape employees.
Machine learning needs full-stack engineers,
someone who is capable of addressing a bigger part of the system.
You need expertise in data collection.
You need expertise in building continuous integration pipes.
You need expertise in effectively deploying models into production.
So being a T-shaped employee means having
a little bit of knowledge about all these different kinds of things
and knowing how everything clicks together.
And that's why companies love T-shape employees,
because it saves them a lot of time and money.
More importantly, it makes you more valuable to the company.
As I've said earlier, T-shaped individuals
have deep vertical skills in a specialized area.
And in our case, your depth is going to come from TinyML, or just ML
at large, since TinyML is really just a specialization of applied machine
learning.
And the breadth itself involves experience and knowledge
about relevant topics and deploying machine learning at scale.
And to that end, MLOps is your breadth.
Being knowledgeable about MLOps and TinyML
will enable you to be flexible, collaborative, and ultimately
make you a good leader.
You'll be flexible in that you'll be able to take on new tasks
and help your team members out so that the project can ultimately
be successful.
You'll be collaborative, because by learning a little bit
about each other's expertise, you'll actually
know how to deal with your team members challenges.
And you'll understand the needs as a whole.
And ultimately, you'll be a good leader, because you see the big picture.
You'll see how you can pull insights and inputs from various members of the team
to make the project successful.
And this is why companies love T-shape employees, especially
in machine learning, because they add a lot of value to them.
So keep this in mind.
And with all that said, I'm going to leave you
in the good hands of Dr. Lara Suzuki from Google
who will introduce you to Google's practical MLOps
approach, which we are adopting as an industry practice in this course.
But that's not to say that this is something very Google specific.
I've made sure that everything we cover here is broadly applicable.
Dr. Lara Suzuki will first paint the big picture of MLOps in the next module,
going one level deeper than what we talked about so far.
She'll explain how it differs, for instance, from [INAUDIBLE] DevOps,
and so forth.
And then I'll come back after you watch those videos
and open up each of these individual stages and do a lot more depth.
And then we'll progressively move through the course.
With that said, have a good time, and I'll see you soon.

\section{Summary of What You’ll Learn}

\subsection{About This Reading}

In this course, we are going to introduce you to MLOps from a TinyML perspective. Lots of ideas relevant to MLOps apply to both larger-scale applications as well as small-scale applications.

We refer to these traditional larger-scale applications as “BigML” in contrast to the TinyML that we have already familiar with in the context of embedded computing systems.

MLOps has become an integral part of machine learning in industry and performs analogous functions to DevOps in software development, except it focuses on both code and data.



MLOps pipeline. This diagram is applicable to both BigML and TinyML. Going from ML Development to ML training to Continuous Training to Model Deployment to Prediction Serving to Continuous Monitoring. Along the way artifacts are produced at each step including: Code and Configurations, Training Pipelines, Registered Models, Serving Packages, and Serving Logs. All steps tie into the overarching themes of Data and Model Management.

\GRAPHICSC{1.0}{1.0}{TinyML/3 Development/12.png}


\subsection{Overview}

At a high level, the differences between BigML and TinyML are threefold: hardware capabilities, software capabilities, and network capabilities. 

Hardware: TinyML works with resource-constrained devices that have limited hardware resources to perform functions. This means that we have to be very efficient with our implementations and minimize resource usage in terms of memory, compute, and communication. In the cloud, we have access to very large amounts of computing resources and data storage, with large caches available on-chip for performing computations. In contrast, mobile devices using processors such as the Cortex-A are more constrained than the cloud (as shown in the figure below), but they still contain relatively large amounts of data storage along with a memory management unit (MMU) for storing memory off-chip. TinyML is even more constrained than mobile ML, having all the memory (SRAM, cache, and eFlash) on-chip, no MMU, and a very small amount of resources available. This presents additional challenges for realizing MLOps.

A figure showing the difference between cloud, mobile, and tiny ML systems as pyramids. Cloud systems have Flash (1TB), DRAM (80GB) off chip and L2 (40MB), L1 (192kb), and a GPU processor on chip. Mobile systems scale that down and have Flash (64GB), DRAM (4GB) off chip and L2 (4MB), L1 (48kb), and a Cortex-A processor on chip. Embedded tinyML systems instead rely on Cloud as the basis of the pyramid followed by on chip base of SRAM(320kb) and eFlash (1MB) with Cache (4KB) and a Cortex-M processor built on top.

\GRAPHICSC{1.0}{1.0}{TinyML/3 Development/11.png}


Software: A corollary of the hardware constraints is the software constraints - the minimal resources available within TinyML means that we are unable to run any software other than things like extremely lightweight clients. This means that traditional software used for MLOps within the realm of BigML cannot feasibly be used, like Docker, Ansible, Kubernetes, Git, and TensorFlow. This brings about constraints that suggest we will probably need a whole new set of tools to perform MLOps in TinyML. TinyML systems are also more heterogeneous than their BigML counterparts, leading to more heterogeneous software operating environments.

Network: Because communication between devices requires a lot of power, the network capabilities of TinyML devices are also more constrained than in BigML and must be treated more conservatively or use more power-efficient wireless technologies for communication. Wireless technologies such as NB-IoT and LoRa are more preferable to the standards used in mobile devices such as 3G and 4G.

\GRAPHICSC{1.0}{1.0}{TinyML/3 Development/comms.png}

A figure showing the difference between cloud, mobile, and tiny ML systems as pyramids. Cloud systems have Flash (1TB), DRAM (80GB) off chip and L2 (40MB), L1 (192kb), and a GPU processor on chip. Mobile systems scale that down and have Flash (64GB), DRAM (4GB) off chip and L2 (4MB), L1 (48kb), and a Cortex-A processor on chip. Embedded tinyML systems instead rely on Cloud as the basis of the pyramid followed by on chip base of SRAM(320kb) and eFlash (1MB) with Cache (4KB) and a Cortex-M processor built on top.

The table below summarizes the differences between BigML and TinyML in the context of (1) compute, (2) software, and (3) network capabilities.

\begin{tabular}{ll}
    BigML & TinyML \\
    Big compute hardware & and memory resource availability \\
    Small compute hardware  & and memory resource availability \\
    Homogenous software operating environments & Heterogeneous software operating environments \\
    Stable network and communication capabilities & 
    Unstable network and communication capabilities \\
\end{tabular}

The Ops corresponding to each of these paradigms (BigML vs. TinyML) is correspondingly different in their characteristics, although the overall procedure remains largely unchanged. TinyMLOps is associated with resource-constrained edge devices, small models, power consumption, and fewer capabilities than BigML systems, such as logging and debugging. Conversely, traditional MLOps is associated with powerful devices, large models, and a large number of capabilities. The differences are summarized in the table below.

\begin{tabular}{ll}
    MLOps & TinyMLOps \\
    Models are trained and deployed on powerful devices & Models are deployed on resource-constrained edge devices \\
    Large models with multiple supported hardwares and Ops & 
    Small models with few supported hardwares and Ops \\
    Accuracy and availability are the important metrics that matter
    & Latency, throughput, power consumption are more important \\
    Containerization, CI/CD, logging, and monitoring is required
    &
    No containers; logging and monitoring is difficult to do
    \\
    Performance checks and model updates are done regularly
    &
    Performance checks and model updates are difficult to do
    \\
    Robustness and autoscaling to workload traffic
    &
    Scaling is difficult after deployment; Robustness is important
    \\
\end{tabular}

\subsection{What You Will Learn}

In this course, we will build upon the topics discussed above with the aim of answering the following questions:

\begin{itemize}
    \item What is (Tiny) MLOps?
    \item Why is (Tiny) MLOps necessary?
    \item How do we automate (Tiny) MLOps?
    \item What are (Tiny) MLOps platform features?
    \item What are some examples of (Tiny) MLOps platforms?
\end{itemize}

By the end of this course, you will understand what TinyMLOps is, why it is necessary, how it works, examples, and the stages involved in developing a TinyMLOps workflow. Let’s get started!

% 1.2: MLOps: The Big Picture

\section{Overview of MLOps Objectives}

\subsection{MLOps: The Big Picture - Module Objectives}

\GRAPHICSC{1.0}{1.0}{TinyML/3 Development/MLOps_Overview.png}

\subsection{Objectives}

The purpose of this section is to give you an overview of MLOps. First, we will start off talking about what are the differences between AIOps, DevOps and alas MLOps. There is often a lot of confusion between these terms so let’s get the definitions right. We are focused only on MLOPs in this course. Briefly, MLOps sits at the intersection of data science, data engineering and software engineering and it is extremely important to successfully make use of ML in practice.

The main topics we will cover in this section include the following:

\begin{itemize}
    \item What are AIOps, DevOps and MLOps?
    \item Why is MLOps so important to get right to be successful?
    \item What are the roles and responsibilities of people involved in MLOps?
    \item What are the key MLOps activities?
\end{itemize}


Keep in mind that in this section we want you to get the core concepts of MLOps right. We focus on the big picture. Subsequently, the later modules will get deep into the details of each section. Also, we will then start to focus on issues that are specific to embedded machine learning, effectively using it as a case study for deploying ML at the edge/endpoints.

\section{What is MLOps, DevOps, and AI Ops}

LARA SUZUKI: Hello.
Welcome to the MLOps for a TinyML course.
In this introductory session, we will learn the concepts and main differences
among AIOps, DevOps, and MLOps.
We will also examine why MLOps processes are
important to machine learning and the personas involved in its lifecycle.
Let's take a look at why we need MLOps ops in machine learning systems.
AIOps represent the intersection of IT operations and artificial intelligence.
It defines the application of AI to enhance IT operations.
More specifically, AIOps use machine learning, big data, and analytics
to ensure that we can collect, analyze, and aggregate data from vast amount
of information generated by multiple ICT infrastructure
components, applications, and performance monitoring systems.
It also intelligently extract signal out of noise to identify and predict events
and also patterns related to system performance, security,
and availability issues.
It's diagnoses the root cause of those and report
them to IT teams to ensure rapid response and remediation.
Or in some cases, it automatically resolve these issues
without having to rely on human intervention.
There are several advantages with replacing multiple fragmented
and manually operated IT operations systems
with a single intelligent and automated and sometimes autonomous AI platform.
It equipped teams to respond timely and even
proactively to slowdowns, security threats, and outages
with a lot less effort.
AIOps also bridges the gap between an ever increasing diverse dynamic
and difficult to monitor IT landscape and user expectations.
User expectations are for little or no interruption in application performance
and availability.
Most IT experts consider AIOps to be the future of IT operations management.
In the case of software engineering, the intersection
between IT operations, software development, and application delivery
forms what we call DevOps.
To start from the beginning, let's take a look
at the life cycle of traditional software.
It's arguably quite straightforward.
At its simplest, you develop, you test, and deploy the software,
and then release new versions with features, updates, and/or
fixes as needed.
To facilitate the process of development, testing, deployment,
and further releases, the traditional software development
can rely on DevOps.
And traditional software development can then
be automated and rely less on manual processes which
are time-consuming and prone to errors.
Amongst the process of DevOps, we highlight the concepts
of continuous integration, also known as CI,
which is all about automating the integration of code changes created
by multiple contributors into a single software project.
We also have continuous delivery, also known as CD.
CD is all about delivery and deployment.
It means that we have got the software that
has been going through the various processes
and now we are ready to push it out at the door and deliver it to end users.
And we also have continuous testing, also known as CT.
It is the process of executing automated test as part of the software delivery
pipeline.
These are employed to reduce development time
while continuously delivering new releases
and maintaining software quality.
In machine learning, we have a different process called MLOps.
It is the integration of software engineering, data engineering,
and data science processes.
Let's take a step back and understand what
happens in ML systems which is the subject of our course.
When it comes to machine learning systems,
it's easy to assume that the ML workflow will follow a similar process
as in software engineering.
After all, it should be just a matter of developing, training an ML model,
deploying it, and releasing a new version as required, right?
Not quite.
The environment in which ML systems operate complicate things.
And that is why we need a different approach
as to traditional software engineering.
For starters, the ML systems themselves are
very different from traditional software due to their data-driven
non-deterministic behavior.
In software engineering, the code drives the behavior of the system.
In machine learning system, the data will drive the behavior of the system.
As such, any change in the behavior of the data, for instance,
a changing its distribution, will lead to machine learning system
to operate differently.
In an ever changing world, ML practitioners
must anticipate that real-world data on which production ML models refer
will inevitably change.
To support the complexities involved in production [INAUDIBLE] ML system,
a special kind of DevOps for ML has emerged, hence MLOps, which
stands for machine learning operations.
MLOps you manage the constant changes in ML systems
and the subsequent need for model redeployments.
MLOps embraces DevOps continuous integration and continuous delivery.
But it replaces the continuous testing phase with continuous training.
Now, let's take a look at some of the personas involved
in a machine learning project.
In an organization, we might have machine
learning engineers, machine learning researchers, data scientists, data
engineers, software engineers, DevOps business analysts, and also product
managers.
It is important to note that MLOps is also about building scalable teams.
As such, MLOps culture should be shared within the entire team.
And also not only with the technical team,
but any stakeholders involved in the ML system development and design,
including CXO with sponsors.
In the next lesson, we we'll be taking a closer look
at some of the personas involved in the MLOps lifecycle.
You might be asking yourself, why is MLOps so
important if I build a very good model?
Let's take a look into this.
It might be surprising to the academic and industry community
to know that only a tiny fraction of the coding many ML systems
is actually devoted to what is described so far, machine learning, training,
and learning, and prediction.
Much of the remainder efforts in ML assistance
might be described as plumbing.
It involves all the processes and capabilities for us
to ensure scalability, availability, portability, reproducibility,
modularity, monitoring and alerting, security and explainability,
and also hosted or serverless systems.
It is very important to remember that in ML systems,
continuous integration is no longer only about testing and validation
code and components, but also testing and validating
data, data schemas, and models.
Continuous deployment is no longer a single software package or a service,
but an ML system that should automatically deploy other services,
for instance model prediction service.
Continuous training is a new property specific to ML systems concerning
to automatically retrain and serve machine learning models.

\section{MLOps Personas}

\GRAPHICSC{1.0}{1.0}{TinyML/3 Development/5.22.png}

It should be becoming evident to you that ML development and deployment at scale is a complex problem. MLOps is a multi-faceted problem. MLOps is best performed by a coalition of experts who have specific responsibilities. At the same time, it is important to have a high-level vision of the whole pipeline to enable cross-discipline improvements. Below is a description of the generic ML Personas. In practice, these roles are often blurred and specific to the company and application. Nonetheless, these job descriptions are pervasive in any real-world application.

As we step through this course, we will highlight the personas that play a vital role in certain and specific roles. We hope that through this you will have a deeper understanding and appreciation for the big picture or 30,000 feet overview of MLOps.

Title

Description

ML Engineer

Designs and implements ML Systems. Responsible for writing the code to train and test the model before deployment. Usually a jack of multiple trades.

ML Researcher

Conducts research experiments to improve the underlying model and/or the training scheme. Usually transfers new developments to the ML Engineer before a new technology is implemented in production.

Data Scientist

Primarily focused on statistical methods and the analysis of the underlying data. A data scientist might be responsible for analyzing the models’ impact on higher-level performance metrics like user satisfaction or user behaviors.

Data Engineer

Designs, implements and maintains the data pipelines, including the creation and processing of the original training dataset and the implementation of the continuous learning system. Often data engineers pass the curated datasets to Data Scientists.

Software Engineer

Writes the software (non-ML) component of the application or development pipeline. Usually integral in the deployment of the model because they are responsible for the surrounding application code.

DevOps

Develops and manages the systems that enable the continuous testing and improvement of the model. Ensures mechanisms are in place to support the work and insights generated by the other personas.

Business Analyst

Analyses market demand and potential revenue sources to provide insight on what ML applications are useful and profitable for the organization

\section{MLOps: Key activities and Lifecycle}

LARA SUZUKI: Hello.
In the last lesson, we looked at ML challenges
on the productionization of machine learning systems.
They included governance of data, features, model, pipelines,
and experiments, the continuous training and deployment, changes
in data sets and features, data validation process,
and also modal analysis.
We also looked at production system challenges,
such as scalability, availability, portability, reproducibility,
modularity, monitoring and alerting, security, hosted or serverless systems.
Now, let's have an overview of the key activities in the MLOps lifecycle.
In the next sessions, we will be looking at each one of these activities
in a little more detail.
Here, we present an overall view of an MLOps process.
We have seven different processes ranging
from writing ML code to continuously monitoring a team production.
We have machine learning development, machine learning training or training
operationalization, continuous training, model deployment, prediction service,
continuous monitoring, data and model management.
Before we dig into each one of these processes,
we need to understand that, for TinyML, this process is incomplete.
We need to add a process between training and deploying our machine
learning models in production.
This step is known as "model conversion process."
It enable our machine learning models to be used for inference on Tiny devices.
We will take a closer look at it in the next lessons.
Now, let's take a look at each one of these steps
and the key outputs they produce.
The first process is known as "machine learning development"
or "ML development."
It concerns with experimenting and developing
a robust and reproducible model training procedure, which
is the training pipeline code.
It consists of multiple tasks from data preparation and transformation
to model training and evaluation.
The typical assets produced are code and configuration files.
Amongst the assets produced; we have notebooks for experimentation
and visualization; metadata and artifacts of the experiments;
data schema; query scripts for training data;
source code; and configurations for data validation and transformation
for creating, training, and evaluating models
for the training pipeline workflow and for unit tests and integration tests.
Once you are done with your experimentation, you build source code
and run various tests.
The outputs of the steps are pipeline components,
such as packages, executables, and artifacts
to be deployed at a later stage.
As such, machine learning training or training or personalization
concerns with automating the process of packaging,
testing, and deploying a repeatable and reliable training pipeline.
Amongst the typical assets produced, we have training pipeline executable
containers; training pipeline and runtime representation, which
references the components stored in artifacts or repository;
continuous training concerned with repeatedly
executing the training pipeline in response to the environment.
It could be new data or code changes or a retraining,
which is triggered on a schedule but then treated
with the new training settings.
This continuous training of new models, which
includes redeployment of those models and all the technical efforts
that goes along with it, aims to address three notable aspects
of machine learning projects.
The first one is the need for explainability.
We need to understand how and why a model makes certain predictions.
This is especially important for auditing purposes
to meet regulations and/or certain levels of predictive performance.
Second, model decay, which is the reduction of the production
model predictive performance over time.
Such model decay is due to new and changing real-world data encountered
by the model.
The continuous development and enhancements to the model
are often driven by business requirements.
Typical assets produced include a trained and validated model
stored in the model registry, training metadata
and artifacts stored in the ML metadata and artifacts repository.
That includes pipeline execution parameters, data statistics,
data validation results, transformed data files, evaluation metrics,
model evaluation results, training checkpoints, and logs.
In the case of machine learning models that
will be executed by resource constrained devices, ML for Tiny devices, hence
TinyML, models should be converted to a format that
is suitable for those devices.
In this a step, we can convert to TensorFlow model,
for instance, into a compressed flat buffer with the TensorFlow Lite
converter.
Amongst the typical assets produced, we have training pipeline executable
containers, training pipeline and runtime representation,
which references the components stored in an artifacts repository.
Modern deployment concerns with packaging, testing,
and deploying a model to a serving environment for online experimentation
and production serving.
Typical assets produced are model serving executable applications.
For example, a container image is stored in a container registry
where a Java package is stored in an artifact repository.
Online experimentation evaluation metrics
are also stored in the machine learning metadata and artifact repository.
The prediction serving is all about serving
the model that is deployed in production for inference.
Typical assets produced are request response payloads
stored in a serving logs store and feature
attributions of the predictions.
Continuous monitoring concerns with monitoring
the effectiveness and efficiency of the deployed model in production.
Typical assets produced are anomalies detected in serving data during a data
drift detection process, evaluation metrics
produced from continuous evaluation.
Metrics will depend on the data, features, model selected,
and runtime environment.
Data and model management is central cross-cutting function
for governing ML artifacts to support auditability, traceability,
and compliance.
Data model management can also promote shareability, reusability,
and discoverability of ML assets across an organization.
Typical assets produced are a unified repository for data sets and features;
data sets; features; metadata and annotations; information
about process executions; information for reproducibility, tracing, lineage,
and discoverability of data and models; those for analyzing, comparing,
and visualizing the metadata and artifacts; a process
to register, review, validate, and approve models for deployment;
reporting on the performance of the deployed model
and also explainable reports.

\section{How can you apply MLOps thinking to your future projects?}

Now that you have explored MLOps at a high level, we thought it might be useful to reflect on what you have learned so far. We’ve proved a couple of possible prompts to consider below:

\begin{itemize}
    \item How can you apply MLOps thinking to your future projects?
    \item What changes do you think need to be made to industry today to make these processes move more smoothly?
    \item Does your new knowledge of MLOps make you more or less likely to use TinyML?
    \item Does it make TinyML seem harder or easier? More or less exciting?
\end{itemize}

% 1.3: ML Development

\section{Overview of ML Development}

\subsection{ML Development}

\GRAPHICSC{1.0}{1.0}{TinyML/4 MLOps/pasted_image_0.png}  

\subsection{Objectives}

    The first stage of the MLOps pipeline is ML development. This stage should be somewhat familiar to learners who have taken TinyML Course 1--3 (since we have been doing this repeatedly since Course 1)! But don’t worry if you are taking Course 4 alone, you aren’t missing out. ML development involves the development of a model that is able to meet the performance criteria of a well-defined ML use case. This requires answering several questions such as:
    
    \begin{enumerate}
        \item What is the task?
          \begin{itemize}
              \item E.g., keyword spotting, wake word detection, object recognition
          \end{itemize}
      \item  How can we measure business impact?
    
          \begin{itemize}
           \item Will our model increase sales, reduce inefficiencies in business processes, increase security?
          \end{itemize}
    \item  What is the evaluation metric? And are these metrics measurable before deployment?
          \begin{itemize}
             \item Sales might require measuring performance before and after deployment, while image classification can be readily evaluated through a model’s test accuracy
          \end{itemize}
    \item  What data is needed? And does this data already exist?
          \begin{itemize}
    
             \item  Some features may be fundamental for developing an effective model, and in some cases, we may have no or minimal data and may need to procure data or simulate it in some cases.
          \end{itemize}
    (5)  What are the training and serving requirements?

          \begin{itemize}
    
             \item Some models might need to be updated daily, weekly, or monthly
             \item Some models may require high inference rates (e.g., 10 inferences per second)
             \item For TinyML, a key serving requirement we have is that the model must be small enough in memory to fit on our resource-constrained device.
          \end{itemize}

    \end{enumerate}

    In this section, we will look at these questions all structured in a systematic way, organized as follows: Problem Definition → Data Selection → Data Exploration → Feature Engineering → Model Prototyping → Model Validation → Model Evaluation → (repeat)
    
\subsection{Experimentation Is Key}

    A key aspect of ML development is experimentation. Experimentation is split up into the following tasks: data discovery, data preparation, and model prototyping. From the model perspective, this process involves trying different machine learning models (not just neural networks!), different model architectures, hyperparameters. From the data perspective, this involves deciding which features are most relevant, performing feature engineering, and determining whether the marginal gain in performance of including a specific feature outweighs the costs to onboard memory and compute operation.
    
    Keeping track of the impact of all of these items can be overwhelming, so often experiment tracking is used to simplify the process. This involves keeping track of all of the features, hyperparameters, and models being used for a given task and comparing the relevant performance metrics. Experiment tracking is very useful during the prototyping phase and determining which model should be selected for the final project that will be sent for model validation and potential deployment.
    
\subsection{Where Do MLOps Fit In?}

    MLOps can help with the model development process in various ways. For experiment tracking, we can save our models and all of their relevant information in a model registry so that we can easily reference and compare them. If a model needs to be updated or retrained on a regular basis, this can be set up using a continuous integration/continuous delivery pipeline. This procedure is commonly done in the software development world. For example, TensorFlow runs nightly builds every evening (tf-nightly) that include all-new feature additions and ensure that the key functionality of the package remains intact using a cron-like scheduler. This feature requires the model pipeline to be integrated with a version control system such as GitHub but allows multiple individuals to work on and update the model infrastructure simultaneously, making the development process more efficient and collaborative.
    
    We will delve into the ML development process in more detail for the rest of this chapter and we will use Keyword Spotting as the TinyML application use case to present the MLOps details.
    
\section{ML Development}

    0:00 / 0:00Klicken Sie auf die Pfeiltaste nach oben, um das Geschwindigkeitsmenu aufzurufen. Regeln Sie dann mit den Hoch und Runter Pfeiltasten die Geschwindigkeit. Klicken Sie Enter, um die jeweilige Geschwindigkeit festzulegen.Geschwindigkeit1.0xKlicke auf diesen Knopf und diess Video stumm oder nicht stumm zu schalten, oder drücke die HOCH oder RUNTER Knöpfe um die Lautstärke zu erhöhen oder abzusenken.Maximal Lautstärke.Videotranskript
    Beginn des Transkriptes, zum Ende springen.
    LARA SUZUKI: Hello.
    In your last lesson, we had a high-level overview
    of the each processes of our MLOps framework.
    Now let's take a closer look at each one of them in detail.
    In this session, we will cover ML development,
    which concerns with deploying a robust and reproducible model procedure.
    In the MLOps ops lifecycle, the first stage
    concerns with experimenting with data and models
    and developing a robust and reproducible model training procedure, also known
    as training pipeline code.
    This training pipeline code consists of source code and configurations for data
    validation and transformation, for creating, training, and evaluating
    models, for the training pipeline workflow,
    for unit tests and integration tests.
    An ML environment should allow data scientists
    to write modular, reusable, and testable source
    code that is virtually controlled.
    To achieve this efficiently, the steps of the ML experiment
    can be orchestrated using ML pipelines.
    A machine-learning pipeline orchestration
    defines the order that each one of the aforementioned tasks should be executed
    and when, for instance, the arrival of new data
    should automatically trigger the early steps of the machine learning pipeline.
    Such as data validation and transformation,
    to model training and evaluation, to model redeployment.
    Machine learning pipeline orchestration leads to rapid iteration and better
    readiness to move to production.
    The core MLOps capabilities employed in this stage
    are data sets and feature repository, data processing, experimentation,
    model training, model registry machine learning metadata,
    and artifact repository.
    This diagram illustrates the process of ML development.
    As we can note, the primary source of development data
    is the data set and feature repository.
    The experimentation and prototyping of ML model
    consists of data discovery, selection, exploration, preparation
    and engineering, model prototyping and validation, and also evaluation.
    We can use what-if, fairness, bias and explainability.
    The configurations for each experiment is persisted in the ML metadata
    repository.
    There's a pointer to the version of the training code, model architecture,
    pre-trained models, and also hyperparameters.
    There is also information about training, validation, and testing
    the splits.
    And also about model evaluation metrics and the validation procedure.
    The formalization of the ML training procedure for retraining
    represents the end-to-end pipeline.
    The key success here are experiment tracking, reproducibility
    and collaboration.
    Data scientists should be able to find previous experiments
    and compare experiments and run to understand models behavior.
    Ende des Transkriptes, zum Anfang springen
    Downloads und Transskripte
    Video
    Videodatei herunterladen
    Transskripte
    Download SubRip (.srt) file
    Download Text (.txt) file
    
    
\section{Problem Definition}

    0:00 / 0:00Klicken Sie auf die Pfeiltaste nach oben, um das Geschwindigkeitsmenu aufzurufen. Regeln Sie dann mit den Hoch und Runter Pfeiltasten die Geschwindigkeit. Klicken Sie Enter, um die jeweilige Geschwindigkeit festzulegen.Geschwindigkeit1.0xKlicke auf diesen Knopf und diess Video stumm oder nicht stumm zu schalten, oder drücke die HOCH oder RUNTER Knöpfe um die Lautstärke zu erhöhen oder abzusenken.Maximal Lautstärke.Videotranskript
    Beginn des Transkriptes, zum Ende springen.
    VIJAY JANAPA REDDI: So Laura just introduced us
    to the big picture of machine learning development.
    In this section, I'm going to break this complicated stage down further
    into smaller bits and pieces, and we'll step through each
    and every single one of those stages.
    The first step in any machine learning project
    is to understand the project objectives and requirements from a domain
    perspective.
    Then we have to convert this knowledge into a machine learning
    problem definition with a plan designed to achieve the necessary objectives.
    Now, and our projects are often structured
    around specific needs of an industry use case or something
    that's tailored for a particular problem that we're trying to get to.
    Now a successful ML project will start from a clean definition of what's
    actually needed to solve a problem.
    In this video, I'm going to walk us through this process using
    keyword spotting as the example of TinyML application.
    And for those of you who are not familiar with keyword spotting,
    that's like saying OK Google or Alexa.
    Those are keywords that trigger the consumer device.
    Now, in any organization, we have ML engineers,
    researchers, data scientists, data engineers, software engineers, DevOps,
    business analysts, project managers, and so forth.
    So as we step through the different stages, what I'm going to do
    is I'm going to include a persona that is most likely going to be
    responsible for this particular task.
    In this particular context, it might be the business analyst
    whose job is really to clearly define what the problem definition is.
    Now a clear and precise problem definition is important for any machine
    learning project's success.
    Domain experts can often be extremely critical down here.
    For example, let's say we want to design a recommendation
    system for an e-commerce website.
    We must first study how people will visit the online store
    and how do they actually interact with the store
    and whether that actually leads to an actual sale at the end of the day.
    Now if you're using something like Netflix
    or some other online software like Amazon as an example in your head,
    you'll need to understand how the user's behavior patterns are actually going
    to translate to purchasing patterns.
    Now, without domain expertise, the problem statement
    might actually be stated rather vaguely, like let's build
    a better recommendation system that will enhance net revenue.
    Well, however, a domain expert won't think that way.
    A domain expert will likely think about what's
    going to lead to user interest being increased that are specifically
    being caused by a certain set of recommendations that are being made.
    Now, apart from recommendations, there might
    be other sort of explanations for why certain sales increases, like maybe
    it's the Christmas holidays or maybe it's new years so there's
    a lot of shopping is going on.
    So those kinds of nuances are the things that a domain expert would typically
    know that you and I might otherwise miss.
    So the bottom line is that we have to have
    a strong and clear understanding of the challenge
    that we're trying to solve here.
    Now at this stage, you should really be asking yourself,
    if you understand what's the clear definition,
    then is machine learning really the right approach to solve this problem?
    Often, people will get things conflated.
    They hear about the excitement around machine learning and AI
    and they want to use it everywhere.
    So ask yourself, is machine learning really the right approach
    for the problem at hand?
    I can tell you for a fact that even at large enterprises like Google,
    Microsoft, Alibaba, and so forth, these engineers who have years of experience
    are repeatedly told to do that, which is ask yourself and justify that it
    is indeed a machine learning problem.
    You want to think hard about this, because if this
    is the wrong approach-- that is, machine learning is not the right approach--
    then you will waste a lot of time and effort, as you will see,
    because building an ML apps flow is non-trivial, though it is absolutely
    necessary for ML projects.
    Assuming that ML is indeed the right solution for the problem at hand,
    the next question you have to ask yourself
    is, what data do you really need?
    Data is king in machine learning projects.
    I can't emphasize this enough.
    Think hard and carefully about where you are going to get your data from.
    What are the sources you're going to gather it from?
    How are you going to scrub it?
    You know, what's the process you're going
    to be going through in terms of labeling the data for supervised machine
    learning models?
    As all of that is extremely costly.
    Later we'll talk a lot about this in the context of a keyword spotting example
    and I will go into more details.
    Next, what is the desired outcome from your machine learning project?
    Think about the metrics that are actually involved.
    Often, people focus on the wrong metrics as they're
    starting to build their machine learning projects.
    For example, it's often common to think about what's
    your machine learning accuracy.
    Well accuracy alone is not sufficient.
    Let's say, for example, if you're dealing with TinyML,
    you have a lot more to worry about, such as latency, energy efficiency,
    and so forth.
    And all those are part of the desired outcomes.
    Also, you want to think about what kind of machine learning algorithms
    are really the right approach to solve the problem.
    There are many different kinds of algorithms involved.
    Obviously, the one that is causing the most amount of ruckus in the community
    are deep neural networks.
    Now we'll be focusing on supervised deep learning approaches in this course,
    but you also want to keep in mind that there
    are many unsupervised learning-based methods, like decision
    trees, random forest, k-means clustering, and so forth, that are also
    very powerful and especially in the context of TinyML
    because they tend to be much more computationally efficient.
    So that is something you really want to keep in mind.
    So it's not always about having deep learning-based methods that
    are computationally expensive just because of the jazz around them.
    It's about simplicity, because simplicity is really the key nugget
    in Tiny machine learning.
    Next, you want to keep in mind the way a network behaves
    is really a function of its past experiences, right?
    Because that training data that you have,
    which is going to influence how the network works,
    has a lot in terms of how it'll actually perform downstream.
    What I'm trying to say is that the training data might not necessarily
    reflect what the real-world data is.
    So keep this in mind as we'll come back to it over and over,
    talking about skews and talking about drifts in the data distributions
    as time goes by.
    All right, so I've kind of loosely talked
    about problem definition in a very abstract sense,
    so let's try and ground this using keyword spotting.
    Now I always go back to keyword spotting because it's a bread and butter
    application in TinyML, because in the future,
    I honestly believe we'll have voice controlled interfaces for literally
    everything.
    Your toaster might be able to understand a few key words,
    like on, off, toast for one minute, for two minutes, three minutes,
    and so forth.
    Or your laundry machine might be able to understand gentle cycle, dry cycle, on
    or off, and so forth.
    So this sort of limited keyword spotting capability
    has a potential to really change how we interact with our devices around us.
    So that's why I'm pretty excited about it.
    Now let's go back to the problem definition around keyword spotting.
    What are the first things that come to mind?
    For me, I've just listed out the main ones that jump out, which are like,
    what are the keywords?
    What environment is this particular solution going to be operating in?
    What are the data sets that I'm going to be using?
    And what are the right metrics?
    Let's go into each one of these into a little bit more detail.
    For keywords, you really want to understand
    what task are you trying to perform.
    There are bound to be some common keywords like no, yes, or ignore,
    and so forth.
    Man, if you're trying to control like a small robot, in that case,
    you need the right keywords like left, right, go, stop, halt, something
    to that effect.
    Or perhaps you're building a device that needs to actually understand
    temperature and time.
    And, in that case, you need the machine to understand numbers--
    0, 1, 2, 3, 4, 5, 6, 7, 8, 9, 10, 11, 12, and so forth.
    These things seem so simple.
    I'm sure you're probably thinking, why is Vijay
    going on talking about these things?
    But I can't emphasize enough the getting this right during the initial problem
    definition stage is extremely crucial, because it
    has a massive downstream impact.
    For example, in the context of keyword spotting,
    what are the demographics where you actually going
    to be running your keyword spotting application?
    There are always 7,000 different languages spoken across the world,
    so you need to make sure you understand the user base
    that you're targeting for your keyword spotting application.
    Keep in mind even for a single language like English
    there are many different accents.
    So you really want to understand the user base,
    not just think about, Oh, I got my keywords.
    Now I know what to do.
    Because understanding these nuances and aspects
    of where things are going to be deployed is
    going to have an impact as to how the system will perform.
    Then you want to think about what the operating
    conditions are for this device.
    For instance, if you are building a keyword
    spotting an algorithm into a washing machine,
    then you need to understand the background noise
    that it will experience as that's going to be vastly
    different from the keyword spotting model that's going to be sitting
    inside a toaster in the kitchen.
    Now you can't assume that just the raw keyword data--
    the yes, no, and so forth-- are enough.
    The environment or the operating conditions noise is extremely critical.
    It has a huge impact in terms of how these models actually
    perform in the real world.
    Now this is true for even things that extend beyond keyword spotting.
    Now we'll go more into detail soon.
    Ultimately though, when you're running this keyword spotting model,
    you've got to think about what are the right metrics or KPIs involved.
    Someone has to sit down and write these things
    down so that everything that is done from the problem definition phase
    down to model deployment can be thought about correctly.
    For instance, if you're building a TinyML use case,
    and you don't talk about energy consumption,
    oh, that's a big oversight.
    I can tell you that this actually happened in real world use
    case with the customer where a startup company failed
    to capture the requirements clearly where the customer was expecting
    a 10-year lifespan on the TinyML device, but in practice, they
    were only able to deliver a device that had a two-year lifespan.
    And this was because the ML architects were actually not told the specifics
    about the deployment scenario.
    They were simply told, OK, this is the particular task,
    and I need you to build a model for this.
    And they simply chose the wrong device and deploy
    the model that fit on the device and pushed it into production.
    And they said, well, this is the best we have right now.
    Well that was not what the customer wanted.
    So this happens in practice is my point.
    So as a quick recap, look through the problem definition stage carefully.
    There are many more questions than what I can actually fit in a slide,
    but these are some of the important ones.
    And do not overlook this problem definition stage,
    because it can come back to bite you down the path.
    That said, I'll see you at the next video.
    Ende des Transkriptes, zum Anfang springen
    Downloads und Transskripte
    Video
    Videodatei herunterladen
    Transskripte
    Download SubRip (.srt) file
    Download Text (.txt) file
    
\section{How Might You Use KWS}

Now that you’ve learned a little bit about KWS (or been reminded about it if you took the previous courses), how might you scope out the “problem definition” for KWS?

Here is a quick recap of the problem definition process:

\begin{itemize}
  \item What challenge do you wish to solve?
  \item How might data science and machine learning approaches be used to solve this problem?
  \item What data is required for a successful machine learning prototype?
  \item What could the desired outcomes look like?
  \item What algorithms do you wish to put to the test?
  \item What do you make of the model's outputs?   
\end{itemize}
         
\section{Data Selection for KWS}

    0:00 / 0:00Klicken Sie auf die Pfeiltaste nach oben, um das Geschwindigkeitsmenu aufzurufen. Regeln Sie dann mit den Hoch und Runter Pfeiltasten die Geschwindigkeit. Klicken Sie Enter, um die jeweilige Geschwindigkeit festzulegen.Geschwindigkeit1.0xKlicke auf diesen Knopf und diess Video stumm oder nicht stumm zu schalten, oder drücke die HOCH oder RUNTER Knöpfe um die Lautstärke zu erhöhen oder abzusenken.Maximal Lautstärke.Videotranskript
    Beginn des Transkriptes, zum Ende springen.
    VIJAY JANAPA REDDI: Once you have the problem definition stated clearly,
    it is time to bring together subject matter experts and the data scientists
    to begin the journey of selecting the right data.
    It all starts with the search for suitable data.
    And finding this data might seem simple, but in practice, it's
    actually an arduous task.
    So let's dive into it.
    Data selection is about constructing a data set from one or more data sources
    to be used for exploring and the design of model.
    It is a solid practice to start with an initial data set
    to get familiar with the data, to discover first insights into the data,
    and have a good understanding of what might possibly be some data quality
    issues that might creep up.
    Now data preparation is often time consuming
    and is heavily prone to errors.
    The old saying of garbage in, garbage out
    is particularly applicable to data science projects or machine
    learning projects at large.
    Because the data has been gathered from many invalid and potentially
    out-of-range related sources and might have things like missing values
    and so forth.
    So analyzing the data is extremely key.
    If you do not analyze the data and screen it properly,
    then you're going to have a model that can likely
    be mis-predicting downstream, and have severe penalties.
    So the success of many machine learning projects really
    depends on the quality of the data preparation process.
    Typically, this is the kind of thing that a data engineer or a data
    scientist is concerned with at large.
    Recall that machine learning is all about data, right?
    It reminds me of my mom always talking about location, location, location
    when I was buying my first home.
    So when it comes to machine learning, I like to say data, data, data.
    You take a big and large original data set that you got your hands on,
    and then you chop it up into small pieces for the training
    set then test set and then you usually just kind of focus
    on that little piece.
    For most of the time when you're building a machine learning project.
    But really, you should be spending just as much time asking the key question
    about what is the source of this original data set?
    Often, we just assume that this data is available to us,
    and we tend to diminish the value of actually asking the tough question of,
    is this data really set up correctly for me to use in a commercial flow?
    More likely than not, if you worked with data sets in the past,
    especially as you're learning machine learning,
    you probably use data sets on the internet,
    or you probably got them from something like TensorFlow, the data catalog.
    There are a large number of public data sets
    that are available there across like a range of machine
    learning tasks for audio, image, NLP, text and so forth.
    Now more often than not, these data sets are typically
    built by researchers, people like me, that are made publicly available so we
    can foster research.
    But are they really useful for commercial purposes?
    On the other hand, you have data sets that
    are being created through crowdsourcing based approaches for example.
    Mozilla common voice is a crowdsourced data set.
    It's built by individual contributions, by people like you and me.
    That's a free data set that exists for speech recognition related software
    and has over 50 different languages that can be
    used to train machine learning models.
    All very exciting stuff.
    But there are also sources like ML Commons, which is actually
    a non-profit organization that is aimed at accelerating machine learning
    for everyone.
    The focus is more exclusively on commercially useful data sets.
    As a side note, I'm one of the founders of ML Commons,
    and we strongly are committed to widening access to applied machine
    learning.
    And what we mean by that is that we want to actually have data
    sets that will help mom and pop stores kind of come up to life real quickly,
    rather than having to invest very heavily in costly data
    engineering practices to build new data sets.
    Originally, ML Commons started off benchmarking machine
    learning systems, something that we'll talk
    about way down later in the course.
    But more recently, we've started focusing
    on building large public data sets available to train ML models,
    so that everybody can actually get their hands on data, high-quality data
    easily.
    Now the focus really again is not much on research but much more
    on the commercialization side of the path,
    and we're focused on really getting the right licenses in place
    so that the commercialized models can be used with no concerns about legality.
    So talking about legality, one of the big questions you got to think about
    are what's the quality of the data?
    Really, building a large data set is often not that hard a problem.
    It's really hard to have a high-quality data set.
    This is a crucial difference, because someone really
    has to vet that data set.
    If someone is not checking every entry in the data set, which is undoubtedly
    a tedious and costly task, then you might end up
    having a poorly performing model, because your data is actually
    of not high quality.
    So therefore, understanding the source of the data
    is extremely crucial and maintaining that high quality is also
    extremely costly.
    Now in fact, Andrew Ng has lately been talking
    about pushing the community to think more about data quality,
    and something that I will come back to later on, as we go into model valuation
    later on.
    Now going back to a keyword-spotting example,
    which we discussed earlier during the problem definition stage.
    In the context of keyword spotting, one of the most widely used data sets
    is really Google speech commands.
    Or really, it should be called Pete Wharton's data set.
    Pete is one of the pioneers of TinyML movement.
    He created this data set specifically for keyword spotting.
    Now it was developed using crowdsourcing-based approaches,
    so it's got a fairly limited number of samples and so forth.
    Now despite this awesome effort, and despite the fact
    that this data set is something that's actually
    commercially visible for a lot of people,
    many startups do not actually use it in practice
    when they're actually building their MLOps stage.
    And the reason is really simple.
    It's because it's not really set up for commercial purposes.
    It's really kind of just set up to do initial prototyping and exploration.
    Now don't get me wrong.
    I'm not bashing on Pete's effort.
    Instead, I'm simply making the point that some data sets
    that are really publicly available are not ripe
    enough for commercial use cases, because they might lack the right background
    noise.
    They might lack the right environment variability.
    The data distribution might not be of the right set so
    that it meets your customer use cases.
    So keep these things in mind.
    The next question is how easily can your data scientists and engineers
    get their hands on the data?
    Do you have the right data pipeline built
    so that it's easy to pull the data from a repository?
    For example, if you take the speech commands data set, by default
    it's available in the TensorFlow ecosystem,
    unsurprisingly because it's one of Google's artifacts.
    But typically, what you have to do is, you
    have to create a data lake or a data stream
    and you have to facilitate some sort of easy access for your data scientists
    to be looking at the data.
    Really, it's the data engineer's job to make sure that this happens.
    More on this when we talk about the data engineering stage.
    It's also extremely important that what's inside the data.
    One of the emerging themes is the notion of data sheets for data sets.
    It's sort of like a nutrition label on a food package,
    so you know what's actually inside the data.
    This way, you know what the ingredients are,
    like what's the data distribution look like for this particular data,
    before you train models.
    In the context of keyword spotting, it might be like,
    how many male speakers are there?
    How many female speakers are there?
    You want to know this way before you start.
    Now this is a very simple checklist.
    It's not rocket science, but you'll be surprised
    if you look at the state of what our data sets look like today.
    These are still best practices that are just starting to emerge.
    And these best practices are extremely important to put in place for MLOps.
    One of the major issues or concerns with data sources
    is actually having properly labeled data.
    Just because you have a machine learning data set
    does not automatically mean that you have a high-quality data
    set that has the right set of labels.
    For example, here, when I'm showing you the spectrograms,
    are you sure that yes is really a yes and that the no is really a no,
    and that there's no mislabeling here?
    Well, the only way you can actually verify this is you
    have to have someone physically listen to the sound.
    This is one of the biggest problems that we actually
    have is making sure that the data is actually nice and clean.
    Also, when you think of data sets, you often
    tend to think of them as like some sort of static entity.
    Here's my data set.
    Well, in the real world, the data is constantly changing.
    Whereas your data does not really change.
    Your data set is static.
    So the question you have to be thinking about when you're
    thinking about your data set sources are,
    how are you going to keep this data set that you are going to work on updated?
    Because if you don't think about it, you might
    create a model that's kind of good for this particular snapshot
    or version of the data, but as the data changes in the real world,
    then this model is not going to perform well, because you are not
    updating your data set.
    So this is, again, one of those things that a data engineer should really
    be sweating quite a lot on.
    More often than not, your data engineer is
    thinking about this much more than your data scientists,
    because the data scientist is kind of really focused
    on just analyzing what's inside the data, exploring the data.
    And your data engineer's job is really to try and make
    sure he or she can keep it updated.
    Next question.
    One of the key things to keep in mind is the legality around this data.
    While you might be able to arbitrarily select
    data from what you get access to on the internet or so forth,
    or from these different data catalogs, you
    have to ask yourself this key question as to whether you can actually
    train a machine learning model on that data for commercial purposes.
    Because sometimes, you will have data sets
    that are completely fine for research purposes,
    but they are not fine for commercial entity purposes.
    And this can be a big snag.
    Because again, if you go through the effort of building a model,
    and then later you can't use it, it's a waste of time and effort.
    The whole reason you have to worry about the licensing and so forth
    is because your data might have personally identifiable information.
    Now this is a big problem in certain countries
    which actually take privacy-related aspects very seriously, like in Europe
    for instance.
    Now this PII, or Personally Identifiable Information,
    is basically any information in the data set
    that might be used to identify a person.
    You cannot have this.
    And if you get this wrong, you are done.
    Because your model is not going to be useful.
    So again, I'm re-emphasizing these points,
    because these are all parts of best practices
    that need to be implemented inside the MLOps.
    One of the more nuanced issues is whether they need restricted features.
    And I already kind of hinted towards this earlier,
    which is sometimes you might be able to use it for research,
    but you cannot use it for non-commercial purposes.
    So just because the data is available or you might be able to slice
    from a particular data, does not mean that data set is actually viable.
    All right, so just kind of keep that in mind.
    So keep all of these questions in mind as you're
    selecting the data set source.
    Getting this wrong can have consequences downstream in the MLOps flow,
    because, as you'll see, you'll be working quite a bit with the data set
    that you've actually chosen.
    So I'll see you in the next video.
    Ende des Transkriptes, zum Anfang springen
    Downloads und Transskripte
    Video
    Videodatei herunterladen
    Transskripte
    Download SubRip (.srt) file
    Download Text (.txt) file         
    
\section{Why Real Data Matters}


\subsection{Garbage In, Garbage Out}

The phrase “Garbage in, garbage out” is an old computer science adage that is particularly relevant to machine learning systems. The implication of this statement is that the data in our training set needs to match the data our system will observe during deployment as closely as possible, otherwise we risk a loss of performance and other undesirable effects. The mismatch between the data we use to train our models and that which is seen during deployment can manifest in a number of different ways.
    
    Firstly, a scarcity of data is typically our most prevalent issue. It takes a lot of time and effort to procure and manage a large and high-quality dataset. Fortunately for us, machine learning models can often function with as little as a few hundred or thousand data samples, but this much data can still be challenging to obtain. The more complex our data distribution is, the more samples we need for our machine learning model to learn an accurate approximation of this distribution. If we are trying to build an algorithm to understand voice commands, we need to take into account different types of accents, pitches, inflections and intonations. This can be particularly problematic if you are part of a small team, and you may need to get creative and do some crowdsourcing of data (as Pete Warden did with his Speech Commands dataset!).
    
    A separate but related issue is how representative our data is of the real-world use case. As we just highlighted, a model to respond to voice commands should work for a broad range of accents and voice types. If we built a model based purely on data of our own voice, it might work great when you give a voice command, but work terribly for everyone else (even if they try to imitate you!). The reason this example fails is our model is not representative of our data distribution, only a small subset of it. To improve our model’s performance, we would have to obtain more data samples that have a greater diversity of accents and voice types.
    
\subsection{Legal Challenges}

    So, we have established that curating a dataset is difficult. However, you may have had a thought while reading the last section - why don’t we just use datasets that already exist? This is an excellent question. In our example, can’t we just go ahead and use Pete Warden’s Speech Commands dataset? The answer is that it depends. Some datasets are shared under specific licenses, and if we plan to use them, we must make sure to abide by the stipulations of those licenses, especially if we plan to use the dataset in any commercial capacity. However, whether commercial models can be deployed that are trained on proprietary datasets is a bit of a gray area, highlighted nicely by Authors Guild v. Google. In this case, Google was sued for training a search algorithm on copyrighted books. However, Google won the case because it was deemed that the algorithm actually benefited the copyright holders since the search algorithm made it easier for the books to be found via a search engine. While this is an isolated case, this sets an interesting and note-worthy precedent.
    
\subsection{Technical Challenges}

    Beyond legal limitations, there are also technical limitations associated with using external datasets. Once again, they might still not be representative of our end user use case. As an example, audio data that was recorded in a recording studio might not represent the real environment, where there may be background noise such as traffic or people talking. The difference between these two environments would add bias to our data that might impact our model’s performance for the end user. Even using different compression techniques can impact the performance of our model. In the examples below, we see spectrograms of different encodings for two examples, (1) a baby crying, and (2) a dog barking, before and after different audio compression techniques. These techniques leave characteristic markers in our data that are not obvious to the listener but would be obvious to a machine learning algorithm. Thus, sometimes merely transferring between data types can impact the usability of our data!
 
 \GRAPHICSC{1.0}{1.0}{TinyML/4 MLOps/why_real_data_matters1.png}   
 
    Spectrograms of high-quality baby cry recording (left) and MP3 encoded version (right). Visually the top of the spectrogram is missing on the encoded version and there are other areas of the image that are also missing (the colors of the spectrogram are black). This is due to losses from the compression of the audio.

 \GRAPHICSC{1.0}{1.0}{TinyML/4 MLOps/why_real_data_maters.png}   
    
    Spectrograms of high-quality dog bark recording (left) and AAC encoded version (right). Again just like with the MP3 encoding above, a lot of imformation is last as there are now errors throughout the encoded spectrogram where sections of the audio are now the wrong color and the whole top of the spectrum again is completely lost and shown in black.
    
    Similarly, training an algorithm on images that are always horizontal might result in a performance loss if our camera is placed at an angle during deployment, and poor lighting may cause similar problems if this was not accounted for in our training data. Naturally, it can be difficult to consider all of the possible nuances of our deployment environment, which is why curating high-quality datasets is challenging. One possible way around this is to fine-tune our model on a small set of data from our deployment environment, after broader training using a larger dataset. This provides the benefit of generalizability, while also taking into consideration the particularities of our particular deployment environment.
    
    Hopefully, you now have a greater appreciation for how important, but also how difficult, it is to produce a high-quality dataset to train our machine learning models. We will delve deeper into this in the following sections, but the key takeaway here is to always keep the end-use case in mind when developing your training dataset.    
    
\subsection{Additional Resources}


    \href{https://keras.io/examples/keras_recipes/sample_size_estimate/}{Estimating required sample size for model training}
    
    \href{https://datacentricai.org/}{Data-centric AI Resource Hub}
    

\section{Data Exploration}


    0:00 / 0:00Klicken Sie auf die Pfeiltaste nach oben, um das Geschwindigkeitsmenu aufzurufen. Regeln Sie dann mit den Hoch und Runter Pfeiltasten die Geschwindigkeit. Klicken Sie Enter, um die jeweilige Geschwindigkeit festzulegen.Geschwindigkeit1.0xKlicke auf diesen Knopf und diess Video stumm oder nicht stumm zu schalten, oder drücke die HOCH oder RUNTER Knöpfe um die Lautstärke zu erhöhen oder abzusenken.Maximal Lautstärke.Videotranskript
    Beginn des Transkriptes, zum Ende springen.
    VIJAY JANAPA REDDI: Once you've identified your data set sources,
    the next step is to explore that data.
    Data exploration involves understanding what's really inside the data set,
    even if you have various data set sources from which you can select data.
    Your data needs to have the right characteristics that
    make it useful to train machine learning models.
    Now, yet again, this is the kind of a thing
    that a data scientist would take a strong liking
    towards in terms of examining the data.
    Sometimes, even ML researchers get into the mix.
    It really just depends on the organization--
    how it's kind of structured.
    But here, you can see that the focus is now shifted away from the data engineer
    to the machine learning researcher and data scientists.
    So keep an eye out for these kinds of shifts,
    because this is part of that breadth that I mentioned in the t-shaped skills
    that you need to know.
    Because you want to understand who's responsible for what.
    Now, the problem with ML models is that they're highly
    susceptible to data distribution.
    So simply put, the data quality.
    If you have bad data, then you end up with a poorly performing model.
    Let's take a real-world example where bad data
    has resulted in a serious problem.
    Now, Amazon once created an AR hiring tool to screen resumes.
    They used the past 10 years' worth of Amazon applicant resume data to train
    the machine learning model.
    However, the recommendation of the model was sort of biased heavily towards men,
    and it even, in fact, penalized any resumes
    that had words related to women, such as women's chess club captain,
    and so forth.
    The reason for such bias is because of the unbalanced number of males
    and female applicants that were there in the past 10 years,
    as the figure shows on the left.
    There are more males than females in the technical roles in this figure.
    Now, does this really mean that females are less qualified than males
    for these jobs?
    No, absolutely freaking not.
    The moral of the story here, from this real world example,
    is if you had actually explored the data first before training the model,
    you would have easily caught this issue.
    It would have been so easy to uncover this imbalance by simply
    looking at the distribution between the male and female applicants,
    for instance.
    You don't need a rocket science degree for this.
    But many companies do fail at simple things like this,
    and that's the key message down here.
    So what are some ways to look at data?
    You can look at individual data points and consider how noisy the labels are.
    You can do reservoir sampling, which is really
    a fancy way of saying that you just pick a random sample without replacement
    of K items from an unknown population size in a single pass.
    Or, you can do a human test.
    If a human can't figure out the distinguishing factor between any two
    samples that are given, then it would be difficult to train
    a machine learning model to accomplish pretty much the same objective.
    You can listen to the keywords, in the context of keyword spotting.
    For instance, are they distinct?
    Is there too much background noise?
    Too little background noise?
    Is the labeling actually correct or incorrect?
    Bounding boxes, if you're talking about things like image classification
    or image detection or object detection kind of things.
    Like, are the bounding boxes actually precise or are they imprecise?
    This is the kind of stuff that a data scientist is typically exploring
    once a data set is given to them.
    So if you were the data scientist, you might be thinking,
    how can I handle all of this data that's coming at me?
    Well, typically, you don't have to look at the large corpus of the entire data
    set.
    What you try and do is you try and be clever about it.
    For instance, you would select a small portion of that data--
    of, let's say, 20 to 30 samples--
    where you have, for example, a golden thing.
    Where it's like, golden means that you know
    that the data labels are 100% correct and that they are clean.
    Then you overfit your model to those examples of, let's say, 10 examples.
    And then you'll evaluate on those exact same 10 examples,
    and you should get 100% accuracy at that point.
    Makes sense, right?
    Because you trained it on exactly 10 examples and you're
    evaluating on exactly those 10 examples.
    Then, what you do is you take that model and then
    you evaluate it on those other 20 examples
    or so that you have in your bucket.
    And then you try and determine how big is that performance delta there.
    So through overfitting on small data examples,
    you might discover problems with your architecture, for instance.
    Maybe the autoencoder bottleneck you're using is too small.
    So things like that will become quite obvious.
    So you don't really have to look at the whole data set.
    You can actually be very clever by just looking at small data.
    Now, there are a rich array of different ways
    you can slice and dice and examine data.
    I'm just going to give you a little sampling,
    so you sort of know what your colleague, who's a data scientist,
    might actually be doing.
    For example, they'll be looking at outliers-- observations that
    are relatively large or small value, compared to the majority of the samples
    that you have.
    Which can have a huge impact in terms of your model's ability to quickly learn.
    This can be caused, for instance, by a glitch in a sensor data stream.
    Maybe a sensor, just for a little bit, turned off, but then it came back on.
    These things do happen in production scenarios.
    There are many different techniques to deal with them,
    but that's really not important.
    For example, you have a box plot or a Cleveland dot plot
    or something that you'd use to visualize and quickly detect
    this sort of outlier characteristics.
    But this is what a data scientist would do.
    If you're interested in these kinds of things,
    I strongly encourage you to take a look at the Harvard Data Science
    class, which goes really deep into some of these analysis things.
    I'm not trying to do that here.
    I'm trying to just give you a flavor for the things that a data scientist might
    be doing, so you appreciate his or her role in the company.
    They might also be looking at what's the homogeneity of the variance,
    because it's often a critical assumption that we make.
    The homogeneity of variance is the assumption
    that the variance between different groups is relatively even.
    That is to say, all the groups have a similar variation between them.
    Now this is quite useful, because it's actually
    an integral part of just making sure you can actually train effectively to this.
    And there are different ways to test this homogeneity.
    Like, you can actually just look visually.
    You can plot the residuals or do a statistical test,
    and you can actually tell.
    This is the kind of thing, again, your data scientist colleague
    might actually be doing.
    One of the key underlying backbone assumptions in machine learning
    is that your data is normally distributed.
    Histograms are an awesome tool for visualizing what these data
    distributions look like, and it's something
    that we should all be doing before we start training the model.
    Because you want to understand the characteristics that you
    want inside the data before you start training the model,
    and then figuring out, hey, why is my model not performing well
    for this particular group of people?
    Or something to that effect.
    Well, it wouldn't perform well if your data doesn't actually
    have that sort of a distribution profile inside.
    So these are mechanisms that you want your data scientist colleague
    to actually be thinking about.
    There are also simple things like just cleaning up your data.
    Having zero or null values inside your data can really throw havoc,
    and so very often, data scientists need to think about cleaning up the data.
    Just to wrap things up a little bit more,
    correlation analysis is also one of those key checks
    that you need to do across various stages of a project,
    like, before you get into feature engineering
    or feature selection and so forth, which we'll be doing next.
    Now, if you ignore this particular stage,
    one is likely to end up with a highly confusing statistical analysis
    in which nothing seems significant to the machine learning model.
    However, if you drop one covariate out of it,
    then other variables might actually start to show more prominence.
    So there are many different ways to detect
    colinearity, that include things like VIF, which
    is like Variance Inflation Theory.
    Pairwise scatter plots, comparing different covariates,
    doing correlation coefficients, and so forth.
    A whole bunch of techniques exist in data science principles.
    Anyway, a data scientist is actually looking at these kinds of things,
    and they're really responsive for this.
    What I'm trying to get to, fundamentally,
    is that there are a lot of this sort of data cleaning stuff
    that you have to do.
    After you have sourced the data-- you've got all these labels
    and so forth-- then you've got to look into the data
    and make sure it's actually got the attributes that you want it to have,
    and that there are no wacky distributions inside or wacky outliers
    and so forth inside.
    Because that will really skew your model.
    You don't want to do feature engineering and things-- which
    we'll talk about next-- until you get this portion right.
    So that's the key thing.
    And this list of things that I'm talking about can go on and on and on,
    and I've just listed a few.
    But there's one very important point, which
    is you need to have the right tools in order to help you understand the data.
    More often than not, what I've seen is that people
    will build a lot of infrastructure and tooling
    for training machine learning models, but when it comes to actually examining
    the data that they're using to train the models, there aren't really any tools.
    My experience has been that you often end up
    having a data scientist that specializes in advanced scientific tools like R
    or scikit-learn and these kinds of things.
    That's great, but I think increasingly, we're
    going to need to see tools in the MLOps flow
    where it's also more easily accessible for everybody.
    One of the reasons that I'm talking about this
    is I want to give you a flavor for the kind of things
    that need to be put in place so that you can efficiently and effectively
    make sense of the data when you want to train the models.
    I have a reading coming up that talks a little bit more in order
    to give you pointers around useful resources
    to explore data through visualization mechanisms,
    especially in the context of MLOps platforms.
    So make sure you actually check that out,
    because it's got a lot more specifics involved.
    Taking a step back, what I want you to take away from this particular video
    is that, before you start engineering the models
    and feeding it into the training pot, make
    sure you have actually identified the right data with the right setup,
    and that you have a data scientist who is actually
    examining that data to try and understand what
    are the characteristics of the data.
    Because once you do this, you're going to be iterating on this data
    over and over and over, so you want to get this right, is the message.
    Ende des Transkriptes, zum Anfang springen
    Downloads und Transskripte
    Video
    Videodatei herunterladen
    Transskripte
    Download SubRip (.srt) file
    Download Text (.txt) file    
    
\section{Data Visualization Tools}

\subsection{Overview}

    One of the greatest strengths of machine learning models is that they are able to pick up associations from data that humans are unable to parse. However, the lack of interpretability inherent in some of these models has led them to be referred to as “black boxes”. Neural networks are a prime example of this; the non-linear mapping inherent in neural networks makes it difficult to explain how the network came to a conclusion based on its data inputs.
    
\GRAPHICSC{1.0}{1.0}{TinyML/4 MLOps/data_vis.png}    
    
    As machine learning models become increasingly complex, there is a trade-off between accuracy and interpretability which are plotted on the x and y axis. The linear regression model is the most interpretable but the least accurate. The neural network is the most accurate but the least interpretable. In between those fall other machine learning approaches like decision trees, SVMs, and random forests. To learn more see: https://towardsdatascience.com/guide-to-interpretable-machine-learning-d40e8a64b6cf
    <span id=" />
    
\subsection{Interpretability Questions}

    One of the ways we can ameliorate this lack of interpretability is through understanding our data, specifically by asking questions like:

\begin{itemize}
  \item Are there any outliers present? (How outliers can be handled is a complex topic, but the interested reader can delve into extreme value theory)
  \item Is the data homoscedastic or heteroscedastic (i.e., does variance in a variable grow with the size of the variable, or is it constant)
  \item Are the data normally distributed? Or do they follow other distributions (e.g., Poisson, Exponential)?
  \item Are missing values present? If so, will data be filled in via imputation or omitted?
  \item Is there any collinearity in covariates? (e.g., two separate variables with one representing temperature in Fahrenheit and the other temperature in Celsius)
  \item Are there any interactions between variables in our data? (e.g., relative humidity influences temperature, and vice versa)
  \item Is each data point independent? (i.e., is there any autocorrelation between data points)
\end{itemize}

    These are very technical questions, but by better understanding our data we can usually build better models and gain additional insights that were previously hidden from view. 
    
    One way of understanding our data is through the use of visualization tools and techniques. The strength of visualizations as a tool is that they replace cognition with perception. That is, instead of having to read through and interpret tables to highlight associations between data, visualizations help to uncover those associations with great ease, and often uncover new perspectives.
    
    The cubist paintings of Pablo Picasso provide a helpful analogy. When most people look at a painting, they see a single perspective. However, Picasso specifically designed his paintings so that they had multiple interpretations, and by continuing to look at the painting you would start to see alternative perspectives within the same image. This is similar for our dataset; a single bar chart or scatter plot of our data set will only give us a low-dimensional representation of our dataset that has been filtered by our perceptive field, and only by seeing our data from multiple perspectives can we really start to understand and build higher-dimensional insights. This is where interactive visualization tools come in handy.
    
\subsection{Visualization Techniques: Dimensional Reduction}

    There are multiple techniques available to visualize high-dimensional data. These typically work by compressing the dimensionality of our data into a small representation that can be visualized on a two- or three-dimensional plot, as in Edge Impulse’s Feature Explorer. Hence, these techniques are called dimensionality reduction techniques. These techniques can be linear (e.g., principal component analysis, non-negative matrix factorization) or non-linear (e.g., tSNE, UMAP).
    
    As an example, we will look at Edge Impulse’s Feature Explorer to see how misclassified examples can be spotted within a dataset. The Feature Explorer visualization shown below allows us to compress our data features into a three-dimensional image, and then highlight these images by their classification output. In this case, the classifier output is one of four different motions: “wave”, “updown”, “snake” and an “idle” value if there is no input. The visualization allows us to click on each of the samples and see the corresponding accelerometer waveforms of the sample. You could apply the same techniques to compressed image or audio features.
    
\GRAPHICSC{1.0}{1.0}{TinyML/4 MLOps/feature.png}    
    
    Feature exploration of 6,458 samples of data by Edge Impulse. The image shows a 3d point cloud visualization of the x, y and z axes of the root mean square of acceleration data. There four clusters of datapoints that correspond to different actions: idle (in blue and centered around 0,0,0), wave (in orange and distributed widely with large accelerations in the x and y axis), updown (in green with large accelerations in the z axis), and snake (in red with small accelerations in all three axis). There is also a view of the timeseries of a particular selected datapoint shown at the bottom of the screen where the data is plotted as line graphs showing the x, y, and z axis acceleration over time.
    
    To highlight how the Feature Explorer can be used, we will use an example of highlighting misclassified data samples. In the below image, we have filtered the classifier outputs to pick up only samples which correspond to the keyword “yes”. Upon clicking on the sample, we can listen to the audio sample to see what it sounds like.
    
\GRAPHICSC{1.0}{1.0}{TinyML/4 MLOps/feature_explorer.png}    
    
    The image shows a 3d point cloud visualization of some 1,526 data points of “yes”, “no”, “noise”, data showing only the “yes” data that is clustered into a circular pattern. There is an outlier data point which does not fall in the pattern with the rest of the data and is highlighted by a red arrow. This may correspond to a misclassified data sample.
    
    Upon listening to the audio sample, it becomes clear to us that this audio sample is mislabeled, as the sample clearly sounds like “no” more than “yes”. This can be done with other samples to prevent incorrect classifications from reducing the performance of our machine learning algorithm.
   
\GRAPHICSC{1.0}{1.0}{TinyML/4 MLOps/explorer.png}    
    
    A time series view of the outlier audio sample mentioned above. There is also an area on the bottom of the screen that enables the user to play the audio sample and adjust the volume and listen to the particular sample to determine if this is just an outlier or a mislabeled sample.
    
\subsection{Additional Resources}

\begin{itemize}
  \item \href{https://netron.app/}{Netron} (model visualization)
  \item \href{https://www.edgeimpulse.com/blog/visualizing-complex-datasets-in-edge-impulse}{Visualizing complex datasets in Edge Impulse}
  \item \href{https://projector.tensorflow.org/}{Tensorflow Embedding Projector}
  \item \href{https://www.tensorflow.org/tensorboard}{Tensorboard}
\end{itemize}
  
\section{Feature Engineering}

    0:00 / 0:00Klicken Sie auf die Pfeiltaste nach oben, um das Geschwindigkeitsmenu aufzurufen. Regeln Sie dann mit den Hoch und Runter Pfeiltasten die Geschwindigkeit. Klicken Sie Enter, um die jeweilige Geschwindigkeit festzulegen.Geschwindigkeit1.0xKlicke auf diesen Knopf und diess Video stumm oder nicht stumm zu schalten, oder drücke die HOCH oder RUNTER Knöpfe um die Lautstärke zu erhöhen oder abzusenken.Maximal Lautstärke.Videotranskript
    Beginn des Transkriptes, zum Ende springen.
    SPEAKER 1: Now that you have your data set ready,
    it's time to get into engineering the model.
    Now feature engineering is a technique of extracting features
    from the raw data.
    The objective of feature engineering is to employ these extra attributes
    in order to improve the machine learning results
    over just delivering a raw data.
    You're trying to help it identify the salient characteristics.
    Now this is the kind of thing that is more closer to an ML researcher or ML
    engineer who's more responsible for building the model.
    He or she is effectively taking a clean or cleaned up
    data set that's coming in from the data scientist for instance,
    who's already explored it.
    So here I'm trying to give you a flavor how the roles have shifted
    from problem definition, which was something the business analyst was
    doing to, something that the data engineer was responsible for going
    and finding the right data, handing it to the data scientist, who is then
    responsible for actually doing exploration,
    and then the ML researcher who's now trying
    to build a model based on that data set which has supposedly
    got the right characteristics.
    With that said, let's go back to our keyword spotting application
    where we want consumer-facing devices to be triggered
    on and off using voice commands.
    In keyword spotting, there are three main stages:
    those are the input, pre-processing step, and the neural network stage.
    In the input stage, we take signals that are
    coming in from the audio microphone on a device.
    One thing that is seemingly simple here, because it's just
    converting the raw audio signals into a digital signal.
    Looks fairly simple, right?
    However, what most people often don't realize,
    is that when you're capturing this input signal,
    there are a number of different hyperparameters involved.
    Examples include: what's the length of the sampling window, how
    big or small is each of those steps that you're looking at,
    what's the down sampling rate of the signal from the microphone?
    These might just seem like high level parameters,
    but they have a big impact on the machine learning models
    when you're actually training them.
    Because the characteristics we capture here
    are going to influence how the model actually
    learns to predict a certain outcome.
    As such, it is extremely important that in embedded devices
    that feature engineering of machine learning models
    is strongly tethered to the input sensor characteristics.
    Don't assume that you can capture these signals independently of the model,
    as that will affect your model's accuracy downstream when it is actually
    listening to a different microphone which
    has different characteristics that are then feeding it.
    So, a lot of nuances but very, very important.
    Now the next major step is pre-processing that captured input
    signal.
    In the context of keyword spotting, we typically convert the signals
    into an image, a spectrogram.
    A spectrogram is a visual representation of the spectrum of frequencies
    of a signal as it varies over time.
    Now when applied to an audio signal, spectrograms,
    which are also known as sonographs, voice prints, or voice grams,
    basically it's the image that can be fed into a convolutional neural network
    to make predictions.
    At which point it's really no different than training
    a model that's able to tell the difference between a cat and a dog.
    But even here, we have a rich and wide array of different hyperparameters
    in the pre-processing step.
    For instance, with spectrograms, you have various parameters
    like what's the length of each frame.
    How big is a stride across continuous audio signal inputs?
    These parameters determine how we generate the spectrogram image
    on the right, which is then fed into the network.
    So, just because you get an audio signal and you convert into a spectrogram
    does not mean the spectrograms are always the same.
    It depends on how you generate the spectrogram itself.
    It's got input characteristics.
    So that's a hyperparameter search in itself.
    Now even within the realm of spectrograms,
    you can have multiple types of spectrograms.
    You can have simple spectrograms that are just frequency representations,
    or you can have spectrograms where we enhance
    the way our ears really perceive sound.
    So on the left, I'm showing you spectrograms and on the right,
    I'm showing you MFCCs.
    Hopefully MFCCs will be clearer for the machine learning model.
    MFCCs basically stands for Mel-Frequency Cepstral Coefficients.
    These are coefficients that collectively make up a MFC.
    They tend to have better accuracy because they're
    kind of more closely reflecting how human ears behave.
    Certain frequencies are emphasized a bit more than other frequencies.
    Just the similar way that our ear actually
    differentiates different types of frequencies.
    Long story short, the spectrogram on the left, for instance,
    will have an accuracy of 80%.
    The MFCC for this particular set of keywords
    actually has an accuracy of 90%.
    It's almost a 10% difference.
    So clearly, the type of spectrogram in itself has a huge impact.
    Once again, there are a huge number of different types of parameters here,
    and each hyperparameter has a range of values.
    So the bottom line is that there's a large space for feature engineering.
    And as an ML researcher, one has to be busy twiddling these different things
    because it has a clear impact on the accuracy of the model.
    Recall that I said that the MFCC is about 10% better than the spectrogram
    that we're looking at.
    The spectrogram on the left versus the spectrogram on the right
    are two totally different things.
    Well not purely from an accuracy standpoint,
    it seems like it's a no-brainer that you should be doing an MFCC.
    But that's not always the case.
    Now if you focus only on accuracy, you will tend to overlook other key issues.
    For example, if you look in the bottom row, what I'm showing
    you is on device performance on a tiny ML device.
    You can clearly see that there's a trade off between MFCC on the right
    and the spectrogram on the left.
    While the spectrogram on the left only takes 123 milliseconds to run,
    the MFCC takes 229 milliseconds to run.
    Now this is because the MFCC needs you to do a bit more pre-processing
    on the embedded device, which has to happen in real time.
    In other words, it's on the critical path
    from the sensor to the neural network.
    The MFCC processing is being done in the middle.
    So depending on your deployment scenario,
    you're going to have to be very conscious about this.
    Now, the reason I'm emphasizing this is that your ML researcher needs
    to understand this, and that they cannot just focus on accuracy alone.
    This I cannot emphasize enough.
    I have seen this happen multiple times with companies
    where the ML researcher just thinks he or she can focus only on the accuracy
    aspect, and then they toss this model over to the people who are actually
    deploying it on device, and you will find that when they deployed on device,
    it doesn't work as expected because either the ML researcher did not
    think about some sort of key performance indicator metrics
    that were not given to him or her, and so this
    leads to a lot of back and forth which effectively is not a good thing
    to have.
    And especially with ML apps where you're trying to automate this process,
    you need to capture these kinds of characteristics.
    So at least you need to be really aware about this.
    So in general, feature engineering is extremely important.
    It's important for your ML researcher to understand
    that there are alternative methods to costly pre-processing of MFCCs.
    MFCCs might be the first place a lot of people will generally go.
    And it's true.
    People do go there.
    But it's not necessarily the best thing.
    Instead, we can in fact apply a variety of different classical techniques
    in order to improve the quality of the audio signal which
    is a huge impact, especially in the context of tiny ML models.
    Now this is something that most people often overlook.
    Because when you're traditionally dealing with big ML,
    you're not thinking about the computational resource implications.
    It's only in the context of deeply embedded machine
    learning that you're often worried about the computational overheads associated
    with simple things like just doing MFCC processing on it.
    And it's because these issues are so particular to the deployment
    of consumer-facing use cases and embedded systems
    that we have to worry about these kinds of different ways of solving problems.
    Again, it all comes back to the fact that we're
    dealing with resource constrained systems, right.
    It's unsurprising.
    So given how critical feature engineering is, often, what you end up
    doing in my ops pipe is that you will set up what is called a feature store.
    Simply put, feature stores are a data management layer for machine learning
    that allows one to easily discover, share, and reuse features.
    Ultimately, they help create a more effective machine learning pipeline.
    One interesting tidbit here.
    While I've seen feature engineering stores built up
    for large scale machine learning deployments like places like Google
    and so forth, I'm yet to discover feature stores for small embedded ML
    use cases where you're specifically dealing with sensor data.
    So hey, this might be an opportunity for a startup
    or a key feature in your product.
    Keep that in mind.
    Feature stores can be this big behemoth behind the scenes,
    as I'm showing you here in this figure.
    They're used throughout the life of an MLOps pipe,
    both from the engineering of the original model
    to the deployment and maintenance of that model.
    As you can see from the front part of this image,
    feature stores is especially useful for iteratively training and evaluating
    models.
    So investing them as part of the ML pipeline is useful for scaling
    out tiny ML.
    With that said, I'll see you in the next video.
    Ende des Transkriptes, zum Anfang springen
    Downloads und Transskripte
    Video
    Videodatei herunterladen
    Transskripte
    Download SubRip (.srt) file
    Download Text (.txt) file    
   
\section{Feature Engineering for KWS

\subsection{About This Reading}

    Keyword spotting (KWS) is one of the archetypal examples of TinyML in action. As we have discussed previously, the workflow consists of three stages: input, preprocessing, and output. Each of these stages is chained together in a data pipeline to output a label prediction for a specific audio signal. We will talk about each of the stages involved in the KWS workflow and delve into some of the specifics.
    
\subsection{Input Signal Capture}
    The input signal is obtained from an audio file or in real-time from an audio microphone. The conversion of sound from an audio microphone into a digital signal is relatively simple. The diaphragm in an audio microphone vibrates in the presence of sound, converting mechanical wave energy into an AC voltage signal (the method depends on the type of microphone, e.g., condenser vs. dynamic microphones). This AC signal is then converted into a digital signal. Thus, an audio microphone is essentially a transducer to convert sound into a voltage.
    
\GRAPHICSC{1.0}{1.0}{TinyML/4 MLOps/pasted image 1.png}  
    
    (adapted from course 2): Input sound waves cause the diaphram of the microphone to move causing the magnet to move in the coil. This creates a continuos time analog electrical singal in the form of a wave. This raw smooth analog signal is converted into a digital signal and passed into a neural network to compute the output.
    
    Beyond the microphone itself, there is a lot more going on in the input stage. Aspects such as the window length, window step, and the downsampling rate are important hyperparameters. Once determined, data is selected by “dragging” a window across the audio file and sending the data within this frame to the pre-processing stage.
    
\subsection{Pre-processing Blocks}

    During the pre-processing stage for KWS, the audio inputs are translated into a visual data format known as a spectrogram (as shown below). This is an image that reflects which frequencies were contained in the audio for each segment of time. Working with image data is easier than time-series data, and so converting the audio time-series to the spectrogram allows for more convenient analysis using a convolutional neural network.
    
\GRAPHICSC{1.0}{1.0}{TinyML/4 MLOps/Untitled.png}  

    (same as course 2): A visual depiction of the sliding window of an FFT and the resulting spectogram. On the left there is a time-domain digital signal (time on the x-axis vs. amplitude on the y-axis) overlayed with many red boxes showing the slices of sliding windows that are used to construct the spectogram. On the right is the spectogram which is shown to again have time along the x-axis but now frequency on the y-axis and then the color is the amplitude. It now looks much more like a paint splatter as compared to custered bars.
    
    As you might have expected, there are a wealth of hyperparameters that come with this preprocessing stage. Most notably, the frame length, frame stride, and frequency bands. These parameters tell us how long each segment of audio that we analyze will be, how many bands our frequency range will be split into, and how big the gaps between each window segment will be. Different values for these parameters will result in different spectrograms, which may be either easier or harder for our convolutional neural network to differentiate between. Ideally, we want to select values for these hyperparameters so that our network can easily differentiate between the different classes we wish to classify.
    
    An alternative visual representation of audio data is the Mel-frequency Cepstral Coefficients (MFCCs) – as shown below. These coefficients are often used as features in speech recognition algorithms, and generally result in better performance compared to spectrograms. In human speech, a single MFCC usually corresponds with a particular phoneme, which is the smallest unit of sound that we can create. Syllables are created using a combination of phonemes, words and a combination of syllables. MFCCs are slightly less interpretable than a spectrogram, since they correspond with specific phonemes, but are more useful for extracting features by a convolutional neural network. Hyperparameters for MFCCs include the number of coefficients, frame length, frame stride, and FFT length.
    
\GRAPHICSC{1.0}{1.0}{TinyML/4 MLOps/Untitled2.png}  

    (adapted from course 2): Spectrograms for the words “yes” and “no” are shown on the left and their respective MFCCs are shown on the right. Visually one can see how the MFCCs are much sharper in their color gradients and contrasts. This is because the MFCC clustering makes the differences more stark.
    
    The last audio-visual representation we will discuss for KWS is the Mel-filter banks. While MFCCs were popular for many years, nowadays Mel-filter banks are becoming increasingly popular. The two techniques are related, and, in fact, MFCCs can be derived from Mel-filter banks by the application of a discrete cosine transform. While the procedure is somewhat complex, the gist is that the audio signal is sent through a pre-emphasis filter, is sliced into overlapping frames, and then a window function is applied to each frame where a Fourier transform is then applied. From this, we can compute a power spectrum and subsequently compute the filter banks. For the interested reader, I recommend looking here for more information about filter banks.
    
\subsection{Summary}

    In this reading, we have talked a bit about some of the feature engineering techniques used in KWS, mainly, the depiction of audio in a visual representation that can then be parsed by a convolutional neural network. The basis of these techniques: the spectrogram, MFCCs, and Mel-filter banks, lies in digital signal processing theory, and while not strictly necessary to understand to implement these for performing KWS, it is certainly helpful to understand the workings of each to determine which is the most suitable for a given application.   
 
 
\section{Model Prototyping}

    0:00 / 0:00Klicken Sie auf die Pfeiltaste nach oben, um das Geschwindigkeitsmenu aufzurufen. Regeln Sie dann mit den Hoch und Runter Pfeiltasten die Geschwindigkeit. Klicken Sie Enter, um die jeweilige Geschwindigkeit festzulegen.Geschwindigkeit1.0xKlicke auf diesen Knopf und diess Video stumm oder nicht stumm zu schalten, oder drücke die HOCH oder RUNTER Knöpfe um die Lautstärke zu erhöhen oder abzusenken.Maximal Lautstärke.Videotranskript
    Beginn des Transkriptes, zum Ende springen.
    VIJAY JANAPA REDDI: Once you've got the features in place,
    then you get into model prototyping.
    Doing model prototyping is where you're experimenting
    with many different kinds of machine learning models
    in order to get a model that gives you good accuracy.
    This is the kind of thing that a machine learning researcher
    is going to be quite interested in, and she or he
    is kind of taking that cleaned up data that we talked about earlier
    and starting to engineer a very specific machine learning algorithm.
    Now typically, when you think about model prototyping,
    there are two ways to go about this.
    You can focus more on the model or the network architecture itself,
    or you can take a very data-centric approach.
    In this video, I'm going to focus more on the model-centric approach,
    and later on we'll talk about the data-centric approach in the next two
    or three videos.
    So in the model-centric approach, when you think about training models,
    you typically think about starting from scratch.
    Now, one of the key challenges when you think about training models
    from scratch is that you need a lot of computational horsepower to do this.
    And it's often not cheap because as an ML researcher what you want to do
    is you want to iterate through the training of models
    continuously until you find a model that you're really happy with.
    In terms of computational requirements, I
    think this chart does a rather good storytelling
    about the growing demands for compute.
    Since 2012, when AlexNet was introduced, which
    set the whole deep learning and demotion,
    the amount of compute used for training in our models
    has been increasing exponentially with the 3.5 month doubling time.
    That means every 3 and 1/2 months the computational capability
    has had to double in order to keep up with the general trend.
    If you can pair this with basic Moore's law, which
    has been the driving force behind CMOS electronics,
    that doubling happens roughly about every 18 months or two years.
    So how have we then been able to bridge this?
    Well, we've been building a lot of custom hardware for this, right?
    Google has been building GPUs.
    Nvidia has built custom GPUs dedicated for running
    machine learning very efficiently.
    So this costs a lot of money, and it's not something
    that everybody can just take for granted.
    Just to give you an example, you can look at what
    it costs in order to train XLNet.
    XLNet and that is something that's actually
    used for language understanding.
    Or here, take a look at GPT-2, which has over 1.5 billion parameters.
    And it costs over \$200 to get this model trained.
    Of course, it's a very awesome model as you can see on the right.
    You type something, and it actually generates text
    for you, which is really, really interesting and really cool.
    But today, GPT-2 is not even the hottest model in the market.
    So you might be arguing and say like, OK, well,
    hey, BJ, that is all in the context of this big model.
    I thought we're talking about TinyML here.
    Well, you could potentially argue we're not training such large models.
    But even for a five word keyword spotting model,
    it can take up to two hours.
    Ah, now, that's just for training in a one shot flow.
    Now imagine if you couple this with all the hyperparameters
    that I mentioned earlier from feature engineering.
    Every time you want to experiment with one hyperparameter
    change, it's going to need you to re-trigger the whole training
    process again, which means from an MLOps perspective,
    you're going to be spending a lot of money.
    So you need to be efficient about how you do your training.
    This is where we're increasingly seeing the trend towards pre-trained models.
    For me, this is like saying we don't all have to know
    how to build a 747 or a 787 plane.
    How does the economy work today?
    Few companies like Boeing or Airbus will go ahead and build these costly planes.
    And most of the companies will just go ahead and buy the basic plane
    and then customize it for their customers' needs.
    And the idea here of pre-trained models is very similar.
    You don't want to have to do everything from scratch on a brand new objective
    loss.
    It's much more difficult to do this than to try and take
    a backbone that's already got some sort of pre-trained characteristics,
    and then fine tune it a little bit.
    And that's really where this idea of transfer learning really comes in.
    Transfer learning is this process that what you want to do
    is that you want to get a network that you are interested in, which
    is capable of doing the kind of classifications or detections
    that you're interested in but not having to retrain everything
    every single weight in the network.
    So the first portion of any network generally
    ends up extracting critical features.
    And then the back portion of the network typically ends up being task-specific.
    It's what learns to say OK, well, a cat is a cat,
    or a dog is a dog and so forth.
    So what this means is that you can generally
    rely on fixing that front portion by locking it or freezing
    its layers and its weights and not changing it after you've
    trained the first version of it.
    And then only train the back end to be task specific.
    And this way, you can dramatically cut down the amount of training time
    you have.
    So it's very important from an MLOps standpoint
    to think about efficient ways of training networks.
    Now, embeddings are another way of looking at transfer learning.
    It's a slightly different variant of transfer learning.
    The general feature extraction can move or embed the input in a feature space
    like the space that I'm showing you here where
    inputs with rather similar features are likely to be
    grouped close to one another.
    And here, for example, I'm showing you that grouping in a t-SNE plot.
    T-SNE stands for a t distributed stochastic neighbor embedding.
    And it's a statistical method that is often
    used in machine learning for visualizing high dimensional data.
    Basically given a data point, it's got a certain location.
    It can map that data point into a certain 2-D or 3D sort of a map.
    And so you can sort of see the clusters.
    In this particular example, we're looking
    at different clusters for keywords in the TinyML use case
    where each cluster sort of corresponds to a different word.
    And each point within the cluster is a different recording or
    from different speakers basically, so it's really interesting.
    You can see from different words, there are these core latent characteristics
    that allow you to group words together.
    OK, so now, let's say you wanted to train a keyword spotting model real
    efficiently.
    This way, if you use embeddings, your ML researcher
    doesn't have to start from scratch.
    Instead, he or she can start with an embedding model that's
    been trained on a large prior set of samples
    as I'm showing them the left side in the green box.
    And then you have an embedding at the end.
    Then you can fine tune the model easily for exactly what you
    want to be able to do.
    So on the right, I'm showing you an example
    for how the embedding model has been fine tuned for three categories, which
    is five training examples, OK?
    So the three categories are keyword, something that's unknown,
    or it's background or silence.
    Now, this is awesome and wickedly useful not only
    to cut down the training time of the model and the cost associated with it,
    but also when you typically don't tend to have a very large data set.
    So it's an extremely powerful kind of concept.
    The take away or in summary, what I'm trying to emphasize
    is that you have different approaches to model prototyping.
    Often when we're talking about building models,
    we tend to talk about them in terms of training from scratch.
    But when you're really operationalizing things,
    you need to think about reducing the overheads of starting from scratch.
    So as I showed you, it can cost a lot of money to train models.
    Given how heterogeneous embedded systems are
    and given that you have a rich number of feature engineering knobs,
    you want to be efficient in terms of very quickly getting
    to a model that actually works.
    Now, when we get into a later MLOps stage called continuous training,
    I will dig much deeper into methods like AutoML and neural architecture search
    and so forth.
    I'm kind of deferred them because it would
    be too much to cover all that in the ML development section.
    So I've kind of spread some of these topics around.
    All right, with all that said, let's move on to the next section, which
    is model validation.
    Ende des Transkriptes, zum Anfang springen
    Downloads und Transskripte
    Video
    Videodatei herunterladen
    Transskripte
    Download SubRip (.srt) file
    Download Text (.txt) file
 
\section{Model Prototyping: Research vs. Production

\subsection{About This Reading}

    Once an ML model has been developed by a research team, there are a lot of steps required to convert it into a production-ready model. This is because many of the characteristics of an ML model are different in a production environment. These include characteristics related to the data, model training, overall goal, priority of metrics, and additional constraints that are not present in the research environment. In this reading, we will explore some of these differences.
    
    \begin{tabular}{ccc}
       & Academic/Research ML &   Production ML \\
        Data &  Static &    Dynamic - shifting \\
     Priority for Design & Highest overall accuracy &  Fast inference, good interpretability \\
     Model Training & Optimal tuning and training &  Continuously assess and retrain \\
        Fairness & Very important &     Crucial \\
  \end{tabular}        
    
\subsection{ifferences}
    
    \textbf{Data.} The most obvious difference in the production environment is the dynamicity of data. Data is often being collected continuously, which also means the data distribution will be continuously changing. Furthermore, the overall data size will get larger over time. Thus, our ML models have to be configured to take in and take advantage of new data, while ensuring it is able to handle changes in the data distribution which might affect the served predictions.
    
    \textbf{Priority for Design.} The goals in the research environment and production environment are often misaligned. The priority in an academic context is often a single metric, which is typically accuracy. Skim through half a dozen ML papers and you’ll typically see accuracy being compared between a variety of models on a particular task. However, this single-mindedness often does not translate into the production environment, where things like speed, efficiency, and interpretability become important. After all, what good is an algorithm that gives 100\% accuracy if it is impossibly slow and also violates the GDPR rule on the right to explanation? If a model is planning to be productized, these additional metrics should be considered alongside accuracy so that the goals in the research and production environments are appropriately aligned.
    
    \textbf{Model Training. }In a research context, the models are often only being developed to obtain a specific metric to be placed in a journal article, conference paper, or PowerPoint presentation. This is fundamentally different from a production ML environment, where the goal is high performance at a minimal cost. This means that instead of excessively fine-tuning a single model, we care about continuously assessing a model’s performance over time and retraining it once a significant amount of new data is available or the performance metrics begin to dip below an acceptable limit. Setting up an environment for continuous monitoring and retraining requires a lot of additional patience, expertise, and resources on top of the initial development of the model.
    
    \textbf{Fairness.} We already touched on this slightly, but in some jurisdictions - most notably the European Union - there are stipulations that ML models must be able to “explain” how they arrived at their predictions. This is often used in circumstances such as when deciding if an individual is given a bank loan for transparency purposes, giving the affected party the opportunity to see on what basis their bank loan application was rejected. For models in production, ensuring algorithmic fairness is crucial for legal compliance as well as ensuring equal opportunities to all potential algorithmic users. In a research context, less emphasis is generally placed on algorithmic fairness and interpretability since this often comes at the expense of performance. For example, the performance of a neural network is typically far superior to linear regression, but while linear regression is highly interpretable, neural networks are decidedly less interpretable.
    
\subsection{Summary}

    As we have just discussed, research and production environments for ML are quite different in their characteristics. Typically, data scientists and machine learning engineers are responsible for the research aspect, while MLOps/DevOps teams and production engineers are responsible for the production environment. In an ideal world, these two sides would work together and align their practices with the end-user in mind. Next time you are developing a model, try to think more critically about how it will be implemented in a real-world application, whether the way you are designing and evaluating the performance of your model is correctly aligned with this, and how the production environment will vary from your current research environment. 
    
    
\section{Model Validation}

    0:00 / 0:00Klicken Sie auf die Pfeiltaste nach oben, um das Geschwindigkeitsmenu aufzurufen. Regeln Sie dann mit den Hoch und Runter Pfeiltasten die Geschwindigkeit. Klicken Sie Enter, um die jeweilige Geschwindigkeit festzulegen.Geschwindigkeit1.0xKlicke auf diesen Knopf und diess Video stumm oder nicht stumm zu schalten, oder drücke die HOCH oder RUNTER Knöpfe um die Lautstärke zu erhöhen oder abzusenken.Maximal Lautstärke.Videotranskript
    Beginn des Transkriptes, zum Ende springen.
    VIJAY JANAPA REDDI: At this point, we have a model, which we have protyped,
    and we want to validate.
    Now, model validation is a phase in the ML space
    where you want to quantify how well the model's actually going to perform
    in terms of producing predictions.
    And you want to check whether it produces predictions
    with enough fidelity that they can used reliably
    to achieve the business objectives that we set out to achieve in the problem
    definition phase.
    It basically quantifies the performance of the model
    on unseen data, data that it's actually not seen during the training phase.
    And this is the kind of thing that an ML researcher's
    is going to be very interested in.
    Now, typically, when you're trying to do some sort of error analysis
    on the model, you can think about it either from a data-centric perspective
    or a model-centric perspective.
    What do I mean by a model-centric perspective?
    What I mean is that you tend to look at the characteristics
    of the network itself.
    In other words, you'll try and understand, well,
    can I change some activation function [INAUDIBLE] loss function, optimizers,
    or maybe add some dropout layers and whatnot in order to figure out
    if the model can be improved?
    But really, in a way, it is actually quite limiting.
    So I'd rather think about it from the data
    a perspective to try and take a network that you prototyped, and then
    really subjected to various characteristics of the data.
    Typically, when you're optimizing the hyperparameters of your model,
    you might end up overfitting if you were to optimize for a given test and train
    set.
    So that's because your model's naturally searching for hyperparameters
    that best fit that test and train set that you made.
    So to solve this issue, you generally create a holdout set.
    Now, this is about 10% of the data that you would typically not
    use during the validation stages.
    Now, after you optimize the model on the test and train
    splits, then you introduce it to the holdout set that is actually not seen.
    And this is fairly typical.
    But to minimize a sampling bias as to when
    you take that one holdout, what you want to do is you actually
    want to create multiple splits and validate the model
    on all of those different splits.
    And this is called K-fold cross-validation,
    where K is roughly about 5 or 10 given.
    And that's fairly representative.
    Now, a variant of that K-fold cross-validation
    is leave-one-out cross-validation, where each sample in the data
    is sort of treated as a separate test set,
    while the remaining samples form the training set.
    This variant is effectively equal to saying
    when K is equal to the number of different observations you have.
    It is computationally much more intensive to perform,
    although the results are much more reliable,
    and the unbiased estimation that you get from the model
    is actually much more robust.
    This is a technique that is actually really good
    when you have small data sets, which is kind of typical in TinyML cases.
    And that's why I'm specifically bringing this up.
    You got to make the caveat that this can be
    computationally intensive, because you got to churn through the data.
    But assuming that that's not as much an issue
    and you really are focused on the fact that you
    want to have a good model with a limited data set size,
    then this is a good approach to actually consider for your ML researcher.
    Now, the issue with the K-fold cross-validation
    is that you might want each foal to only contain a single group.
    For example, let's say you have a data set of 20 different keywords
    and you want to predict the success rate for--
    of keyword spotting on those words.
    Then to keep those folds, quote, unquote, "pure,"
    you create it in such a way that every one of those folds only
    contains a single keyword.
    So you basically create a fold for each word,
    and that way you create a version of K-fold cross-validation
    where you leave one of those words out.
    This gives you a much better and robust way
    of doing validation to study the effectiveness of your model.
    These are simple techniques, but they're also
    very powerful techniques when you don't tend to have large data
    sets available for you.
    And of course, in TinyML, anomaly detection's a very big thing,
    and loads of people want to use anomaly detection for various things,
    like predictive maintenance and manufacturing.
    It's actually considered one of the biggest market spaces for TinyML.
    Now, in this case, we definitely want to be
    careful with how you handle the data.
    The reason is because you're going to end up overfitting it.
    Overfitting is a major concern here, since your training data could actually
    contain information from the future.
    It's important that all your training data
    happens before your test data in the context anomaly detection,
    because anomaly detection is a time series kind of data,
    so you can't be looking into the future.
    Just because you have the data doesn't mean
    you can really be looking into the future and making--
    using some of that data to make the predictions, because that's just bogus.
    So one way of validating time series data
    is by using the K-fold cross-validation approach
    and making sure that each of the fold contains the training data only
    before the test data.
    So this is a way of separating out.
    So please be careful, when you're doing time series and TinyML
    that you don't fall into this pitfall--
    that is, if you're an ML researcher or a data scientist-- something
    to keep in mind.
    So far we've been talking about how an ML researcher might
    go around shuffling the data to make sure that the model's robust enough,
    but error analysis can also be more systematic than just doing that.
    You can break things down more into a data-centric sort of an approach.
    This is a chart where I have to give credit to Andrew Ng,
    because it's a chart that beautifully and very simply
    explains a data-centric approach.
    What do you do in a model-centric approach?
    In a model-centric approach, you are constantly
    tweaking the model to try and make it better
    in order to improve the accuracy.
    In a data-centric approach, which you really do is you try and improve--
    take that time that you have and improve the quality of the data,
    and finding clean labels, for instance.
    So you can clearly see that, when you have a noisy data
    set versus a clean data set where, for example, you looked closely at,
    why is your model performing poorly--
    OK, well, it turns out it's performing poorly
    on these particular types of samples.
    And then you focus your effort on going and getting
    good samples put in there so your model gets exposed to those kinds of samples.
    That's a much more effective use of your time
    than trying to change the hyperparameters in your neural network
    to try and get it to learn some of these rare cases in the data.
    So as you can clearly see, the clean data set
    definitely has a higher accuracy for any given number of samples
    that you would pick.
    And again, I want to give a shout-out for Andrew,
    because he's the one who's popularized this notion
    of data-centric view of machine learning to improve model quality.
    Now, when you think about this, you want to analyze the model's error
    during development.
    A data-centric approach basically looks at how
    to improve the model quality by looking at characteristics within the data.
    Now, obviously, when you're training the model,
    you don't want the model to overfit or to underfit.
    You want to stop at that sweet spot where
    the validation error and the training error are not too widely apart.
    And one of the ways you can do this, if you're
    looking at it from a data-centric approach in order
    to improve model accuracy, is you can use a variety of data augmentation
    techniques.
    And these data augmentation techniques, like blurring, cropping,
    scaling, rotating, creating different barrel rolls, and so forth of images
    can drastically improve the quality of your model.
    And this is without really going in collecting new data.
    So what I'm talking about is data augmentation
    is you have your existing data set, and you're twisting and changing
    the data in different ways in order to effectively artificially increase
    the scale of that data.
    And it's a really powerful technique.
    In fact, again, talking about what Andrew had done,
    Andrew had launched a data-centric competition in '21--
    in the year 2021-- to foster more error analysis on the data,
    rather than trying to fix the model.
    Effectively, the competition was set up such
    that it inverts the traditional machine learning approach.
    Instead, it basically asks you to improve
    the data set for a fixed neural network architecture.
    As it turned out, the leading winners on this leaderboard had actually used
    a rich variety of data augmentation techniques to actually boost
    the model's performance on the [INAUDIBLE] data challenge that they
    had here--
    so really interesting the point that I'm landing hard on
    is let's focus on the data.
    Focus on specific goals within the issues you have seen.
    For example, let's say you have a lot of false positives.
    Then try and focus specifically on false positives.
    If you have false negatives, where you're doing poorly,
    then reduce the false negatives.
    You can effectively try and change decision margins
    in order to be able to deal with these sorts of issues.
    You can also apply a wide range of different data augmentation techniques,
    like blurring and so forth, random cropping and so forth.
    And if your issue falls in the high false positive rate for,
    say, detecting some sort of defect, then increase the confidence threshold
    for detecting the defect.
    These are some things that are actually being
    used in practice at a large number of startups,
    and they actually prove much more effective because startups don't really
    have a wide amount of resources to go out and build complex models.
    But they can be much more efficient in handling the data.
    And this is a general theme that we're actually
    learning from startups that are actually exploring the data much more
    efficiently.
    The bottom line point that I'm trying to get to
    is that, when you're talking about error analysis, that classical approach has
    always been to look at the model, and blame the model,
    and say the model's not doing well.
    Instead, the approach now that's increasingly becoming popular,
    especially for MLOps, is hey, let's look at the data
    and let's try and do better error analysis on the data itself.
    So you effectively want to be able to bin the data, bin the characteristics
    that you find into various kinds of buckets of problems
    that you're spotting on how the model is performing or not performing,
    and then you take those particular buckets, the ones with the most
    problematic scenarios-- you take those buckets,
    and then you focus on those buckets by systematically trying to figure out,
    what is the model not performing well on?
    And then try and improve the quality of the data for it,
    because that definitely helps you improve the accuracy of the model
    much faster than trying to come up with better hyperparameters
    for the neural network architecture.
    That's the key takeaway here.
    Ende des Transkriptes, zum Anfang springen
    Downloads und Transskripte
    Video
    Videodatei herunterladen
    Transskripte
    Download SubRip (.srt) file
    Download Text (.txt) file    
    
\section{Model Validation}


    0:00 / 0:00Klicken Sie auf die Pfeiltaste nach oben, um das Geschwindigkeitsmenu aufzurufen. Regeln Sie dann mit den Hoch und Runter Pfeiltasten die Geschwindigkeit. Klicken Sie Enter, um die jeweilige Geschwindigkeit festzulegen.Geschwindigkeit1.0xKlicke auf diesen Knopf und diess Video stumm oder nicht stumm zu schalten, oder drücke die HOCH oder RUNTER Knöpfe um die Lautstärke zu erhöhen oder abzusenken.Maximal Lautstärke.Videotranskript
    Beginn des Transkriptes, zum Ende springen.
    VIJAY JANAPA REDDI: At this point, we have a model, which we have protyped,
    and we want to validate.
    Now, model validation is a phase in the ML space
    where you want to quantify how well the model's actually going to perform
    in terms of producing predictions.
    And you want to check whether it produces predictions
    with enough fidelity that they can used reliably
    to achieve the business objectives that we set out to achieve in the problem
    definition phase.
    It basically quantifies the performance of the model
    on unseen data, data that it's actually not seen during the training phase.
    And this is the kind of thing that an ML researcher's
    is going to be very interested in.
    Now, typically, when you're trying to do some sort of error analysis
    on the model, you can think about it either from a data-centric perspective
    or a model-centric perspective.
    What do I mean by a model-centric perspective?
    What I mean is that you tend to look at the characteristics
    of the network itself.
    In other words, you'll try and understand, well,
    can I change some activation function [INAUDIBLE] loss function, optimizers,
    or maybe add some dropout layers and whatnot in order to figure out
    if the model can be improved?
    But really, in a way, it is actually quite limiting.
    So I'd rather think about it from the data
    a perspective to try and take a network that you prototyped, and then
    really subjected to various characteristics of the data.
    Typically, when you're optimizing the hyperparameters of your model,
    you might end up overfitting if you were to optimize for a given test and train
    set.
    So that's because your model's naturally searching for hyperparameters
    that best fit that test and train set that you made.
    So to solve this issue, you generally create a holdout set.
    Now, this is about 10% of the data that you would typically not
    use during the validation stages.
    Now, after you optimize the model on the test and train
    splits, then you introduce it to the holdout set that is actually not seen.
    And this is fairly typical.
    But to minimize a sampling bias as to when
    you take that one holdout, what you want to do is you actually
    want to create multiple splits and validate the model
    on all of those different splits.
    And this is called K-fold cross-validation,
    where K is roughly about 5 or 10 given.
    And that's fairly representative.
    Now, a variant of that K-fold cross-validation
    is leave-one-out cross-validation, where each sample in the data
    is sort of treated as a separate test set,
    while the remaining samples form the training set.
    This variant is effectively equal to saying
    when K is equal to the number of different observations you have.
    It is computationally much more intensive to perform,
    although the results are much more reliable,
    and the unbiased estimation that you get from the model
    is actually much more robust.
    This is a technique that is actually really good
    when you have small data sets, which is kind of typical in TinyML cases.
    And that's why I'm specifically bringing this up.
    You got to make the caveat that this can be
    computationally intensive, because you got to churn through the data.
    But assuming that that's not as much an issue
    and you really are focused on the fact that you
    want to have a good model with a limited data set size,
    then this is a good approach to actually consider for your ML researcher.
    Now, the issue with the K-fold cross-validation
    is that you might want each foal to only contain a single group.
    For example, let's say you have a data set of 20 different keywords
    and you want to predict the success rate for--
    of keyword spotting on those words.
    Then to keep those folds, quote, unquote, "pure,"
    you create it in such a way that every one of those folds only
    contains a single keyword.
    So you basically create a fold for each word,
    and that way you create a version of K-fold cross-validation
    where you leave one of those words out.
    This gives you a much better and robust way
    of doing validation to study the effectiveness of your model.
    These are simple techniques, but they're also
    very powerful techniques when you don't tend to have large data
    sets available for you.
    And of course, in TinyML, anomaly detection's a very big thing,
    and loads of people want to use anomaly detection for various things,
    like predictive maintenance and manufacturing.
    It's actually considered one of the biggest market spaces for TinyML.
    Now, in this case, we definitely want to be
    careful with how you handle the data.
    The reason is because you're going to end up overfitting it.
    Overfitting is a major concern here, since your training data could actually
    contain information from the future.
    It's important that all your training data
    happens before your test data in the context anomaly detection,
    because anomaly detection is a time series kind of data,
    so you can't be looking into the future.
    Just because you have the data doesn't mean
    you can really be looking into the future and making--
    using some of that data to make the predictions, because that's just bogus.
    So one way of validating time series data
    is by using the K-fold cross-validation approach
    and making sure that each of the fold contains the training data only
    before the test data.
    So this is a way of separating out.
    So please be careful, when you're doing time series and TinyML
    that you don't fall into this pitfall--
    that is, if you're an ML researcher or a data scientist-- something
    to keep in mind.
    So far we've been talking about how an ML researcher might
    go around shuffling the data to make sure that the model's robust enough,
    but error analysis can also be more systematic than just doing that.
    You can break things down more into a data-centric sort of an approach.
    This is a chart where I have to give credit to Andrew Ng,
    because it's a chart that beautifully and very simply
    explains a data-centric approach.
    What do you do in a model-centric approach?
    In a model-centric approach, you are constantly
    tweaking the model to try and make it better
    in order to improve the accuracy.
    In a data-centric approach, which you really do is you try and improve--
    take that time that you have and improve the quality of the data,
    and finding clean labels, for instance.
    So you can clearly see that, when you have a noisy data
    set versus a clean data set where, for example, you looked closely at,
    why is your model performing poorly--
    OK, well, it turns out it's performing poorly
    on these particular types of samples.
    And then you focus your effort on going and getting
    good samples put in there so your model gets exposed to those kinds of samples.
    That's a much more effective use of your time
    than trying to change the hyperparameters in your neural network
    to try and get it to learn some of these rare cases in the data.
    So as you can clearly see, the clean data set
    definitely has a higher accuracy for any given number of samples
    that you would pick.
    And again, I want to give a shout-out for Andrew,
    because he's the one who's popularized this notion
    of data-centric view of machine learning to improve model quality.
    Now, when you think about this, you want to analyze the model's error
    during development.
    A data-centric approach basically looks at how
    to improve the model quality by looking at characteristics within the data.
    Now, obviously, when you're training the model,
    you don't want the model to overfit or to underfit.
    You want to stop at that sweet spot where
    the validation error and the training error are not too widely apart.
    And one of the ways you can do this, if you're
    looking at it from a data-centric approach in order
    to improve model accuracy, is you can use a variety of data augmentation
    techniques.
    And these data augmentation techniques, like blurring, cropping,
    scaling, rotating, creating different barrel rolls, and so forth of images
    can drastically improve the quality of your model.
    And this is without really going in collecting new data.
    So what I'm talking about is data augmentation
    is you have your existing data set, and you're twisting and changing
    the data in different ways in order to effectively artificially increase
    the scale of that data.
    And it's a really powerful technique.
    In fact, again, talking about what Andrew had done,
    Andrew had launched a data-centric competition in '21--
    in the year 2021-- to foster more error analysis on the data,
    rather than trying to fix the model.
    Effectively, the competition was set up such
    that it inverts the traditional machine learning approach.
    Instead, it basically asks you to improve
    the data set for a fixed neural network architecture.
    As it turned out, the leading winners on this leaderboard had actually used
    a rich variety of data augmentation techniques to actually boost
    the model's performance on the [INAUDIBLE] data challenge that they
    had here--
    so really interesting the point that I'm landing hard on
    is let's focus on the data.
    Focus on specific goals within the issues you have seen.
    For example, let's say you have a lot of false positives.
    Then try and focus specifically on false positives.
    If you have false negatives, where you're doing poorly,
    then reduce the false negatives.
    You can effectively try and change decision margins
    in order to be able to deal with these sorts of issues.
    You can also apply a wide range of different data augmentation techniques,
    like blurring and so forth, random cropping and so forth.
    And if your issue falls in the high false positive rate for,
    say, detecting some sort of defect, then increase the confidence threshold
    for detecting the defect.
    These are some things that are actually being
    used in practice at a large number of startups,
    and they actually prove much more effective because startups don't really
    have a wide amount of resources to go out and build complex models.
    But they can be much more efficient in handling the data.
    And this is a general theme that we're actually
    learning from startups that are actually exploring the data much more
    efficiently.
    The bottom line point that I'm trying to get to
    is that, when you're talking about error analysis, that classical approach has
    always been to look at the model, and blame the model,
    and say the model's not doing well.
    Instead, the approach now that's increasingly becoming popular,
    especially for MLOps, is hey, let's look at the data
    and let's try and do better error analysis on the data itself.
    So you effectively want to be able to bin the data, bin the characteristics
    that you find into various kinds of buckets of problems
    that you're spotting on how the model is performing or not performing,
    and then you take those particular buckets, the ones with the most
    problematic scenarios-- you take those buckets,
    and then you focus on those buckets by systematically trying to figure out,
    what is the model not performing well on?
    And then try and improve the quality of the data for it,
    because that definitely helps you improve the accuracy of the model
    much faster than trying to come up with better hyperparameters
    for the neural network architecture.
    That's the key takeaway here.
    Ende des Transkriptes, zum Anfang springen
    Downloads und Transskripte
    Video
    Videodatei herunterladen
    Transskripte
    Download SubRip (.srt) file
    Download Text (.txt) file
    
\section{Data Engineering}

    0:00 / 0:00Klicken Sie auf die Pfeiltaste nach oben, um das Geschwindigkeitsmenu aufzurufen. Regeln Sie dann mit den Hoch und Runter Pfeiltasten die Geschwindigkeit. Klicken Sie Enter, um die jeweilige Geschwindigkeit festzulegen.Geschwindigkeit1.0xKlicke auf diesen Knopf und diess Video stumm oder nicht stumm zu schalten, oder drücke die HOCH oder RUNTER Knöpfe um die Lautstärke zu erhöhen oder abzusenken.Maximal Lautstärke.Videotranskript
    Beginn des Transkriptes, zum Ende springen.
    SPEAKER: Before wrapping up ML development,
    I want to touch on one key thing which is data engineering, because it's
    an extremely important portion and I don't want to ignore the ML development
    stage from the data engineering component.
    We have talked about data engineering previously
    when we talked about data selection and I didn't find the right sources.
    But really, when we're talking about it from the perspective of ML Ops,
    we need to have someone just dedicated or responsible
    for dealing with data engineering, and that's the data engineer's job.
    Their goal is really to help identify what the requirements are for the data,
    go ahead and try and figure out how to gather it,
    then make sure it's being refined, and then sustain it.
    In other words, keep it updated.
    So from a requirement standpoint they got
    to identify what the problem is, figure out what
    are the right sets of permissions and rights that need to be there.
    Then they need to build up pipelines that
    allow us to quickly go out and gather the data,
    figure out how the collection labeling is going to be done.
    Then it needs to be validated or vetted to make sure the quality of the labels
    is good, and then continue to improve this over time.
    I've mentioned all of this in some capacity,
    but here I'm trying to focus only on the data engineer's role
    and try and put this into context.
    Let's take for example, the speech commands
    data set that Pete actually put together.
    And in order to do this, if you look at the data engineering effort that
    went in, the first thing was I didn't find the requirements that it's
    an English data set specifically focused on keywords
    and it's got to have the right licenses.
    Then had to go through the effort of gathering this data by crowdsourcing
    various kinds of volunteers in order to contribute
    the keywords that are needed.
    Luckily, we have over 1,000 examples per keyword,
    which is great with a variety of different accents.
    Then there was a process of actually manually removing
    poor data labels in there.
    In other words, one had to actually physically
    go and listen to the audio file and then remove the quiet or the low volume
    recordings or mislabeled recordings.
    That's a tedious process.
    And then of course, speech commands has improved from version 1 to version 2
    with the inclusion of additional keywords.
    Now the point that I'm trying to make, though, is that all of this effort
    is really labor-intensive, because someone is manually
    labeling the data set.
    And that is problematic.
    Why?
    Because it's expensive.
    Because it actually includes effort.
    And if you're trying to do this from a commercial standpoint,
    it's going to have to scale really well--
    which point you're going to have to explicitly pay for data sets.
    You can't really just have crowd sourced data,
    because you can't control a crowd sourced data set.
    Here, I'm showing you images of cats and dogs.
    But you get the idea that it's speech commands
    data set it really is about just having keywords like yes, no, up,
    down, left, right, and so forth.
    And that has been manually done, but ideally what we want to do
    is replace that manual effort with an automated effort.
    Later on in other parts of the ML Ops pipe I will talk about automated data
    flow engines that have been created in order to expedite this whole process.
    This is an extremely important piece of the ML development stage,
    and that's the only reason I'm touching on it.
    But we already covered a fair bit of material
    in the ML development ops stage.
    So I'm going to refrain from cramming in much more.
    But I just wanted to give you a heads up that data engineering is
    a critical portion of the ML development cycle,
    and we'll come back to it later on.
    Keep that in mind.
    And it's the data engineer's job to make sure
    that there are these nice smooth data flows that
    are automated to create the data sets that are of interest to you.
    And that data can actually be taken and then
    used to train learned representations where you can either
    start training them all from scratch or you
    can use the transfer learning based methods
    that are talked about earlier in the model prototyping stage.
    Ende des Transkriptes, zum Anfang springen
    Downloads und Transskripte
    Video
    Videodatei herunterladen
    Transskripte
    Download SubRip (.srt) file
    Download Text (.txt) file    
        
\section{ML Development Impact on MLOps}


    0:00 / 0:00Klicken Sie auf die Pfeiltaste nach oben, um das Geschwindigkeitsmenu aufzurufen. Regeln Sie dann mit den Hoch und Runter Pfeiltasten die Geschwindigkeit. Klicken Sie Enter, um die jeweilige Geschwindigkeit festzulegen.Geschwindigkeit1.0xKlicke auf diesen Knopf und diess Video stumm oder nicht stumm zu schalten, oder drücke die HOCH oder RUNTER Knöpfe um die Lautstärke zu erhöhen oder abzusenken.Maximal Lautstärke.Videotranskript
    Beginn des Transkriptes, zum Ende springen.
    VIJAY JANAPA REDDI: Congratulations, you're
    just about to wrap up the ML development stage.
    But before we close shop, one of the things that I want to talk about
    is how these various components that we have seen here,
    that are all part of the ML development stage,
    actually can have a downstream impact effect on the ML DevOps pipeline.
    Remember that the ML development is one of the first stages in the MLOps
    pipeline.
    That means you should be asking yourself this question,
    what's the impact of the ML development stage
    on the rest of the MLOps pipeline?
    Now, I'm not going to be able to step through every single one
    of the potential connections that we have here,
    so I'm going to cherry pick a few just to kind of give you a sampling.
    And as we go through the course, I will give you
    more and more of these connections, and you'll hopefully
    be able to build a bigger picture of how this really
    requires you to kind of think about MLOps as an integrated flow rather
    than discrete items.
    We kind of broke things out into these discrete boxes so it's easy for you
    to gather the concepts.
    So, if we talk about ML development, let's try and understand
    how it has an effect on model deployment.
    Typically model deployment concerns packaging, testing,
    and deploying a model to a serving environment for online experimentation
    and production serving.
    That's what model deployment is all about.
    That's where the pedal really meets the metal.
    Now in the context of embedded systems and tinyML,
    it's all about knowing the kind of fleet you're going to be deploying into.
    The heterogeneity of that fleet.
    What are the resource constraints?
    What are the limitations of those embedded devices?
    Now, as I've said earlier in the model evaluation phase of ML development,
    if you are not careful and you ignore the model evaluation phase,
    you're going to have a serious problem when
    it comes to deploying the models into the production pipes.
    I have seen this happen literally firsthand at some major organizations,
    where the ML researchers ignore the exact target device, and they go ahead
    and build a model, and when they deploy it, it does not work.
    And this is one of the reasons I am picking intentionally
    this particular connection to emphasize.
    It happens and I've seen it, and I have seen it happen multiple times.
    So what this is telling us is that we need to be aware of the decisions
    we make ahead and what the downstream impact is.
    Let me pick another example.
    Typically when you deploy models into the field,
    you tend to do continuous monitoring.
    Continuous monitoring is about monitoring
    the effectiveness and efficiency of a deployed model in the wild.
    In a traditional cloud-scale ML system, what you end up doing is
    you end up getting continuous feedback about how
    the model is doing on the various inputs that are coming in
    and what are the predictions that it's making.
    And you gather that feedback continuously
    because network connectivity is something
    that's obviously going to be guaranteed in a data-center environment.
    However, think about this in the context of embedded systems.
    When I was talking about embedded systems in the intro module,
    I talked about some of the limitations around communication.
    I talked about the cost of communication versus computation.
    What that means is that when you create an ML deployment for tinyML and you try
    and scale it out, you cannot assume you're going to have continuous
    feedback from these embedded systems that are actually performing these
    inferences.
    Which means you cannot rely on continuous monitoring in the same
    traditional way that you would assume continuous monitoring in a big ML
    system.
    And that is another place where the ML engineers need to be aware.
    Because the assumption you might be making in the ML development stage
    is that I might actually be able to look at historical data that's continuously
    coming in, and then based on that I can do agile decisions for how
    I'm going to tune the model, be it for model prototyping
    or be it for error analysis and so forth.
    But if you can't have that sort of a tight continuous monitoring feedback
    loop, then you need to make alternative decisions.
    Now what those are is certainly going to depend on the particular use case
    you're dealing with.
    But nonetheless, I hope it conveys the point that you have to be aware
    of these feedback loops, or these interactions between various stages
    of the ML lifecycle.
    Last but not least, we have talked about the various people that
    are involved in the MLOps lifecycle.
    Just in the ML development stage, we went from the business analyst
    helping us define the problem statement or creating the key question,
    to then talking about the data engineer whose role really
    is to help us source the right data and be able to build data engineering
    pipelines that will allow us to have the data scientists explore that data,
    find interesting trends, and let the ML researcher then go ahead
    and actually build the models, and validate them and then evaluate them.
    And my friends, with that said, I want to congratulate you again
    on getting past the first stage of MLOps.
    You've seen a lot of new things in this particular module.
    Specifically, you have seen how Larissa Suzuki introduces the MLOps stage
    and then you see how I open up the MLOps box.
    And then I go in, and then I talk about some
    of the issues and opportunities and challenges
    that we face when we are talking about MLOps in the context of tinyML.
    Keep this general pattern of principle in mind,
    because it's going to repeat over and over, because that's
    how I have structured the course.
    And hopefully that is working out well for you.
    With that said hopefully you take a nice break
    and come back, do the next stage of MLOps.
    Ende des Transkriptes, zum Anfang springen
    Downloads und Transskripte
    Video
    Videodatei herunterladen
    Transskripte
    Download SubRip (.srt) file
    Download Text (.txt) file    

\section{Overview of Training Operationalization}

\GRAPHICSC{1.0}{1.0}{TinyML/4 MLOps/003.png}


\subsection{Objective
    ML Training operationalization sounds like an intimidating term at first, but it is essentially the process of automating and improving upon aspects of the standard ML workflow that we have come to know and love. The procedure is shown above. The topics we will cover include:

\begin{itemize}
  \item What is ML training operationalization and why is it important?
    \begin{itemize}
      \item Continuous integration
      \item What is continuous integration?
      \item What are the sources and various roles that CI plays in ML? 
      \item What are some challenges for TinyML with CI?
      \item Why is testing TinyML challenging compared to traditional ML systems?
    \end{itemize}    
  \item Continuous delivery
    \begin{itemize}
      \item What is continuous delivery?
      \item What are the various stages of continuous delivery?
      \item What is staging and why is it helpful?
      \item Who does acceptance testing?
    \end{itemize}    
  \item Production deployment
    \begin{itemize}
      \item What is it about?
      \item What are smoke tests?
      \item What are commonly used ML deployment strategies for releases?
    \end{itemize}    
  \end{itemize}    


\subsection{Motivation}

    Take, for example, a medium-sized company with around a hundred employees. At this organizational level, we may have dozens or even hundreds of models being developed. All of these models will be at different stages; some in development, some ready to be tested, and some already deployed. These models may have been developed by different individuals, perhaps using different packages and software versions. As you might suspect, from all of this heterogeneity in the model development process, problems are likely to ensue. These problems may just be inefficiencies, with individuals spending time poring over an unfamiliar PyTorch implementation while everyone else in the company is using TensorFlow. However, sometimes these can be more severe problems, giving rise to errors in production that might have been missed during the testing stage due to an inconsistent testing regime, or perhaps a model that has not been tested yet being put into deployment, or even losing a model because it was not stored somewhere central for the rest of the team to access!
    
\subsection{Training Operationalization}
    How can we tackle these kinds of problems in an efficient way? Training operationalization! The idea of training operationalization is that the training and deployment processes are streamlined and automated in such a way that consistent testing and monitoring procedures are applied to each and every model. For example, a model registry can be created (e.g., using MLflow) to store model versions and relevant artifacts (e.g., model parameters, dependencies) that can be compared later. A model registry is very similar to a GitHub or DockerHub repository for machine learning models. Continuous integration/delivery capabilities (e.g., Jenkins, GitLab, Ansible) can be introduced in order to dynamically update models and pass them through automated testing procedures (e.g., unit tests, acceptance tests, smoke tests) to ensure all of their functionalities remain intact after they have been edited. Not only can these be tested for a single system, but they can be tested across a range of hardware platforms, operating systems, and package versions. This is very similar to how DevOps is used to streamline software development, but reimagined for streamlining the training, development, and deployment of machine learning models.
    
    The specific procedure for ML training operationalization that an organization might take is highly variable. Smaller companies may only have a single testing stage, whereas a larger organization may have multiple stages and a procedure more complex than that outlined above. Over the remainder of this chapter, we will become more familiar with training operationalization and how it can help us to mitigate errors and improve the efficiency of our ML workflow.

\section{Training Operationalization}

    0:00 / 0:00Klicken Sie auf die Pfeiltaste nach oben, um das Geschwindigkeitsmenu aufzurufen. Regeln Sie dann mit den Hoch und Runter Pfeiltasten die Geschwindigkeit. Klicken Sie Enter, um die jeweilige Geschwindigkeit festzulegen.Geschwindigkeit1.0xKlicke auf diesen Knopf und diess Video stumm oder nicht stumm zu schalten, oder drücke die HOCH oder RUNTER Knöpfe um die Lautstärke zu erhöhen oder abzusenken.Maximal Lautstärke.Videotranskript
    Beginn des Transkriptes, zum Ende springen.
    LARA SUZUKI: Hello.
    In our last lesson, we had a closer look at the ML development
    stage of our MLOps framework.
    It's now time to move to the next stage, that is,
    of training operationalisation, or simply model training.
    In this session, we will cover deeper details of model training
    which concerns to creating a reliable training
    pipeline for a machine learning model.
    As in the last lesson, once data scientists
    are done with experimentation, they build source code
    and run various tests.
    The output of this step are pipeline components,
    such as packages, executables, and artifacts
    to be deployed at a later stage.
    The training operationalisation stage concerns
    with automating the process of packaging,
    for example, container images stored in a container
    registry, the training-pipeline runtime representation, which references
    the components stored in the artifacts repository,
    and also the testing and deployment of repeatable and reliable training
    pipelines.
    In summary, it mainly involves the process
    of continuous integration and continuous delivery
    that we also discussed in our first lesson.
    The source code for components must be modularized and ideally containerized.
    This will help data scientists to decouple the execution environment
    from the custom code runtime.
    This will make the code reproducible between different environments
    and isolate each component.
    The output of this phase is a validated training model.
    Throughout the entire run, data scientists
    track the metrics and metadata for the pipeline alongside the model,
    and they retrain model provenance and they ensure
    model traceability and auditability.
    The core MLOps capabilities employed at this stage are ML pipeline--
    target environment, data sources, service account, containers,
    and others--
    the CI/CD pipeline, version control, software artifacts, training data sets.
    In training operationalisation, the machine learning pipelines
    consists of configurations to deploy the pipelines, the environment, data
    source, service account, amongst others, and the deployment of a pipeline
    using standard CI/CD processes and tools.
    In this diagram, I would like to highlight the two
    main stages of the CI/CD pipelines.
    In the continuous integration-- the CI stage--
    the source code is unit tested and the training pipeline is filled
    and integration tested.
    Any artifacts that are created by the field
    are stored in the artifact repository.
    In the continuous delivery-- the CD stage--
    the tested training pipeline artifacts are deployed to a target environment.
    In this target environment, the pipeline goes to end-to-end testing
    before being released to production environment.
    Typically pipelines are tested on known production environments
    on a subset of the production data, while the full-scale training
    is performed only in a production environment.
    The newly-deployed training pipeline is then mock tested.
    If the new pipeline fails to produce a deployable model,
    the model serving system can fall back to the mode
    that was previously produced by the current training pipeline.
    Ende des Transkriptes, zum Anfang springen
    Downloads und Transskripte
    Video
    Videodatei herunterladen
    Transskripte
    Download SubRip (.srt) file
    Download Text (.txt) file
    
\section{CI/CD Triggers}

    0:00 / 0:00Klicken Sie auf die Pfeiltaste nach oben, um das Geschwindigkeitsmenu aufzurufen. Regeln Sie dann mit den Hoch und Runter Pfeiltasten die Geschwindigkeit. Klicken Sie Enter, um die jeweilige Geschwindigkeit festzulegen.Geschwindigkeit1.0xKlicke auf diesen Knopf und diess Video stumm oder nicht stumm zu schalten, oder drücke die HOCH oder RUNTER Knöpfe um die Lautstärke zu erhöhen oder abzusenken.Maximal Lautstärke.Videotranskript
    Beginn des Transkriptes, zum Ende springen.
    VIJAY JANAPA REDDI: So Laura just introduced
    us the concept of ML training operationalization where
    the goal is to move from prototyping and experimentation
    from the model development stage to actually pushing
    the code out into the field.
    Now, to do that there are two core parts.
    One is you actually have to do some serious testing
    and the second part is that you actually want
    to automate the whole process of packaging, testing, and deploying
    so that things are repeatable and you result in a reliable training pipeline.
    That's the really key thing about this stage.
    It's about repeatability and it's about reliability.
    In this module, I'm going to peel into each one of these three main stages--
    the continuous integration, the continuous delivery,
    and the production deployment.
    If you aren't familiar with CI/CD terms, don't worry.
    I will explain it as I go through the course,
    and I'll explain why it's actually quite important for your ML project.
    We'll first step through each of these at a high level.
    And then in subsequent videos, I'll actually
    dig deep into each one of these.
    Specifically, my goal is to highlight some of the unique issues that
    come about in the context of a TinyML that don't typically show up in big ML,
    things like cloud machine learning.
    Typically, this is the kind of thing that a software engineer or a DevOps
    team is responsible for in your ML project.
    The software engineer is going to be responsible for packaging up
    the prototype models that come from the ML researchers or the data scientists.
    And the DevOps team is generally responsible
    for the continuous integration, and continuous delivery,
    and production deployment.
    Before I go into the details of CI/CD and PD stuff,
    I want to first and foremost talk about how you actually get here
    from the normal development stage.
    What triggers it?
    As you may recall, your researcher or a data scientist
    is the one who's actually responsible for doing the experimentation
    and prototyping.
    Once she or he is done with the model validation and evaluation,
    they'll push their artifacts into a central repository
    that is shared by everyone.
    Now, these artifacts generally include the training code
    that they have created, the data set that they have used,
    what is the accuracy, and so forth, some amount of metadata about the model.
    Now, in a mature system, this push will basically happen automatically.
    And that'll trigger the CI/CD flow.
    And there's really nothing new here in terms of machine learning specifically.
    It's more so that you have to have this process in place.
    Perhaps the data element is a new element when machine learning comes in.
    You have to have the data set and so forth added in.
    So it's not just your GitHub code, for instance.
    It's the actual data set that you're actually going to be interested in,
    and the model, of course.
    Now, in a typical system, all of this happens
    when the developer pushes the code through something like git push.
    And then automatically, this will trigger the build system,
    which is responsible for the code.
    It will actually build it, and typically it does a Docker build.
    And if you're not familiar with Docker, it's
    a containerization system that makes it really easy to ship code so that it can
    run robustly on any system in the wild.
    I'll talk much more about this later on when
    we talk about scaling the TinyML deployment down in the MLOps pipe
    later on.
    But for now, let's just treat it as some kind of shrink wrap in software
    where you put your software under a special box
    and then it's guaranteed to magically work wherever you deploy it.
    Now, once you have this build, you really want to test it thoroughly.
    You're testing it for any kind of errors that
    show how the model is going to behave.
    Is the model really robust?
    In the traditional sense, as in non-ML systems,
    think of it as like the unit test that you would do.
    For example, you might test for divide by 0, something like that.
    Now, in the context of the ML-based system,
    you might have to do a little bit more, obviously,
    because you're trying to understand how your model is actually
    going to respond to data.
    You'll also want to do integration tests.
    Integration tests are different from unit tests.
    Unit tests tend to focus on isolated testing.
    In integration tests, the testing is really involved around the modules
    that the ML is going to encompass in a way.
    So remember, an ML project is more than just the machine learning model itself.
    There's a lot of code that is non-ML as well that has to work in conjunction
    with the machine learning model.
    So you need to make sure all of those components fit together.
    And that's what the integration tests are all about.
    If all that clears out, you'll then prepare the model into production.
    This isn't as simple as simply turning the light switch on,
    because things can go kaboom if you're not
    careful about how you have done things.
    So there's some very clever strategies that
    are best practices in deploying machine learning
    systems that we'll learn about.
    Yet again, before we dive into this box, I
    want to quickly talk about what the output of this box is,
    as in what are some of the artifacts.
    Laura hinted at this.
    Let me do a quick recap.
    First and foremost is a code of the trained model.
    Most rookies make a very naive mistake here.
    We often tend to think, hey, we've got the model
    and the model is the most important thing.
    Well, the reality, my friend, is that it's not the model that's
    actually really important.
    It's the training code that generates the model.
    You want that model prominence, because you
    need to be able to have a repeatable pipeline here
    that you can test very robustly over and over.
    So you need to ensure that one of the artifacts is the training code itself.
    And all of this has to be packaged up nicely and make sure it runs safely
    and robustly.
    All right, all good up to that point.
    At this point, we've figured out the input and the output
    for the CI/CD pipes.
    Now let's get into the actual box, which involves three main stages--
    the continuous integration, the continuous delivery, and the production
    deployment.
    The goal of all three of these pieces is fundamentally
    to ensure that you have a repeatable pipeline
    that you can safely deploy into your target execution environment.
    That seems straightforward, but there's some gotchas
    when you're trying to specifically tackle
    this in the context of scaling out for TinyML.
    Let's talk about the continuous integration piece first.
    Continuous integration is the practice of automating
    the integration of code changes from multiple contributors
    into a single coherent software project.
    It involves automatically checking the code
    for testing through unit tests and integration tests.
    In the continuous delivery stage, we want
    to get ready to deploy to our target environment.
    And continuous delivery is a process of building a fully packaged and validated
    version of the model.
    And the end result is a package being ready for deployment into production.
    To do this, we have to package up the validated version of the model
    along with the necessary hyperparameters and so forth
    that the network would actually run with.
    So it looks, from everything that I've done,
    so far we're confident that things will work.
    Let's just make that assumption.
    But will they really is a key question.
    Well, that's where you get into the production deployment cycle, which
    is like the separate stage.
    Now, here is where you do this non-exhaustive set of tests
    that ensure that the important functions really
    work before you completely go into online serving
    with the models you have created.
    What's challenging down here is that you would actually have to understand what
    are some of the nuances around CI/CD-related issues,
    specifically in the context of a TinyML.
    So OK, so let's take this CI flow, for instance.
    And we have this notion of we have to understand
    what the build environment looks like.
    Now, recall the heterogeneity that I talked about
    in the introduction module.
    Well, that's going to really come back and bite us right now.
    Because when you talk about continuous integration in the build environment,
    you're not building for just one particular embedded device.
    We're talking about many devices here.
    It's all about scaling tiny amounts.
    So then suddenly, you have to understand the heterogeneity that is there.
    So we'll have to peel the layers of the onion a little bit down here.
    There's a wide array of heterogeneous devices.
    And manually testing them, now that's a real challenge.
    And you can't possibly assume that you're
    going to plug in every single embedded device and test.
    It you need to come up with some sort of creative ways.
    So hold on to this.
    I'll come back and fill in the holes when
    we talk about the continuous integration in much more detail in the next video.
    Now, in the context of actually doing continuous delivery,
    some of the things that we take for granted, which is like we
    have to stage the environment.
    Well, typically staging means that before you push things into production,
    there's a step where you push things into a production stage
    but it's actually not visible to all the end users.
    And its main purpose really is to ensure that any new changes that
    are being deployed from previous environments
    are still working as intended before you hit the end user.
    All right, so here's some food for thought.
    When we're talking about TinyML, we want to put deployments
    into some remote areas.
    Like for example, a forest--
    talk about wildlife conservation, for example.
    We want this device to be operating remotely
    under some really harsh conditions.
    The lighting might be very, very different.
    The sounds and so forth will be very different, all of which
    might affect how your model really performs.
    So the question is, how do you stage that in your testing environment?
    How do you replicate that sort of an environment?
    So you need to think very creatively about this.
    And there are some ways to emulate those environments that I'll talk about next.
    And finally, when you're talking about doing production releases,
    this is where the team is really responsible for signing off
    that everything works well.
    And the QA team specifically writes its own tests,
    which are independent of the developers' tests.
    Because they understand the big picture.
    And they really have to give the sign off that if they turn the switch on,
    the model will work as expected.
    Now, there are some issues here as to why you don't just
    simply push the model into production, why the QA
    team tends to be careful down here.
    Because you can imagine that you've trained the model on some data set.
    Well, that data set is a static data set.
    And these data sets evolve, right?
    The real world is continuously evolving the data set,
    and the real world is continuously evolving.
    So you don't want to just simply throw this model that
    has been trained on some data set, which might, in fact, be six months old,
    for instance, and then put it out in production.
    Because then your model won't perform as well.
    And then you've got a serious issue in the field.
    So you end up having to come up with very clever strategies of,
    how do you stage this deployment?
    You don't want to go from 0 to 1.
    Instead, you want to progressively turn it on and make
    sure it's actually running well.
    There's some textbook-style strategies here that I'll talk about.
    All right, so far I've touched on the big picture of CI/CD and PD.
    In the next few videos, I will dive into the details
    of each one of these individually.
    And my goal is to help explain what specifically makes them
    challenging in the context of TinyML.
    So first, we'll learn about the concepts of CI/CD and PD by themselves.
    And then we'll look at the challenges and some potential solutions.
    So I'll see you in the next few videos.
    Ende des Transkriptes, zum Anfang springen
    Downloads und Transskripte
    Video
    Videodatei herunterladen
    Transskripte
    Download SubRip (.srt) file
    Download Text (.txt) file    

\section{Continuous Integration Challenges for TinyML}


    0:00 / 0:00Klicken Sie auf die Pfeiltaste nach oben, um das Geschwindigkeitsmenu aufzurufen. Regeln Sie dann mit den Hoch und Runter Pfeiltasten die Geschwindigkeit. Klicken Sie Enter, um die jeweilige Geschwindigkeit festzulegen.Geschwindigkeit1.0xKlicke auf diesen Knopf und diess Video stumm oder nicht stumm zu schalten, oder drücke die HOCH oder RUNTER Knöpfe um die Lautstärke zu erhöhen oder abzusenken.Maximal Lautstärke.Videotranskript
    Beginn des Transkriptes, zum Ende springen.
    VIJAY JANAPA REDDI: Welcome back.
    In this video, we're going to pick up on continuous integration.
    Recall that continuous integration is the practice
    of automating the integration of code changes from multiple contributors
    into a single project.
    It's also the process of automatically checking the code through testing.
    Now, when I talk about continuous integration,
    I'm talking about it through the lens of an application perspective,
    not just the machine learning model.
    Keep that in mind.
    Being an MLOps course, you might be thinking,
    oh, this is all going to be about machine learning modeling only.
    No, it's not.
    MLOps is really about taking the model and making
    something effective out of it.
    You want to get some effective use out of it.
    And so that means you have to think about more than just the model.
    Definitely keep that in mind, because that
    will help you understand some of the challenges
    around continuous integration.
    So earlier on, we raised some of these questions
    about what the build environment looks like
    and what are some of the assets we might actually
    want to produce during this continuous integration phase.
    Let me try and think them through, as these are issues
    that we're facing today with TinyML.
    And I'm going to try and post some of these as questions to you,
    so as we're going through, you'll be able to reflect on them along with me.
    Let's start off first by discussing about the building moments.
    What does a typical build environment look like for embedded devices?
    Well, as you may recall, first and foremost,
    we have lots of different kinds of embedded systems.
    Typically, when we are programming a device, we tend to think of one device.
    My go to system often tends to be the Arduino Nano 33
    BLE, because it's a really good platform that
    comes with a rich array of sensors, and it's got a fairly good development
    environment.
    That said, though, these devices all have very different characteristics.
    For instance, I hope you recall that they have different clock rates,
    they have different memory system capabilities,
    they have different sensors, different communication capabilities.
    All of this is sort of a real pain for us
    when we actually want to test things at scale.
    What's not captured here, also, is the fact
    that the underlying software frameworks are completely different
    based on the embedded system you're talking about.
    So for example, writing code for the Arduino
    is totally different than writing code for the ESPI
    at the bottom of the table.
    Or it's very different than writing code for the SparkFun.
    Each of these systems have their own infrastructure
    for actually targeting their particular embedded system.
    Pick your favorite embedded system, and that's
    going to be the exact same problem.
    So a lot of the embedded systems have their own custom hardware
    build environments, and that is really where
    the portability and continuous integration and continuous sort
    of delivery becomes a real nightmare for us.
    Let's take the Arduino device as an example.
    Now, reflect on this for a moment when we
    deal with the device on a daily basis.
    Assume you want to build an image that does keyword spotting or person
    detection or something like that.
    You start off writing the code in TensorFlow,
    TensorFlow Lite Micro, or something that is your favorite framework.
    Then you build a binary, and then you get it running on the device
    by pushing it through the firmware.
    The first challenge is the fact that the toolchain
    that you're using when you go from writing your code in the Arduino device
    to pushing it through is it's a proprietary surfer tool chain that
    only applies to the Arduino device.
    Now, think for a moment that if you have this problem with a traditional desktop
    or server system-- in other words, if I do a build environment,
    I compile a program to run on my Intel processor.
    Will it run pretty much anywhere?
    By and large, I would argue that it's going
    to run anywhere, because environments are very homogeneous in the big ML
    scenario.
    Why do I say that?
    Well, this is because really, the only two operating
    systems that largely dominate the world, you could say it's Windows and Linux.
    OK, fine, so maybe there are a couple of other variants
    and there are differences between Windows and Linux.
    But really, that is almost negligible.
    More importantly, at the lowest level-- at the processor level--
    you have homogeneous instructions at architecture.
    The x86 instruction set architecture has been around since the 80s.
    It's the bedrock foundation for computing to be portable, by and large.
    One could argue that.
    So when you compile code down to run on an x86 processor,
    it doesn't matter whether it's Intel who built a processor
    or whether it's AMD who built the processor
    or it's Centaur who built the x86 processor.
    It doesn't matter.
    It's an x86 processor.
    And that is a beautiful thing, and that's
    the whole reason we've been able to scale out
    computing across multiple domains across the whole planet.
    But in the context of embedded systems, that is most definitely not the case.
    The processors are highly specialized.
    You have processors from ARM, Cadence, Tensilica-- pick your favorite vendor
    out there and you'll find that the processors are highly specialized even
    within families of the processors.
    In other words, even if I take a single family of the processor,
    depending on the version that comes out it'll have a different ISA.
    The reason we have had to do this in the embedded system world
    is because we wanted extreme efficiency, and this
    leads to serious portability problems when you try and scale out
    from doing your CI/CD flows.
    And the second thing here to note is that when you're pushing images,
    like when you're building the code and you're pushing it,
    you often use a physical connection and you push the form right through.
    That's sort of the default approach, by and large, with any embedded system.
    Now, can you imagine trying to do this for a whole fleet of different TinyML
    devices?
    How do you think we're going to be able to push code to all of them?
    And certainly in some rare cases, you can do over the air update,
    but I would argue that this really depends
    on the particular device you have.
    This is not something that is generally applicable across all kinds of devices.
    And also, there are protocol differences in doing over the air updates.
    So it's a very limited thing even though, yes, it is possible to do this.
    So this is yet another one of those places
    where we end up having problems with scaling.
    Now there are some workarounds.
    For example, when uploading things over the air,
    you can potentially come up with an integration architecture like this.
    For example, you can compile a particular build that you want
    and that compiled build can be sitting in some sort of AWS S3
    bucket or something like that, and then you can have a URL associated to it.
    And then you can have the device basically issue
    a request for that particular device, and then you
    can do an over-the-air update.
    We can pull the build.
    And all this kind of stuff is hypothetically possible
    and there's some systems that are out there,
    but nothing that's honestly at the scale of what it needs to be today in order
    to really productionize TinyML.
    And so this is both-- a challenge and a potential opportunity for someone
    who's looking to do something new here.
    Conceptually, we can come up with these things,
    but we don't quite have a lot of robust frameworks.
    There are a couple of systems that I know
    of like Azure Percept, for instance.
    But often, when you're dealing with some of these systems,
    they only support a handful of devices.
    Which I guess is better than nothing, but nonetheless, it's
    not as flexible as a traditional big ML situation.
    Now hopefully this gives you an idea of some
    of the challenges of the build environment.
    So let's move on towards the testing side.
    Generally, when you're thinking about TinyML,
    there are multiple components involved that create a complete system.
    You have the embedded system board, you have different sensors,
    you have peripherals, you have interfaces that bind them together.
    Now, all of these components really have to click together
    in order to enable you to have a smooth CI/CD flow, because that's
    what is a real system.
    All these components click together to be able to do something effective.
    And so TinyML is more than just the ML.
    You have all the embedded components.
    And as a result of that, there's a lot you
    have to actually test when you're doing this continuous integration testing.
    You have to test the code, the model, the data, the sensors.
    And so let's step through each one of these real quickly.
    On the code front, you of course have to remember that your model is
    sitting inside an application.
    And you have to really test that application itself, not just the model.
    Because remember, there are implications about what the application's code
    size is, because you could test a model and everything looks great.
    But when they actually try and deploy it,
    if the application takes up too much memory-- more than the model,
    as in the model plus the application takes up too much memory--
    then you've got a serious problem, because that actually
    won't fit in the flash or in your srun during execution.
    So you have to worry about such things.
    Now, on the model things, things are a little bit more interesting.
    Recall that the model is part of the solution.
    So if your researcher is only focused on the accuracy
    and they skip the evaluation phase that I mentioned in the ML development
    stage, then you've got a situation here.
    So for instance, if you're thinking about KWS, like keyword spotting,
    and then you want to test it, you have to run the model effectively
    through the input and the pre-processing stages.
    You can surely bypass the input stage by saying, hey,
    I'm not going to physically have the microphone.
    Instead, I'm just going to give the audio file directly to the model
    and I'm going to test it that way.
    Yes, that is what most people do.
    But stop and think for a moment about what
    are the pitfalls of that particular approach.
    How will you be able to validate that the deployment scenario is actually
    robust enough?
    Because if you don't measure the timing for the code that
    wraps the model where you're bringing in the audio signal
    and doing the pre-processing, then that's
    going to affect what the end user's experience is going to be,
    because that's what really ends up happening.
    So you've got to think about the fact that you actually
    have to test the whole thing.
    Now, as you're testing the model, the code portion-- for example,
    the neural network-- is a piece of the puzzle.
    So you also have to test the rest of the application pipe.
    So there are two things you're really trying to test,
    and that's where the data becomes especially interesting, because you
    have to pay attention to what kind of data are you
    running through the network.
    And perhaps the hardest part of all of this
    is really going to be testing the sensors, of which you have quite a few.
    How do you test the sensors because that sensor
    is going to be tightly bound to how the model is going to perform,
    because the input for the model is often coming from the sensor itself.
    So if you're talking about one device, surely you
    can deploy it onto your device and test that particular sensor.
    But if the intent is that you're going to be deploying this
    to a fleet of embedded systems, then you've got a problem.
    How do you test them all?
    How do you physically test them all and test all the code
    that's associated with it, too?
    So some of the issues here are that the commodity sensors are often cheap,
    so they tend to have variability amongst them.
    So you can't assume that just because something runs on one device,
    the whole fleet will naturally behave the same.
    So you have to have some sort of statistical confidence.
    This is the sort of thing that embedded machine learning scenarios
    that are specific in the context of TinyML bring about new challenges
    that we have to address.
    A common way to kind of address this problem is to use emulators.
    Emulators are useful to develop, debug, and test multi-node devices
    in a very reliable and then very scalable and efficient
    sort of a manner, because it's an emulator.
    So one such example is a Renode emulator.
    A Renode is a development framework which
    accelerates IoT and embedded systems development
    by letting you simulate the physical hardware components altogether,
    including the CPU, the peripherals, the sensors, and the whole environment
    if you need to, be it in a wired or sort of a wireless setting.
    So you have some flexibility around this.
    But still, even when you talk about these kinds of rich emulators,
    they have limitations in how faithfully they
    can mimic the variations of sensors or how faithfully they
    can mimic the true quality of what a sensor input would look like.
    Nonetheless, it's best effort.
    So I'm trying to just give you a sample here.
    But hopefully, you're starting to get an appreciation for some of the issues
    that you can imagine facing, just because of the fact
    that your ML is being deployed into a physical embedded system.
    Now I want to give you an example of how this is really handled in smartphones.
    The reason I'm coming back to smartphones
    is because they're really just glorified embedded systems.
    A typical smartphone is a traditional embedded system
    with a rich array of sensors, like a light sensor, a proximity sensor,
    accelerometer, a gyroscope, GPS, a magnetometer, and so forth.
    So smartphones come in various shapes and sizes.
    There are loads and loads of different smartphones that are out there,
    and we have had to do CI/CD testing for these things as well.
    It's not like we have managed to skip continuous integration
    testing for mobile devices.
    We have kind of figured a way to do this.
    This has been a real problem.
    And the way this has been done in the past
    is that people have built device farms.
    For example, the Amazon AWS device farm is one particular example.
    The device farm is--
    it's an application testing service.
    It lets you improve the quality of your mobile apps
    by testing them on an extensive range of different, real mobile devices
    without having to you, yourself, physically provision
    or buy and manage all of that testing infrastructure that is needed.
    You can think of it as a remote Android having 700 different phones
    and testing them all remotely.
    And the service enables folks like you and me
    to run all kinds of different unit tests and integration tests
    on physical, real devices that are remotely located.
    Now, you can think of it as a one stop shop
    for testing mobile applications without ever having
    to write your own testing scripts.
    We can directly use one of the supported frameworks that Amazon provides.
    Now, I've used similar things myself at Google when
    I used to work on the Pixel 6 phone.
    We would have an Android farm of about 700 different devices that
    are in a particular location, and I would basically
    be sitting at my computer doing all kinds of tests.
    And they were actually physically running
    on the device, which was really, really cool.
    What's nice about this is that you can actually not only just kick off
    a script that actually does some sort of tests,
    you can actually simulate characteristics
    like gestures, swiping, and all those kinds of behaviors
    all remotely, which is really cool.
    But unfortunately, I don't know of anything
    personally where we have something like this for embedded ecosystems yet.
    Again, might be an opportunity, and that's the point.
    But nonetheless, even though we don't quite have them yet,
    there are prior known solutions that we can adapt from,
    so that's why I'm trying to bring these up here.
    All right, so hopefully you see that it's not
    as simple to do continuous integration and testing.
    Remember, MLOps is more than just a ML model,
    so you have to think about the big picture and the end-to-end flow.
    And debugging and building embedded systems at scale
    is a real challenge that we are yet to overcome.
    With that said, next, let's go into continuous delivery.
    Ende des Transkriptes, zum Anfang springen
    Downloads und Transskripte
    Video
    Videodatei herunterladen
    Transskripte
    Download SubRip (.srt) file
    Download Text (.txt) file


\section{Continuous Integration Tools

\subsection{What is Continuous Integration?}

    To anyone not familiar with DevOps, the term continuous integration sounds a little daunting. In reality, though, continuous integration is a simple but powerful concept. Imagine you are part of a team that is developing an open-source machine learning application that allows users to upload pictures of mushrooms that then tells them if they are poisonous or not. This seemingly simple application has a lot of moving parts to consider: data, models, and a web interface.
    
    \GRAPHICSC{1.0}{1.0}{TinyML/4 MLOps/mushroom.png}    
    
    The web interface for our mushroom identifier application showing the picture of a mushroom and the evaluation – fortunately, it looks like our mushroom is not poisonous!
    
    Most likely, your team will all have to work on different aspects of the same code base, using some form of decentralized source code repository such as GitHub or GitLab. After a hard day of work, you finally get the application up and running. You decide to push the code changes to the group repository so that everyone now has access to this working version, and then go out to celebrate. The next day, you find that another colleague has pushed some additional code to add some extra functionality. The problem? It broke your application! Now nobody can learn whether their mushrooms are poisonous. The real problem here is not really that the application broke, but the fact that nobody seemed to notice. This is where continuous integration comes in to save the day.
    
    Continuous integration is a software development technique of automating the integration of changes to our coding environment. This might entail performing a suite of tests whenever our code is changed, to ensure that everything is still functioning correctly, like our mushroom application, as well as things like revision control, and build automation. This technique is helpful in highlighting problems that may not have been picked up on, and automating how these changes are dealt with - for example, do we decide to reject changes if they fail? Is it okay if certain tests fail as long as the major functionality remains?
    
\subsection{Benefits}

    There are several benefits to continuous integration:
    
    \textbf{Reduced risk.} Defects and bugs are detected faster and more frequently, which means they are less likely to sneak through into a production environment and are dealt with more efficiently.
    
    Improved communication. By working in a more collaborative, automated, and regulated environment, communication is faster and more transparent.
    
    \textbf{Improved product quality.} Because errors are more likely to be found during the development stage, the overall quality of the product will be improved.
    
    \textbf{Reduced waiting time.} When using continuous integration, we do not need to wait for someone to send us code or to perform manual checks which may take up valuable time.
    
    \subsection{Continuous Integration Tools}
    
    \textbf{Renode.} Renode is a development framework which accelerates IoT and embedded systems development by letting you simulate physical hardware systems - including both the CPU, peripherals, sensors, environment and wired or wireless medium between nodes. It lets you run, debug and test unmodified embedded software on your PC - from bare System-on-Chips, through complete devices to multi-node systems. It is specifically useful for TinyML! 
    
    Renode strengths:
    
    \begin{itemize}
      \item full determinism of execution, shared virtual time
      \item transparent \& robust debugging, tracing, analysis, even in multi-node setups
      \item easy integration with your everyday tools, plugins
      \item rich model abstractions with additional functionality "for free" + modular platform description format
      \item automated tests and CI integrations, other collaboration features
    \end{itemize}
    
    \textbf{Jenkins.} Jenkins is a more general purpose continuous integration/continuous delivery (CI/CD) server and is one of the most commonly used CI tools. The architecture works by:
    
    \begin{itemize}
      \item Files are committed to a code repository
      \item Jenkins checks the repository at regular intervals to pull recent code changes
      \item A build server compiles the code into an executable file; if this procedure fails, feedback is sent to developers
      \item If successful, the executable is sent to a test server; if this procedure fails, feedback is sent to developers
      \item If the code runs successfully on the test server, it is deployed on the production server
    \end{itemize}
    
    \textbf{GitHub Actions.} GitHub Actions works similarly to Jenkins, but is built-in to GitHub. This works by detailing tests to run based on certain events, such as when a new branch is merged or new code is pushed to the repository. These details are listed in YAML format in a “.github/workflows” directory. Some examples of GitHub Actions can be found here.
    
   \textbf{TravisCI.} TravisCI is a popular open-source CI tool that can be easily integrated with GitHub repositories. It works very similarly to GitHub Actions, except the YAML-format commands are stored in a “.travis.yml” file. This tool is popular in the Python community for testing that applications run correctly on multiple versions of Python.
    
    Other examples also exist, such as Amazon SageMaker, Google TFX, CircleCI, Bamboo, and Buddy. We will discuss these later in the course as they become more relevant.
    
\section{Continuous Delivery Challenges with TinyML}

    0:00 / 0:00Klicken Sie auf die Pfeiltaste nach oben, um das Geschwindigkeitsmenu aufzurufen. Regeln Sie dann mit den Hoch und Runter Pfeiltasten die Geschwindigkeit. Klicken Sie Enter, um die jeweilige Geschwindigkeit festzulegen.Geschwindigkeit1.0xKlicke auf diesen Knopf und diess Video stumm oder nicht stumm zu schalten, oder drücke die HOCH oder RUNTER Knöpfe um die Lautstärke zu erhöhen oder abzusenken.Maximal Lautstärke.Videotranskript
    Beginn des Transkriptes, zum Ende springen.
    VIJAY JANAPA REDDI: In this video, we're going
    to talk about continuous delivery.
    Recall that continuous delivery is the process
    of building a fully packaged and validated version of the model where
    the end result is a package that's ready for deployment into the real world.
    The key is development in short cycles.
    You want very little manual intervention here so
    that you can actually expedite this whole complex process where reliability
    and robustness is extremely key.
    Now, continuous delivery, ideally, tries to automate the entire software release
    process.
    I just said software, which is where, I think, when it comes to TinyML,
    things get interesting.
    So let's go on.
    Every revision that's typically committed
    triggers an automated flow that builds, tests, and then
    stages things for an update, right, and the final decision
    is to deploy things into a production environment
    where the whole thing can be automated or someone will manually
    sign off before a release happens.
    OK, now, let's talk about continuous delivery
    in the context of web development, because it's a fairly mature space.
    So, this way, we have some sort of a baseline to compare against.
    So this is a website of MLCommons, which I talked about earlier in the model
    development stage.
    Now, let's say the web developer wants to make some changes to the website,
    as shown here.
    Now, simply making this update looks fairly harmless and simple,
    but the reality is that it's not that simple, even in the web context.
    The reason is because we have many different browsers,
    many different browser versions.
    We have devices like laptops, desktops, mobile phones.
    And they all render things a little differently,
    so you do have to test them.
    Therefore, typically, you don't just push the website live
    without doing some thorough checking.
    Normally, developers that are working on the site
    will push things into a staging site.
    This is a protected site that is not broadly
    accessible to everyone in the public.
    In other words, if you go to the main MLCommons URL,
    you won't see the prototype changes that are being made.
    You will, instead, see the changes on a staging site, as shown in the top URL,
    here.
    With that staging URL, what you can then do
    is you can emulate a mobile view of the web page,
    for example, to check how the rendering actually happens
    on a mobile device like an iPhone.
    You can go to the extent of even picking the particular device
    and then having the device kind of show you what it'll look like.
    So this is kind of an emulation-based mechanisms.
    And we have a rich array of different kind of tools and infrastructures
    in place in order to do this kind of testing before we broadly
    scale the web applications out.
    All right.
    At this point, you've done all the testing.
    You got a nice little confident view of the website
    as to what it'll look like on a rich array of devices,
    and then you publish it.
    This looks awesome from a web development standpoint.
    Now, let's talk about this in the context of embedded systems in TinyML.
    As I mentioned earlier, some of the key questions that you
    want to be thinking about are, what does the staging environment look like?
    Think about it for a second.
    These are embedded systems that are deployed in the field.
    So how can you create a staging environment for this?
    OK.
    Next thing is acceptance testing.
    Acceptance testing is when you put the system through all kinds
    of different staging testing scenarios.
    And then you examine how the whole system behaves,
    and then you would effectively sign off on that you've got acceptance testing,
    and then you push things out into production.
    But, if you don't have acceptance testing,
    then how do you actually get confidence the model
    will behave as expected in the field?
    All right, let's take a case study here.
    The example that I'm going to use is a smart power line monitoring system.
    It's actually a real-world case study.
    There's a reading about this that'll come up later on.
    But, here, let me kind of focus on some of the high-level nuggets.
    As I mentioned, one of TinyML's strengths
    is actually industrial applications, so we'll
    be looking at power line monitoring in this particular case study.
    Power line monitoring is important because a simple failure
    in the infrastructure can cause all kinds of various problems
    like power outages or even cause severe damages.
    So in the, quote, unquote, RAM-1 project, the goal for the company
    was to provide remote monitoring, advanced analytics of power grids
    and surging activities in the power grid.
    And the reason for this is because you can do predictive maintenance.
    Like, if you're able to detect any kind of outages or surges
    that are happening, then you might be able to do intelligent load balancing.
    And you might also be able to do very energy-efficient conservation by doing
    voltage reduction and so forth.
    So they're both predictive maintenance opportunities as well as
    efficiency mechanisms that are capable.
    So, therefore, this company went out and built
    this module called the RAM-1, which effectively plugs into the power grid
    and actually monitors the live current that's going through this.
    Now, let's talk about this.
    If we take this example, what does a staging environment look like?
    How do we create a staging environment for an electrical grid measurement
    system which is trying to use a TinyML device?
    Do we test it on synthetic data or previously gathered real-world data?
    But that's very limiting, right?
    It's very hard to get confidence around whether the device will really work
    or not.
    The next thing is run acceptance testing where the key goal, generally,
    is that you want to make sure that everything functionally works
    in this particular application scenario.
    All right?
    For example, if it's a website instead of a power grid,
    right, an acceptance test might actually ensure
    that a user can, in fact, add a product into their cart and then check out.
    You could totally imagine being able to do that on a website testing setup.
    How do you do that here, where you have this device that needs to physically
    be subjected to the high voltages?
    This is a real challenge that they actually faced.
    So, somehow, you wish you could automate these things, but you can't.
    And automation is key because you want these things to move fast.
    Right, the ecosystem is moving fast.
    But you can't, so it's a very manual process.
    Now, because it's a physical gizmo, we can't really
    have fast acceptance testing methods.
    In fact, if we have to do that, the only way we can do it is we
    have to actually go into the deployment stage.
    And only by physically deploying it, by having electricians install
    this device into the power grid, can we actually get
    some sort of acceptance testing verification
    as to whether the device will really perform as we expect or not.
    And this is, in fact, what they actually had
    to do in order to really get confidence, and this is not even the part
    where they're actually putting it into production.
    This is just the part about actually testing the device.
    So the point that I'm trying to land hard, here, on
    is the fact that when it-- because we're dealing with embedded devices,
    it's really hard for us to have smooth flows for accelerated testing, which
    is key.
    So, in an ideal world, continuous delivery
    lets your team sort of automatically build, test, and prepare
    code changes for release into production.
    Now, in a typical software environment where
    you have a machine learning model, sure, you
    might be able to do that based on the things that just kind of talked
    to you about, and your software delivery will be really fast and efficient.
    But because we're dealing with an embedded device, where
    we have physical devices in our loop, it really slows the whole process down.
    Now, I'm not saying that these are problems
    that are completely insurmountable, but these challenges do exist
    and these are some of the challenges that I've
    heard personally about from companies who
    are trying to take TinyML technologies and put them into production in order
    to provide new value.
    And these are things that, if we knew about ahead of time, then
    we'd be able to actively think about how to address them efficiently.
    With that said, I'll see you in the next stage,
    where we'll talk about how to actually deploy TinyML into production stages.
    Ende des Transkriptes, zum Anfang springen
    Downloads und Transskripte
    Video
    Videodatei herunterladen
    Transskripte
    Download SubRip (.srt) file
    Download Text (.txt) file

\section{Production Deployment Challenges with TinyML (Part 1)}

    0:00 / 0:00Klicken Sie auf die Pfeiltaste nach oben, um das Geschwindigkeitsmenu aufzurufen. Regeln Sie dann mit den Hoch und Runter Pfeiltasten die Geschwindigkeit. Klicken Sie Enter, um die jeweilige Geschwindigkeit festzulegen.Geschwindigkeit1.0xKlicke auf diesen Knopf und diess Video stumm oder nicht stumm zu schalten, oder drücke die HOCH oder RUNTER Knöpfe um die Lautstärke zu erhöhen oder abzusenken.Maximal Lautstärke.Videotranskript
    Beginn des Transkriptes, zum Ende springen.
    VIJAY JANAPA REDDI: All right, this is where
    we're going to be talking about production deployment.
    This is the phase where we do a non-exhaustive set of tests
    that ensure that all the important functions are working before we go
    into production release with our model.
    This is really where, for the first time,
    we're talking about the model getting released into the wild.
    And it's the kind of thing that your DevOps team has really
    got to be on top of.
    I've already talked about testing in the context of continuous integration
    and continuous delivery, so you might be wondering,
    what can go wrong at this point?
    Ah, well, my friend, we took a lot of different things for assumption here.
    First and foremost, we assumed that the testing environments
    are almost the same as a production environment.
    I already kind alluded to some of the challenges
    around this in the last video.
    Second, it's hard to make the test environment match the production
    environment-- also kind of alluded towards this with the case study
    for power line monitoring.
    Third, we said it's important to kind of test exhaustively,
    but it's really hard to test everything.
    It's virtually impossible to have every kind of corner case tested, right?
    Now, this is true whether it's a machine learning model
    or it's not a machine learning model, but the reality
    is with machine learning, it's an added level of complexity, if we recall.
    Because it's not only about the code, it's also about the data.
    So it's not only about just making sure the model actually runs properly,
    but it's also making sure the way the model would actually
    run given a set of data is kind of what you want it to do.
    And this is why it's become critical to test in live production environments
    with real traffic.
    [SIGHS] It sounds very scary, and it's going
    to be that much more fun with TinyML.
    There are several pieces to this part.
    So in this particular video, I'm going to talk about the first part,
    and then we'll kind go on to the next part.
    So the first and foremost thing is--
    it's really nothing specific to TinyML.
    This is generally how we would do a production release.
    So I'll talk about it a little broadly and then kind of talk about it
    in more detail with respect to TinyML.
    Specifically, it's going to be around the production deployment strategies.
    So assuming we've done all the testing we can,
    and we finally get the sign off to say, OK, now go ahead with the release.
    Then comes a question of, do you just put the model
    straight into live production in one sudden shot,
    or do you kind of hold back?
    Well, going back to that ML Commons website that I showed you earlier,
    think if you want to suddenly publish a website live,
    what if you haven't caught all the different issues in the field that
    might happen, especially with ML?
    Again, remember, though, we have data, right?
    So the data is going to make things a lot more complicated for us.
    And this is where you need to be careful about how you actually push things out
    into deployment.
    So there's some well-established strategies for this.
    They basically fall into three main buckets,
    and I've chosen these three main things to talk about because I
    think this is what you need to know.
    The first is canary deployment.
    The second is shadow deployment.
    And the third is what we call blue/green deployment.
    Canary deployment is one of the most popular strategies,
    where a small portion of the end users, for example, are exposed to the model.
    So not everybody gets to see the model.
    And if it turns out that there's any kind of problem
    with that small set of users that are actually getting the traffic routed
    towards this new model that you've actually pushed into production,
    then you might be able to quickly look at what's the issue.
    And then you can fix it, or you can either revert it back.
    So you hopefully you get the idea.
    So what you're doing is you keep the old one, whatever is working,
    in practice, if you have an old one.
    And then you got this new model that you want to put in place.
    So when you're pushing the new model in, you
    don't want to have everybody suddenly just start using the new model.
    So maybe in a fleet of TinyML devices, what you really want to do
    is you want to subject only 5% of the inferences queries
    coming through into this.
    It sounds great, but there are issues with it.
    We'll talk about that after we understand
    canary deployment a bit more.
    Some of the pitfalls here you have to worry about
    is that even if you're not directing all the traffic here,
    you have to worry about what region or what sort of users
    are really getting the traffic into this particular embedded fleet, right?
    So if this embedded system, the one that you're
    doing a cannery deployment on with the new model--
    maybe it doesn't see the kind of data that the rest of the fleet
    is actually seeing.
    Then what are you going to do?
    Because now you have this notion of representativeness concern, right?
    The new model might only be seeing a small percentage
    of the actual characteristics, and so it might actually be working well
    for this particular deployment.
    But when you actually spread it across the whole fleet,
    then the model has actually crashed.
    The model actually doesn't perform well because, well,
    the behavior in this particular region is quite different from what
    this model is getting exposed to.
    So that's a serious problem you have to worry about with canary deployment.
    And therefore you have another method called shadow mode
    or what's called dark launch at Google, where
    the technique is you would basically route some of the traffic--
    some of the production traffic through to the actual new one as well as
    the old one.
    But the idea here is that you have a TinyML device, for instance,
    which has two models now because you effectively
    are trying to mimic the second shadow model behavior.
    The new model is called the shadow.
    So you still have the original model that's doing all the things
    that it's doing, and then you're giving it the behavior,
    and you're loading in a new model that's also going
    to be getting the input reference.
    So for example, if I'm a Ring doorbell and I'm
    doing a shadow deployment in my Ring doorbell,
    there are two models that are simultaneously
    looking at the camera input.
    And so one of them is looking and actually performing inferences,
    and maybe it is then sending me a text message.
    But the other one is also looking at the same image.
    It's just that it's actually not acting in the live scenario.
    So it doesn't actually contact me, for instance,
    if that's what is supposed to happen.
    And that's a typical shadow deployment.
    Broadly speaking, there are two key reasons you would do a shadow mode.
    Because one is it ensures that the model is actually handling inputs the way
    you intended for it to do, and then you can ensure that your system is actually
    working as a whole.
    So it's a really nice thing.
    Beautiful thing about this is you're not affecting anything
    on the primary device, right?
    Because the primary model is still doing what it's supposed to be doing,
    and you get all the benefits of checking the stability and so forth.
    And it's very nice because when you have a shadow deployment, what you can do
    is you can actually analyze how the model is really behaving.
    So you can collect the inferences that it's getting
    and the inference decisions that it's predicting.
    And then you can compare it to what the old model is doing to really see, hey,
    are you actually performing better, right?
    And so you can do a lot of quality analysis that way.
    But there's some problems down here when we
    talk about this in the context of embedded machine
    learning because now the cost goes up.
    Do you have the memory to simultaneously run two models?
    Do you have the compute bandwidth?
    Or are you going to slow the main model down
    because you are running on microcontrollers that can't keep up?
    Do you have the energy to spare, especially if it's based on a battery?
    So these are some of the complications that you really have to worry about.
    And also, remember, here when I said you do a shadow launch,
    as in you push the new model in with the old model, well, that's
    just kind of assuming a traditional computing system
    where you have network connectivity and all those kinds of good things.
    But that's not guaranteed when you're talking about over-the-air updates
    with these embedded systems, as I've hinted many, many times already.
    So honestly, I'm not even sure we can actually do a shadow of deployment
    in practice because I have not come across any companies that
    are trying to do this sort of a thing, especially in the context of TinyML.
    But nonetheless, it's something to keep in mind.
    Now, the third and final approach is what we call blue/green deployment.
    Now, as the slide here clearly says, in its purest form,
    blue/green asks us basically to duplicate every resource in the system,
    in the application, as a separate thing.
    And then what you do is-- in practice, you can't always duplicate everything.
    For example, if it's a website and it has to access a database,
    the database can't also be duplicated because there is just
    complications with trying to just duplicate all the data because then you
    have to synchronize the two databases, and that causes
    all kinds of havoc in the background.
    So you can't always do duplication.
    Nonetheless, what this means is that if you have a blue/green deployment, what
    you can do is you can have some of the traffic
    be routed over to the blue and some of the traffic
    be routed over to the green.
    And so you need some sort of a load balance or a router
    in the middle of that kind of manages all of this.
    And then what it allows you to do is to effectively
    be able to test how the system is actually working.
    So the blue is the original system, and the green
    is a new system to which you actually want to move to.
    This is what the blue/green deployment does.
    Blue is nice and stable.
    Green is the one where you're temporarily routing traffic to.
    So for example, if it's a website, you can
    imagine you have this new website, the staging
    website, that you push into production.
    You can watch its traffic, right?
    And so that's one way of doing it.
    And then slowly, if everything looks nice and clear
    on the green one, the new one, then you would push all the traffic there.
    And when that is the case, then you would effectively
    route all traffic over the green one.
    Nice thing about this again is that there's no downtime.
    The system is continuously alive.
    You get to do a lot of testing effectively because you are kind
    of looking at the real world traffic.
    And so you get a lot of confidence in your deployment.
    But again, in the context of TinyML, you've
    got to ask yourself this question, what if like
    you can't make continuous updates?
    In other words, there might not be able to quickly route the traffic from one
    model to another model because, again, all of this
    is assuming connectivity and stability in the ecosystem
    for you to be able to quickly do these things.
    What if you can't support two identical sort of environments that are running?
    Are we saying here that now we need two, let's say,
    embedded TinyML boxes that are sitting next to each other that
    are looking at the same image or processing the same sensor
    or inputs that are coming in?
    What does this really mean?
    So this is where you start realizing that some of the traditional deployment
    practices that we tend to have often become
    challenging in the context of a TinyML.
    And also, there's a situation where perhaps the infrastructure simply
    does not allow us to have some way of routing this sort of traffic, right?
    You need mechanisms in place.
    Now, if it's just a web server, obviously, these kinds of things
    are easy to do because you have all the network infrastructure in place.
    So I hope you have two key takeaways from this video.
    First and foremost is how you go into production deployment.
    In other words, you don't want to flip the switch
    and then just toss in your new model and have it handle everything
    because it's a potential risk that things might break catastrophically.
    Following that, The second thing is that you have
    learned about three different methods--
    that canary deployment, the shadow deployment,
    and the blue/green deployment.
    Each one has some pros and cons, and keeping those in mind
    is going to be key because those are bread and butter
    sort of design principles that are well exercised in practice.
    However, as it comes to TinyML, I hope you
    realize that those deployments don't necessarily translate immediately well.
    So even though you might have experience having doing those kinds of things,
    or now that about those things you might want to do them,
    I want you to keep in mind, that in the context of TinyML,
    it might not be very easy to do.
    And this is where your DevOps and your software engineering team
    needs to be aware of these kinds of things-- and limitations
    and challenges-- so they know how to go into production deployment well.
    Next, I'm going to talk a little bit about online experimentation
    in terms of how do you sample from the space of different users
    that you might have access to when you actually want to do online testing.
    I'll see you then.
    Ende des Transkriptes, zum Anfang springen
    Downloads und Transskripte
    Video
    Videodatei herunterladen
    Transskripte
    Download SubRip (.srt) file
    Download Text (.txt) file

\section{Online Experimentation}

    0:00 / 0:00Klicken Sie auf die Pfeiltaste nach oben, um das Geschwindigkeitsmenu aufzurufen. Regeln Sie dann mit den Hoch und Runter Pfeiltasten die Geschwindigkeit. Klicken Sie Enter, um die jeweilige Geschwindigkeit festzulegen.Geschwindigkeit1.0xKlicke auf diesen Knopf und diess Video stumm oder nicht stumm zu schalten, oder drücke die HOCH oder RUNTER Knöpfe um die Lautstärke zu erhöhen oder abzusenken.Maximal Lautstärke.Videotranskript
    Beginn des Transkriptes, zum Ende springen.
    VIJAY JANAPA REDDI: Welcome back to part two of production deployment,
    which is all about online experimentation where
    we are making the fundamental assumption that we have somehow managed
    to push our model out into production.
    And the question now is, what can we learn from online experimentation
    of the model?
    Now, the two parts to it which I'll talk about, one
    is smoke testing and the other one is called A/B testing.
    Smoke testing is a term that actually came out of the plumbing profession,
    where the intent was to try and see if there were any leaks in pipes.
    And the way you would do that is by basically pumping in smoke in one end,
    smoke comes out of the other end, everything was good.
    But if there were any leaks, you'd be able to spot the smoke coming out
    elsewhere.
    My other favorite version of smoke testing
    is really when you do hardware testing, when you turn the switch on.
    And if it goes kaboom, then, well, and smoke starts coming off,
    then you've got an issue.
    But not really very useful in practice, of course.
    Anyway on a serious note, smoke testing is a software testing technique
    that is used after a program has been built
    to ensure that the critical functions of the product
    are all operational when it's actually live.
    Now, we've talked about this briefly in the context of ML Commons website.
    Remember, how we talked about we want to make updates to the website,
    then we want to understand whether the website actually renders well
    on a mobile device, for instance.
    And we can use staging mechanisms, and so forth, and test things out.
    But when it comes to smoke testing the whole intent
    is to actually do some of this testing in a live scenario, where we're not
    looking at an emulation of the web page.
    We've already published the website to be live,
    and end users are actually using it on the devices.
    And we're going to get data back from that.
    You've got to think about if this was a context for an embedded system
    deployment, and you're trying to do smoke testing, where you're
    doing online experimentation, I want you to stop for a moment
    and think about what does this really mean for TinyML.
    I've already kind of talked about a lot of things
    earlier on, which are with respect to the code, the model, the data,
    the sensors, and so forth.
    So I want you to kind of stop for a moment
    and reflect about what are some of the challenges
    that you would have if you want to do smoke testing?
    And what are some of the benefits that you would have
    if you were able to do smoke testing?
    Now, I don't want to repeat a lot of the things again,
    but nonetheless, I'll leave room for you to discuss this in the forum
    with your colleagues.
    The next thing is run A/B testing.
    Now A/B testing is really simple.
    It's a widely used online experiment, where
    you have two versions of a live website, for instance, version A and version B.
    And what you really want to do is you want to try and root users
    to both versions of the websites, and then
    you want to be able to understand what the user's experience is, and maybe
    then rate which version of the website is actually performing better,
    and then use that as a main website.
    Now, again, if you think about this in the context of an embedded deployment,
    do you really think we'd be able to do this?
    Think about this, you have a fleet of embedded systems
    that are out there all running different kinds of models,
    and you want to do A/B testing.
    What are some of the issues that you think you're going to get into
    are what I want you to reflect in the upcoming forum.
    Next, there's a version called MAB testing, which is a Multi-Arm Banded
    testing, which is effectively a souped up version of the core A/B testing
    principles.
    And I'm not going to get into it, because the underlying principles are
    largely the same.
    It's just about doing A/B testing much more intelligently.
    I'll encourage you to kind of read about it,
    but I won't personally be getting into the details of it as much.
    Fundamentally, when you're doing online experimentation,
    by and large, what I want you to think about
    as you discuss with your colleagues is what are some
    of the challenges with connectivity.
    What are some of the latency implications of these kinds of things?
    You have very finite resources in terms of memory and battery lifetime.
    So if you were to do some of this online experimentation,
    be it smoke testing or be it the A/B testing,
    what are some challenges that you're going to face?
    And how would you try and address them if you still
    wanted to get the data back in order to learn from how the model's actually
    performing in the world?
    That's what I want you to really reflect on.
    So bottom line, what I'm trying to say, though,
    is that be it continuous integration, or continuous delivery,
    or production deployment, you've got to be very careful about how
    you're doing releases in some capacity.
    You don't want to just flip the switch and toss these models overboard,
    because they can be serious mechanisms in which they will break,
    and that can cause real havoc for your use case.
    So you want to be coming up with strategic means of actually doing
    product releases, even in the context of embedded machine learning,
    despite the fact that we have a whole set of challenges.
    Ende des Transkriptes, zum Anfang springen
    Downloads und Transskripte
    Video
    Videodatei herunterladen
    Transskripte
    Download SubRip (.srt) file
    Download Text (.txt) file
    
    
\section{Production Deployment in ML Deployment}

\subsection{About}

    We have previously discussed deployment in the context of ML Training as it was also relevant to the issues of testing. More often than not it is important to test in production as it is impossible to know if you have covered all your bases despite having integration, smoke, as well as acceptance testing and so forth. So here we will do a quick recap on the deployment methods.
    
\subsection{Deployment Methods}

    First, we have the basic deployment. This is the riskiest technique, and it involves simply replacing our current model (if one exists) with our new model. If everything works fine, then everything runs smoothly and there are no issues. This type of deployment is also very fast since there is little to no delay in predicting serving. However, if there are any issues, the lack of a failover mechanism can cause problems. The issues may be time-consuming to fix and even require the entire production environment to undergo downtime.
    
    Next is the multi-service deployment. In this type of deployment, we do not “overwrite” a saved model or application, but merely add an updated version of it. This way, if there is an error in the production environment, we can always failover to an earlier version. This method is safer than the basic deployment but still suffers from issues in terms of maintenance since extensive testing might be necessary to make sure any infrastructure changes are compatible with all current and previous versions.
    
    Rolling deployment offers some new advantages. With this method, all of our deployment environments will be updated with our new model or application, but this update will be done in batches. The advantage of this method is that only a few deployed devices are updated at a time making it easier to roll back a release and making it safer than a basic deployment as only a few users at a time will be affected by the update.
    
    A similar approach is a canary deployment. This allows us to put our model into production, but only targets a subset of users which may be selected randomly or through some form of filter (e.g., geographically). This allows us to test our model in a true production environment and to compare the real-time performance to the existing deployment. This deployment strategy is essentially an implementation of A/B testing.
    
\GRAPHICSC{1.0}{1.0}{TinyML/4 MLOps/5-4-10 PD Part 1.jpg}    
    
    A graphical depection of a canary deployment showing that most users are given the old version of a system while a few users are given a new version.
    
    Another popular strategy is blue-green deployment. This is a common deployment strategy and involves two identical environments, one called the “staging” environment, and the second called the “production” environment. Within the staging environment, user testing and quality assurance can be performed before our model or application is deployed to the production environment. Once we are satisfied with our model’s performance in the staging area, the model is then deployed into the production environment. This method is more expensive than rolling or canary deployment since it requires both a production and staging environment which each may be extremely large. 
    
\GRAPHICSC{1.0}{1.0}{TinyML/4 MLOps/5-4-10 PD Part 1 (2).jpg}    
    
    A graphical depiction of a blue-green deployment showing that users are split into a blue and green environment. Each environment has a seperate set of servers and services but both environments share a database to analyze and compare the results.
    
    The final deployment strategy we will discuss is shadow deployment. This involves having both our current and tentative deployments running simultaneously in the production environment. Real-time traffic is sent to the “shadow” deployment (which represents our tentative model or application), but predictions are not sent to the user. Instead, they are monitored to test for performance and stability. Once some threshold values are reached, the model can then be solely deployed in the production environment. This procedure is complex to set up compared to several of the other methods.
    
\GRAPHICSC{1.0}{1.0}{TinyML/4 MLOps/5-4-10 PD Part 1 (1).jpg}    
    
    A graphical depiction of a shadow deployment showing that the calling code triggers both a live API that uses the live model as well as a shadow API that uses a shadow model.
    
    As you can now see, there are many ways for us to introduce a new model into our production environment. Some of these are easy to implement, but relatively risky, while others are complex but relatively safe. The most suitable deployment strategy depends strongly on the application, but also things like cost, existing infrastructure, and technical expertise available. We just talked about production deployment from the lens of all the things that can go wrong even after we have done numerous amounts of testing. Therefore, we want to have some sort of a staged deployment to ensure that deploying our models into production won’t cause things to crash. This is why we learned about various rolling deployment strategies such as the canary, shadow, and blue-green deployment strategies.     
 
 \section{Case Study Discussion}   
 
Now that you’ve explored CI/CD, it's time to explore a case study about deploying CI/CD in the real world. In the document below you will explore a case study of using tinyML to monitor power lines with the RAM-1 device. After you finish reading the case study  please reflect on the following prompts in the discussion forum below:

\begin{itemize}
  \item What particular CI/CD challenges would you expect to have using the device? 
  \item What features does the device have that would minimize those issues? 
\end{itemize}

\section{Training Operationalization Impact on MLOps}

    0:00 / 0:00Klicken Sie auf die Pfeiltaste nach oben, um das Geschwindigkeitsmenu aufzurufen. Regeln Sie dann mit den Hoch und Runter Pfeiltasten die Geschwindigkeit. Klicken Sie Enter, um die jeweilige Geschwindigkeit festzulegen.Geschwindigkeit1.0xKlicke auf diesen Knopf und diess Video stumm oder nicht stumm zu schalten, oder drücke die HOCH oder RUNTER Knöpfe um die Lautstärke zu erhöhen oder abzusenken.Maximal Lautstärke.Videotranskript
    Beginn des Transkriptes, zum Ende springen.
    SPEAKER: Congratulations on almost completing the second stage of MLOps.
    We've covered a lot of ground around continuous integration,
    continuous delivery, and production deployment.
    Hopefully you've gotten two key takeaways at this point.
    Generally, you should know what are the best
    practices around testing and integration for machine learning systems.
    And second, hopefully you understood some
    of the unique challenges that emerge in the context of embedded machine
    learning as you try and test them and as you try and deploy them.
    So yet again, we have to ask ourselves this meta question, what
    is the impact of ML training on the other stages in the ML development
    pipeline?
    For one, you can clearly see that it's actually quite important
    to go from experimentation and prototyping
    to actually having a systematized method for robustly testing
    and robustly deploying machine learning models.
    This is extremely crucial, especially because researchers such as myself
    often tend to just kind of focus on getting the solution with the best
    accuracy, latency, and so forth.
    But we don't necessarily think about it in terms of how
    do we robustly test this continuously?
    How to be version this code?
    How do we version does data?
    And how do we actually progressively deploy this into the field?
    So hopefully, some of the things that we've talked about here
    give you an appreciation for why this is in fact
    a separate stage away from the machine learning development stage.
    And why the training operationalization as a whole bunch of goop
    that needs to be appreciated all by itself.
    And also as you will see later, training is not a one shot thing.
    We'll be re-triggering this training pipeline over and over.
    So it's actually important, therefore, to have a repeatable pipeline where
    everything is streamlined for continuous integration, continuous delivery.
    And it's extremely important because you don't
    want to waste time and effort doing a poor job in the operationalization
    phase.
    Because if you do that, then you might not
    be able to quickly fix a machine learning model that has gone rogue
    in the field, or you might not be able to quickly pull it
    back and be able to replace it with something that's old, that's mature.
    So that's another example where continuous training, which
    is downstream, is actually going to rely on the fact
    that you have a good operationalization pipeline built up.
    Last but not least, I just want to give you one more example.
    Prediction serving is when you are deploying
    your machine learning solution as a service,
    like machine learning as a service.
    When you do this, you really have to worry about heterogeneous deployment
    situations because you're talking about typically
    deploying into a fleet of different embedded devices.
    Now, some of these things have already crept up
    when we're talking about continuous integration for testing, building,
    and so forth.
    So hopefully, you will have an appreciation when
    you go down for prediction serving.
    Like, oh well I've already kind of tackle
    some of these difficult problems in my operationalization phase.
    So now when I come to actually deploying it at scale,
    I can actually be deploying it with some level of confidence.
    Last but not least, I hope you see how software engineers will
    be critical for packaging up the machine learning model into the application,
    as well as doing all the necessary things that are needed for testing
    the machine learning model robustly.
    And I also hope you see how the DevOps team becomes
    really essential in carefully deploying the models into production.
    Remember, it's not like you flip the switch and everything just goes online.
    You need a team of people who are dedicated
    for actually carefully testing and deploying the model in phases
    into production.
    So that everything goes on smoothly, and if there are any gotchas
    they're able to quickly retract back and keep your business intact.
    With that said, we'll wrap up this particular stage
    and move on to continuous training.
    I'll see you next.
    Ende des Transkriptes, zum Anfang springen
    Downloads und Transskripte
    Video
    Videodatei herunterladen
    Transskripte
    Download SubRip (.srt) file
    Download Text (.txt) file
    
    
% 1.5: Continuous Training    
  
\section{Overview of Continuous Training}
\subsection{Module Objectives}

\GRAPHICSC{1.0}{1.0}{TinyML/4 MLOps/continues_Training.png}    
  
  The MLOps pipeline. Going from ML Development to ML training to Continuous Training to Model Deployment to Prediction Serving to Continuous Monitoring. Along the way artifacts are produced at each step including: Code and Configurations, Training Pipelines, Registered Models, Serving Packages, and Serving Logs. All steps tie into the overarching themes of Data and Model Management.
  
\subsection{Objective}

  In most respects, our world is in a constant state of flux. These constant changes transfer over to our data, which we can consider a simple “snapshot” of reality that our machine learning algorithms use to build models that describe the physical world. This presents us with a problem: if the world, and hence our data, is always changing, won’t our models always be wrong? In some respects, the answer to our question is yes. To that end, we will first start with a broad overview of Continuous Training and then dive into the specifics of each of these:
  
  \begin{itemize}
    \item  Data Processing Engine
      \begin{itemize}
        \item Data ingestion
        \item Data validation
        \item Data transformation
        \item Keyword spotting datasets
      \end{itemize}
  \item   Model Training Engine
    \begin{itemize}
      \item AutoML
      \item Neural Architecture Search
      item Hyperparameter tuning
      \item Why these are important and how they can be helpful
    \end{itemize}
  \item Evaluation
    \begin{itemize}
      \item Use-case driven metrics
      \item Continuous Training pipeline metrics
    \end{itemize}
  \end{itemize}
  
  We learn these concepts through the lens of Keyword Spotting as a TinyML use case.
  
\subsection{Motivation}

  Our models will likely never perfectly reflect reality - just think, when was the last time you trained a model that was 100\% accurate? However, for most applications, our models do not need to be perfect, just consistently accurate. Accurate implies that our model performs well for a given task; for example, we might decide that a malware detection algorithm with 90\% accuracy is sufficiently accurate to implement as part of an intrusion detection system. Consistent implies that our model performance does not change significantly over time. In the case of malware detection, it is clear that malware changes quite rapidly in response to different software vulnerabilities that are discovered (see the CVE List for examples). If our algorithm’s accuracy falls to 30\% after three months of being in use, this is not great, and you will probably end up with all manner of software viruses unless you make sure to periodically update the model using more recent data. So far, we have motivated the need for periodically updating our models, which we call retraining. While this will help to resolve our consistency problem, it presents us with several more questions. How often should I retrain my model? How can I tell if my model is underperforming? How can I translate this into monetary costs?
  
  \subsection{Continuous Training}
  
  The answers to these questions are highly context-dependent. For example, a model trying to model customer preferences will change significantly over time in response to various stimuli, including socioeconomic and political factors. These changes may be chronic, occurring over a long time frame such as years or decades, or immediate, occurring perhaps on the order of just minutes or hours. Clearly, these situations would have to be approached radically differently. In contrast, modeling defects within a manufacturing process might only change if the process is altered. Hence, retraining this model over time would largely be unnecessary (other than to improve the model’s accuracy by training on a larger set of data). Therefore, the frequency of retraining is largely data-dependent. In the case of modeling consumer preferences, a more advanced form of retraining might be implemented, called continuous training. 
  
  Continuous training, as the name suggests, refers to the procedure of retraining our model on a continuous basis. This might be achieved in different ways. Training can be scheduled, as in the case of periodic retraining, using a cron-like scheduler. This might tell a system to retrain a model using the most recent data once every night, similar to the nightly build of TensorFlow. Alternatively, training might be event-driven, by what we call retraining triggers. These triggers might be a metric that is monitored, and once the model has fallen below a specific threshold, the model is retrained, or when a new threshold amount of data has been collected. The metrics or thresholds selected for this purpose are, again, context-dependent. For example, a binary classification algorithm might care more about monitoring the F1 score or true positive rate as opposed to raw accuracy. Training might also be manually triggered, such as in the case of a new system update.
  
  Retraining models is very important for some applications, and less so for others. Taking COVID-19 as an example, a large number of commercial models in the finance, e-commerce, and business sectors failed during this period due to rigid models that were unable to accommodate rapid and profound changes in socioeconomic conditions. Machine learning models are great when they work, but can cause a lot of headaches when their assumptions are violated! In this section, we will delve deeper into the concept of continuous training and how it can be added to our arsenal of MLOps tools.
  
  \subsection{Module Structure}
  In this section, we will first start with a broad overview of Continuous Training and then dive into the specifics of each of these:
  
  \begin{description}
    \item[Data Processing Engine:] Importance of building robust data engineering pipelines so that we can easily keep improving the quality of our datasets
    \item [Model Training Engine:] Benefits of AutoML, NAS and other automated methods for easing Continuous Training
    \item [Evaluation:] How to consider use-case driven metrics to ensure we are improving models systematically over time driven by the end user value as well as metrics that matter in terms of the efficiency of the Continuous Training pipeline itself      
  \end{description}
        
\section{Continuous Training}


    0:00 / 0:00Klicken Sie auf die Pfeiltaste nach oben, um das Geschwindigkeitsmenu aufzurufen. Regeln Sie dann mit den Hoch und Runter Pfeiltasten die Geschwindigkeit. Klicken Sie Enter, um die jeweilige Geschwindigkeit festzulegen.Geschwindigkeit1.0xKlicke auf diesen Knopf und diess Video stumm oder nicht stumm zu schalten, oder drücke die HOCH oder RUNTER Knöpfe um die Lautstärke zu erhöhen oder abzusenken.Maximal Lautstärke.Videotranskript
    Beginn des Transkriptes, zum Ende springen.
    LARA SUZUKI: Hi.
    In your last lesson, we had a closer look at the training operationalization
    stage of our MLOps framework.
    It's now time to move to the next stage that is of continuous training.
    In this lesson, we will cover deeper details
    of creating a responsive training pipeline through continuous training.
    Continuous training refers to the repeated execution
    of the training pipeline in response to data changes, code changes,
    or on a schedule, potentially with new training settings.
    It utilizes a trained and validated model is stored in the model registry.
    Training metadata and artifacts stored in the ML metadata and artifacts
    repository includes pipeline execution parameters, data statistics,
    data validation results, transformed data files, evaluation metrics,
    model validation results, and training checkpoints and logs.
    The core MLOps capabilities used in this stage
    are dataset and feature repository, ML metadata and artifact repository,
    data processing, model training, model evaluation, ML pipelines,
    and model registry.
    The MLOps stage continuous training and to orchestrate and automate
    the execution of training pipelines.
    It mirrors the step of a typical data science process.
    The data validation and model validation plays a critical role here.
    They are the gatekeepers of our model.
    Any deviation from what is expected might result in a new pipeline
    to not be deployed.
    The continuous training takes into account
    the business value versus the cost of retraining models.
    It can be considered to run on demand or on a schedule,
    on availability of new training data or new training code,
    or on model performance degradation.
    The continuous training involves the typical machine
    learning training pipeline workflow such as data ingestion, validation
    and transformation, model training and tuning, evaluation, validation,
    and registration.
    The ML pipeline can register a model candidate in the model registry.
    Runtime information in the artifacts are tracked in the metadata and artifact
    repositories, it enabled debugging, reproducibility, and lineage analysis.
    It also enables retrieving the hyper-parameters, evaluations, data
    after transformation, data summaries, a schema, and also feature distributions.
    Ende des Transkriptes, zum Anfang springen
    Downloads und Transskripte
    Video
    Videodatei herunterladen
    Transskripte
    Download SubRip (.srt) file
    Download Text (.txt) file        
  
\section{Retraining Triggers}

    0:00 / 0:00Klicken Sie auf die Pfeiltaste nach oben, um das Geschwindigkeitsmenu aufzurufen. Regeln Sie dann mit den Hoch und Runter Pfeiltasten die Geschwindigkeit. Klicken Sie Enter, um die jeweilige Geschwindigkeit festzulegen.Geschwindigkeit1.0xKlicke auf diesen Knopf und diess Video stumm oder nicht stumm zu schalten, oder drücke die HOCH oder RUNTER Knöpfe um die Lautstärke zu erhöhen oder abzusenken.Maximal Lautstärke.Videotranskript
    Beginn des Transkriptes, zum Ende springen.
    VIJAY JANAPA REDDI: We at large say that continuous training orchestrates
    and automates the execution of the training pipeline, which
    we put together in the training operationalization stage of the MROps
    pipe.
    Now, in the context of continuous training,
    there are lots of moving parts here.
    We'll be specifically focusing on this main section,
    and I'll allude to some of the peripheral components as needed.
    More specifically, we will be focusing on these three core engines--
    the data processing engine, the training engine, and the evaluation engine.
    We'll be talking about all three components using
    the keyword "spotting example" as a case study,
    because that'll help ground these concepts into the real world.
    Now as usual, multiple folks are going to be involved.
    The continuous training pipe is rather a big piece.
    So there are going to be many different kinds of individuals
    in the organization playing a role.
    For example, you're going to have the data
    engineer who's typically responsible for the data processing engine.
    The ML researcher will be responsible for the model training engine,
    especially in the lens that we're going to look at it at, which is the auto ML,
    or auto machine learning.
    And the ML engineer is generally going to be responsible for maintaining
    this continuous training pipe.
    Once again, remember this is just a rough guideline.
    Depending on the size and velocity of your company,
    the characteristics of who's involved might change.
    All right, before we dive into the details of the continuous training
    pipeline, let's start by understanding why we need continuous training.
    In other words, why is it important to retrain models from time to time?
    Well, it all starts with drift--
    model drift, to be more specific.
    Model drift, which is also sometimes known as model decay,
    refers to the fundamental point that your model's prediction capabilities
    will degrade over time.
    Now, this can happen due to a variety of different reasons.
    It might be a change in the environment, change
    in the relationship between the different variables
    that you have to trained your model on, and so forth.
    In this video, I'm going to touch primarily on the notion of drift
    lightly, but I promise you I'll talk a lot more about it
    as we get into the continuous monitoring stage, which is where drift
    is extremely critical to watch out for.
    That's going to come way at the end of the section.
    So let's hang on and let's continue.
    Broadly, model drift can be broken down into concept drift and data drift.
    Concept drift is when the statistical features of the target variables,
    which the model is attempting to predict,
    vary in unexpected ways over time.
    Now this presents a major issue, because naturally, your model
    is going to be predicting poorly, which you don't want to happen.
    Let's take, for example, what happened with the COVID-19 era.
    As that virus has come about, the way we have been behaving
    is completely changed, which means the way models might have been predicting
    in the past is not what is representative of today's real world
    characteristics.
    For example, we end up washing our hands much more.
    We end up engaging in much more social distancing.
    We tend to pray more.
    Much to my mom's disappointment, I have not gotten any better at this.
    Nonetheless, I appreciate the fact that we are praying.
    I believe in faith.
    Getting back on point, we tend to avoid restaurants.
    We have been hoarding a lot of food and water, just to shelter ourselves,
    and more importantly, we've been wearing masks,
    which we have not done in the past.
    So imagine that.
    So all these are things that actually affect
    how your model is going to perform in the real world,
    because these are things that have only started happening after the virus has
    come about.
    Previously, we didn't have any of this.
    Also, this issue of data drift.
    Now, when you train a model in the case of supervised learning,
    the training data is usually label data.
    So when you deploy the model into production, there is actually no label,
    so no matter how accurate your model is in the testing phases,
    it's only valuable and right when it's actually
    working well with the data that's actually in the field.
    And what if it doesn't, right?
    And that's when you call it data drift.
    Data drift is fundamentally defined as a variation in the production data
    from the data that was used to test and validate the model before deploying it
    into the real world.
    Now, as an example of natural data drift, is seasons.
    That's why I'm showing you that picture here.
    You can see it's the exact same tree, but as the seasons change, as
    the temperature changes, the environment changes, the look and feel
    of that particular tree is remarkably different.
    All while it's still exactly in the same place.
    So that's an example of natural data drift.
    From a technical standpoint, you tend to have data skews or data
    changes in the form of either schema skew or a value skew.
    So at a high level, schema skews are when
    there are anomalous input data, which means the downstream processing
    steps may fail effectively as it receives this data,
    because it is not something that it's actually expecting.
    When the schema skew basically happens, the pipeline execution must be stopped,
    and then you've got to basically go back and look at the schema of the data--
    the structure of the data that you're ingesting.
    Now, examples of data schema skew will include things like unexpected schema.
    Like in other words, the structure of the data fundamentally looks different.
    Are not receiving all of the schema and receiving only features
    with unexpected values.
    I'll give you more specific examples in the context of the TinyML use cases
    next.
    All right, so that's in the schema itself.
    Now there's a value skew, right?
    This is when the value of the features changes significantly
    that the statistical properties of the data has fundamentally changed.
    And this indicates that the model needs to be retrained,
    in order to capture the change in the statistical properties in the data.
    All right, so basically the two big data skews-- one is the data schema skew
    and one is the data value skew.
    All right, let's talk about this a bit more specifically.
    In the TinyML examples, like often I say that, well, manufacturing
    is a big use case for TinyML, right?
    All right, well, in manufacturing you're often
    dealing with a lot of sensor data.
    So imagine if a sensor in the manufacturing plant is replaced
    and it generates values differently.
    For example, this sensor, instead of reporting data in centimeters,
    is now starting to report data in meters.
    And that's a fundamental structural change
    inside how the data is actually getting reported,
    and that can basically break the model, because you're not
    ready to handle data that's coming in meters.
    Instead, you were trained on centimeter sort of behavior.
    Similarly, it's also possible that, for example, your sensor physically
    breaks, which is something I alluded to in the previous section.
    When you're trying to deploy these devices into the real world,
    because they're physical systems, the sensor itself might physically break.
    In which case you might be getting null data or just bogus data coming through.
    And so you've got to catch that issue and then address it.
    And as I mentioned earlier, you might have covariate shifts,
    which is basically the change in the relationship between different features
    that you would actually train your model on.
    All right, so later on, when we get into the continuous monitoring pipe--
    which is way down--
    you will see that when we build pipelines,
    we'll build them in such a way that we can actually
    monitor the effectiveness and efficiency of a deployed model in production.
    We'll rely on generally what I call data set monitors,
    in order to tell us if any of the skews or things are breaking in the field.
    These data set monitors do a couple of things in the continuous monitoring
    stage.
    They will detect and alert us when some data drift
    or skew's happening, when the concept is kind of seeming to drift.
    They'll analyze historical data and then compare it
    with what's actually happening, in order to tell us
    if there is actually a drift happening.
    So it's not just kind of looking at only the raw, real time data.
    It's actually comparing it to historical data.
    And it will actually collect and profile a lot of information
    and then trigger flags.
    Now, these kinds of systems exist in big ML systems,
    but not so much in the TinyML use case where
    I'm trying to hint that it's actually quite important,
    because you're dealing with sensor data.
    And all of this is going to be tied back to the continuous training pipeline,
    right?
    So whatever the continuous monitoring stage is doing,
    it's actually going to influence the continuous training pipeline, which
    is why it's extremely critical for us to understand how we actually
    engineer this pipeline.
    For now, though, let's just assume that these data
    set monitors are in place down in the continuous monitoring stage,
    and that somehow, they know when to retrigger the continuous training
    pipeline.
    And then, we can just continue to dive into the continuous training
    pipeline for now to just be focused.
    Ultimately, the goal of the continuous training pipeline
    is to prevent model decay, as shown on the left.
    As you can see, the model accuracy is dropping over time.
    Instead, we want to keep retraining the model so
    that we can keep fixing any model inaccuracies that appear.
    The figure on the right is also showing something quite important--
    that model decay is effectively inevitable.
    In other words, you will detect the need to retrain your model,
    and you will retrain the model.
    But that's not a one-shot thing.
    In other words, you have to keep doing it,
    because the real world data is continuously changing.
    So you're going to have to keep repeatedly updating the model.
    And this is why, my friends, the continuous training stage
    relies so heavily on the training operationalization stage.
    You need to have a operationalization stage that is really nice and robustly
    built, so that you can rely on the continuous training stage
    to be able to effectively have a short cycle
    and being able to update the model and push it out robustly into the field.
    While doing continuous integration, while doing continuous delivery,
    and while doing all the things we have to do for production deployment.
    All right.
    With all that said, let's now dive into the details of the continuous training
    pipeline, and I'll see you in the next video.
    Ende des Transkriptes, zum Anfang springen
    Downloads und Transskripte
    Video
    Videodatei herunterladen
    Transskripte
    Download SubRip (.srt) file
    Download Text (.txt) file        

\section{Data Processing Overview}

    0:00 / 0:00Klicken Sie auf die Pfeiltaste nach oben, um das Geschwindigkeitsmenu aufzurufen. Regeln Sie dann mit den Hoch und Runter Pfeiltasten die Geschwindigkeit. Klicken Sie Enter, um die jeweilige Geschwindigkeit festzulegen.Geschwindigkeit1.0xKlicke auf diesen Knopf und diess Video stumm oder nicht stumm zu schalten, oder drücke die HOCH oder RUNTER Knöpfe um die Lautstärke zu erhöhen oder abzusenken.Maximal Lautstärke.Videotranskript
    Beginn des Transkriptes, zum Ende springen.
    VIJAY JANAPA REDDI: The first of the main continuous training tasks
    is data processing.
    The data processing engine is responsible for dealing
    with three main tasks, data ingestion, validation, and transformation.
    Data ingestion is where the training data
    is extracted from the source data set and feature repository
    using some filtering criteria.
    Data validation is when we actually take the training data
    that we have extracted and validate it to make sure
    that the model is trained on good data, not skewed or corrupted data.
    And data transformation is when we actually
    split the data up into training evaluation and testing splits.
    And the data is then usually transformed into some form that's
    actually useful for training the data, where we can actually engineer
    the features and go on as usual.
    Now this is the thing that your data engineer is typically dealing with.
    I'm going to use keyword spotting as an example yet again.
    One of the interesting things with keyword spotting
    is that we actually want to make it ubiquitous.
    Already, we see it everywhere.
    However, by and large, keyword spotting is limited to a handful of languages.
    And this is because training a keyword spotting model
    usually requires thousands of examples from many different speakers, which
    presents a serious challenge in terms of data collection.
    So one of the grand objectives is to try and make
    keyword spotting-based interfaces broadly available in any language.
    Now let's compare this with the gold standard,
    which is Google Speech Commands.
    Now this data set, as you may recall, has 35 different keywords
    and has over 100,000 different samples in it.
    But in reality, it only works on English when, in fact, there
    are over 7,000 different languages that are being spoken across the globe.
    So the question is this, how do we bring about equity in the system
    so that we can enable machine learning technologies
    such as keyword spotting for everyone?
    The key challenge, as I mentioned earlier,
    is that as you try and cater for the long tail,
    we don't tend to have data sets for a wide number of languages.
    So recall why this is the case.
    It's because data engineering is typically very costly.
    Gathering the requirements, going and collecting the data,
    then refining it and sustaining that whole process is extremely costly.
    It's also highly error-prone.
    Yet it is very critical that we get it right.
    Now in the context of speech commands, as you can see here,
    it took Pete about two years to go through this process just
    to create that small, little data set that he can make broadly available.
    So the key question is, how can we scale this part up so that we can actually
    support a wide range of languages?
    And this is what I'm going to use in order
    to help you understand the value of continuous training pipes.
    Now remember, the key challenge is really
    that labeling is really expensive for us.
    What we want to do is try and figure out if there's
    a way to automatically label data sets.
    To land on that, I'm going to use a concrete example called
    the multilingual spoken words corpus, which is a keyword spotting data set.
    I'm going to use this to help establish the connection through the data
    ingestion, validation, and transformation,
    all done automatically through a continuous training pipe.
    Let me first give you a high-level overview of MSWC
    so you understand what it is.
    And then in subsequent videos, I will actually dive in deeper.
    At a very high-level, MSWC has one-second
    spoken word examples in 50 different languages.
    The data is basically categorized for each different language
    based on the number of hours that's actually present in the data set.
    And as you can see, there are three main areas
    that you can break them down into.
    This is just an arbitrary breakdown just to help you understand things.
    For instance, the data set has over 100 hours
    of keyword data in 12 different languages.
    And these 50 languages are currently estimated
    to be spoken by over five billion people,
    representing a significant potential reach for this particular data set.
    Now an average for low-resource languages,
    there are over 500 unique keywords and a total of 10,000 audio samples
    inside this data set.
    Now for many languages, this data set is the first publicly available
    keyword data set.
    Now for high-level resource languages, there's
    a substantial amount of corpus in here.
    On average, there are over 23,000 keywords available
    and over 1.8 million audio clips in the data set.
    In total, this behemoth data set contains 340,000 keywords
    with over 23 million samples, which totals to about 6,000
    plus hours of speech.
    If you're curious, you can go check it out at ML Commons.
    And you can actually download it for free.
    So what I've done is set up the context for you in this video.
    Next, what I'm going to do is we'll actually
    learn the lessons behind data ingestion, validation, and transformation
    using this data set as the example.
    And the reason is because the MSWC data set was created
    using the continuous training pipe.
    And so hopefully, you will be able to make the connections as
    to how such a large data set was created thanks to all
    the operationalization that was done.
    I'll see you in the next video.
    Ende des Transkriptes, zum Anfang springen
    Downloads und Transskripte
    Video
    Videodatei herunterladen
    Transskripte
    Download SubRip (.srt) file
    Download Text (.txt) file

\section{Data Engineering For Everyone}

\subsection{Background}

    During the early days of software development, many programs, tools, and architectures were closed-source (i.e., individuals could not freely peruse the codebase). The advent of open-source software to help overcome problems faced by large distributed teams of software engineers propelled software development forward massively, enhancing productivity and broadening access to resources. Today, we face a largely similar challenge with respect to data.
    
    \subsection{The Importance of Public Data}
    
    Machine learning is becoming increasingly ubiquitous, with tens of thousands of research papers now published on the subject every year. The majority of these algorithms require data to achieve their objectives, meaning that the appetite for data is concomitantly increasing. Expertise in data engineering is becoming increasingly necessary to architect, build, and manage large and high-quality datasets used for training machine learning algorithms. However, a substantial proportion of these datasets are built and managed by large corporate entities, and may or may not be open-source (approximately 1 in 5 published works by Google, Facebook, and Microsoft utilized private datasets). You can find the source of the details here.
    
        \GRAPHICSC{1.0}{1.0}{TinyML/4 MLOps/Screen_Shot_2021-11-30_at_12.50.07_PM.png}   
    
    A pie chart showing that 78.3\% of datasets are open datasets while 21.7\% are closed private datasets. Total responses to the survey were 429 for private and 1548 for public.
    
    \subsection{The State of Public Data}
    
    Within the machine learning community, there has been a primary focus on architectural or algorithmic problems as opposed to a focus on data. Consequently, the development of efficient architectures to crowd-source and democratize large datasets for use in machine learning algorithms has been largely sidelined. Many of the data benchmarks that exist today, such as the ImageNet and COCO datasets, were developed within academia by a small group of researchers. However, since most machine learning algorithms today require more data than can be readily produced by such groups, there is a growing need for such massive datasets to be produced through open-source contributions.
    
    Although there has been a minor improvement in the prevalence of open-source dataset usage within academic publications over the years, this number should ideally be 100\% to allow results to be readily checked and reproduced. When comparing publications by Google, Facebook, and Microsoft, the prevalence of private datasets has remained fairly consistent. This suggests that even for the giants of industry, improving the use and development of open-source datasets is not a primary goal.
    
    \GRAPHICSC{1.0}{1.0}{TinyML/4 MLOps/Screen_Shot_2021-11-30_at_1.21.20_PM.png}   
    
    A series of bar charts showing the use of private vs. public datasets for Facebook, Google, and Microsoft. We see that for all three companies the number of total publications is growing and the breakdown of private vs. public is different for each company. For Facebook the data is from 2015 to 2020 and shows (public-private) by year of (15-7), (19-3), (34-10), (124-33), (170-66), (209-47). For Google its (23-4), (30-12), (87-26), (93-27), (101-29), (157-43). For Microsoft the data starts in 2017 and its (98-38), (85-26), (137-33), (189-39).
    
    \subsection{Progress}
    
    Initiatives such as MLCommons have been developed to help encourage assessment using open-source benchmarks, datasets, and best practices, but there is still a long way to go before data engineering is on a similar level of open-source parity as software engineering.        
        
\section{Data Ingestion}

    0:00 / 0:00Klicken Sie auf die Pfeiltaste nach oben, um das Geschwindigkeitsmenu aufzurufen. Regeln Sie dann mit den Hoch und Runter Pfeiltasten die Geschwindigkeit. Klicken Sie Enter, um die jeweilige Geschwindigkeit festzulegen.Geschwindigkeit1.0xKlicke auf diesen Knopf und diess Video stumm oder nicht stumm zu schalten, oder drücke die HOCH oder RUNTER Knöpfe um die Lautstärke zu erhöhen oder abzusenken.Maximal Lautstärke.Videotranskript
    Beginn des Transkriptes, zum Ende springen.
    VIJAY JANAPA REDDI: In this section, I'm going
    to talk about the data ingestion stage, which
    is the very first stage of the data processing engine.
    Recall that we'll be using the MSWC corpus as our case study.
    In total, MSWC is huge because it's got over 340,000 keywords,
    with over 23 million samples, totaling for over 6,000 hours of speech.
    It's important to note that MSWC is a generated data set which
    is assembled automatically by extracting individual words
    from continuous speech.
    That seems powerful so let's actually learn what the data ingestion
    pipeline would look like in order to create a massive corpus like this
    automatically.
    In MSWC, the keyword extraction pipeline ingests a pair
    of sentence audio and text transcripts.
    This could be sourced manually, for instance, when you just go out
    and buy a data set from someone.
    And by the way, that's actually quite common.
    Folks do go and buy data sets, and they sell data sets.
    And the value of the data really just kind of depends
    on the quality of the data set.
    So bottom line, it basically means that you
    have various kinds of data sets you might be getting your hands on.
    The simplest case would be that you end up
    with a static data set that someone just has this pair,
    audio and transcript data set, and then they just give it to you.
    Or you can actually source it from audio video films.
    Let me explain what I mean by that.
    There's a wealth of information out there
    that you might actually be able to leverage.
    Here, I'm showing, for example, how you might
    be able to use YouTube data as a source of information.
    You see captions.
    Well, those captions are actually the source for keywords.
    For example, when Dylan is saying concepts to keep in mind
    latency and throughput in computing is the transcript.
    And there is an audio file obviously associated with it
    because it's a YouTube video file.
    And that becomes the input.
    And then you can use a system like forced alignment
    to estimate per word timing boundaries.
    That's a well established technique in the speech community.
    It's effectively a way to match the audio with the transcript.
    Now, using the estimated timings, we apply some sort of inclusion criteria.
    Just kind of depends on what we're looking for, right.
    You might say OK, I'm interested in the work time, ladder, yes, no,
    and so forth.
    In this particular example here, we're basically
    saying whatever you extract, let's just sink
    that into the data set, which is kind of what they did for the MSWC data set.
    Now, you can also apply some sort of intelligent inclusion criteria.
    For example, you might want to insist on having
    a sufficient number of characters, at least for certain types of languages,
    like Chinese and so forth, that are actually matters,
    or that you have at least five samples per word.
    The reason for imposing some sort of inclusion criteria
    generally is that you want to have sufficient data.
    Now, I just said five samples, for example.
    You might say well, is five samples really any good?
    Well, it really depends on how you are actually training your model.
    We'll talk about that once we actually get into the model training stage.
    Once you have done all of this, you will end up
    with a large data set, like the MSWC, and then you
    can go ahead and train the model by handing it to your ML researcher.
    Now, the cool thing about this sort of a pipeline, for instance,
    is the fact that you can retrigger this pipeline whenever you want.
    You might want to target new languages.
    Instead of just English, German, Spanish,
    you might want to expand into Telugu, Hindi, Tamil, and so forth.
    In which case, all you have to do is point it towards a new source,
    and then the system automatically will generate new languages for you.
    Now, there are, of course, some key challenges
    that you have to address during the data ingestion stage.
    For data sets that are continuously updated, you
    need to track the metadata, where each sample was collected,
    so that you can actually backtrack in case something is needed to be done.
    For example occasionally, you'll have to automate
    the removal of samples that were collected from a set of faulty sources.
    Ultimately, you might realize some categories
    that are undersampled in your data set.
    And having provenance information can actually
    help you collect more samples for those categories.
    It's like having a data sheet for your data set.
    So in a way, as much as you want to automate things,
    you also want to be able to actually collect metadata about what you're
    actually putting into the data set.
    And as your model's use cases evolve in production,
    you might also experience data drift, where some training
    data is no longer relevant.
    So in those kinds of cases, yet again, you might want to throw some data out.
    But again, you need to have metadata for all of that.
    Now, another thing to keep in mind is that when supporting updates
    to the training data in a continuous training pipe,
    you should always ensure that any training data which stays the same
    is not accidentally repartitioned into validation or testing data,
    or vise versa.
    You will often simultaneously be working on improving the model performance.
    And this requires holding fixed which samples
    are used for training in order to compare accuracy metrics for a given
    category.
    Otherwise, if you're randomly shuffling your data in the training, validation,
    and testing sets as you're improving your data set,
    then you might get into a situation where
    it becomes impossible to truly compare the performance of the new model
    against the previous version of the model
    because you're effectively mixed up the data.
    And so keep that in mind because this is a critical part of the data ingestion
    pipeline.
    That said, I'll see you in the next step.
    Ende des Transkriptes, zum Anfang springen
    Downloads und Transskripte
    Video
    Videodatei herunterladen
    Transskripte
    Download SubRip (.srt) file
    Download Text (.txt) file        
        
\section{Data Validation}

    0:00 / 0:00Klicken Sie auf die Pfeiltaste nach oben, um das Geschwindigkeitsmenu aufzurufen. Regeln Sie dann mit den Hoch und Runter Pfeiltasten die Geschwindigkeit. Klicken Sie Enter, um die jeweilige Geschwindigkeit festzulegen.Geschwindigkeit1.0xKlicke auf diesen Knopf und diess Video stumm oder nicht stumm zu schalten, oder drücke die HOCH oder RUNTER Knöpfe um die Lautstärke zu erhöhen oder abzusenken.Maximal Lautstärke.Videotranskript
    Beginn des Transkriptes, zum Ende springen.
    VIJAY JANAPA REDDI: In this video, I'm going
    to talk about what comes after the data ingestion stage.
    We're going to move on to the data validation
    stage, which is the second stage.
    Now for a moment before we talk about MSWC, let's go back
    and look at the Google speech commands data set.
    This data set, as you may recall, was actually
    manually validated which took a long time to actually get right.
    Clearly doing some manual validation for the MSWC corpus is going to be
    challenging-- don't you think--
    because the corpus is huge.
    Now automatic extraction pipelines, like the one that generated the MSWC,
    are fantastic because you get a large number of keywords,
    a large number of samples, and multiple languages and all that good stuff.
    Let's go ahead and actually listen to two sets of extractions
    to see how well the system actually does because that
    will help us understand how critical data validation is actually
    going to be for us.
    For example, let's go ahead and listen to these good audio samples
    that were from the data set.
    [AUDIO PLAYBACK]
    - Learning.
    - Learning.
    - Learning.
    [END PLAYBACK]
    VIJAY JANAPA REDDI: Great.
    Now, let's go ahead and listen to these other audio samples,
    which might not sound so great.
    [AUDIO PLAYBACK]
    - [INAUDIBLE]
    - Learning.
    - [INAUDIBLE]
    [END PLAYBACK]
    VIJAY JANAPA REDDI: As you can see here, there
    are several sources from where you can end up having errors
    inside the extract of the keyword.
    For example, this might happen because you have transcript mismatches
    or you have incorrectly labeled word boundaries.
    And all of this will lead to a poor extraction.
    And that's why in these particular samples,
    you are not able to figure out that the word still is indeed learning.
    So we need to somehow clearly validate this data set before we start using it.
    Well unlike Google speech commands, which is a much smaller data set,
    it's incredibly hard to validate an extremely large data
    set, which is kind of what you would actually
    have for a production use case.
    So if you have 350,000 words with 6,000 hours of speech,
    well that would take forever to be able to validate
    if we're having a human actually going and listening
    to every single one of them.
    So we need to cook up some automatic way of doing data validation.
    Now, the way you would go about doing data validation
    is, of course, going to be very dependent upon the type of data
    you're actually looking at.
    For instance, in the context of keyword spotting data,
    we can use some clever techniques like self-supervised outlier detection.
    For instance, you can use a self-supervised nearest
    neighbor or anomaly detection algorithm to automatically filter out
    erroneous samples using a distance metric.
    This works by effectively having a small number of clusters.
    Let's say 5 or 7 or something like that with a handful of random samples drawn
    from the full data set.
    And then, you treat any of those samples that are far away from these clusters
    as anomalous.
    Now just taking the MSWC corpus and running this
    through with about 75 keywords for English, Spanish, German,
    and so forth, what you end up getting is that on average,
    23% of the farthest 50 clips that were sorted by distance
    are actually erroneous.
    So the users of the data set might actually
    come up with various other kinds of metrics,
    but you can imagine that you want to have some base level of quality
    checking going on for this data.
    In the context of an evolving and continuously updating data set,
    now consistency is going to be important.
    Consistency in validation is critical in order
    to be able to accurately compare the new models you might end up creating
    using this with the older ones.
    Conversely, you might also find issues in the data validation process itself.
    And then, you will have to revisit some of the old previously validated data.
    For example, you might increase the input resolution
    on an object detection model.
    And then suddenly, you find that the bounding boxes that you
    have on the previously labeled data are actually now imprecise
    at that higher resolution, which immediately
    means that you would have to go back to relabeling
    which is one of the first steps in your data collection pipe.
    So the point that I'm trying to make here
    is that it's extremely important for us to have a data validation, an automated
    data validation setup, as much as it is important
    in order to have an automated data ingestion sort of a pipeline.
    So in the continuous training pipe, everything
    is supposed to be automated so that way you can effectively continuously
    be updating things robustly.
    So keep that in mind.
    I'll come back and talk to you about the last stage of the data processing
    engine.
    Ende des Transkriptes, zum Anfang springen
    Downloads und Transskripte
    Video
    Videodatei herunterladen
    Transskripte
    Download SubRip (.srt) file
    Download Text (.txt) file        
 
\section{Data Transformation}

    0:00 / 0:00Klicken Sie auf die Pfeiltaste nach oben, um das Geschwindigkeitsmenu aufzurufen. Regeln Sie dann mit den Hoch und Runter Pfeiltasten die Geschwindigkeit. Klicken Sie Enter, um die jeweilige Geschwindigkeit festzulegen.Geschwindigkeit1.0xKlicke auf diesen Knopf und diess Video stumm oder nicht stumm zu schalten, oder drücke die HOCH oder RUNTER Knöpfe um die Lautstärke zu erhöhen oder abzusenken.Maximal Lautstärke.Videotranskript
    Beginn des Transkriptes, zum Ende springen.
    SPEAKER 1: In this section, I'm going to talk about the data transformation
    stage.
    The data transformation stage concerns of how
    you go from raw training samples and your disk,
    such as the audio waveform or an image file,
    and some associated label information, and convert that into a format that
    can be used to train a model.
    More likely than not, data transformation
    is going to involve some sort of feature pre-processing for model input.
    For example, speech models often rely on spectrogram data rather than
    raw waveform samples.
    Similarly, image classification models might
    require uncompressed JPEG data and resizing the image to the network input
    shape.
    You might also want to normalize the data in each color channel
    so that there are zero-centered.
    These are pretty textbook style standard approaches.
    But pre-processing can also be more involved depending on the use case.
    For example, you might have to support volumetric convolutions
    for medical imaging data.
    In such a case, you have to extract the feature volume from an input image
    first using a pre-trained network.
    Now, later on I have a Colab if you're interested in checking out
    about how some of these steps work.
    Basically, it works on the multilingual spoken words corpus,
    and it basically shows you how we go from taking
    a raw audio sample into converting that into a spectrogram.
    More specifically, it converts the audio sample
    into frequency bands that more closely approximate the human hearing
    perception characteristics for speech.
    So it leads to a much better keyword spotting model.
    Data transformation can also apply to labels, not just the training
    samples or the features themselves.
    For example, one challenge you might face in supporting continuous training
    is having to combine and ingest multiple data sets or data sources.
    So for example, let's say you're doing person detection or object
    detection, which is a fairly common tiny ML use case.
    There might be many data sets available for this particular task.
    But they're often coming in different formats.
    For example, data set one versus data set
    two might specify the bounding boxes in totally different ways.
    Thus, your data transformation pipeline will actually
    need to convert these different sources into a more standard format
    to create a single baseline that you can use for training the model.
    But data transformation can also involve more complex things
    like relabeling the data to support more advanced machine learning capabilities.
    For example, you might want to choose to relabel object detection data
    that you already have access to in order to better support a segmentation model
    where you'll need to convert all the bounding boxes from your object
    detection data set into segmentation maps for your segmentation model.
    A continuous training pipe might naturally therefore
    have to support various types of training jobs,
    in order to kick one off, for instance, to do object detection and another
    to do segmentation, and so forth.
    It really depends on what your customer's or end user's needs are.
    Now current research on repurposing training data
    for different kinds of applications like,
    has been based on bootstrapping the feature extractors by using things
    like self supervised learning.
    Now, these may become common elements of data transformation steps
    in future continuous training pipelines, but it remains to be seen.
    At this point, though, I hope that you have
    a rough idea about what the data processing engine does,
    starting from data ingestion to data validation and the need
    to actually do data transformation.
    I'm going to wrap this up, and next we'll
    be moving on into the model training/tuning stage.
    I'll see you there.
    Ende des Transkriptes, zum Anfang springen
    Downloads und Transskripte
    Video
    Videodatei herunterladen
    Transskripte
    Download SubRip (.srt) file
    Download Text (.txt) file 
 
 \section{raining vs. Tuning}
 
     \GRAPHICSC{1.0}{1.0}{TinyML/4 MLOps/pasted_image_0__1_.png}

     The ML model development process can be broken down into two possible parallel paths. The first is to train a model from scratch. The second is to fine-tune a pre-trained model.
     
     
     \subsection{About This Reading}
     
     Training a model from scratch might seem to be the norm. Certainly, this is almost always feasible if you have wealth access to a large corpus of data and training resources. But quite a few times, we are unable to train machine learning models from scratch. This might be due to the lack of computational resources, time available, or a paucity of available training data.
     
     Ideally, we would like to minimize unnecessary computation to help save us time, money, and also to minimize carbon emissions as a result of our computation, but also to port existing models used for general tasks to our more specific application. To this end, this section will discuss the pros and cons of training a model from scratch vs. fine-tuning a pre-trained model.
     
     \subsection{Examples}
     
     Imagine a scenario where you are developing an image-captioning model to produce a description of a patient’s chest x-ray. You might only have access to a small amount of data to develop this, say, 1000 x-rays and reports, and thus training an image captioning model from scratch would produce poor results. However, we can take an existing, more general, image-captioning model and use this as the starting point for our model. Sometimes, this is called “pre-initializing” or “pre-training” our network, since we import the weights from a previously trained network. Following this, we can train the latter layers of our network (while keeping the early layers frozen, since these encode more primitive image properties) on our smaller dataset of chest x-ray data. The results from this implementation should far exceed those of a model trained from scratch on the chest x-ray data. Clearly, this only works when the two tasks are similar enough that the model can be ported from the more general application to the more specific application. To this end, repositories of models have been developed (e.g., Hugging Face) to store models that are freely available for use. Some examples of popular models are BERT, VGG16, Inception, and GPT-2. These models were developed by organizations with vast amounts of computational resources, and are thus prohibitively expensive for any individual to produce. However, the fact that we can use and adapt these models to our own applications is a triumph of open-source machine learning.
 
\section{Training with AutoML}

    0:00 / 0:00Klicken Sie auf die Pfeiltaste nach oben, um das Geschwindigkeitsmenu aufzurufen. Regeln Sie dann mit den Hoch und Runter Pfeiltasten die Geschwindigkeit. Klicken Sie Enter, um die jeweilige Geschwindigkeit festzulegen.Geschwindigkeit1.0xKlicke auf diesen Knopf und diess Video stumm oder nicht stumm zu schalten, oder drücke die HOCH oder RUNTER Knöpfe um die Lautstärke zu erhöhen oder abzusenken.Maximal Lautstärke.Videotranskript
    Beginn des Transkriptes, zum Ende springen.
    VIJAY JANAPA REDDI: In this section, I'm going
    to be talking about the model training engine.
    Specifically, we'll be talking about training models from scratch
    efficiently, as well as tuning models that have already been pre-trained.
    Now if you recall, by and large, I've talked
    about training earlier already in the ML development
    and training operationalization stages.
    But here, I want to talk about making that process really efficient
    and exploring the vast design space that we have.
    Now, this is something that's typically done
    when you have a really good operationalized pipeline to start with.
    That is the MLOps stage that came right before continuous training.
    And it's the sort of thing that your ML researcher
    is going to be wickedly excited about.
    Specifically, I'm going to talk about AutoML,
    which is the whole process of automating the design and optimization
    decisions in the traditional ML workflow you have seen before,
    where we typically go from collecting data to training models,
    and we do that whole flow manually, and then we go on
    to deploying things into the field.
    Now, the goal of AutoML is to iterate through this flow over
    and over and over, not just do it once, with the intent that we
    can arrive at a potentially really good model at the end of the day.
    Now, if you remember keyword spotting, the three main stages
    to it, right, the input processing, the preprocessing, and the output stage.
    Well, even if we just look at the first part of the input stage,
    you might recall that there are a large number of different hyperparameters
    you can actually tune, the window length, the window size, and so forth.
    And when you look at the preprocessing, again,
    I told you previously in the ML development stage,
    there are a large number of different methods
    you can use for your preprocessing.
    You can try and end up with a spectrogram,
    or you can end up with the MFE, or you can end up with the MFCC,
    each one of which in itself has a whole different number of hyperparameters
    that you can tune.
    And all of these ultimately end up impacting the accuracy.
    They can also end up impacting the memory footprint of your model
    and how long the network might actually take to run on the end of the device.
    Taking a step back, what I'm trying to say
    is that there's a large cross product of things that can really
    affect the quality of the model.
    And ideally, you want to explore this design space in order
    to arrive at the best possible model for your customer.
    Now, as I've hinted earlier, you have input features.
    Then input features lead into preprocessing hyperparameters
    that you can actually play with.
    Of course, in addition to that, you actually
    have the neural network itself.
    Then you actually have the target device and its capabilities.
    You have various metrics you can go after, and the quality of the data
    set itself could also affect how well your network actually
    ends up performing.
    So given that there are so many different pieces in it,
    how do we explore the space?
    It's to cover this rich design space is why AutoML has actually
    emerged as a main topic.
    It was originally a research idea but has definitely
    matured into production pipelines.
    AutoML tries to automate the design and optimization decisions.
    And the key thing that is tries to do, especially
    in the context of embedded systems, is to try and get you
    to the best set of characteristics that you can arrive at,
    where in the context of embedded systems,
    you care more than just accuracy.
    You care about that latency.
    You care about that memory footprint.
    You care about that energy consumption and so forth.
    And so AutoML can actually help us with all these things.
    By navigating the complex space, it can try and figure out
    which of these things it can optimize the best.
    Ideally, in each case, it will maintain that accuracy as high as possible
    and minimize the latency, minimize the memory footprint,
    and minimize the energy consumption.
    AutoML is this really big thing that goes all the way
    from data all the way to deployment.
    Well, you don't necessarily have to think of AutoML in that one big loop,
    you can actually break it down.
    You can actually try and look at AutoML in its sub pieces, for example,
    you can try and focus on feature engineering.
    Alternatively, you can just focus on improving the network capabilities.
    For example, you might want to try different kinds of network topologies.
    You might want to change the learning rate.
    You might want to try different optimizers.
    You might want to try and play around with the weight initialization.
    You might want to change the loss functions.
    There are so many different ways of actually doing optimizations
    to the network itself.
    And then, of course, you can also train to target specific characteristics,
    such as optimizing the embedded system's specific capabilities,
    for example, targeting devices like cortex M0, or M4, M7, and M33,
    and beyond.
    That's because all of these different processors
    come with different capabilities.
    For instance, the M0 does not have a floating point in it.
    But the M4 has a floating point in it, but it doesn't have any specialized DSP
    or SIMD instructions, which the M7 has.
    And so as you step through these kinds of different devices
    and different processing engines, you have a rich design space that opens up.
    And you want to try and optimize a network to be as good as possible
    for that particular platform.
    Doing this manually is virtually impossible.
    So you want to try and automate this.
    So to ground all of this discussion on AutoML,
    I've chosen to talk about neural architecture
    search, which is a technique for automating
    the design of neural networks.
    It's a subfield of AutoML so we won't learn about everything in AutoML.
    Nonetheless, it'll give us one level of depth for how these systems will
    actually work in practice.
    Keep in mind, however, that there is no such thing as a free lunch.
    There's an inherent balance between the amount of time
    it takes for an AutoML or a NAS engine to run and the amount of search
    you actually can perform.
    So AutoML can take extremely long amount of time
    to finish the scope of the search, and we'll talk about some
    of the pros and cons around here.
    But bottom line, I do want to emphasize that many companies are actively
    specializing in just AutoML capabilities.
    And therefore, it is actually quite important
    to know that it actually exists and how you
    can take advantage of that as part of your continuous training pipe.
    Ende des Transkriptes, zum Anfang springen
    Downloads und Transskripte
    Video
    Videodatei herunterladen
    Transskripte
    Download SubRip (.srt) file
    Download Text (.txt) file 
 
 
\section{Neural Architecture Search (NAS) - Part 1}

    0:00 / 0:00Klicken Sie auf die Pfeiltaste nach oben, um das Geschwindigkeitsmenu aufzurufen. Regeln Sie dann mit den Hoch und Runter Pfeiltasten die Geschwindigkeit. Klicken Sie Enter, um die jeweilige Geschwindigkeit festzulegen.Geschwindigkeit1.0xKlicke auf diesen Knopf und diess Video stumm oder nicht stumm zu schalten, oder drücke die HOCH oder RUNTER Knöpfe um die Lautstärke zu erhöhen oder abzusenken.Maximal Lautstärke.Videotranskript
    Beginn des Transkriptes, zum Ende springen.
    VIJAY JANAPA REDDI: In this video, I'm going
    to give you an overview of neural architecture search.
    At its most basic form, neural architecture search or NAS
    automates the process of designing a neural network architecture in order
    to improve the model's accuracy.
    The way a model is designed can have a significant impact
    on the accuracy, often in difficult to predict ways.
    Now, designing and optimizing a model structure
    can actually be a very time consuming and manual effort.
    Especially when the model size is extremely constrained,
    you worry very much about the model design
    because you are very susceptible to accuracy changes.
    However, keep in mind that when you're talking about TinyML,
    we have extremely tight constraints on latency and memory as well as energy.
    So therefore, when we're using a neural architecture search algorithm,
    we must optimize for latency and memory consumption of the model in addition
    to accuracy.
    This is what is often called as a multi-objective search.
    Multi-objective search is often easier said than done,
    as any changes you make to a model's architecture
    often tend to have an impact on different metrics in different ways,
    be it accuracy, or latency, or memory, or even energy consumption.
    So for example, let's say you make a model larger
    in order to get better accuracy.
    Well, you may also end up increasing the latency of the model.
    Now, multi-objective NAS aims to find an optimal point
    where the accuracy, latency, and memory constraints
    are all cohesively met together.
    And it's important to remember that the optimal point depends entirely
    on the constraints of the application.
    Another important consideration is that the specific model configuration
    must match the architecture of the target hardware platform
    and therefore execute more efficiently with respect to a target platform.
    This means that the performance of a model compared
    to other possible configurations is highly dependent on the deployment
    platform's capabilities.
    If you recall earlier, I said that if you look at the ARM Cortex
    class of microcontrollers, where you have the M0 all the way up to M33,
    they have very different capabilities.
    The M0 for instance just provides you base capability
    to execute very simple instructions.
    The M4 steps it up by giving you a floating point
    unit, which allows you to achieve better accuracy because you're
    using floating point arithmetic.
    When you go to the M7, you can actually run the floating point model
    more effectively, and you might also be able to use many more computations
    because you have a DSP capable or a SIMD capable
    engine in there that allows you to run instructions much faster, specifically
    in the context of neural network executions.
    And as you scale it up M33, you might be able to run
    even more complex instructions.
    So you really want to try and get the NAS engine
    to be able to understand the characteristics of the hardware
    platform.
    So bottom line, you have to think about the hardware in the loop.
    And this is the hardware where NAS is a little different than traditional NAS
    where you're just kind of exploring the architecture independently
    of these hardware specific platforms.
    But given that in TinyML these things matter a lot,
    it's extremely important that we focus on hardware of our NAS.
    In summary, NAS basically aims to optimize your architecture
    to meet certain constraints of the application, given a target platform,
    in our case.
    There are many platforms out there like AutoKeras cameras
    and Google AutoML that abstract a lot of the details and make it super easy.
    So the point here is that it's clearly become an important concept to learn.
    Especially in the context of continuous training
    where you're trying to continuously keep models updated,
    you want the engine to not just generate one model.
    You want it to actually search a design space of models
    in order to generate a really good model.
    Now, coming up next, what I want to do is
    I want to open up the hood a little bit and tell you
    about some of the core building blocks that are there inside the NAS pipeline.
    This helps you understand some of the fundamental concepts.
    Now, the exact implementations within each of these boxes
    might vary across different implementations,
    but by and large, the community has kind of settled on this sort of a framework
    in terms of how NAS is actually deployed and executed.
    And we'll talk about more on this in the next video.
    Ende des Transkriptes, zum Anfang springen
    Downloads und Transskripte
    Video
    Videodatei herunterladen
    Transskripte
    Download SubRip (.srt) file
    Download Text (.txt) file 
 
\section{Neural Architecture Search (NAS) - Part 2}

    0:00 / 0:00Klicken Sie auf die Pfeiltaste nach oben, um das Geschwindigkeitsmenu aufzurufen. Regeln Sie dann mit den Hoch und Runter Pfeiltasten die Geschwindigkeit. Klicken Sie Enter, um die jeweilige Geschwindigkeit festzulegen.Geschwindigkeit1.0xKlicke auf diesen Knopf und diess Video stumm oder nicht stumm zu schalten, oder drücke die HOCH oder RUNTER Knöpfe um die Lautstärke zu erhöhen oder abzusenken.Maximal Lautstärke.Videotranskript
    Beginn des Transkriptes, zum Ende springen.
    VIJAY JANAPA REDDI: In this section, we'll
    be discussing the underlying components of the neural architecture search.
    Neural architecture search can broadly be broken down
    into these three core, big pieces.
    One is the search space, the search strategy, and the performance
    estimation strategy.
    First, we'll look into the search space, which
    defines the possible range of models your NAS algorithms can actually
    explore.
    Now, the search space is a set of different architectural points
    where you might have some that are good, some that are bad,
    but they also offer various kinds of trade-offs
    that only you can decide upon.
    For example, if you want a small model, you
    might get that really small model, which is great
    because you might be able to fit it onto a very small, little, embedded device.
    You can also have a really fast model if latency is critical for you.
    Or you might have a really highly accurate model
    when accuracy is paramount, perhaps even at the expense
    of some amount of latency.
    Now, the search space size is a critical decision.
    If the search is very large, you will spend far too much time
    looking for the optimal model.
    And the key concern here is that it takes a lot of time,
    and it can also cost you a lot of money in order
    to be able to navigate the search space.
    Later on, I have a reading that's coming up
    that talk about some of the pitfalls of doing too much exploration
    in order to try and find that optimal model.
    If you make the search space too small, you
    might actually never end up finding that good model that you're looking for.
    Ideally, we'd want to have some sort of prior knowledge
    that we can apply in order to pick a search space for the models that
    are more likely to be accurate or meet the device constraints.
    Now, there are a number of different types of search spaces
    that constrain the architecture of the model in many different types of ways.
    One common type is to use a chain structure
    that linearly structures the layers.
    The size, the shape, and the type of the layers can all be constrained.
    Another type is to use multi-branch where multi-branch search often
    defines a much larger search space because it allows
    for more complex and potentially optimal models to be found,
    but, yet again, you have the cost of long search times here.
    Now, many NAS algorithms will use cell-based structures
    as well that repeat a chunk of the model a certain number of times.
    The inner structure of the cell is often the target of the search in addition
    to how many times you repeat that cell.
    This search type is an example of using some prior knowledge in order
    to improve the search time for the NAS algorithms.
    In this case, the prior knowledge is that the repeated structures often
    produce good results.
    The next component of NAS is really the search strategy.
    The search strategy defines how you pick the model
    that you're going to test next.
    It involves sending the candidate model to the performance estimation strategy
    and then receiving some result that comes back that tells you
    how well the model performed.
    And the core purpose of this search strategy
    is, how do you use the feedback from the performance estimation strategy
    in order to guide the search to pick the next model that you're actually
    going to evaluate?
    And that's how you kind of iterate over and over.
    And that, hopefully, allows you to arrive at a good model.
    The question is, what type of search strategy do you use?
    Well, there are a lot of different kinds.
    You can use Bayesian methods.
    You can use evolutionary methods.
    You can use reinforcement learning.
    You can use gradient-based methods.
    NAS is still a very active area of research,
    even though you are seeing it in production.
    So it's kind of hard to say which one of these strategies
    really works well, which also makes it problematic because that
    is why people want to explore a rich search design space.
    Now, the best strategy, more often than not,
    is really going to be task dependent.
    So, for example, if you're doing it for keyword spotting versus person
    detection, that's going to have a totally different conversion
    sort of point.
    And the method can also depend on what level of restrictions or flexibility
    you have.
    For example, something called DNAS, or differential NAS,
    is a direct and fast method.
    It's really good, but it assumes you have continuous functions.
    Now, there are other methods, like Blackbox-based testing
    methods, that are often slower, but they're much more diverse and flexible.
    So you have to kind of explore the right mechanism that works for you.
    And, finally, we have this performance estimation strategy,
    which we basically said actually tells us how well the model performed.
    Well, yes, that's good, but then there are challenges
    in the context of embedded systems.
    Remember, this performance estimation strategy
    has to try and profile and understand how
    the model will behave for latency, for memory consumption, for energy,
    and so forth.
    Now, the performance of the model can either be measured directly on device,
    or you can try and estimate it through proxy metrics.
    Now, the reason you would be concerned with measuring things directly
    on device with TinyML is because it takes a long time.
    Remember, when we were talking about continuous integration
    and continuous delivery, we talked a lot about staging environments
    and the challenges with having a whole fleet of different embedded devices.
    Well, that kind of comes back and bites us here
    because here we ideally want to be able to farm out and explore
    a whole different range of these devices,
    but then, if you try and rely on running the model on the physical device,
    it's going to really slow you down.
    So, often, what folks end up doing is they
    end up with proxy metrics based on the number of ops,
    for instance, that the network has.
    We estimate how long or how much energy will the network really consume.
    And, based on that, we use the estimation strategy
    to give feedback back.
    So NAS is still an active area of research.
    It is constantly evolving, and, therefore,
    people are working really hard to improve the process.
    Nonetheless, it is being deployed in the field today.
    Especially in TinyML, I've already seen deployments in commercial products
    where they actually want to explore this design space because energy and latency
    and so forth are extremely tight constraints in embedded devices.
    And, therefore, by incorporating this into the continuous training pipe,
    you can actually try and design models that are really well matched or well
    suited for the target platform.
    Now, NAS can give us very good results.
    More often than not, you can get as much as an order of magnitude or more
    in performance improvement, depending on your device.
    But techniques like NAS that are part of AutoML are definitely not free.
    In the next reading that's coming up, you'll
    learn about the carbon footprint or the carbon price of AutoML systems
    and the amount of energy it takes in order
    to try and arrive at optimal models.
    And it is something that we actually have
    to keep in mind because sustainability of AI
    is now starting to become a very important topic.
    So please do make sure you read that reading carefully.
    Ende des Transkriptes, zum Anfang springen
    Downloads und Transskripte
    Video
    Videodatei herunterladen
    Transskripte
    Download SubRip (.srt) file
    Download Text (.txt) file  

\section{The Carbon Price of AutoML: CO2}

    \subsection{Carbon Cost of ML Training}
    
    Training In recent years, machine learning has undergone a bifurcation, a splitting into two scales: big and small. At the small scale, we have TinyML, where model size and efficient computation are crucial for models to be feasibly deployed onto resource-constrained microcontrollers. At the opposite end, we have large-scale neural networks trained on cloud computing clusters, using dozens or even hundreds of GPUs. At smaller scales, such as those in TinyML or on an individual laptop, the carbon footprint of models does not have a significant footprint. However, at a larger scale, these models can have carbon footprints that rival typical emission sources such as planes, trains, and automobiles.
    
    Models developed for natural language processing (NLP) are of particular concern since these seem to exhibit a performance improvement that scales proportionally with model size. This means that NLP researchers will continue to produce larger and larger natural language models, since currently no limit has been reached where performance starts to degrade with size. In principle, this could lead to networks with billions or even trillions of parameters. Computational resources also scale proportionally with the number of parameters, meaning that larger models will begin to have increasingly large carbon footprints.
    
\GRAPHICSC{1.0}{1.0}{TinyML/4 MLOps/005.png} 
    
    
    A graph showing model size in terms of billions of parameters in a log scale revealing the exponential increase in model size from 2018 to 2022. Starting with ELMo (early 2018 94M), BERT-Large (late 2018 340M), GPT-2 (early 2019 1.5B), Meatron-LM (mid 2019 8.3B), T5 (late 2019 11B), Turing-NLG (early 2020 17.2 B), GPT-3 (mid 2020 175B), and Megatron-Turing-NLG (late 2021 530B).
    
    \subsection{AutoML Can Exponentially Increases The Carbon Cost}
    
    This is quite concerning, especially since these models already account for a substantial portion of carbon emissions from data centers. However, the incorporation of AutoML with large-scale language models such as these adds even more fuel to the fire, so to speak. As we have seen,  AutoML allows us to automate the optimization of a model in terms of feature engineering, hyperparameter tuning, neural architecture search (NAS), and auto-optimizing compilation. Since large-scale language models already take a significant amount of resources to train, they take an even more significant amount of resources to optimize, in some cases orders of magnitude more resources.
    
\GRAPHICSC{1.0}{1.0}{TinyML/4 MLOps/006.png} 
    
    Our standard ML workflow: Collect Data, Preprocess Data, Design a Model, Train a Model, Evaluate and Optimize a Model, Convert a Model, Deploy a Model, and Make Inferences. Overlayed are the possible optimizations that can require steps to be run a second time, or more. Feature Engineering moves backward from model design to preprocessing. Hyper-parameter tuning goes from Evaluate and Optimize to training. Neural architecture search goes from evaluate and optimize to model design. Auto-optimizing compilers go from inferences to evaluate and optimize. And finally, AutoML goes all the way from inferences back to data collection, possibly re-running the entire pipeline.
    
    A study by Strubell, Ganesh, and McCallum (2019) estimates that training a large-scale transformer with NAS has a carbon footprint equivalent to that of five cars over their working life. Performing NAS on this transformer required over 3000x the resources compared to training the original transformer model. Since NAS is only a subsection of AutoML, one could easily imagine the resources being required for AutoML being even more significant. Not only that, but training this model cost upward of \$1 million! Clearly, the cost and carbon footprint of training large-scale models such as these need to be considered in their development
    
    Chances are, you will not be running NAS on a large-scale language model with billions of parameters. However, it is important to understand the environmental impact that such models are beginning to have due to their inherent scale, and to be cognizant of this when building your own models.

\section{Continuous Training with Transfer Learning}

    0:00 / 0:00Klicken Sie auf die Pfeiltaste nach oben, um das Geschwindigkeitsmenu aufzurufen. Regeln Sie dann mit den Hoch und Runter Pfeiltasten die Geschwindigkeit. Klicken Sie Enter, um die jeweilige Geschwindigkeit festzulegen.Geschwindigkeit1.0xKlicke auf diesen Knopf und diess Video stumm oder nicht stumm zu schalten, oder drücke die HOCH oder RUNTER Knöpfe um die Lautstärke zu erhöhen oder abzusenken.Maximal Lautstärke.Videotranskript
    Beginn des Transkriptes, zum Ende springen.
    VIJAY JANAPA REDDI: In this video, I'm going
    to focus on transfer learning in the context of continuous training.
    I've already talked about transfer learning a while back
    in the ML Ops development stage, but I want
    to close the loop here and just show you how powerful transfer learning can
    be when you integrate it into the continuous training pipe.
    Training a model normally requires a large amount
    of data for any given task.
    So for instance, if you want to do an image classification task,
    people would normally train it on ImageNet with over 1 million samples.
    Now, the cost of collecting that data for specific tasks is really high.
    It's almost prohibitively expensive.
    But our intent is that we want to enable AI for everybody, which
    means everybody's got to have efficient ways of training models
    that meet their specific use cases.
    Furthermore, remember, training resources
    can be quite costly, especially because you often tend to iterate over,
    and over, and over when you're trying to train models.
    It's not a one shot thing, you tend to repeat this process
    many times over until you're happy with the final model that you get.
    However, we don't always have to train from scratch, if you remember.
    Instead of training neural networks from scratch,
    we can instead try and use transfer learning, where
    we freeze the general feature extractor intact,
    and then only train the last few layers.
    This would allow us to maintain the general pattern recognition
    capabilities, while learning new high level features in the relationships
    on the final task side of the world.
    So this allows us to effectively be able to cut down on training time
    quite dramatically.
    And it's called transfer learning for an obvious reason, right?
    It's the process by which we learn to transfer knowledge
    of a pre-trained model to a new task.
    Now, as an example of the usefulness of transfer
    learning in the context of continuous training,
    let's consider the previously discussed application
    of keyword spotting from a multilingual sort of perspective.
    Let's consider this scenario.
    Let's say, hypothetically, you're a startup company that's
    invested in keyword spotting models.
    That's what you do.
    And you have a customer base that is constantly evolving and growing,
    as in, the demographics are continuously changing.
    Now, this is both an opportunity and a challenge for you.
    The opportunity is that you have a growing customer base,
    and they all want to be able to use your keyword spotting model
    in their language.
    However, the challenge that you have is that you
    don't have a large volume of data.
    Recall that in the data engineering section,
    I said how many languages have very little data?
    Now, this is where you need to stop and ask yourself,
    what if I had a little bit of data?
    Can I use transfer learning to do something here?
    Well, let me show you how that can actually be a powerful concept.
    We can create a model, for starters, that classifies the audio spectrograms
    into many different classes.
    Then what we can do is, we can try and train this model.
    In order to just get this original model,
    we can try and train this model on high resource languages
    where there's a lot of data already accessible.
    So, for instance, with the multilingual spoken words corpus,
    in that particular case, you can use a high resource languages,
    like English, German, Spanish, and so forth,
    for which there's a lot of big data already available.
    And then, once the model is trained, we can use the penultimate layer
    as the embedding representation, which embeds the high level feature
    information about the input speech.
    Now, this representation has the benefit of embedding
    multiple utterances of the same word close to one another,
    regardless of the speaker.
    And when you visualize this in the context of the t-SNE plot, which
    I mentioned in the ML development stage, we
    can see that the features learned by the front portion of the model
    can be easily clustered into these different word labels,
    meaning that the embedding representation does
    a good job of encoding the human speech in there.
    This is really powerful, because now we can take this embedding model,
    and then we can use transfer learning to create keyword spotting models
    for arbitrary keywords.
    By using the pre-trained embedding, we can reduce the usual requirement
    of having a huge volume of data in a new language,
    to just a few utterances of the target keywords that you care about.
    For example, here, we can just take five words
    and get a whole new keyword spotting model in that new language,
    with just five examples.
    Literally, I could just stand there and give five examples, and that's it.
    This model's immediately trained for a new language.
    And that's a remarkably powerful concept.
    You don't need thousands of words.
    So, putting this back into the context of continuous training,
    imagine how efficiently you'll be able to support
    new languages using the embeddings and transfer learning approach.
    These are some creative ways that we actually
    want to be thinking about doing things as we move into the future,
    so that we can enable machine learning at scale without always
    having to start training from scratch.
    So in summary, transfer learning allows us
    to reuse the features learned from a previously trained model,
    and then fine tune it to a new task, thereby saving
    an insane amount of time, and reducing the computational requirements,
    as well as reducing the data requirements for new ML applications.
    I really want you to take this video to heart,
    because it tells you about a really powerful concept, which
    in the context of TinyML, is actually getting a lot of traction,
    and it's starting to become the bread and butter way of very quickly training
    small ML models that are highly specialized for different types
    of tasks.
    So keep that in mind, and I'll see you in the next video.
    Ende des Transkriptes, zum Anfang springen
    Downloads und Transskripte
    Video
    Videodatei herunterladen
    Transskripte
    Download SubRip (.srt) file
    Download Text (.txt) file

\section{Pros and Cons of Transfer Learning}


    
    \subsection{Motivation}
    
    Sometimes, we are unable to train machine learning models from scratch. This might be due to the lack of computational resources, time available, or a limited amount of available training data. Ideally, we would like to minimize unnecessary computation to help save us time, money, and also to minimize carbon emissions as a result of our computation, but also to port existing models used for general tasks to our more specific application.
    
    Imagine a scenario where you are developing an image-captioning model to produce a description of a patient’s chest x-ray. You might only have access to a small amount of data to develop this, say, 1000 x-rays and reports, and thus training an image captioning model from scratch would produce poor results. However, we can take an existing, more general, image-captioning model and use this as the starting point for our model. Sometimes, this is called “pre-initializing” or “pre-training” our network, since we import the weights from a previously trained network. Following this, we can train the latter layers of our network (while keeping the early layers frozen, since these encode more primitive image properties) on our smaller dataset of chest x-ray data. The results from this implementation should far exceed those of a model trained from scratch on the chest x-ray data. Clearly, this only works when the two tasks are similar enough that the model can be ported from the more general application to the more specific application.
    
    \GRAPHICSC{1.0}{1.0}{TinyML/4 MLOps/UntitledReuse}  
    
    A high level diagram describing transfer learning showing re-training of the last few layers and freezing/reusing pre-trained initial layers. This is shown visually by graying out the first set of layers, which take in the input and have already learned to produce general features from the input, and then highlighting the last few layers that take those features and produce the output layers and noting that those are re-trained.
    
    To this end, repositories of models have been developed (e.g., Hugging Face) to store models that are freely available for use. Some examples of popular models are BERT, VGG16, Inception, and GPT-2. These models were developed by organizations with vast amounts of computational resources, and are thus prohibitively expensive for any individual to produce. The fact that we can use and adapt these models to our own applications is a triumph of open-source machine learning. However, despite all of the available resources, it is often difficult to tell whether transfer learning is suitable (or feasible) for a particular application.
    
    Transfer learning is an extremely useful technique in scenarios where we have (1) limited computational resources, (2) a relatively small dataset, and (3) a specialized application of a more general task, but there are limitations to its usage as well.
    
    \subsection{Advantages}
    
    \textbf{Resource Saving and Efficiency.}     Transfer learning allows us to take existing models that require vast amounts of computational resources, and reuse these for similar applications. This helps to save resources, reducing the associated cost, time, and power consumption of training models. As a result, efficiency is improved because resources can be diverted towards more productive activities
    
    \textbf{Synergistic Model Performance.} In situations where the initial and target model have a high degree of similarity, model performance can exceed the performance of a model trained from scratch on the transfer learning dataset. However, it should be noted that the improvement may only be marginal in some cases (or negative in the case of negative transfer) if the initial and target models are not strongly related.
    
    \textbf{Small Datasets.} Transfer learning allows large neural network models to be trained with relatively little data since we are more so tuning the model as opposed to training it directly. However, generally, it is true that better performance will be achieved if a larger number of data points is used since standard error scales inversely with the number of data points.
    
    \subsection{Limitations}
    
    \textbf{Negative Transfer.} A major limitation of transfer learning is the problem of negative transfer. If the tasks of our initial and target models are sufficiently similar, transfer learning should provide results superior to a model that is initialized with random weights instead of the weights from a network already trained on a similar task. However, in some cases, transfer learning can lead to worse results than a model trained using random weights, which is often referred to as negative transfer. This tends to occur if the tasks are sufficiently dissimilar that the structural information encoded in the initial network is detrimental to the target model’s predictions. Currently, there is no clear way to determine what the conditions are to ensure that negative transfer does not occur.
    
    \textbf{Data Suitability.} While one of the benefits of transfer learning is that we can build a large machine learning model with only a small amount of training data, this is also one of the disadvantages. If our dataset is too small, and our model is too large, we run the risk of overfitting our model to the small amount of data, meaning we will achieve a poor performance at inference time. The number of data samples needed to minimize this possibility is context-dependent, but a helpful rule of thumb is that several hundred samples should be used in a transfer learning example at a minimum.
    
    \textbf{Constrained Architecture.} Perhaps the most important limitation of transfer learning is that of a constrained architecture. Since we are basically importing model weights from a previously trained network, we are constrained to use the same architecture as the previous model. If we start to add or remove layers, we will likely distort the network performance to its detriment.
    
    Transfer learning is a powerful resource to have at your disposal, but only if it makes sense for a given application and the architectural constraints are not prohibitive. Otherwise, you may have to resort to an alternative technique, such as training your model from scratch!

 
 \section{Continuous Training Metrics}
 
     0:00 / 0:00Klicken Sie auf die Pfeiltaste nach oben, um das Geschwindigkeitsmenu aufzurufen. Regeln Sie dann mit den Hoch und Runter Pfeiltasten die Geschwindigkeit. Klicken Sie Enter, um die jeweilige Geschwindigkeit festzulegen.Geschwindigkeit1.0xKlicke auf diesen Knopf und diess Video stumm oder nicht stumm zu schalten, oder drücke die HOCH oder RUNTER Knöpfe um die Lautstärke zu erhöhen oder abzusenken.Maximal Lautstärke.Videotranskript
     Beginn des Transkriptes, zum Ende springen.
     VIJAY JANAPA REDDI: An important component of continuous training
     are the metrics that are used to evaluate the models.
     These metrics can tell us when a model is
     ready to be deployed into production, and they can also
     be used to tell us when to retrigger another round of training
     when the metrics fall below a certain threshold.
     Model evaluation is primarily handled by ML engineers,
     and metrics are something that is universally understood.
     Or you want to make sure that the right metrics are universally
     understood, both, for example, being a ML researcher
     and an ML engineer who's putting things into the continuous training pipe.
     The most basic metric, for instance, is accuracy.
     This metric simplifies the performance in the model
     to a single metric that can be used to show improvements over previous models,
     or as a guide when training is completed.
     Now, test set accuracy is far more important than training set accuracy,
     as test accuracy gives us an idea of how well the model is really performing
     in the wild on unseen data.
     I'm assuming you know this.
     Nonetheless, refreshers are always a good thing.
     Now, having use case appropriate metrics for model evaluation
     can substantially influence the overall ML application design process.
     Consider what we've been talking about with respect to keyword spotting.
     Keyword spotting can actually broken down a little bit further
     into subcategories, like wakeword detection,
     for example, when you just say a wakeword like Alexa
     and then nothing else, and then the machine just wakes up.
     Now, we're assuming that is a very nice and clean input
     signal to the machine learning model.
     But what happens if you're doing keyword search in place of wakewords.
     On Its surface, both of them are looking like wake words,
     but the key challenge in the keyword search
     is that the audio is continuously streaming.
     And that's like trying to find a needle in a haystack.
     So here, you need a slightly different metric
     that will help you faithfully address whether the model is
     capable of keyword search or not.
     A very basic example of a domain specific metric
     here is streaming accuracy, which refers to the accuracy of a model at audio
     that is being streamed in it over time, instead
     of having individual clips that are nicely
     cropped to it, where the model is expecting just a nice one-second clip.
     Instead, the audio streaming case, it's continuously
     seeing an audio signal going through it.
     So streaming accuracy can actually reveal some critical issues
     with the model, where standard accuracy, as in the case of wakewords,
     would totally miss that.
     So for example the model might incorrectly
     trigger on an ending syllable of an unknown word, which
     would be bad because then you have this machine constantly waking up
     and saying what would you like me to do.
     Now, an extension of streaming metrics is really precision recall.
     Precision tells us what percentage of the detections
     were actually the keyword, and recall tells us
     what percentage of the utterances of that keyword did we really detect.
     Now, let me give you a crazy example where it's extremely important
     to get this kind of a thing right.
     Now, let's say, for instance, you're trying to keep us safe,
     and you can have a high recall by calling everyone
     who's taking a flight a terrorist.
     Even if there's only one in a million, you have 100% recall.
     But imagine the consequences of that, despite your well intentions.
     All of the flights will be delayed, and virtually, no one
     will ever take off, as we don't trust anyone to get on a plane.
     In such a case, you also need high precision.
     How accurately are you able to pinpoint someone is
     a terrorist only when they truly are?
     So I strongly encourage you to understand the nuances of these metrics
     because if you don't understand the metrics,
     you're going to be ending up doing bogus evaluations.
     So precision, recall, AUC, accuracy, F1 scores
     are all highly important metrics, and you should know when to use them.
     I'll give you some pointers later on so you can kind of go and explore them
     a little bit further.
     Now coming back to the point, in this slide here,
     I'm showing you some false positives and false negatives
     in a continuous streaming signal.
     As you can see, there are times when we actually correctly detect the word,
     and there are times when we don't correctly detect the word.
     And in order to kind of take a big macro level view of this,
     you need to understand how precise are you being
     and what's your recall capability.
     And that often combines into what we call
     as a F1 score that allows us to have a confidence in both precision
     and recall because the metric combines both of them
     to tell us how well the keyword spotting model is actually doing.
     Now, recall this particular pretrained model
     from the previous video, whereas telling you that we only need five samples.
     So we have a embedding representation, which we then
     use, along with transfer learning, to quickly get the model
     to recognize a new language, or recognize
     a handful of words in new languages.
     Well here, I'm showing you the data of what
     happens when you don't use the right metric to actually assess
     the keyword spotting model.
     You don't have to understand everything in this graph.
     The only thing that you need to know is that the line should ideally
     be closer towards the origin of the graph, which is towards the zero, zero.
     And that would mean that you have low false rejection rates
     and false acceptance rates.
     Now, if you want to improve the accuracy of the keyword spotting model
     for streaming audio content, then you can provide it
     with some sort of context, as shown on the right.
     And this context can help boost the streaming accuracy of the model.
     So assuming you have provided this context to the model,
     now, you can see that the bars have all effectively move closer
     to where it's zero, zero.
     This shows that for many different languages,
     the keyword spotting model is actually performing better with context.
     We can also ensure that wakeword performance is not
     impacted in this particular scenario, using additional contextual data,
     but that's a different story.
     In summary, what I'm trying to get to is that although accuracy
     is a well-understood metric, it often fails
     to connect model's performance to more user centric, or use case specific,
     kind of metrics.
     And often when you're talking about continuous training,
     these kinds of nuances are extremely important because they help you
     determine when to refigure training.
     With that said, this effectively wraps up
     the core part of continuous training.
     And I will summarize this in the context of the rest of the MLOps pipeline next.
     I'll see you then.
     Ende des Transkriptes, zum Anfang springen
     Downloads und Transskripte
     Video
     Videodatei herunterladen
     Transskripte
     Download SubRip (.srt) file
     Download Text (.txt) file
 
\section{Forum discussion on Metrics for Continuous Training}

Is there a tinyML application in which you think continuous training would be useful?

 What metrics would you use to trigger continuous training?
 
  What metrics would you use to evaluate its results? 
  
  
\section{Continuous Training Impact on MLOps}

    0:00 / 0:00Klicken Sie auf die Pfeiltaste nach oben, um das Geschwindigkeitsmenu aufzurufen. Regeln Sie dann mit den Hoch und Runter Pfeiltasten die Geschwindigkeit. Klicken Sie Enter, um die jeweilige Geschwindigkeit festzulegen.Geschwindigkeit1.0xKlicke auf diesen Knopf und diess Video stumm oder nicht stumm zu schalten, oder drücke die HOCH oder RUNTER Knöpfe um die Lautstärke zu erhöhen oder abzusenken.Maximal Lautstärke.Videotranskript
    Beginn des Transkriptes, zum Ende springen.
    VIJAY JANAPA REDDI: Congratulations on wrapping up continuous training.
    As you might recall, continuous training has
    got a large number of moving pieces.
    By and large, you'll tried to break them down into three core components--
    the data processing engine, the model training engine,
    and the model evaluation engine.
    Now, the data processing engine-- the way
    we talked about it was by talking about how we can create an automated data
    engineering pipeline where we took the ability to ingest audio and transcript
    data to automatically create keyword data sets.
    Now, you could imagine that instead of just focusing on keywords,
    you could possibly apply it for vision tasks.
    So there are many different use cases you
    might be able to apply such a fundamental principle to,
    but the intent there was to help you understand
    how you can have an automatic data generation pipeline so that you
    can keep revising your models quickly and efficiently in short light cycles.
    Then we talked about the model training engine.
    Specifically in the model training engine
    I introduced the notion of AutoML and neural architecture search.
    And the reason I did that was because I want
    you to have an appreciation for the fact that there
    are a rich number of hyperparameters that are possible
    and that some configuration of them will lead to better or worse models.
    And in the context of TinyML and embedded use cases,
    you want to really try and achieve that optimal model.
    But, of course, you have to learn how to navigate
    that search space efficiently and, therefore, we dug deep
    into neural architecture search.
    Finally, we talked a little bit about really
    evaluating the metrics that are coming out of the continuous training pipe.
    You want to evaluate both the model and you
    want to evaluate the continuous training pipeline itself.
    Now, taking a step back up, though, hopefully you
    see how the continuous training part is really critical portion of the MLOps
    lifecycle.
    More specifically, you have to think about what
    is the role of continuous training in the context of this MLOps.
    Now, they're clearly interdependent relationships
    between all of these different stages with continuous training.
    However, the one that is of most importance for me
    is really emphasizing that how continuous training
    and continuous monitoring are tethered together.
    If you recall, we started talking about continuous training
    by talking about drift, talking about model drift,
    and talking about data drift, and so forth.
    We talked about changes in the schema and the way the scheme is skewed
    can happen, the way the value skew can actually happen.
    All of that are things that are part of the continuous monitoring pipe.
    In the continuous monitoring pipe, we need
    to build all this infrastructure that will actually
    be monitoring for various kinds of drift so that when it does detect a drift,
    it immediately comes back and tells us that we
    need to re-trigger the continuous training pipe so we can actually push
    a new updated model into the field.
    Sometimes these models are pushed automatically into the field
    and are other times they gated by someone
    who has to sign off to launch the model into production.
    Nonetheless, as always, there are multiple people involved
    in trying to get all of this going.
    The data engineer helped us with the data processing engine.
    The ML researcher is critical in the context
    of building up the AutoML pipeline.
    The ML researcher or the ML engineer is typically responsible
    for making sure all of this work smoothly
    and that the model that's coming out at the end of it
    is indeed ready for production deployment.
    With all of that said, I'll see you in the next video.
    Ende des Transkriptes, zum Anfang springen
    Downloads und Transskripte
    Video
    Videodatei herunterladen
    Transskripte
    Download SubRip (.srt) file
    Download Text (.txt) file  
 
\section{Optional: Multilingual Spoken Words Colab}

Now that we’ve explored continuous training and the multilingual keyword spotting example in detail, it's your turn to get to try out the process. In this Colab you’ll get work through the process used in the MKWS paper to explore embeddings and train your own model on a subset of the dataset. Follow this link to get started:

\url{https://colab.research.google.com/github/harvard-edge/multilingual_kws/blob/main/multilingual_kws_intro_tutorial.ipynb }
 
\section{Overview of Model Conversion
    The MLOps pipeline. Going from ML Development to ML training to Continuous Training to Model Deployment to Prediction Serving to Continuous Monitoring. Along the way artifacts are produced at each step including: Code and Configurations, Training Pipelines, Registered Models, Serving Packages, and Serving Logs. All steps tie into the overarching themes of Data and Model Management.
    
    \subsection{Objectives}
    
    While running standard Tensorflow, Pytorch, or Keras models in Colaboratory is very useful on a desktop or laptop computer, they are not very useful when it comes to microcontrollers - the models are just way too big! These frameworks include many features and as such are very large. When running models only for inference, most components are non-essential and can be stripped out to save on memory and compute overhead. This is exactly what Google did with TensorFlow, culminating in TensorFlow Lite (TFLite) for Microcontrollers.It is important to understand what the differences are between various frameworks and optimizations schemes.
    
    The model conversion procedure is not particularly complicated, but it is one of the most important aspects of TinyML. Certain inference libraries may only be able to work with certain file formats, or may perform inadequately due to poorly optimized operations.
    
    In this section, we will delve into some of the particularities of the model conversion process:
    
    \begin{itemize}
      \item  ML Frameworks
        \begin{itemize}
          \item TensorFlow vs. TensorFlow Lite vs. TensorFlow Lite Micro and others
          \item Understand how inference frameworks differ from one another and why that’s important to keep in mind
        \end{itemize}
      \item Common Optimizations
        \begin{itemize}
          \item Pruning, Quantization, Knowledge Distillation, Joint optimization, etc.
        \end{itemize}
      \item Role of Optimizations
        \begin{itemize}
          \item Why these optimizations are so important for TinyML
        \end{itemize}
    \end{itemize}
    
    \subsection{Motivation}
    
\GRAPHICSC{1.0}{1.0}{TinyML/4 MLOps/covert_model.png}
    
    The core runtime for TFLite Micro is only around 16 kB, several orders of magnitude smaller than that of TensorFlow. However, despite the fact that TFLite Micro is derived from TensorFlow, it cannot natively run TensorFlow models. In order for us to run our model on-device, we must go through a model conversion process.
    
    Model conversion occurs when transferring between representations of a neural network model. For example, we could convert between a Pytorch and TensorFlow model, or between a Keras model and one cast in the ONNX format. Each of these formats has been built independently and optimized for a specific purpose. TFLite Micro was built with resource-constrained microcontrollers in mind, and is optimized for a tiny memory footprint and highly efficient computation, but is stripped of all unnecessary functionality. This includes functionalities that we might be very used to using, such as debugging!
    
    \subsection{Conversion}
    
    To convert our model from a TensorFlow format to TFLite, we must use the TFLite Converter Python API. This API converts our model into a FlatBuffer, which is an efficient serialization format for performance-critical applications. In this data format, computations are a lot more efficient, and it is optimized for statically typed languages (whereas JSON is often the format of choice for dynamically typed languages such as Javascript), making it a good choice for our resource-constrained microcontrollers which often utilize C, which is a statically typed language.
    
    OK, now we have our TFLite model. However, the model conversion does not stop here. Because most microcontrollers do not natively contain filesystem support, we cannot copy and paste our file onto our device, instead we must embed the model in the microcontroller. This is typically done by converting it to a C array (using the xxd command) and then compiling it into our program.
    
    \subsection{Optimizations}
    
    Are we done now? Not if you want the job done properly. We still have to include model quantization and model pruning! Model quantization involves taking the weights of our trained neural network and changing them into a different numeric representation. Typically, this is from 16-bit floating point (FP16) or 32-bit floating point (FP32) to an 8-bit representation. Some applications have even gone so far as to develop binary neural networks - where the neural weights are all represented by a single bit. Using a simplified numeric representation can have a profound effect on the model size. Model pruning is a similar method for reducing the size, but instead of altering the numeric representation, simply cuts out weights that are not that important to us. For example, all weights that are somewhat close to zero, it might just set to zero and ignore them, somewhat similar to how sparse matrix calculations simply ignore zero values and just focus their attention on the non-zero values. 
 
\section{Model Conversion}

    0:00 / 0:00Klicken Sie auf die Pfeiltaste nach oben, um das Geschwindigkeitsmenu aufzurufen. Regeln Sie dann mit den Hoch und Runter Pfeiltasten die Geschwindigkeit. Klicken Sie Enter, um die jeweilige Geschwindigkeit festzulegen.Geschwindigkeit1.0xKlicke auf diesen Knopf und diess Video stumm oder nicht stumm zu schalten, oder drücke die HOCH oder RUNTER Knöpfe um die Lautstärke zu erhöhen oder abzusenken.Maximal Lautstärke.Videotranskript
    Beginn des Transkriptes, zum Ende springen.
    LARA SUZUKI: Hello.
    In your last lesson, we had a closer look at the continuous training
    stage of our MLOps framework.
    We learned how to repeatedly execute the training pipeline in response
    to data changes, code changes, or when a schedule potentially
    with the new training settings.
    It's now time for us to move to the next stage of our MLOps framework
    that is of convert model.
    In this session, we will cover deeper details
    of creating a model that is catered to run on resource-constrained devices.
    Imagine embedding machine learning into appliances
    that can adapt to a daily routine and are
    industrial sensors that will understand the difference between problems
    and normal operation.
    For the purpose of this session, we will cover TensorFlow Lite
    and TensorFlow Lite Micro.
    In summary, TensorFlow Lite is an optimized on-device machine
    learning framework.
    And it addresses five key constraints.
    The first one is latency--
    there's no round trip to a server.
    Second, privacy-- there's no personal data leaving the device.
    Third, connectivity-- internet connectivity is not required.
    Fourth, size-- reduced size model and binary size.
    And lastly, power consumption-- efficient inference
    and a lack of network connections do not impact the system.
    TensorFlow Lite is a set of open source tools
    that enable on-device machine learning by helping developers
    to run their models on mobile-embedded and IoT devices.
    TensorFlow Lite Micro, on the other hand,
    is an open source machine learning inference framework
    for running deep learning models on embedded systems.
    Besides tracking efficiency requirements,
    this framework solves for fragmentation challenges that
    hinder cross-platform interoperability.
    Microcontrollers are typically small low-powered computing devices
    that are embedded in within hardware that requires basic computation.
    By bringing machine learning to tiny microcontrollers,
    we can boost the intelligence of billions of devices
    that we use in our lives, including household appliances
    and the internet of things devices.
    And that is all done without relying on expensive hardware or reliable internet
    connections.
    And that is often subject to bandwidth and power constraints and results
    in high latency.
    It can also help to preserve privacy, since no data leaves the device.
    TensorFlow Lite for microcontrollers is designed for specific constraints
    of microcontroller development.
    If we are working on more powerful devices,
    for example, an embedded Linux device like a Raspberry Pi,
    the standard TensorFlow Lite framework might be easier to integrate.
    To work with TensorFlow Lite, you can directly create a TensorFlow Lite model
    and training workflow, or you can convert a TensorFlow model
    into a TensorFlow Lite model.
    You can simply use the TensorFlow Lite converter
    and convert a TensorFlow model to a TensorFlow Lite model.
    During conversion, you can apply optimizations such as quantization
    to reduce model size and latency with minimal or no loss in accuracy.
    In the case of TensorFlow Lite micro, after converting a model
    to TensorFlow Lite, you must convert to a C byte
    array using standard tools to store it in a read-only program
    memory on device.
    The core MLOps capabilities used in this stage
    are model registry, machine learning metadata, and artifacts repository.
    Ende des Transkriptes, zum Anfang springen
    Downloads und Transskripte
    Video
    Videodatei herunterladen
    Transskripte
    Download SubRip (.srt) file
    Download Text (.txt) file 

\section{ML Frameworks \& The Lay of the Land}
    0:00 / 0:00Klicken Sie auf die Pfeiltaste nach oben, um das Geschwindigkeitsmenu aufzurufen. Regeln Sie dann mit den Hoch und Runter Pfeiltasten die Geschwindigkeit. Klicken Sie Enter, um die jeweilige Geschwindigkeit festzulegen.Geschwindigkeit1.0xKlicke auf diesen Knopf und diess Video stumm oder nicht stumm zu schalten, oder drücke die HOCH oder RUNTER Knöpfe um die Lautstärke zu erhöhen oder abzusenken.Maximal Lautstärke.Videotranskript
    Beginn des Transkriptes, zum Ende springen.
    In this section, I'm going to peel some of the layers
    off the model conversion stage of the MLOps pipe.
    Model conversion and optimization is an extremely important part
    of deploying ML models into resource-constrained devices.
    Now normally, this is part of the ML deployment
    process, which as you can see, is basically the very next stage.
    However, we intentionally chose to pull it out,
    because we want to highlight its significance.
    So in this section, I want to cover the following topics.
    First and foremost, I want to introduce you
    to the rich landscape of ML frameworks, which I'll do in this video.
    Then next, I want to talk a little bit about the differences
    between TensorFlow, TensorFlow Lite, TensorFlow Lite Micro, and then kind
    of get into some of the optimizations that someone like an ML engineer
    would really be responsible for doing.
    Now there are a lot of issues around this.
    And that's why I intentionally want to spend a bit of time talking about this.
    So you actually have an appreciation for what your ML engineer is
    going to have to go through when they're getting ready to actually push
    a model into the embedded ecosystem.
    Let's start by understanding the big picture,
    the lay of the land for different ML frameworks.
    Why, you might be wondering?
    Well, the simple reason is that when you're deploying models,
    you want to understand the different trade
    offs between these various systems.
    Now when it comes to ML frameworks, they're like ice cream flavors.
    They're rich and wide variety of flavors out there.
    So it's important to know how they're similar
    and how they're actually different.
    Now, for instance, we have the wildly popular TensorFlow framework,
    which was developed by Google and released to the public in 2015.
    When Google developed it, the original intent
    was to create an open source machine learning framework
    in order to train models.
    Today, of course, it's evolved into a very mature deep learning
    framework that supports a wide array of both training and inference
    capabilities, as well as a lot of strong visualization capabilities
    in order to enable responsible AI.
    Now, similar to them, there's PyTorch, which
    is another open source machine learning library
    that's based on the Torch library.
    That was originally developed by Facebook's AI research lab.
    For all practical purposes, you can say it's
    a very strong competitor of TensorFlow.
    Compared to TF, however, PyTorch is still a relatively young framework.
    However, it's got a very strong community base out there
    and it is much more friendly to Python in terms of programming.
    It's gaining a lot of popularity primarily because of its simplicity,
    ease of use, dynamic computational graphs, memory efficiency,
    and so forth.
    Some of which, as we know in the context of embedded systems,
    is actually going to be very clean.
    However, PyTorch does not quite yet support an embedded version
    that allows us to target TinyML.
    But I would only expect that it's a matter of time.
    Finally,
    I want to touch a little bit on Keras, because Keras was originally
    created to be a friendly modular and easy-to-extend sort of API
    interface that sits on top of Python.
    So you can actually allow programmers to focus
    on what they're actually trying to do in terms of programming neural network
    graphs.
    And effectively reduce the cognitive load of a lot of the details
    that come about when you're actually programming directly inside TensorFlow.
    That was written in Python and it supports
    a large number of different multiple back ends,
    however TensorFlow is the primary interface for Keras.
    So that's why I'm bringing it up.
    Because very often, more often than not, you're
    actually going to see Keras popping up all the time when you're
    talking about models in TensorFlow.
    Now here's a very high-level summary that puts together the differences
    between Keras, TensorFlow and PyTorch. , Again just to kind of quickly help you
    understand the Delta.
    I'm not going to focus on every single cell.
    I'll encourage you to kind of just pause and take a look at this table.
    But briefly, Keras has got very friendly API, similar to PyTorch,
    which is also known for having good APIs.
    Keras is just a high-level API, while TensorFlow and PyTorch
    are actual frameworks with low-level API capabilities.
    I would say that a very good competitor to Keras plus TensorFlow
    is really just PyTorch, because people often
    say that it's very easy to make changes in the PyTorch environment.
    It's especially popular in the research communities.
    But nonetheless, people do use Keras with TensorFlow.
    Last but not least, the debugging capabilities
    are really crucial when you're kind of building programming frameworks.
    And, of course, the debugging capabilities vary a little bit.
    Nonetheless it's actually quite key for the ML engineer who's actually
    responsible for the deployment process.
    Now often most people will think that during inference,
    the frameworks are virtually identical, right.
    Have already trained the neural network and I'm going to try and deploy this.
    So why should I really worry about the frameworks?
    Well it's actually not true that the frameworks
    are going to run any neural network graph in a completely identical manner.
    For example, let's consider this image down here
    where we're trying to do object detection
    and draw a box around the actual object of interest.
    Now, this needs the support of non-maximal suppression,
    where the input is a list of proposal boxes as shown on the left,
    with some sort of overlap.
    And on the right, you're seeing a filtered list
    of proposals, which clearly tries to identify where exactly the object is.
    Now the way this would normally work is that you
    select some sort of a box with the highest objectiveness score,
    and then you compare the amount of overlap between the boxes,
    and then you use some sort of arbitrary threshold
    in order to remove the bounding boxes with the most amount of overlap.
    And then you repeatedly kind of prune this process.
    And finally, you'll end up with an image on the right.
    As you can see, as you go from the left to the right,
    you basically have iterative pruning happening.
    This is a very common process literally in every object detection task.
    TinyML or otherwise, you're going to end up with something like this.
    And you'll find that when you actually look
    into the code of between like say TensorFlow and PyTorch, both of them
    provide support for NMS as a function call.
    But the reality is that the actual way these frameworks implement
    are not quite identical.
    In other words, NMS operation inside TensorFlow
    runs a little bit differently from the NMS in PyTorch.
    This can lead to actual accuracy differences,
    even though you're running the exact same network.
    Because the code is not just within the network,
    it's also part of the framework when it's calling NMS.
    So that's why you have to kind of have an appreciation for some
    of these subtle deltas, because people often
    think like, you're going to get the exact same result.
    And you don't, just because of the framework you're using.
    Now that's just for two frameworks.
    The reality is there are many other frameworks out that we have.
    Caffe, MXNet, Scikit, Learn and many others.
    When in fact you actually consider all the proprietary in-house
    frameworks that a lot of the companies are developing.
    So one of the major challenges that's kind of emerged in this space
    is how do you handle all these different frameworks
    in an open and interoperable way?
    There's a real need for greater interoperability in AI tools,
    because you have to support all these deployment
    targets as shown on the right.
    So while people are working on great tools,
    developers are often locked into a framework or a certain ecosystem,
    because things are not easy to share.
    So the whole goal of ONYX was to see if there's
    some way we can actually share models.
    In other words, it might be coming from any different framework,
    but as long as the graph itself is an ONYX neutral sort of a format,
    these frameworks should be able to run them on these target devices.
    Put another way, this is the way you should envision it.
    The goal of ONYX is to provide a definition
    of extensible computational graph models as well
    as definitions for built-in operators and standard data types
    that networks will use.
    And these operators are going to be implemented externally to the graph,
    but the set of built-in operators are going
    to be portable across frameworks, which is the key thing.
    Now every framework supporting ONYX will then
    provide some sort of implementations for these operators
    on the applicable data types.
    In this way, it allows operators to support
    a rich array of different hardware solutions,
    such as what we're seeing at the bottom layers here of CPUs, GPUs, FPGAs
    and various other custom processing units.
    In an ideal world, if everything was to work, great.
    One would be able to take a trained network, say, coming out of PyTorch
    then easily export it to ONYX and then import it into the TF Lite Runtime,
    for instance, and then be able to run it without any issues like NMS
    for instance.
    But in practice, that is far from the truth.
    Often, what you find is that while ONYX is a great effort,
    it's still fairly nascent, so you end up with portability issues
    where some of the frameworks are not equally supporting all the ops.
    The frameworks may also implement certain library functions
    that are not standardized upon.
    So this leads to a lot of this questioning about
    whether a certain network is really going to be able to run neutrally
    through the ONYX format.
    So they're always going to be differences.
    And this is why, for better or for worse,
    most developers end up picking a single ecosystem.
    And my recommendation in large, independent of TensorFlow,
    is that you'd want to pick a single open-source ecosystem.
    Because that helps with the portability and neutrality.
    So for instance here, I'm showing you if you stay within the TensorFlow
    ecosystem, then you can reap a large number of the benefits that
    come from a lot of the prior work.
    Like for instance, you can go from leveraging all of that TF training
    toolchain that you have to reuse in the TF Lite Explorer, which was already
    designed for exporting to mobile devices, which was then later extended
    to support embedded devices.
    So you can actually get a lot of the optimizations that were done for free.
    And a lot of the debugging support that comes already pre-built in the existing
    toolchain, you can reuse that.
    The reason I'm mentioning this is because most people will often
    get very excited about having a custom neural network engine
    specialized for deploying to embedded devices,
    but in reality this can become a royal pain.
    Because it's great to support a handful of operators,
    but it's really hard to support a rich and diverse set of models, all of which
    are going to be using very different operators.
    And you might not have all the necessary support built in.
    And that often becomes a major issue when
    you're trying to scale deployment out.
    I can't tell you how important this is, because I've actually seen it
    in practice where it's easy to come up with a small little engine that's
    efficient on a device, but the second you try
    and target more and more devices, things start breaking down.
    So therefore, you want to pay attention to this.
    I actually have a write up that talks a little bit more in detail
    about these kinds of issues.
    And I want you to keep that in mind as we move forward.
    Coming up next is going to be talking about TensorFlow versus TensorFlow Lite
    versus TensorFlow Lite Micro just so you know what your ML engineer is actually
    going to deal with when it comes to deploying the models while they're
    doing the conversion of the models.
    I'll see you there.
    Ende des Transkriptes, zum Anfang springen
    Downloads und Transskripte
    Video
    Videodatei herunterladen
    Transskripte
    Download SubRip (.srt) file
    Download Text (.txt) file
 
\section{TF vs. TFLite vs. TFLite Micro}

    0:00 / 0:00Klicken Sie auf die Pfeiltaste nach oben, um das Geschwindigkeitsmenu aufzurufen. Regeln Sie dann mit den Hoch und Runter Pfeiltasten die Geschwindigkeit. Klicken Sie Enter, um die jeweilige Geschwindigkeit festzulegen.Geschwindigkeit1.0xKlicke auf diesen Knopf und diess Video stumm oder nicht stumm zu schalten, oder drücke die HOCH oder RUNTER Knöpfe um die Lautstärke zu erhöhen oder abzusenken.Maximal Lautstärke.Videotranskript
    Beginn des Transkriptes, zum Ende springen.
    VIJAY JANAPA REDDI: In this video, I really want
    to talk about what the differences are for different machine learning
    frameworks within the same family.
    For example, there's TensorFlow, then there's TensorFlow Lite,
    and then there's TensorFlow Lite Micro.
    And I want to talk about what the primary differences are
    between these different frameworks.
    Now, TensorFlow can run on big beefy machine learning systems
    like Google's CPUs, and NVIDIA's GPUs, or just Intel's CPUs, and so forth.
    But at the same time, there's a slightly lighter version
    that's capable of running models and mobile devices,
    which is what TensorFlow Lite is.
    And there is an even more lighter version
    that can run models on Android devices, which is what TensorFlow Lite Micro is.
    Now, if you're dealing with anything around embedded systems,
    you are hands down using TensorFlow Lite Micro.
    Let's try and understand how we can go all the way from the left to the right,
    and what's the typical process.
    Generally, we want to train a neural network,
    and that's going to start with training networks using TensorFlow, which
    has a lot of support for doing things like back propagation
    or doing a lot of distributed scheduling in order
    to make the most of the resources, both within a single system as well
    as multiple systems.
    Once the network is effectively trained using TensorFlow,
    then you actually want to strip it down--
    strip the network down into its core components, which is often,
    in the case, just a very simple TFLite file.
    This TFLite file is both lean and mean so
    that it can actually be efficiently deployed onto embedded devices.
    Even a smartphone would actually use a TFLite.
    And while a smartphone is just a glorified embedded system,
    it's much more capable than our real TinyML systems.
    And of course, you would typically take this TFLite file,
    and then using some sort of format, like a flat buffer format, which
    is easily mapped into memory, can actually
    be deployed onto embedded devices.
    And then, of course, you end up feeding the network with some data that's
    coming in from sensors and so forth.
    This is the general process.
    Now, let's talk a little bit about what are
    the key differences from a technical standpoint within TensorFlow,
    TensorFlow Lite, and TensorFlow Lite Micro.
    I would say that, as of today, this is largely
    where the primary differences sit.
    TensorFlow is great for training, but a lot of it
    is just blown when it comes to inference-- for instance,
    when you're actually deploying the models to mobile and embedded devices.
    This is why TensorFlow Lite and TensorFlow Lite Micro
    were actually born.
    And in those systems, you don't see any support for training,
    because you won't be doing training.
    You're just doing inference.
    You're never going to be really training a big model on an embedded device.
    So that makes the code really lean and mean when you're actually
    taking out all the training bloat.
    Another main difference between the three different frameworks
    is the number of ops they support.
    As you can see, for instance, TensorFlow supports a lot of operators,
    and that's because researchers are often experimenting with different operators.
    So they might be experimenting with operators just used once,
    or some operators are widely supported.
    In reality, though, many of the operators,
    of the 1,400 operators inside TensorFlow,
    are actually quite rarely used.
    That's why, when you go to mobile and embedded devices
    for TensorFlow Lite and TensorFlow Lite Micro, respectively,
    you see that the number of ops drops down to 130 and 50.
    Now, when it comes to optimizations, the latter two-- that is,
    TensorFlow Lite and TensorFlow Lite Micro--
    support some sort of specialized optimizations.
    I'm actually going to talk about these in a separate video in much
    more detail.
    Now, on the software front, you can see the TensorFlow and TensorFlow
    Lite rely on operating support--
    things like Windows, like Linux.
    Or in the case of TensorFlow Lite, it's going to be Android or iOS.
    Now, you can imagine things like memory allocation
    and so forth are high-level functionalities
    that you can actually take for granted on those kinds of mature operating
    systems, but that's totally not the case when it comes to TensorFlow Lite Micro.
    Sometimes, you don't even have RTOS.
    So you can't even assume base metal, sort of basic operating system
    support on an embedded device.
    So in a way, if you want to think about scaling things out,
    you have to be independent of that real-time operating system support.
    Another major difference is that TensorFlow and TensorFlow Lite support
    accelerated delegation, meaning that they can actually
    schedule the neural network code to different accelerators on the system,
    be it your CPU, GPU, or FPGA, and so forth,
    while TensorFlow Lite Micro doesn't really
    need that kind of specialized scheduling code which takes up memory,
    because in embedded devices, you don't really
    have that much flexibility at the lowest level of the hardware.
    So again, not supporting this functionality is not a bad thing.
    In fact, it's a good thing, because it allows
    us to keep the inference framework really lean
    and mean, and keep the research requirements to a bare minimum.
    So that's what really allows us to have very small binary footprint.
    So TensorFlow Lite Micro, which runs on TinyML systems,
    is at least two orders of magnitude smaller than TensorFlow
    because we don't have all of that training support,
    and we have taken out some other goo.
    Now, if you recall, in TinyML systems, memory constraints
    are important both from a flash perspective as well as
    the runtime random access memory perspective.
    So TensorFlow Lite Micro is very cognizant about those kinds
    of optimizations, and so we actually have to take out a lot.
    Just to wrap things up here, what I'm touching on
    is the fact that, between this video and the previous video,
    there are lots of different frameworks you have,
    both for training and inference.
    And there are actually major differences when
    it comes to deployment, even within a single framework,
    as we saw in this particular video.
    And, while people have been trying to build open standards like ONNX,
    they are not a panacea because they do not yet
    have the level of maturity that allows us to be platform neutral
    and be easily able to port across multiple devices.
    With that said, in the next couple of videos, what I want to do
    is I want to talk about some optimizations that are bread
    and butter for embedded deployments.
    Ende des Transkriptes, zum Anfang springen
    Downloads und Transskripte
    Video
    Videodatei herunterladen
    Transskripte
    Download SubRip (.srt) file
    Download Text (.txt) file  
 
  
\section{TFLite Micro for TinyML}


\GRAPHICSC{1.0}{1.0}{TinyML/4 MLOps/5-6-4_Reading_Image.png} 

Diagram showing a server, phone, and microcontroller for TF, TFLite, and TFLite Micro.

\subsection{Motivation \& Background}

TensorFlow, developed by Google, is a very popular machine learning framework for constructing and running machine learning models, and we have used it several times in this series of modules to construct our machine learning models. However, there are several downsides to TensorFlow that make it unsuitable for use in microcontroller units (MCUs).

Most notably, it is just too darn big! For starters, downloading TensorFlow on your laptop or desktop computer will take up around 1 GB of hard disk space - on microcontrollers, we are lucky if we have 1 MB of space! Next, the RAM required to run even the most basic of models is too high for any MCU, which typically have RAM of 500 kB or less. TensorFlow also works with floating point arithmetic by default (FP16/FP32), which is not supported on most MCUs due to the lack of a floating point unit (FPU). Most MCUs only have access to an arithmetic logic unit (ALU) for performing 8-bit integer operations.

There is a happy ending to this story though. The TensorFlow team at Google recognised that there is a market for running machine learning models on resource-constrained devices that would require a more lightweight machine learning framework. Enter TensorFlow Lite (TFLite). TFLite is a framework for deploying ML models on multiple devices ranging from mobile (iOS and Android), desktop, and edge devices. An even more lightweight version exists specifically for microcontrollers, called TFLite Micro (sometimes called TFMicro).

\subsection{TFLite Micro Overview}

TFLite Micro is a barebones ML inference engine written in C++11 which acts as a stripped out version of the TensorFlow library, removing all unnecessary components such as stored models, debugging capabilities, and unnecessary operations. The team, headed by Pete Warden, managed to compress the core runtime into approximately 16 kB (on an ARM Cortex M3), which is sufficient to run on a large number of resource-constrained devices. This is incredibly impressive given some of the additional challenges associated with resource-constrained devices. Note that TFLite Micro is only an inference engine, and thus does not natively support on-device training.

The heterogeneity of embedded systems makes interoperability across the landscape of embedded devices challenging. As discussed previously, most MCUs cannot handle floating point operations. However, some can. Similarly, some work with 16-bit integers as opposed to 8-bit. Their architectures often vary, as do their peripherals and the presence of other computational components. This makes it incredibly difficult to design a framework that is configurable with all types of heterogeneity environments where devices have no common standard. When we download TensorFlow on our laptop or desktop computer, we have to specify what system we are using. Is it 32-bit or 64-bit? Is it running MacOS, Windows, or some form of Linux? Which version is the operating system? For microcontrollers, we come across the same issues, but the number of options is far greater due to the lack of a common standard. Currently, the only constraint imposed by TFLite Micro is the need for a 32-bit platform. The framework has been tested extensively on the ARM Cortex-M series of microcontrollers, and has also been ported to ESP32 and can be downloaded as an Arduino library. Boards currently supported by the framework include the following: 

\begin{itemize}
  \item Arduino Nano 33 BLE Sense
  \item SparkFun Edge
  \item STM32F746 Discovery kit
  \item Adafruit EdgeBadge
  \item Adafruit TensorFlow Lite for Microcontrollers Kit
  \item Adafruit Circuit Playground Bluefruit
  \item Espressif ESP32-DevKitC
  \item Espressif ESP-EYE
  \item Wio Terminal: ATSAMD51
  \item Himax WE-I Plus EVB Endpoint AI Development Board
  \item Synopsys DesignWare ARC EM Software Development Platform
  \item Sony Spresense
\end{itemize}


\bigskip
There is currently a trade-off between offering interoperability and performance. Because no industry standard currently exist, making TFLite Micro increasingly portable would require more code, which will inevitably make the library larger. Currently, the library only offers support for a limited set of devices (listed above), a limited set of TensorFlow operations (addition, subtraction, matrix multiplication, ReLU, etc.), and requires manual memory management (since the library does not utilize dynamic memory allocation).

Other TinyML frameworks are gradually coming into existence, most notably Pytorch’s GLOW which utilizes the GLOW ahead of time (AOT) compiler. TFLite Micro seems to be the dominant framework currently, but this could change in the future. For a deeper dive into TFLite Micro and its structure, capabilities, and limitations, we refer the reader to the origin TFLite Micro research paper available on arXiv.  


\section{Model Pruning}

    0:00 / 0:00Klicken Sie auf die Pfeiltaste nach oben, um das Geschwindigkeitsmenu aufzurufen. Regeln Sie dann mit den Hoch und Runter Pfeiltasten die Geschwindigkeit. Klicken Sie Enter, um die jeweilige Geschwindigkeit festzulegen.Geschwindigkeit1.0xKlicke auf diesen Knopf und diess Video stumm oder nicht stumm zu schalten, oder drücke die HOCH oder RUNTER Knöpfe um die Lautstärke zu erhöhen oder abzusenken.Maximal Lautstärke.Videotranskript
    Beginn des Transkriptes, zum Ende springen.
    VIJAY JANAPA REDDI: Starting with this video,
    I'm going to be talking about a range of different model optimizations
    that are commonly performed when we're trying to push models
    into deployment, especially at scale.
    Before we dive into the details, let's try
    to understand what are the goals of these optimizations.
    First and foremost, for instance, we might
    want to cut down the latency or the inference cost for embedded devices.
    Another thing is that, given that embedded devices have
    extremely severe resource constraints in terms of processing capabilities,
    memory availability, power consumption, network, and so forth,
    we might just want to squeeze every last bit of juice out of the models
    in trying to shrink them down while trying to keep the accuracy high.
    Now, later on, when we get into continuous monitoring, which
    triggers back into continuous training where we want to update the models,
    we also have to worry about what the payload size is.
    In other words, we're going to have to push
    a model for an over-the-air update, so you need to keep the model small down
    in that size.
    So that's another reason why you might want to, for instance,
    shrink the models down.
    So there's a wide range of reasons as to why we actually
    want to optimize the models to be really small.
    Now, this is the kind of thing that your ML engineer is really
    going to be concerned about.
    He or she is constantly thinking about how
    to do model conversion, so that they can actually push the models out
    efficiently and effectively into the field.
    Now, the area of model optimization can involve
    a variety of different techniques that include, but are not limited to, by
    and large, these three key things.
    One is you can reduce the parameter count.
    You can reduce the representational precision of the network.
    Or you can update the network topology to make
    it more efficient with a reduced number of parameters to make it run faster
    or whatnot.
    And there are a wide array of different techniques you can also use.
    You can use things like pruning, quantization, clustering, knowledge
    distillation, and so forth.
    I picked those to mention just because we're going to be talking about those,
    just so you have a flavor for what people do normally.
    We're going to start off by talking about pruning here.
    Pruning is fundamentally all about just taking a neural network
    and trying to make it sparse by "pruning," quote unquote.
    Cutting out some of the weights that are inside the neural network.
    Now, trying to look at it in a different way, at the end of the day,
    all these neural networks are just some representation of matrices.
    So either you have a dense matrix or you have a sparse matrix.
    A sparse matrix is a matrix with some weight values that are actually zero.
    So usually, in a real world neural network,
    you'll find that a large number of these matrices end up actually being zero.
    So every place you're using a dense matrix, for instance,
    or a dense tensor, you might actually be able to swap it with a sparse one.
    And that's the key thing here, because if you can do that,
    there are a certain amount of benefits.
    What do I mean by swapping?
    So here, I'm showing you those matrices at the bottom,
    and as you can see, as we go from the left to the right,
    between the first one and the third one, we're increasingly adding in zeros.
    Now, the pros of having lots of zeros is that you actually
    don't need to compute on them.
    In other words, you don't need to do the necessary gem operations
    and multiplications for these various sensors and matrices.
    And that means you gain both speed, in terms of processing capability,
    but you can also get size benefits, because you can actually
    shrink the network down.
    Which is really awesome, because you can run compression algorithms on it,
    doing all that to compact everything as I'm showing you here.
    So you take a sparse one and you relay it out so you can actually jam pack it,
    and then that becomes smaller and then you can efficiently
    run that on your target device.
    Now, the key question is, how do you actually
    go about pruning the neural network?
    There are lots of different ways.
    Now, I'm going to talk about things in the context of weight pruning,
    and as you can see, there are many different optimization methods here.
    We'll take a look at a few of them.
    Weight pruning basically means eliminating any unnecessary values
    in the weight sensors.
    We're wanting to set the values or the network parameters
    to zero to remove them so that are any unnecessary connections in the network
    will actually be eliminated.
    And this is done sometimes during the training process
    and sometimes you can also do it offline.
    So let's take some examples.
    Let's start with magnitude-based pruning,
    because this is the most simplest weight pruning algorithm.
    The intent here is simply to look at the salience of your link
    just in terms of its absolute weight.
    Now, through this method, what you want to do
    is that you want to set some sort of arbitrary threshold
    and you say that, OK, if the weight is not above
    or if it is below a certain weight, then you might want to just remove it.
    This technique, as simple as it sounds, is remarkably effective in practice.
    In fact, in the real world, you'll find that this model
    can lead to 6x better improvement in compression itself,
    with a minimum loss in the actual accuracy of the network.
    Remember, as you're taking weights out or you're
    trying to clip them down to zero--
    if it's 0.152 or something like that-- you convert it down to zero,
    you're going to lose some level of computational representation in there.
    But that's OK.
    Let's take a quick look at some other types of weight pruning methods.
    In contrast to magnitude-based pruning, let's take
    a look at sensitivity pruning.
    Sensitivity pruning is really simple.
    The idea is that, rather than arbitrarily
    picking one value that applies throughout the network for all
    the weights, what you're trying to do is understand
    each layer's characteristics.
    And then you can actually try and prune each layer a little differently,
    so you're sort of customizing a little bit based on the distributions
    that you actually see in the network.
    That's really all there is.
    So you end up doing a pruning sensitivity analysis first,
    and then you effectively look at some sort of a threshold
    or a sensitivity parameter that is tuned with respect
    to some standard deviation in terms of the weight distributions,
    and then you pick that arbitrary number and then
    you effectively go ahead and do magnitude-based pruning.
    It's just called sensitivity pruning because you're being more aware
    or it's more granular than magnitude pruning.
    Finally, let's take a look at structured pruning,
    which is actually quite different from the previous two
    that I just hinted at, which are just looking at the raw weights.
    So there's structured and unstructured.
    Let me start off with unstructured.
    So unstructured is effectively where the individual weight connections
    are removed from the network by sending them to zero.
    Now, pruning as we have talked about it with weight pruning or sensitivity
    pruning, is effectively a means of achieving unstructured pruning.
    This is obviously good.
    As I said, you can get up to a 6x compression in the model size, which is
    great for flash storage on the device.
    But you might also have structured pruning.
    With structured pruning, the idea is that you
    take groups of weight connections and you remove them all together.
    This can be done effectively at varying levels of granularity.
    So for example, you can do this at the channel level
    or you can do this with filters, or you can even do this with the filter
    shapes themselves.
    So there's a lot of flexibility down here.
    But in the end, what you're really doing in that is that you're undoubtedly
    improving the latency, because you're taking out big chunks of the processing
    that's required here.
    And this is a great way in order to effectively
    be able to improve performance across a wide range of devices,
    because this is something that's eliminating entire network connections.
    Now, undoubtedly, as you perform these optimizations,
    you're going to have an impact on accuracy.
    So what I'm showing you here is just how much of an impact
    you might have in terms of your accuracy.
    For example, with inception, what you see
    is that, even as you increase the amount of sparsity pruning,
    while your accuracy drops a little bit, it's not a huge drop.
    It's not a catastrophic drop.
    Well, it kind of depends on which use case you're going.
    If you're going after a medical application, well, in that case,
    this can be a catastrophic drop.
    But if it's just being able to tell if this is a cat versus a dog,
    it might not be such a big issue.
    So it's a little subjective.
    But nonetheless, the point is that the drop is not dramatic.
    It doesn't go down to something like 10%.
    And the second thing that I want to emphasize on this particular slide
    is that the amount of accuracy drop you have
    when you're doing these kinds of optimizations
    is really going to vary from one network to another.
    So for example, if you look at MobileNet, which is actually
    used quite widely in TinyML use cases, you'll
    find that the accuracy drop here is really
    quite different from how inception is behaving.
    So you need to keep that in mind.
    It's highly subjective based on the network characteristics.
    And the same is true whether you're dealing with image tasks or language
    tasks.
    So for example here, I'm showing you what the language translation
    models are going to do when you actually do pruning on them.
    What's crazy about the amount of sparsity optimizations you can do here
    is that you are able to push it to as much as 90%
    with a very minimal impact on accuracy.
    Which in this case, is measured in terms of a blue score.
    Blue score is effectively a metric that allows
    you to evaluate machine translation text, similar to how
    you use Top 1 or Top 5 for image classification tasks.
    Finally, here, I'm showing you for keyword spotting.
    For example, depthwise separable convolutional neural networks are great
    state-of-the-art models for actually doing keyword spotting in the field.
    And they had close to 95% accuracy, and you
    can see that you're able to prune the network by as much as 50%
    with a very small drop in the accuracy.
    So in summary, what I'm trying to explain to
    you is that, given that the size of the network
    is actually a critical issue for us, we want
    to be able to prune the network down while we're keeping the accuracy.
    So this is, hands down, one of the most important optimizations
    that all TinyML conversion tools try and do
    as part of the conversion process itself.
    Later on, we'll look at some of the other optimizations.
    We'll go into the complementary version of pruning, which is actually
    called clustering, which is a slightly different, clever way of doing
    compression.
    I'll see you there.
    Ende des Transkriptes, zum Anfang springen
    Downloads und Transskripte
    Video
    Videodatei herunterladen
    Transskripte
    Download SubRip (.srt) file
    Download Text (.txt) file

\section{Model Clustering}

    0:00 / 0:00Klicken Sie auf die Pfeiltaste nach oben, um das Geschwindigkeitsmenu aufzurufen. Regeln Sie dann mit den Hoch und Runter Pfeiltasten die Geschwindigkeit. Klicken Sie Enter, um die jeweilige Geschwindigkeit festzulegen.Geschwindigkeit1.0xKlicke auf diesen Knopf und diess Video stumm oder nicht stumm zu schalten, oder drücke die HOCH oder RUNTER Knöpfe um die Lautstärke zu erhöhen oder abzusenken.Maximal Lautstärke.Videotranskript
    Beginn des Transkriptes, zum Ende springen.
    VIJAY JANAPA REDDI: In this video, I'm going
    to be shifting our attention over to clustering.
    Clustering is an optimization algorithm that
    focuses on reducing the storage and network transfer size requirements
    of your model.
    Let me use a concrete example to help you understand
    how weight clustering really works.
    Imagine, for instance, you have a layer in your neural network
    that has a four by four matrix of weights
    where each weight is stored as a floating point value.
    So when you say if this particular model to disk,
    you are going to save 16 unique floating point values to memory.
    What weight clustering says is that, can we
    cut that down by looking at similarities in the weight distributions
    and by replacing similar weights in a layer with the same value?
    So instead of having some cluster of similar values,
    we just replace it with one value.
    And then the question becomes, how many of these clusters
    do we really want to have?
    You can have arbitrary number of clusters.
    Let's say, for instance, if we choose four,
    then what this means that we'll end up with four cluster centroids, as shown
    here in the centroid table.
    And each centroid has a unique value between zero and three.
    The next step is that we take each of the weights inside the weight matrix,
    and we replace it with the centroid's index which it corresponds to.
    So for example, with the values 0.86, which is shown in the top left matrix,
    that value is effectively going to correspond to a certain index, which is
    the second index in the centroid array.
    And this way, what we see is that every single weight is effectively
    replaced with a new matrix where you only
    have indices, which you are actually going to pull from the centroid matrix
    when you're actually doing the computation.
    As you can see down here, we were actually
    able to get a substantial reduction in size
    in using weight clustering because we were
    able to go from 16 unique floating point values in the top weight matrix
    down to just four float values in the bottom matrix.
    And of course, you still have the centroid,
    but it's still very small compared to just keeping the original matrix.
    Now, just to help you understand just how much of an impact this can have,
    for tiny amount especially, I'm just showing you
    some arbitrary set of networks here for image classification.
    And there's a lot of data on this table, but I just
    want you to focus on two key columns.
    That's basically what the compression effect is, or model size reduction is,
    and what the impact on the accuracy is for that level of model size reduction.
    And as you can see, we're able to get model size reductions as much as 6.3x
    while the accuracy drop is really, really small.
    So this is a fantastic optimization to actually consider
    during the conversion process.
    So bottom line, again, weight clustering helps us
    immediately by reducing the model size, and it's great
    for over the air updates and things like that we
    would do when we're doing continuous monitoring and continuous training
    later on.
    Weight clustering is somewhat similar to the sparse models, where
    we're effectively doing pruning, except that way
    clustering is trying to reduce the number of unique weights we have.
    Whereas pruning, on the other hand, is trying to set the weights down to zero
    when they are below a certain threshold.
    So stop for a moment and just think about the key differences
    between weight clustering and weight pruning because they seem similar,
    but they're actually quite different.
    And I actually have a quiz coming up that's actually going
    to test you on this understanding.
    Ende des Transkriptes, zum Anfang springen
    Downloads und Transskripte
    Video
    Videodatei herunterladen
    Transskripte
    Download SubRip (.srt) file
    Download Text (.txt) file

\section{Model Quantization}

    0:00 / 0:00Klicken Sie auf die Pfeiltaste nach oben, um das Geschwindigkeitsmenu aufzurufen. Regeln Sie dann mit den Hoch und Runter Pfeiltasten die Geschwindigkeit. Klicken Sie Enter, um die jeweilige Geschwindigkeit festzulegen.Geschwindigkeit1.0xKlicke auf diesen Knopf und diess Video stumm oder nicht stumm zu schalten, oder drücke die HOCH oder RUNTER Knöpfe um die Lautstärke zu erhöhen oder abzusenken.Maximal Lautstärke.Videotranskript
    Beginn des Transkriptes, zum Ende springen.
    VIJAY JANAPA REDDI: So far, we have talked
    about optimizations that focus on parameter count reduction,
    be it through weight pruning or clustering.
    In this video, I'm going to focus on quantization,
    which is all about changing the numerical precision that's actually
    available for us.
    Now, quantization can actually be coupled with pruning and clustering.
    In this video, I'm going to focus just on the role of quantization
    and how we go about using it for TinyML deployments.
    Quantization fundamentally is the process
    of transforming a machine learning program
    into an approximated version of itself.
    And you do this by using simpler operations,
    like int8 instead of floating point values.
    Now, one of the key things that you'll notice immediately
    is that floating point allows you to have a very large range
    and wide amount of precision versus int8 does not have that capability.
    So you're losing some representational capability, which may or may not
    affect your accuracy.
    So in this figure down here, I'm showing you
    a very simple example, where on the left you have floating point values,
    and on the right you're taking that exact same network
    and converting it into an int8 representation.
    It's a bit of an oversimplification, but we'll
    talk about it more in detail soon.
    What are the benefits of doing this sort of quantization,
    where you're going from floating point values
    to some representation of integer values?
    First and foremost, it's a reduced memory benefit.
    In any embedded system, we want our model to be small,
    so by going away or shying away from floating point values,
    we can get substantial improvements.
    So for example, if we go from floating point values to int8 values,
    we can get a 4x model size reduction, which is great for our flash storage.
    We can also get faster compute.
    Integer arithmetic is typically faster than floating point arithmetic.
    Moreover, in certain cases, it's important to be compliant
    with the limitations of the devices.
    Certain devices, like Cortex-M4, might be
    able to perform important arithmetic, but not
    so efficiently as compared to, say, an M33.
    So the difference is in there as well.
    And of course, integer arithmetic can be as much as 30 times more efficient
    in terms of energy consumption as compared to floating point arithmetic.
    And obviously, given that our TinyML devices are often battery tethered,
    we might want to think quite hard about reaping the benefits here.
    And similarly, you can reduce bandwidth.
    What this means is that, because you're using integer arithmetic,
    the number of bytes you're taking is smaller as
    compared to a floating point value.
    If you use int8, then you can-- at any given instant when you're moving data,
    you can move more off int8 data into the microprocessor than
    if you were just using floating point values.
    And moving data is extremely energy inefficient,
    and therefore, if you can cut down the bandwidth requirements,
    it actually helps you in energy efficiency as well.
    And hardware compatibility is key.
    Integers are baseline computational building blocks
    in a lot of microprocessor units.
    Let me elaborate this.
    On some devices you don't have an option.
    You have to be int8 compliant.
    So therefore, you have to quantize.
    So for instance, here I'm showing you here an ATmega microcontroller.
    And as you can see from its core size, it's an 8-bit microcontroller.
    This means that you can't run anything that is not
    doing computations in integer base.
    So there are devices out in the wild that
    are simple and very energy efficient, but they demand 8-bit execution.
    So let's talk about, how do we actually do quantization?
    There are three main different ways or types of quantization involved.
    The first is called reduced float, which is incrementally
    trying to get towards int8.
    So normally, everything is done in FP32.
    FP32 is not a single precision floating point,
    while FP16 is half precision floating point.
    From a numerical point of view, you can see here
    that the exponent and fractional bits are much smaller in the FP16 case.
    What this also means that is the FP16 does not
    have the same range and precision as FP32.
    But over time, what we have realized is that we
    don't need the same level of precision and range at large as FP32
    for a wide range of networks.
    Over time, in fact, Google researchers invented a new floating point
    arithmetic format called BrainFlow.
    What the BrainFlow does is that it actually
    tries to cater for the dynamic range of FP32
    while not worrying so much about the precision itself,
    because the precision that you get from FP16 is largely sufficient,
    but the dynamic range is something that is actually becoming
    important in the context of big models specifically.
    So either way what you can see is that, when you go from FP32 down to FP16,
    you get a 2x reduction in model parameters,
    because we're going to have precision, which is great.
    And there are other benefits, of course.
    Now, in the hybrid case, we're progressively
    inching more towards int8 computations, where
    you have a mix of both integer and floating point operations.
    But quantifying the weights down to int8,
    we get a 4x improvement in compression, because floating point
    take up 4 bytes, or 32 bits, whereas int8 only takes 8 bits.
    So you get a 4x compression right there.
    Naturally, you're going to be running faster on CPUs if you're just
    using integer arithmetic.
    You also tend to be more energy efficient for reasons
    I told you about earlier.
    However, in some cases, though, you have to keep the network activations
    in floating point.
    So it's a mix and match sort of a situation.
    Let me elaborate this a little bit further.
    Here you can see how we take float input values,
    we quantize them into integer 8, and then
    we do matrix multiplication in the integer 8 format,
    but when you take two numbers and you multiply them,
    the result ends up being a bit bigger.
    So therefore, the output ends up being an integer 32 base value,
    where we're back to using 4 bytes.
    And then we can quantize this effectively
    before we feed it into the activation function,
    like AReLU or something like that.
    And then we end up with an FP16.
    So you can see down here, for example, where
    we're able to get some amount of benefit by going into int8.
    But in some of the places, like when we're dealing with activation
    functions-- which tends to be very sensitive,
    and they often tend to deal with very small values--
    we cannot be using integer values, that we have to use floating point values
    in order to have the precision and the range that we actually need there
    in order to address for those smaller decimal values that show up often
    around the activation functions.
    So here we're seeing a mixture of integer and floating point arithmetic.
    And so that's why we're seeing both of them being intertwined here.
    Nonetheless, we get good amount of benefits.
    Now, finally, we have purely integer-based arithmetic,
    where everything is done an integer.
    Obviously, you get a great bang for the buck in terms of model size reduction.
    You get faster computations.
    What's obviously not shown here on the right
    is that you also get energy improvements.
    In summary, if you were to put all of this into a matrix,
    as I'm showing you here, there's a lot of information.
    I encourage you to pause the video and take a look at this matrix.
    But that said, I'm going to focus on one particular column, which
    is really quite important.
    It's going to be the ease of use column.
    So for the reduced float and hybrid quantization, what you see
    is that there's no data requirements.
    What this means is that you can work blindly off of the network itself,
    the network that you get.
    But when you go to integer quantization, you
    will see that you need unlabeled data.
    This is data that we have to gather from the weights.
    In other words, we effectively have to get some sort of calibration data
    in order to really know how to do this sort of integer quantization
    effectively.
    And the last straw, which is quantization aware training--
    we actually need labeled training data.
    I have not talked about quantization aware training,
    but real simply, it is this.
    In quantization aware training, you inject the quantization error
    during the training process, and you train
    against that error that actually comes about-- has a result of executing
    your entrance only in intake.
    What this allows your network to do is to overcome that penalty that ends up
    coming up later on, when you go into 8.
    And therefore, it's naturally coping with the error mechanism,
    and therefore, it tries to get the accuracy back
    by going over multiple epochs.
    And that's quantization aware training, but there the key requirement
    is going to be that you need the training data, which is not always
    possible in some cases.
    So just to show you some sort of improvements here,
    I'm comparing intake versus FP32 performance.
    And you can clearly see that int8 models are a much faster.
    Similarly, you can see that the intake models are also much smaller in size.
    These models are still in the order of megabytes,
    but the general trend will hold true also for smaller models.
    And finally, from an accuracy standpoint,
    you see that, with post-training quantization,
    you can see that the accuracy is actually impacted a small amount.
    But you can do that quantization aware training
    that I was hinting at earlier in order to get some of that accuracy back.
    And so effectively, while you're still doing quantization,
    you can actually benefit quite a bit from making the model smaller,
    but then keeping the accuracy still high,
    which, yet again, is extremely important to us.
    To sum things up, quantization is a very simple concept,
    but it's remarkably powerful, especially for TinyML.
    There are different ways you can go about it.
    For instance, you can talk about, how much do you want to quantize?
    You can go all the way down to binary networks, which
    are just really single 1-bit networks.
    By and large, though, 8-bit networks are the ones
    that are being widely deployed in the embedded space.
    Next, when you do quantize, are you doing it
    during training or just inference time?
    This really depends on whether or not you have
    that training data available to you.
    Nonetheless, usually, when you quantized and you retrain,
    you actually can recoup that accuracy.
    So you have a very small model, but you don't end up compromising accuracy.
    And some networks, like MobileNets for instance, are notoriously hard.
    So often, when you do quantization, if
    you don't do the quantization during the training time,
    you're going actually have a dramatic drop in accuracy.
    Then, of course, you can choose how you quantize weights
    or you also quantifies the activations.
    As we talked about it earlier, you can focus only on the weights,
    but activations are sometimes a little bit more finicky.
    They might actually need to be executed in floating point base,
    but it really depends, again, the characteristics of the network,
    and what particular task are you doing?
    Finally, you can quantize offline or you can actually
    use methods online in order to dynamically do the quantization, which
    are called static or dynamic methods.
    And there are also some wacky connections you can do.
    For example, you can draw connections between different schemes.
    So the point I'm trying to make is that it's a really rich space,
    but quantization, by default, is something
    that is extremely important to do for embedded deployments.
    And this is why it's extremely critical to understand
    doing the model conversion stage.
    Finally, going back to the big picture, one
    can couple pruning and clustering with quantization, as shown in this figure
    here, and various different combinations with either post training quantization
    or quantization aware training.
    So these are the kind of things that your ML engineer, who's
    generally responsible for downstream deployment,
    is going to be worrying about.
    So hopefully you appreciate the nuances that are involved, and then can
    understand the connections that need to be made when you're actually
    trying to push a model into deployment.
    Ende des Transkriptes, zum Anfang springen
    Downloads und Transskripte
    Video
    Videodatei herunterladen
    Transskripte
    Download SubRip (.srt) file
    Download Text (.txt) file

\section{Collaborative Optimizations}

    \subsection{About This Reading}
    
    Previously, we have talked about optimizing our models for size by pruning our models, performing weight clustering, quantizing our weights from 32-bit floating point values to 8-bit integer values, and also touched on quantization-aware training (QAT). However, trying to perform all of these in such a way that we are able to optimize our model performance without removing the benefits of past optimizations is difficult. In this reading we will explore Collaborative optimization, a relatively new technique that allows us to get the best of both worlds, optimizing for multiple characteristics  simultaneously.
    
    \subsection{Overview}
    
    Collaborative optimization takes the optimization techniques that we were using previously and chains them together in such a way that the advantages gained by applying one method are preserved upon the application of another method. Previously, we might apply weight pruning (flattening weights that were close to zero) followed by weight clustering (clustering the existing weights into a few weights such as small, medium, and large weights). However, in applying the clustering, we might inadvertently remove some of the advantages of the pruning stage. Collaborative optimization takes this into account by introducing optimization stages that are preservable. The optimizations offered by TensorFlow currently are:
    
    \begin{itemize}
      \item Sparsity-preserving clustering
      \item Sparsity-preserving quantization aware training (PQAT)
      \item Cluster preserving quantization aware training (CQAT)
      \item Sparsity and cluster preserving quantization aware training
    \end{itemize}
    
    \GRAPHICSC{1.0}{1.0}{TinyML/4 MLOps/007.png} 
    Overview of TensorFlow optimization techniques for a Keras model showing all of the options discussed in this reading.
    
    Using these optimization techniques in a specific sequence (outlined in the above image) allows the compression of a machine learning model to be “optimally” optimized. For example, applying sparsity-preserving clustering after pruning a model helps to prevent the loss of performance gained from the pruning stage. The leaf nodes in the above diagram are deployment-ready models, green nodes represent stages that require retraining or fine-tuning before deployment, and the red border represents collaborative optimization steps.
    
\GRAPHICSC{1.0}{1.0}{TinyML/4 MLOps/008.png}   
        
    A graphical representation of collaborative optimization of a model by first implementing weight pruning, weight sparsity-preserving clustering, and then quantization showing how these combined effects can be used careful to maintain all of their benefits. It shows how pruning a model with weights distributed normally about 0 (thing a standard bell-shaped curve) when passed through pruning will result in a 50\% sparse model. The weights now look more like two small bell curves that are very short next to a huge spike at 0. When this is passed through sparsity preserving clustering the data now looks like a series of spikes along a bell shape as it is now still 50\% sparse but is clustered at discrete spiked locations. Then when passed through sparsity and clustering aware quantization those floating point spikes are converted into integer spikes that have the same visual shape.
    
    Using this optimization tree may not be required for a specific problem. For example, where model performance is already sufficient from a pruned and quantized model, we don’t need to consider PCQAT or CQAT as an alternative. A suitable procedure might be to go through the standard optimization pathway, see if this is sufficient and, if not, proceed with collaborative optimization pathways. Alternatively, all pathways could be tried and the output model with optimal characteristics is selected for deployment.
    
    \subsection{Summary}
    
    While collaborative optimizations may not always be necessary, knowing that they exist might be helpful for applications that are particularly demanding or where optimizing metrics is very important, such as if the model is being deployed on a large number of devices.
 
\section{Student Teacher Networks / Knowledge Distillation}

    0:00 / 0:00Klicken Sie auf die Pfeiltaste nach oben, um das Geschwindigkeitsmenu aufzurufen. Regeln Sie dann mit den Hoch und Runter Pfeiltasten die Geschwindigkeit. Klicken Sie Enter, um die jeweilige Geschwindigkeit festzulegen.Geschwindigkeit1.0xKlicke auf diesen Knopf und diess Video stumm oder nicht stumm zu schalten, oder drücke die HOCH oder RUNTER Knöpfe um die Lautstärke zu erhöhen oder abzusenken.Maximal Lautstärke.Videotranskript
    Beginn des Transkriptes, zum Ende springen.
    VIJAY JANAPA REDDI: So far, we have talked
    about removing weights from networks and pruning the network
    to reduce its parameter count.
    We've also dabbled with quantization to make the model boat smaller and faster
    using low-precision arithmetic.
    Now, I want to talk to you about knowledge distillation, which
    is a technique that is really quite different from aspects of quantization
    or pruning or clustering.
    Think of it is an algorithmic sort of a solution as
    compared to taking a network and just modifying its internal characteristics.
    More formally in machine learning, knowledge distillation
    refers to the process of transferring knowledge
    from a large model to a small one.
    Big machine learning models often tend to have a large knowledge
    capacity than smaller models.
    This is primarily because they have many more neurons.
    But those neurons might not be getting fully-utilized,
    as in, their capacity is not being utilized to the maximum.
    And so therefore, it's possible that we might actually
    be able to distill some of the knowledge from the bigger networks
    into a smaller network.
    Fundamentally, it starts with the assumption
    that we have a trained model out there.
    Now, this is actually quite important.
    Because there are, in fact, a large number of big trained models out there.
    And when it comes to TinyML models, we don't tend
    to have a large number of these models.
    And so we don't necessarily want to invent things from scratch.
    What we ideally would like to be able to do
    is that we want to take a big teacher model
    and get that teacher model or pre-trained network
    to actually teach a smaller student network.
    And if we can do this, we can enable TinyML deployments
    much more efficiently.
    Now, let's talk about how we can actually
    go about creating this knowledge distillation environment so
    that a big, pretrained model can actually try and teach a smaller
    student model who has no knowledge yet.
    By and large, there are three big pillars.
    One is what is the knowledge that you're actually trying to teach?
    What is the architecture that you can have
    between the student and the teacher so that the student can actually
    learn effectively from the teacher?
    And third is, what's that quote unquote, "distillation algorithm?"
    Distillation here is the training algorithm.
    This can basically be broken down into further types
    starting with the knowledge types-- the three main types.
    Now, when you take a big network, there are multiple things
    you can actually learn from it.
    The most obvious and most simplest thing is
    to focus on the final output of what the teacher does.
    The student can learn to mimic the predictions of the teacher model
    just to the ending.
    Very simple.
    But the teacher also has much more knowledge to share.
    Think about all the intermediate steps that the teacher actually went through
    in order to get to that output.
    These intermediate layers are able to discriminate specific features.
    Now, if we can actually take that knowledge,
    we can actually try and teach the student
    that exact same sort of capability.
    In other words, we can teach the student to mimic the intermediate layer's
    behavior.
    And that could be another way of actually mimicking
    what the teacher's behavior is.
    Finally, there's the actual relationship-based knowledge
    where the idea is that the knowledge that actually captures
    the relationship between feature maps can also
    be used to train a student model.
    So as you can see, there are many different ways or many different types
    of knowledge we can actually impart to the student network.
    Next, comes a question of, how do we architect the relationship
    in teaching the student network?
    And this comes into the aspect of, what is that architecture's type?
    Now, the most common architecture type for knowledge transfer
    includes a student model that is either shallower compared to the teacher.
    In other words, there are fewer layers fewer number of neurons and so forth.
    Or it could be that the student network is a quantized version of the teacher
    model.
    Because we just learned about quantization.
    Or, perhaps, because we're shooting for TinyML deployments,
    the smaller student network is made up of very simple,
    basic operations as compared to the bigger model which can probably
    perform more complex operations, because it's not
    being deployed on TinyML devices.
    Now, you can also couple all these variations
    with techniques like NAS, which we talked
    about in the AutoML section of the continuous training pipeline
    in order to help discover the right sort of an architecture.
    So there are many different ways of actually structuring the student
    and teacher relationship architecture.
    Now, the third part is really the distillation algorithm, the training
    algorithm.
    Now, certainly, a student can just learn from one teacher model.
    But that does not have to be the case.
    More often than not, you want to have multiple teachers teaching the student.
    So this is true like in any student who comes to class,
    they are often taking classes from multiple teachers.
    So you can learn different kinds of knowledge
    that can be much more beneficial than just having a single teacher.
    So I'm giving you some examples here.
    So in terms of other types of mechanisms for distillation algorithms
    would be outside of the ensemble or the teachers,
    you can have quantized distillation, which
    is used to transfer knowledge from a high-precision teacher
    to a low-precision student network.
    You can use things like NAS to quickly identify a suitable student model
    architecture that is actually really well-suited for that teaching model.
    Or you can also use things like attention-based distillation where
    you're transferring knowledge from feature embeddings,
    which I talked to you way earlier in the ML development stage,
    which typically ML researchers dealing with--
    so lots of different ways of actually doing the distillation itself.
    Given that there are so many different ways,
    I just want to pick one particular implementation which
    comes straight out of Geoffrey Hinton's paper for how
    we can do knowledge distillation.
    Taking this base model down here, you can see this input on the left most.
    And then we want to be able to train a student network based on some teacher
    network
    I've broken it down further into four core components.
    Let's step through each one.
    First and foremost, we assume there's a teacher network who
    knows how to make predictions.
    And then we have a student network who's starting from scratch.
    The second portion of this is, really, how the student network
    can learn from its own behavior.
    The student network, much like a traditionally-trained neural network,
    can make predictions about the input that's given to it.
    And its predictions can be compared against a ground truth
    that we already have access to.
    And we can compute the standard loss, which
    can then be typically fed back through the training process
    to help it learn a bit better.
    Here, the student is only learning about whether the prediction
    is right or wrong.
    This is, therefore, typically referred to as just student loss.
    And the key shortcoming here is that the students are really not learning
    anything about the distribution of class probabilities
    and so forth, which is really where the teacher network becomes extremely
    powerful for us.
    Because the teacher actually knows the distribution of class probabilities.
    So for example, by looking at the Softmax probabilities
    for this image, what we're seeing is that there's a high probability
    that it's actually a one.
    And a low probability that it's a seven.
    It is true that this could actually be a seven too.
    But this sort of probability mixture is actually extremely key for the student
    to actually learn in order to really have a good grasp on pretending
    to be the teacher itself.
    So what we want to do is we actually want
    to compare how well the student actually does
    with respect to this sort of behavior and then compute the delta which
    is known as the distillation loss.
    And what you want to do is you want to take both the distillation loss, which
    affected tells you how the student is doing with respect to the teacher
    as well as the student's own learning ability, which
    is called the student loss.
    And you can couple them together into a total loss function
    using some arbitrary weightage.
    And then use this to try and train the student network.
    And this would allow you to get to a student network that
    is very good at learning while also trying
    to mimic the characteristics of the teacher network.
    Now, while all of that sounds great, there
    are some fundamental limitations to knowledge distillation.
    Now, one of the key challenges ends up just being,
    what's the capacity of the student network with respect to the teacher
    network?
    If the student network is really, really small,
    then it can be really hard for the student network
    to actually learn the capabilities that the teacher network has.
    There's also another pitfall is that a student network who's actually
    learned from the teacher might actually end up
    performing worse than a student network that's
    actually been trained from scratch.
    And this is something you have to keep in mind.
    So this is, again, a trade-off between, do you take a big, pretrained model
    and then try and distill its knowledge down to a small one?
    Or do you have the resources and the time and money
    to actually be able to train from scratch?
    Clearly, this pertains quite tightly to us when it comes to TinyML deployments.
    Now in summary, our goal of doing knowledge distillation
    as an optimization is that we want to be able to convert models for easier
    deployment into embedded systems.
    With that said, I will see you in the next video
    where I'm going to try and summarize all that we have learned
    about model optimizations and model conversion in the stage
    and how this ties to the rest of the pipeline in the MLOps.
    Ende des Transkriptes, zum Anfang springen
    Downloads und Transskripte
    Video
    Videodatei herunterladen
    Transskripte
    Download SubRip (.srt) file
    Download Text (.txt) file 
 
\section{Model Conversion Impact on MLOps}

    0:00 / 0:00Klicken Sie auf die Pfeiltaste nach oben, um das Geschwindigkeitsmenu aufzurufen. Regeln Sie dann mit den Hoch und Runter Pfeiltasten die Geschwindigkeit. Klicken Sie Enter, um die jeweilige Geschwindigkeit festzulegen.Geschwindigkeit1.0xKlicke auf diesen Knopf und diess Video stumm oder nicht stumm zu schalten, oder drücke die HOCH oder RUNTER Knöpfe um die Lautstärke zu erhöhen oder abzusenken.Maximal Lautstärke.Videotranskript
    Beginn des Transkriptes, zum Ende springen.
    VIJAY JANAPA REDDI: Congratulations on getting through the moral conversion
    stage of the MLOps pipe.
    We covered a fair bit of ground in this module.
    We talked about various kinds of ML frameworks
    and what are some of the common pitfalls as we
    try to work across different frameworks and going say,
    for example, from PyTorch to TensorFlow, to TensorFlow to PyTorch, and so forth.
    And I exposed you to a rich flavor of different ML frameworks
    that are out there.
    And even within the single family of ML framework,
    I showed you that there are very different flavors.
    For example, between TensorFlow, TensorFlow Lite, and TensorFlow Lite
    Micro, there are major differences.
    Then we talked about many of the optimizations
    that we should undertake when we're talking
    about things for TinyML deployments, such as quantization, clustering,
    pruning, knowledge distillation, and so forth.
    Now, look, I know I didn't go into a lot of depth
    in terms of the technical details along these optimizations,
    but that's only because these optimizations are actually
    quite mature, and they're widely available on many popular frameworks.
    So it's more important that you actually know
    how to use them and know what the trade-offs are
    and who's responsible for these tasks than to understand the nitty gritty.
    Now, the key matter on this slide in itself
    is the question as to, what is the effect that conversion
    has on the rest of the MLOps pipe?
    Model conversion is tightly coupled to various stages.
    In this particular case, I'm going to talk about the ML development stage.
    Let me tell you about a little story.
    In a major organization that deploys keyword spotting at scale--
    I mean, scale-- models are deployed in the research wing
    and passed over to the deployment team.
    Now, little did the research team know about the extreme heterogeneity and
    rich diversity of embedded systems.
    So when the model deployment team was actually looking at scaling things out,
    they found a fundamental problem, which was they realized that the ML
    researchers had used some fancy operators that actually cannot run well
    or cannot run at all on some devices in the production environment.
    So what this meant is that they actually had to go back to the ML researchers
    and have the ML researchers re-architect and come up
    with an entirely new model that met their accuracy requirements
    and then be able to go into production.
    So I'm trying to hint towards situations where even in big entities,
    you tend to have these kinds of fundamental problems
    where there's a relationship breakdown in terms of the communicating what
    the requirements are really at scale.
    Another thing is that model conversion pipe is intricately
    tied to the continuous training stage.
    Often in the continuous training stage, you're
    going to retrigger training when you get new data.
    And when you do that, more often than not,
    your models are slowly but steadily going
    to get bigger as you get more and more data.
    It's totally possible, not always the case, but it's possible.
    And then you find that you have to obviously
    try and shrink the networks down, especially for TinyML deployments.
    So you're going to be going through this process of doing
    pruning and quantization and so forth iteratively.
    And so the message that I'm trying to point out here
    is that while you're doing continuous training,
    you also need to be thinking about continuous model conversions that
    need to be taking place.
    Needless to say, going from model conversion into model
    deployment is very closely tied together.
    In fact, by and large, in most places, model conversion
    is tied into the model deployment stage.
    But we intentionally chose to separate it out here
    because in the model deployment stage, I'm
    going to talk about a whole bunch of new issues
    that will come up when you're trying to scale for TinyML.
    Nonetheless, the point I'm trying to make here
    is that this is an extremely tightly and critical piece that
    couples with model deployment.
    And finally, this is the kind of work that an ML engineer
    is going to be worried about in creating a production grade MLOps pipeline.
    With that said, I hope you picked up some of the nuances
    around model conversion.
    And I'm looking forward to seeing you in the model deployment stage.
    Ende des Transkriptes, zum Anfang springen
    Downloads und Transskripte
    Video
    Videodatei herunterladen
    Transskripte
    Download SubRip (.srt) file
    Download Text (.txt) file 
    
    
\section{Model Conversion Case Study - Smart Doorbell}

Provided By \href{}{Deeplite}

\subsection{Introduction}

    Smart Doorbells represent an excellent application for TinyML. Ideally, a video doorbell should record and/or upload video footage of important events that occur in the view of the doorbell and alert the owner when an event has happened. In practice, this is a challenging task since a constant stream of video is too much to constantly upload or store locally. Therefore the doorbell needs to be able to determine, on its own, which video clips should be uploaded or stored by identifying key objects in the video. Additionally, many smart doorbell systems are battery powered and need to last for months at a time and must protect the owner's privacy. At this point in the course, you can see how this presents a challenge that requires TinyMLOps. By efficiently running object detection on the device we can greatly improve the usefulness of the product.
    
    \GRAPHICSC{1.0}{1.0}{TinyML/4 MLOps/smart_doorbell_screenshot.png}
    
    A screenshot of a smart doorbell detecting a person stealing a package from a porch. There is a red bounding box around the person, with a label stating person 0.84, indicating the confidence of the model.
    
    \subsection{Challenge}
    
    For many standard computer vision tasks, there are dozens of options for developers and product managers – they can use a pre-trained model, or utilize off-the-shelf CNN model repositories to fine-tune on their dataset during the model training phase of development. However, it’s important to note that in many cases, such models are over-parameterized and under-optimized for the given use case. When considering what the optimized model architecture would be for a given CV application, it’s essential to consider the nature of the data distribution (how many classes, how each class is represented in the training sample, how much training data is available, how much concept drift is expected) and the required accuracy.
    
    In the context of a smart doorbell device, we have the following requirements:
    
    \begin{itemize}
      \item 3-classes (person, pet, vehicle)
      \item $>90\%$ accuracy per-class
      \item Over 1M training samples
    \end{itemize}

    A model pre-trained on the COCO dataset, like Yolov5-l, despite being one of the best models in literature, is not exactly the best fit for this application. COCO has over 80 classes with fewer data samples than the smart doorbell use case. This means the training requirements are different for an application with higher accuracy requirements. Furthermore, it’s critical to consider the inference requirements for model optimization. Latency, available memory, CPU usage are factors in creating the best model architecture. 
    
    The following are realistic requirements for the smart doorbell in this case:
    
    \begin{itemize}
      \item < 250ms inference latency on 1.2GHz quadcore Arm Cortex-A53 CPU
      \item < 2MB runtime footprint
    \end{itemize}
    
    \subsection{Model Introduction}
    
    TinyMLOps tools can automate the process of optimizing models to meet specific constraints. For example, Deeplite’s Neutrino, shown below, is a design space exploration and low-bit precision quantization engine, which can combine both training and inference requirements into the model optimization process.
    
    \GRAPHICSC{1.0}{1.0}{TinyML/4 MLOps/deeplite_mlops_pipeline.png} 
    
    End-to-End illustration of where Deeplite fits in the MLOps pipeline showing how it takes a MB size trained model and compiles it down into a KB size embedded model through quantization and other optimizations. The pipeline starts with code and data, which feed into building and designing DNNs, which are then trained. Next, the large MB sized model is fed into the neutrino platform which is an automated model optimizer, quantizer, and compiler. The output of neutrino is a KB-sized model, which is then fed to testing, deployment, and, finally, the deeplite rt runtime and inference engine. There is a feedback loop that connects this final inference stage back to the data at the beginning of the pipeline. There is text under some of the stages in the pipeline. The build/design dnns block says “data preprocessing, visualization, hp tuning, architecture design, tensorflow, pytoch. The text under the train block says, “tensorflow, pytorch, onnix, distributed training, cloud/on-prem.“ The text under that neutrino stage says, “automated model optimizer, quantizer, and compiler. The text under the test stage says, “functionality test, integration test, performance test”. The text under the deploy stage says “bundling, flask, compiler/RT”.
    
    Unlike naïve weight or channel pruning or post-training quantization to INT8, combining multiple optimization techniques in a single optimization pipeline allows end users to factor training requirements (like accuracy) with inference requirements (like latency) to rapidly converge on an ideal model architecture that is ready for production deployment. In this case, the Neutrino engine found a compressed Yolov5-s model with 2bit quantized weights and activations that met accuracy criteria, inference latency (119ms vs 250ms) and provided sufficient memory savings (15.5x model compression vs FP32 and only 1.8MB overhead with DeepliteRT).
    
        \GRAPHICSC{1.0}{1.0}{TinyML/4 MLOps/deeplite_graphs.png} 
    
    Different performance tradeoffs for model optimization of person, pet, and vehicle detection model on Cortex-A53 CPU (RPi3b+). On the accuracy benchmark Yolo5m performs the best, followed by Yolo5s and Yolo5n with accuracies of 45.2 vs 37.2 and 28.4 respectively on COCO and 90.25 vs. 97.11 and 76.21 respectively on the VOC dataset. We find that the DeepliteRT solution far outperforms both the TFLite and ONNX versions of the model in terms of the execution time, which is much faster with 0.119 vs. 0.255 vs. 0.376 on Yolo5s and 0.326 vs. 0.714 vs. 0.99 on Yolo5m.
    
    Now that the model has been optimized to meet the training and inference requirements, the next step is to convert the model into a format that can be run efficiently on the target device.
    
    \subsection{Model Conversion}
    
    The model conversion process involves two basic inputs – the output from model optimization and the target HW for deployment. The various dependencies, versions and AI frameworks for model development - which largely focus on high-level abstraction and python programmability – must be reconciled with low level system requirements and optimization for the target hardware. When dealing with edge AI and tinyML, the latter can be a significant engineering challenge given the harsh constraints we are dealing with (\$1-5 chips, 100milliwatt power budget, < 1GOPS). Unlike cloud environments common for training, containerization and VMs are less common when we are dealing with edge AI deployments that will also require continuous model update and data collection once in production.
    
    In order to avoid the engineering cost of converting the model, model conversion tools exist that automate and optimize the process. For example, the Deeplite compiler allows the automatic conversion of an optimized model to a format that supports Arm based CPUs and runs with optimal inference speed. In this case, the model needs to be converted in a proprietary model format called dlrt. This format supports the 2 bit precision, runs on the Linux OS and Arm CPU already installed on the smart doorbell and requires the Deeplite runtime (DeepliteRT) to be installed on the end device to run AI model inference.
    
    \subsection{Model Deployment}
    
    Now that the model has been optimized and converted for inference on the Arm Cortex-A SoC it can be deployed. The first step is to ensure DeepliteRT is installed over-the-air using the device IP address and Linux OS to infer the dlrt model on-device. The entire package is provided as a docker image to the end-user for ease-of-use when automating the model conversion process. Then the model can be sent to the device and periodically updated over-the-air as well.
    
    \GRAPHICSC{1.0}{1.0}{TinyML/4 MLOps/Deeplite_deployment_stack.png} 

    Full model optimization, conversion, and deployment stack from Deeplite showing its compression and quantization engine (Deeplite Neutrino) followed by its compiler followed by the Deeplite RT runtime engine for embedded devices. Within the deeplite neutrino block, a reference model is fed into the quantization and compression blocks which in turn are output into the deeplite compiler. The Compression block outputs a FP32 model and the quantization block outputs a 2bit/mixed-precision model. The deeplite neutrino and deeplite compiler blocks are in a larger docker rectangle. The compiler outputs a .dlrt file which is passed to the deeplite RT which runs on an ARM Cortex-A CPU.
    
    \subsection{Model Conversion}
    
    Using a TinyMLOps stack for conversion and deployment can substantially reduce the engineering hours required. The tighter the constraints on the application, the more difficult the conversion process becomes, as optimizations add complexity. However, new TinyMLOps tooling improves the process and enables TinyML developers to scale their prototypes across thousands of endpoint devices. In this case study, what typically takes years of manual development was reduced to only 4 months of prototyping and testing before deploying deep learning-enabled person, pet and vehicle detection models on over 100,000 smart doorbells.
    
    \subsection{Additional Resources}
    
    \begin{itemize}
      \item \href{https://arxiv.org/abs/2101.04073}{AAAI 2021 paper on DeepliteNeutrino}
      \item \href{https://docs.deeplite.ai/neutrino/index.html}{Deeplite Neutrino}
      \item \href{https://info.deeplite.ai/deeplite-pytorch-developer-day-poster?utm_source=hs_email&utm_medium=email&_hsenc=p2ANqtz-9LTa1hhXuEMLXVs5VZjE6ntg_AYQhGyazAqnYYNE3BT9jikvgvBi-FyU_JvHLJYToC7kQJ}{PyTorch Dev Day 2021 poster on DLRT}
      \item \href{https://www.youtube.com/watch?v=0bJs7dHC8o4}{DeepliteRT Video}
      \item \href{https://github.com/Deeplite}{https://github.com/Deeplite}
    \end{itemize}

\section{Forum Questions}

\begin{itemize}
  \item What are the factors that impact model conversion? 
  \item   How do model conversion concerns impact other stages in TinyML development?
  \item   How does TinyMLOps improve model conversion?
  \item   What are the primary constraints of smart doorbell applications? 
  \item   What additional features might be added to improve the usefulness of a smart doorbell?
\end{itemize}


% 1.7: Model Deployment

\section{Overview of Model Deployment}

\GRAPHICSC{1.0}{1.0}{TinyML/4 MLOps/model_deployment.png}   

    The MLOps pipeline. Going from ML Development to ML training to Continuous Training to Model Deployment to Prediction Serving to Continuous Monitoring. Along the way artifacts are produced at each step including: Code and Configurations, Training Pipelines, Registered Models, Serving Packages, and Serving Logs. All steps tie into the overarching themes of Data and Model Management.
    
    \subsection{Objectives}
    
    Once we have developed, trained, and validated our models, we want them to be implemented into our processes. This stage of implementation is what is often called model deployment. The deployment stage involves packaging the model, performing a suite of tests, and then deploying it to our target environment. For more traditional deployments, this could be a web server or Kubernetes cluster. Deployment to a production environment is usually more complex than just saving over our previous model, or adding a new one. Typically, it also involves continuous monitoring of the predictions, and deployment is often performed using one of a number of techniques which we will briefly discuss. For TinyML, this more so refers to deployment on a specific device and that is where several unique challenges arise. To this end, in this section we will cover the following topics to understand how to robustly deploy models:
    
    \begin{itemize}
      \item   Containers and the role that containers play in easing ML deployments into the wild
      \item Challenges with the lack of containerization methods for embedded systems
      \item Tiny Machine Learning as a Service (tinyMLaaS) as a potential solution for the issue
   \end{itemize}

    \subsection{Motivation}
    
    We have previously discussed deployment in the context of “ML Training” or “Training Operationalization” as it was also relevant to the issues of testing. More often than not it is important to test in production as it is impossible to know if you have covered all your bases despite having integration, smoke, as well as acceptance testing, and so forth. So here we will do a quick recap on the deployment methods: basic, multi-service, rolling, canary, blue-green, and shadow. As you can now see, there are many ways for us to introduce a new model into our production environment. Some of these are easy to implement, but relatively risky, while others are complex but relatively safe. The most suitable deployment strategy depends strongly on the application, but also things like cost, existing infrastructure, and technical expertise available.
    
    In this module, we will be focusing on another part of model deployment known ‘as a service.’ This extends well beyond the Continuous Integration and Continuous Deployment requirements. Often when you have to deploy models into production, you have to have the right infrastructure in place. One can choose to build all of the components themselves or try and lease them from others who are providing such services for a fee—leasing alleviates the customer to focus on their problem at hand, rather than build all the components from scratch with little to no experience just to deploy their solution. For instance, you can set up a web server with a push of a button today, thanks in part to all the “LAMP” (Linux, Apache, MySQL, and PHP) infrastructure components that have all matured well and been integrated well with one another. This allows you to focus on the website you are designing, rather than figuring out how to bolt together all these different major software components together just to get your website hosted. 
    
    Similarly, when we want to deploy ML models into production at scale, we need the right underlying software infrastructure in place. Arguably, many of these components are still missing, and if not they are limited in their capabilities. To understand the challenges and opportunities, we will be introducing Tiny Machine Learning as a Service (TinyMLaaS).        
        
 \section{Model Deployment}
 
     0:00 / 0:00Klicken Sie auf die Pfeiltaste nach oben, um das Geschwindigkeitsmenu aufzurufen. Regeln Sie dann mit den Hoch und Runter Pfeiltasten die Geschwindigkeit. Klicken Sie Enter, um die jeweilige Geschwindigkeit festzulegen.Geschwindigkeit1.0xKlicke auf diesen Knopf und diess Video stumm oder nicht stumm zu schalten, oder drücke die HOCH oder RUNTER Knöpfe um die Lautstärke zu erhöhen oder abzusenken.Maximal Lautstärke.Videotranskript
     Beginn des Transkriptes, zum Ende springen.
     LARA SUZUKI: Hi.
     In your last lesson, we had an overview on how
     to convert to models to run on resource-constrained devices
     and on microcontrollers as well using TensorFlow Lite converter.
     Now it's time to move to the next stage of our MLOps framework
     that is of model deployment.
     In this session, we will cover deeper details
     of deploying ML models in a secure way.
     In the MLOps deployment stage, we package, test, and deploy a model
     to a serving environment for online experimentation and production serving.
     The machine learning model is packaged, tested,
     and deployed to the target environment mainly via CI/CD pipelines.
     In practice, several CI/CD processes are executed,
     such as reading the source code of the model serving component
     from the source code repository.
     In the model deployment stage, the model serving service
     is built, tested, and validated, and deployed to a progressive delivery
     process.
     Smoke testing is often executed to measure
     model service efficiency such as latency and throughput
     as well as service errors.
     The MLOps capabilities employed in this stage
     are model serving, model registry, online experimentation,
     machine learning metadata and artifact repository.
     Let's take a look at the entire model deployment process.
     As mentioned previously, the machine learning model
     is packaged, tested, and deployed to a target environment.
     During the deployment phase, a number of testing steps
     take place and might involve a number of testing environments.
     In the development phase, several CI/CD processes take place.
     A CI/CD pipeline, which reads the source code of the model serving component
     from source and repository.
     It fetches the model from the model registry.
     It integrates, builds, tests, and validates the model serving service.
     It then deploys a service to a progressive delivery process.
     In the continuous integration pipeline, there's a test of the model interface
     to verify whether it accept the expected input format
     and produces the expected output.
     It also validates compatibility of the model with the target infrastructure.
     In the continuous deployment pipeline, there's a progressive delivery process.
     Smoking tests are focused on model service efficiency,
     like latency and throughput, as well as on server errors.
     Another test that takes place is testing the model effectiveness in production.
     The continuous deployment pipeline also enable
     data scientists to conduct online experiments, such as A/B testing.
     Ende des Transkriptes, zum Anfang springen
     Downloads und Transskripte
     Video
     Videodatei herunterladen
     Transskripte
     Download SubRip (.srt) file
     Download Text (.txt) file
 
 \section{Scaling ML into Production Deployment}
 
     0:00 / 0:00Klicken Sie auf die Pfeiltaste nach oben, um das Geschwindigkeitsmenu aufzurufen. Regeln Sie dann mit den Hoch und Runter Pfeiltasten die Geschwindigkeit. Klicken Sie Enter, um die jeweilige Geschwindigkeit festzulegen.Geschwindigkeit1.0xKlicke auf diesen Knopf und diess Video stumm oder nicht stumm zu schalten, oder drücke die HOCH oder RUNTER Knöpfe um die Lautstärke zu erhöhen oder abzusenken.Maximal Lautstärke.Videotranskript
     Beginn des Transkriptes, zum Ende springen.
     VIJAY JANAPA REDDI: In this module, I'm going
     to build on what Lara said about how we want to focus on continuous integration
     and continuous delivery flows for ensuring
     the robust deployment of our models into production.
     But in addition, I'm also going to talk to you
     about some of the unique challenges that come about
     as a result of dealing with TinyML, specifically
     because of the embedded ecosystems fragmentation.
     Now this is a topic that's often not talked about,
     but it comes up real strongly in the thematic area
     of production deployments.
     But before we get there, let's start off talking about continuous integration
     and continuous delivery.
     When you look at this diagram, a few bells
     should really be going off in your head at this point.
     This should look quite familiar to you.
     Why?
     If you recall, I actually talked about continuous integration
     and continuous delivery in the training operationalisation stage.
     In fact, we talked about several aspects around continuous integration,
     continuous delivery, as well as production deployment in the training
     operationalisation stage.
     Now there's a lot going on here.
     Let's quickly zoom in the specific bits as they
     correspond to production deployment.
     As you can see here, in the operationalisation stage,
     when we talked about production deployment,
     specifically talked about how we go into deploying models into the field,
     recall how I said, if you flip the switch
     and you instantaneously swap your old machine learning model out
     with the new one in a jiffy, how things can go wrong.
     There are many reasons why things can go wrong.
     Well, for one durin the operationalisation
     stage we assume the testing environment is
     the same as the production environment.
     And I've talked to you a lot about the issues around skews and data drifts
     and so forth as to why that might not be the situation.
     It's really hard to match the testing environment
     to the production environment.
     There, I talked to you about some of the challenges
     around dealing with TinyML small devices which actually
     might be operating in remote situations that are actually
     hard to mimic in a testing environment.
     Bottom line, the fact is that it's highly impossible
     to test all the different conditions and scenarios in practice.
     So the key challenge that emerges with ML
     is the fact that there's a lot of susceptibility
     around data, which makes it really hard to be
     able to test everything to perfection in the testing environment.
     And this is often why we test live in production environment
     with real traffic.
     That said, there's certainly some careful approaches
     that we have to take when it comes to the deployment strategies.
     And this is why we had learned about canary deployment, shadow deployment,
     and blue/green deployments.
     And all of these actually hold true when we're talking about the model
     deployment stage down here.
     So even though I talked about them earlier, they're coming back down here.
     As I've told you right at the beginning of this course,
     there are a lot of aspects of MLOps that tend
     to repeat up over and over in different aspects of the pipeline.
     So this is one of those instances.
     So is that all to it?
     Well, no, actually there's a lot more to thinking about production deployment.
     I'll tell you what's new, but in order to help you really understand this
     systematically, I'm going to use some very basic examples.
     The fundamental assumption that we have been making so far
     is that when we use those fancy deployment
     strategies, like the canary deployment or the blue/green
     deployment or the shadow deployment, we made the fundamental assumption
     that it is actually extremely easy to push things into deployment.
     Let's take this WWWR world wide web scenario,
     and let's try and understand how websites on the internet really push
     and how we're able to thrive in an ecosystem
     where we can very easily, with the single push of a button,
     pop a website up.
     Because I think it will help us understand
     all the intricate infrastructure that's already existing that we're
     taking for granted, which may or may not exist for TinyML.
     That's why it's important to understand from a basic example.
     Now today, for example, when you look at all of the internet websites that
     are out there, you might find the staggering to here, but about 80%
     of them all rely on what's called the LAMP stack.
     LAMP is basically the combination of Linux, Apache, MySQL, and PHP.
     These are distinct tools which, when combined together,
     effectively run the vast majority of internet sites.
     In fact, LAMP is so well-baked today that this
     is why you can actually very quickly pop up a website
     and then go live in less than five minutes.
     All right, so some fanatics out there might say,
     well, there are lots of issues with the LAMP stack.
     OK.
     Yes, my friends, I agree.
     You might have a certain religious debate with me
     about what are the things that you're most passionate about.
     Nonetheless, the simple point I'm making is
     that this coupling of Linux, Apache, MySQL, and PHP
     as fundamental building blocks in terms of the services they offer,
     allow us to very quickly build up websites.
     So let's take a quick look at these different pieces
     and how they fit together.
     So that we can actually understand how we
     might be able to build a TinyML ecosystem better
     using this as an analogy.
     First and foremost, the LAMP stack is a classic layered architecture
     with Linux operating system at the lowest level,
     the next layer being the Apache, followed by MySQL
     layered on top with the scripting language, which is PHP.
     Linux is the underlying operating system, more than likely,
     you know about Linux or you've heard about it.
     It's a free and open-source operating system
     that has been around since the mid-90s.
     Today it is an extensive worldwide community
     that works across many different components of its code base,
     and it has contributors from many different industries.
     It's wildly popular in part because it offers a lot of flexibility
     and configuration options.
     Than most of the other existing operating systems
     that you might actually know about.
     The next component is the Apache web server.
     The Apache web server processes requests and serves
     web assets via the HTTP protocol.
     So the application is easily accessible to anyone in the public domain
     over a very simple URL.
     Apache is mature.
     It's feature rich.
     And it's a large [? double ?] burr community.
     And there are many websites that fundamentally rely on it.
     From an execution standpoint, the process usually
     starts with Apache receiving a request and then,
     if the request needs a PHP file for instance,
     then Apache will pass the request over to the PHP scripting layer, which
     will then load the file, and execute the code that's
     contained within that particular file.
     PHP might also communicate with the MySQL database.
     Now the MySQL database is specifically there
     because it's an open source sort of relational database management system
     that allows you to store application data.
     With MySQL, you can throw all kinds of information
     in a format that's easily queried through a sequel language, which
     PHP is going to be able to orchestrate or interact with.
     So taking a step back up, the reason I'm talking about these four core
     components is really because LAMP Stack is extremely efficient and handling
     not only static content that's out there, but also all
     the dynamic content, where things might change based on the interactions
     that the end user's actually having.
     OK, so now you might be thinking, all right
     I kind of understand the LAMP Stack.
     And I understand that the LAMP Stack runs a good portion of the internet
     and all these websites are built in that.
     So who cares?
     Who cares?
     My gosh the LAMP infrastructure's huge.
     It's the reason that there are over one billion websites on the internet today.
     And over 80% of them are all relying on the LAMP stack.
     So what I'm trying to get you to understand here in a subtle way
     is that we need mature infrastructure in place to scale deployments.
     It's extremely key that you get this point.
     We take web development for granted, because we
     can scale easily on existing infrastructures that are actually
     built up.
     And that is crucial if we want to be able to scale a TinyML.
     Now the web is one particular example.
     Next I'm going to try and get you one step closer towards understanding
     some of the technologies that we need to have in place
     in order to be able to enable TinyML scaling.
     Specifically what I want to talk about is
     just the way that the web is actually scaled out based on lab infrastructure.
     I want to talk about how big ML or cloud ML has been scaling out
     and how it's been easily deployed widely across many different industries,
     based on mature infrastructures that are actually underneath it.
     There are a set of core pillars that are enabling us
     to kind of deploy machine learning at scale very efficiently for big systems.
     Once you understand that, then I'm going to take you over into the TinyML world,
     and I'm going to talk about what is and is not possible today.
     And what needs to be put in place in order
     to be able to robustly scale TinyML deployments.
     That said, I'll see you in the next video where you'll learn about Dockers.
     Ende des Transkriptes, zum Anfang springen
     Downloads und Transskripte
     Video
     Videodatei herunterladen
     Transskripte
     Download SubRip (.srt) file
     Download Text (.txt) file
 
\section{Containers for Scaling ML Deployment}

    0:00 / 0:00Klicken Sie auf die Pfeiltaste nach oben, um das Geschwindigkeitsmenu aufzurufen. Regeln Sie dann mit den Hoch und Runter Pfeiltasten die Geschwindigkeit. Klicken Sie Enter, um die jeweilige Geschwindigkeit festzulegen.Geschwindigkeit1.0xKlicke auf diesen Knopf und diess Video stumm oder nicht stumm zu schalten, oder drücke die HOCH oder RUNTER Knöpfe um die Lautstärke zu erhöhen oder abzusenken.Maximal Lautstärke.Videotranskript
    Beginn des Transkriptes, zum Ende springen.
    VIJAY JANAPA REDDI: In this video, I want
    to introduce you to how BigML deployments take place today
    on the notion of what we call as "containers."
    Now, it's going to be important to understand containers
    so that we can understand how this actually
    relates to TinyML and its scaling.
    Simply put, containers are crucial for scaling deployments at large.
    Containerization is the process of distributing and deploying applications
    in a portable and predictable way.
    It's typically accomplished by packaging up
    application code and its dependencies in some sort of isolated, lightweight kind
    of environments called "containers."
    Now, this is true even for scaling of traditional web code, like LAMP,
    or it can also be true for BigML deployments.
    Now, let's get started and try and understand
    what we mean by shrink-wrapping code.
    When we're talking about any software that is running in the wild,
    you have to realize that there are two key components.
    There is the software, the applications, the services, and so forth,
    but there is also the hardware side of the world.
    On the software side, you have a rich array
    of different environments designed specifically
    to deal with the frontend web services or the backend services.
    For example, when you're dealing with a web service, you've got a GUI frontend.
    And at the same time, you've got a backend database,
    as I alluded in the previous video.
    So as part of having this environment, you also
    tend to have a lot of libraries and dependencies.
    I'm sure you must have come across a time when you're
    trying to run a very simple Python package
    and then you find that there are all these dependencies that
    need to be installed.
    And often, you might also run into the situation
    where some of these dependencies, depending on the application
    you're running, they conflict with one another.
    They don't want to play well together.
    So these are the kind of issues that often creep up
    when we're talking about lots and big amount of code.
    Now, on the hardware side, you also have a lot of rich deployment capabilities.
    You might be serving your site on a public or private cloud.
    In addition to just running the web server,
    you'd often have to worry about your production environment, which
    might include things like reliability, scalability,
    dealing with performance capabilities or performance capacity reasons,
    and so forth.
    There are all these kinds of nuanced things
    that you really have to worry about.
    So given this hodgepodge mix of different software and hardware
    environments, you ought to have some serious thought, something
    that keeps you up at night.
    And if it doesn't, it's a really big concern.
    What do I mean by that?
    What I mean is that, first and foremost, you
    have to have all of the software components
    actually play really well with one another.
    Then on the hardware side, you don't want
    to get locked into any one particular cloud or on-premise service provider.
    You want to be able to move easily, without breaking
    all of your service infrastructure.
    For instance, you might be wondering, why do I need to move?
    Well, what if your current service provider
    raises the prices dramatically?
    Or if, for instance, their security protocols
    are not up to code with what is actually expected for your application.
    So there are many different reasons as to why you might actually
    want to move service providers, be it on the software side
    or on the hardware side, so you want that sort of a portability.
    Now, if you had that sort of a portability,
    then you have flexibility that comes in.
    If you are able to pick up your bags and move easily,
    it also means that you might be able to scale easily across different services
    and providers.
    And that's a very crucial thing.
    Now, how do you ensure that all of these components really work well
    and play well together across the hardware and software realms,
    across different service providers, and so forth?
    This is a really hard and tricky problem to solve.
    This question ultimately leads us to this crazy matrix
    of possible combinations that you have to consider somehow
    if you want to be able to scale easily or you want
    to be able to move deployments easily.
    Poof.
    My mind just basically explodes just looking at this matrix,
    and I hope yours does, too.
    Because it is a serious challenge that one
    has to address in order to scale well.
    But let's not fret.
    These things have crept up in the past in one form or another.
    I often believe that, if we look back on history,
    we'll often find a lot of really good solutions.
    We might have to tweak them a little bit.
    But nonetheless, we'll find the blueprint.
    And I'm going to show you one blueprint here.
    Let's go back to the pre-1960s era and talk about transportation.
    We faced a very similar problem back then.
    We had a rich set of goods that we actually had
    to transport across various domains.
    And we had a rich number of different methods
    as well to store and distribute these goods,
    be it through railroads, shipping trucks, or whatnot.
    Now, so, on one hand, you had to make sure
    that these different goods, when they're all packaged together,
    they can actually be packaged well.
    For instance, you don't want your coffee beans and your spices sitting next
    to each other, because that strong aroma will actually
    creep up one way or the other.
    And you also want to understand, what is the most effective or best way
    to move things so that things are done efficiently and reliably?
    So going back to that matrix, we're back at the same matrix
    again, just different entries.
    We have spices, coffee beans, oil, and so forth.
    And do does software.
    We have railroads.
    We have trucks.
    We have shipping opportunities.
    We have the exact same kind of thing on the hardware front
    in place of some different kinds of computing technology.
    So how did we handle all of this in the past?
    Clearly, we must have dealt with this issue before.
    Well, yes, we have.
    This is where amazing engineers came up with the notion of standard containers.
    These standard containers, which are more typically known
    as "intermodal containers," they have a fixed size in terms of the dimensions.
    So they can easily fit on top of each other.
    They can easily fit next to each other.
    They can also fit easily across various types of transportation,
    be it your railroad, train, or your container ship and so forth.
    Having that sort of standardization really
    facilitates interoperability for a wide range of different products.
    So a standard container is typically loaded
    with virtually some sort of goods or some mixture of goods,
    and it stays sealed until it finally reaches its destination.
    And in between, it can be loaded, unloaded,
    stacked, transported across miles and miles,
    huge distances all very safely and robustly.
    So this is pretty amazing.
    Right?
    So now, going back to our software and hardware dilemma,
    you can imagine that, if we have some sort of standardized container,
    we can address the disparity within the hardware and software ecosystems,
    similar to how our earlier engineers dealt with this problem.
    We can enable interoperability across these systems
    if we have a standard container.
    It's a mechanism that will allow us to have any kind of payload encapsulated
    in a lightweight, portable, and self-sufficient container
    that can be manipulated using standard operations
    and run consistently on virtually any hardware platform.
    And that is awesome.
    To that end, say hello to Docker, who is one particular type
    of an implementation of a intermodal container for computing systems.
    Simply put, Docker is a lightweight, open-source platform.
    It simplifies building, shipping, and running applications.
    It runs natively on Linux and Windows servers.
    It also supports development environments, not just
    server environments.
    So it's a really rich and portable sort of an environment, which is amazing.
    And it solves the really hard issues that we were concerned with as to,
    how do we ensure all these different components work well
    and play well together?
    Docker fundamentally relies on images and containers
    to enable that interoperability.
    Docker images are read-only templates that are used to build containers.
    Let me explain this a little further.
    So you can think of a Docker image like a cooking recipe,
    and the container is like the dish that you've actually
    made following those recipes.
    So what this means is that you can have some sort of a recipe,
    and you can have multiple instances of a container that's been created based off
    of that original recipe.
    And the Docker itself is a shrink-wrapped bubble that allows you
    to kind of package everything that is relevant all closed up together
    so you can easily push that into any kind of heterogeneous environment,
    but that self-contained execution environment has all the dependencies
    it needs and it does not rely on anything external.
    So if you're installing libraries and so forth, you can package all of that
    up real nicely and then ship it all out.
    And you're guaranteed that it'll actually
    work in the production environment.
    Now, typically, this is the kind of thing
    that your software engineer or your DevOps team is really responsible for.
    The software engineer is typically dealing
    with packaging up all the prototype models
    and whatnot you might get from the researchers
    into something like a Docker container.
    And the DevOps team is really responsible for the deployment
    of this Docker container at scale.
    Now, let's make this a little bit more concrete.
    The software developer typically worries about what's
    inside the container, like the code, the libraries, the app, the data,
    and so forth.
    The ops people, they really worry about what's
    outside the container, the logging, the access, the monitoring, the network
    configurations, the starting, the stopping, the migrating
    of the container, and so forth.
    The idea is that you configure this Docker environment
    once so that it can run anywhere and makes
    the whole entire lifecycle of pushing things into production deployment,
    be it for ML or otherwise, very efficient, consistent, and repeatable.
    In a way, it increases the quality of the code
    that is being produced by the developers because things are going to likely run
    robustly, primarily because it eliminates all kinds of inconsistencies
    between development, tests, production environments, including the customer
    environment, which is often very hard to predict for developers.
    And so because of this, Docker has kind of
    become the sort of bread-and-butter platform
    in order to enable widespread deployment very easily and effectively.
    Now, once you have that core container, you
    can load up these different containers with different kinds of items
    without worrying about cross-contamination
    of any kind of dependencies or libraries.
    In our software world, you can imagine packing up
    all the relevant code that needs to be run together
    into respective containers.
    So when these containers are all put next to each other,
    they're not going to interact because the Docker container ensures that all
    of its contents is well self-contained.
    In other words, it's allowing you to have robust guarantees
    when you actually deploy things into the field, which is awesome.
    This is awesome in the context of ML deployment.
    It's awesome both in terms of the CI/CD flows
    to get consistency and repeatability as well
    as when you're actually pushing it out into production deployment.
    And last but not least, recall how I said
    you can actually have different service providers
    and I said that you don't want to get locked into any one of them.
    Well, similarly, down here, you can see how a fixed format in size and shape
    for the containers has allowed ports to actually thrive
    and ships to do business at different kinds of ports.
    Ships could basically come in, unload some containers, move on to other ports
    without themselves being affected by the port.
    And the ports themselves can actually operate independently,
    because the containers give them a fixed structure
    and they can design their cranes and all that good stuff independent
    of how things might have been loaded onto the container ship itself.
    So this really frees up the design space for both the developers and the service
    providers, and everybody wins in this world.
    And that's the beauty about having an ecosystem like containers
    in place, because it enables really smooth deployment of your application
    code, be it ML or otherwise.
    Now, I'm harping on this particular point
    because, as we transition over into the TinyML landscape,
    you'll see how things get really sticky.
    So hang in there, and I'll teach you more about this
    in the next couple of videos.
    Ende des Transkriptes, zum Anfang springen
    Downloads und Transskripte
    Video
    Videodatei herunterladen
    Transskripte
    Download SubRip (.srt) file
    Download Text (.txt) file

\section{Docker vs. Virtual Machines (vs. Virtual Environments)}
    
    \GRAPHICSC{1.0}{1.0}{TinyML/4 MLOps/Dockers_vs._VMs.png}  
    
    Both virtual machines and container software (like Docker) run applications on top of the base operating system. The difference is that virtual machines use a hypervisor to run Guest operating systems that run applications while container software runs applications in containers on top of the main host operating system.
    
    \subsection{About this Reading}
    
    In this reading, we get a quick overview of dockers, virtualization and virtual environments. We want to learn about these things so that we can understand how they are useful for ML model deployments. And this ultimately leads us to the question. Do we have such things for TinyML?
    
    Docker, virtual machines (VMs), and virtual environments are all powerful computing resources at our disposal, but are often confused or poorly understood. This is likely because the difference between them is nuanced and non-obvious even when you have used them. In this reading, we will try to get a clearer understanding of Docker, VMs, and virtual environments, and the pros and cons of each technology.
    
    Docker and VMs are virtualization technologies that allow a user to run separately running systems within a single host system, allowing the segmentation of computational resources. However, the way Docker and VMs achieve this virtualization is slightly different.
    
    \subsection{Virtual Machines}
    
    VMs achieve this at the operating system level, allowing multiple different operating systems to be run on a single machine's hardware. The VM operating system instances are run on a hypervisor, which manages resources across the different operating systems. VMs allow all of the compute resources given to a particular operating system instance to be used across all applications running on the virtual machine. If we have ten apps, they will all share the resources as if we were running it on our desktop computer. We can essentially think of VMs as multiple physical computers compressed into one, instead of needing to buy both a Mac and a Windows, we need only buy one and run both operating systems on top of it.
    
    \subsection{Docker}
    
    Docker achieves a similar goal but at a higher level of abstraction, often called containerization, segmenting computational resources within a single operating system. This is similar to the way that we can have multiple Jupyter Notebooks running on our computer, each with their own separate kernel and some resources attached to the runtime. However, Docker goes one step further than the Jupyter notebook, allowing us to run different operating systems on our main system through something called a container engine (kind of like the hypervisor, but it runs on the operating system and manages the containers). As an example, let’s imagine that I have a MacBook Pro and am interested in running an Apache Web server as well as an Ubuntu instance. The Ubuntu instance will be used to run a complex machine learning model which requires a lot of special setup and is not compatible with MacOS. In this case, we can take a Dockerfile with all of the configuration details, outlining how to configure system-level dependencies (not just Python dependencies), and run it on our system. This Dockerfile will allow us to create a barebones Ubuntu instance which is perfectly configured to run our code, despite the fact that we are running on a MacOS.
    
    \subsection{What’s The Difference?}
    
    Virtual machines have been around for quite a while, and are the basis for most cloud computing, but containers have become very popular in recent years, especially in the academic community as they provide a useful way of quickly verifying publication results and building upon code related to publications. Containers are great in a number of ways: they are very fast to spin up compared to virtual machines, and are smaller in size since only necessary functionality is added to the container. They are also easier to manage and configure, as they can be done using Infrastructure as Code tools such as Ansible. There are also online repositories with thousands of useful Dockerfiles that can be downloaded to run various operating systems, servers, and models. A popular online repository is DockerHub. However, one downside to containers is that applications cannot share resources in the same way as on a virtual machine. If I have ten containers, I must give each a designated amount of resources, and if one needs more resources than I have available, it will crash. Thus, there are load balancing advantages to using virtual machines, as well as some security benefits. However, for most non-specialized uses, containers are often preferable.
    
    \subsection{Where Do Virtual Environments Fit In?}
    
    In your education journey so far, you may have also come across the term virtual environment, which is also related to Docker and VMs and thus worth explaining. A virtual environment is like a container for a particular Python environment. It manages all of our dependencies and ensures that we do not get stuck in “dependency hell” where none of our code will run due to incompatible libraries. The most popular virtual environment today is Pipenv, which comes in the form of a “Pipfile” and a “Pipfile.lock”. The former stores a list of all the dependencies, while the latter stores the dependencies along with hashes of their libraries.
    
    You can think of a virtual environment as a Docker container which only manages Python dependencies. In fact, many Docker containers used for machine learning often utilize Pipenv within them to manage their Python dependencies!
    
    The difference between these three is a common source of confusion among students of ML, and thus we recommend spending some time becoming familiar with the differences between them to make sure you understand the differences. If you are feeling confident, maybe try running a virtual environment within a couple of Docker containers, where one of each is placed on a separate virtual machine - good luck!
    
    \subsection{Why Dockers for ML Deployment?}
    
    Now that we have a basic understanding of Dockers, lets see why they are useful for machine learning deployment? Creating a machine learning model for our computer is simple. But deploying the model at scale, on various types of servers globally, is more challenging and difficult. It is highly possible that if you design a model and then move it to production or another server, it may not operate properly on other systems. Many things can go wrong. You can run into performance issues, crashes, and even poor optimization. Even while our machine learning model can be implemented in a single programming language (e.g. Python), the application will need to connect with apps written in other programming languages. With Docker, you get granular updates and reproducibility. With Docker, you can shrink wrap things as “microservices,” allowing for scalability and easy addition or deletion of independent services.
    
    When a model is complete, the ML engineer or the data scientist responsible often fears that it will not reproduce real-life findings or that it will be shared with teammates. Sometimes it's not the model, but the need to recreate the entire stack. Docker allows you to simply replicate the machine learning model training and testing environment anywhere.  A training model can be built locally and quickly moved to external clusters with greater GPUs, memory, or CPU power. Moreover, it is easy to distribute the model globally by packaging it as an API in a container and using technology like OpenShift, a Kubernetes distribution (which we will learn about a bit later).
    
    Containerization of ML applications is also simple since we can develop containers using templates and use an open-source registry of user-contributed containers. Docker allows developers to track container image versions, see who produced them, and rollback to prior versions. Finally, your machine learning program can execute even if one of its services is updating, fixing, or unavailable. For example, updating an output message integrated in the entire solution does not necessitate updating the entire program and disrupting other services.
    
    Hopefully now you see how Dockers are powerful constructs for ML deployment. But do such constructs exist for TinyML deployments that don’t have even a bare metal operating system? We’ll dig into that next.
    
\section{Challenges for Scaling TinyML Deployment (Part 1)}

    0:00 / 0:00Klicken Sie auf die Pfeiltaste nach oben, um das Geschwindigkeitsmenu aufzurufen. Regeln Sie dann mit den Hoch und Runter Pfeiltasten die Geschwindigkeit. Klicken Sie Enter, um die jeweilige Geschwindigkeit festzulegen.Geschwindigkeit1.0xKlicke auf diesen Knopf und diess Video stumm oder nicht stumm zu schalten, oder drücke die HOCH oder RUNTER Knöpfe um die Lautstärke zu erhöhen oder abzusenken.Maximal Lautstärke.Videotranskript
    Beginn des Transkriptes, zum Ende springen.
    VIJAY JANAPA REDDI: In this video, I'm going
    to start off by talking about some of the challenges
    with scaling TinyML deployments.
    Let's pick up where we left off.
    Earlier, we talked about containers as being extremely important
    for scalability of BigML models.
    In this video, we're going to try and understand
    what makes scaling in the context of TinyML hard.
    In the next few videos, I'll give you a couple of suggestions or solutions.
    Let's start with the traditional computing paradigm.
    When we have the traditional computing paradigm,
    where we also have endpoint devices, the endpoint device
    is really just there to get the data.
    And we send the raw data up to the cloud to do the inference.
    Now, this endpoint device is a dummy device--
    not really very, very powerful.
    And by and large, it's really relying on the massive computational capabilities
    that are up in the cloud to run the machine learning models.
    Now, in such a model, when we want to scale, it's very easy.
    Why?
    Because all we would have to do from the endpoint device perspective
    is just simply replicate the device.
    So for example, if you're talking about a Google Assistant
    or you're talking about an Alexa device, you could just
    imagine having more of those devices, and they're all
    talking to a main server.
    And, of course, this leads into the question about scalability
    on the cloud computing side of things.
    Now, if you have things like Docker containers,
    then it becomes very easy to scale out.
    And that is why I was actually talking about Docker containers earlier
    and containerization, because in order to be able to scale,
    you need to have the right building blocks in place.
    And Docker was a beautiful way of being able to easily
    package things up and then push things out into the field.
    However, we cannot always rely on just cloud computing capability,
    because for instance, let's say latency is a critical issue for us.
    In that case, we can't rely on the cloud,
    because it's going to be taking quite a long time to get the data up and bring
    the data back.
    Think about this in the context of some sort of manufacturing plant, where
    you're trying to actually use TinyML sensors and the computation
    nearby to try and figure out how to control that little widget.
    If that's the case, you're going to have some sort of real-time constraints.
    That makes it extremely hard to be able to rely on the cloud continuously.
    So therefore, this paradigm of edge computing
    has emerged, where in order to balance that latency requirements,
    we try and bring the computing closer towards the data.
    It's a step in the right direction.
    It's basically a mechanism or a paradigm where,
    rather than bringing all of the data to the compute, we flip the model
    and bring some of that computational intelligent processing capabilities
    closer to the data so we can lower things like the energy consumption.
    We can guarantee better privacy around the data.
    We can try and improve latency, reliability, bandwidth,
    and all those good things.
    Now, edge computing is definitely a step in the right direction.
    However, it too can actually have some critical shortcomings.
    One of the key things that we're assuming still
    when we have an edge computing system is that the endpoint device would still
    have to be communicating to the edge server.
    Yes, it is not communicating to the cloud.
    It's communicating to the edge, which is undoubtedly better in terms
    of latency, performance, and so forth.
    However, if you remember what we say in terms of TinyML,
    we say that communication costs are extremely costly with respect
    to computation costs.
    Communicating with any kind of device can
    be as much as two, three, or four orders of magnitude more costly than actually
    just doing the computation locally.
    So therefore, if we're assuming that, again,
    these endpoint devices that we're trying to convert into intelligent devices
    are all relying on the edge server, it still
    presents a real problem for us because of the energy consumption argument.
    So that's really where the embedded computing paradigm comes in,
    where the entire idea is that we want to have these devices effectively
    be operating locally with one another.
    If needed, they try and communicate with each other locally
    under some minimal requirements, and then we'll
    talk about communication protocols and so forth
    in another stage in the MLOps pipeline.
    But for now, let's just assume that they're ways to communicate.
    And that way, they only communicate when needed.
    And more often than not, they're just doing
    the computation for the machine learning data right at the endpoint.
    Now, at this point, it should be clear as
    to what the key differences are between cloud computing paradigms,
    and edge computing paradigms, and the embedded computing paradigms.
    And this is going to be important for us to understand the scalability aspects
    as we move forward.
    As you can see on the cloud side, scalability really isn't an issue,
    as you only have to worry about scalability of those kinds of container
    environments.
    As more endpoints show up, you can easily add them in,
    as long as you've got the right infrastructure on the cloud side
    to scale.
    On the edge side, you see the benefit of better latency and so forth,
    but the devices still have to talk to the edge server.
    If you think about TinyML, since communicating with the cloud
    or with the edge server is costly, we actually
    have to scale that back as well.
    That's where the embedded model kind of comes in, where most of the inference
    is really happening at the endpoint locally, rather than relying
    in any kind of edge or cloud servers.
    Now, conceptually, everything should be pretty sound and solid here
    because there are very clear and distinctive differences
    between the cloud, the edge, and the embedded [? DRAM. ?]
    But there are issues in trying to realize this embedded paradigm
    at heart, because there are a lot of things that are brushed under the rug,
    assuming that you can easily get these devices to communicate,
    assuming that you can easily push the models down to these devices,
    and so forth.
    Let's be a bit more specific about that.
    On the cloud side, we generally tend to have rich resources, right?
    We tend to have really nice programming languages around.
    We tend to have nice, big model zoos and infrastructure to support that.
    We have stable frameworks, like containers
    that allow us to actually easily package up software and deploy them robustly.
    On the embedded side of things, a lot of things
    are still in the bare bones stage.
    We often don't have highly portable machine learning frameworks.
    There is no concept of containers really.
    I'll be getting more into that soon.
    And all of this means that the lack of these kinds of basic building blocks
    makes deployment really hard.
    Even if a cloud service wants to push a model to an embedded device
    so that the embedded computing paradigm can be the predominant way of how
    things execute, the issue becomes, how does this cloud actually
    communicate with this endpoint device?
    There are also many issues around the fact
    that a lot of the devices in the embedded ecosystem are heterogeneous.
    I'm not going to go back into the details of the heterogeneity,
    because at this point, I've talked about it many times over, from memory,
    from clock speeds, from connectivity, and so forth.
    There are lot of these different constraints
    that really limit us from being able to effectively deploy things.
    Now, if we were to look at things on the hardware side of the world
    and compare how the cloud and the embedded side compare,
    you would see that, on the cloud side, we
    tend to have large machine learning resources where they're
    able to run on big beefy machines.
    We have huge amount of compute, memory, energy storage, and so forth.
    All of those make the deployment of machine learning much more easier.
    However, yet again, on the embedded side, we're resource-constrained.
    Compared to megabytes and gigahertz, we're
    dealing with kilobytes and megahertz ranges.
    So there's definitely an order of magnitude difference
    in the capabilities that we have on the embedded side, the TinyML side,
    versus what we actually see on the cloud ML side.
    For starters, we don't even have simple things,
    like a homogeneous operating system environment
    that we can build our services on.
    And this leads into a lot of challenges.
    So let me close this video up by summarizing the key differences
    for the various paradigms.
    Cloud computing has lots of resources.
    As we move things closer to the data, the edge and embedded ecosystems
    become much more interesting and capable.
    Now, that said, though, given the extreme severe constraints
    of the embedded ecosystem, there is still much to be desired.
    So this begs the question now, how do I actually push models when I'm saying,
    quote unquote, "I want to deploy models at scale in TinyML?"
    Then how do I actually orchestrate the execution?
    Are the models actually running well?
    Is the model even running?
    Has the system crashed?
    How do I do this very basic kind of a thing, which
    we take for granted, thanks to containers, and Dockers, and Kubernetes
    clusters, and so forth-- all of which we have built up
    over many years in the BigML system?
    So what I'm trying to really get across here
    is that, when we want to scale ML deployments,
    we build on top of all the prior work.
    In the cloud ML side, we have done a lot of work
    to build up a lot of this infrastructure on the edge, and more specifically,
    in the embedded ML side, we don't have that infrastructure.
    This is why scaling becomes really hard and problematic for us.
    So let's try and understand what are the nuanced issues that will actually
    creep up in the next couple of videos.
    Ende des Transkriptes, zum Anfang springen
    Downloads und Transskripte
    Video
    Videodatei herunterladen
    Transskripte
    Download SubRip (.srt) file
    Download Text (.txt) file    
    
 \section{Challenges for Scaling TinyML Deployment (Part 2)}
 
     0:00 / 0:00Klicken Sie auf die Pfeiltaste nach oben, um das Geschwindigkeitsmenu aufzurufen. Regeln Sie dann mit den Hoch und Runter Pfeiltasten die Geschwindigkeit. Klicken Sie Enter, um die jeweilige Geschwindigkeit festzulegen.Geschwindigkeit1.0xKlicke auf diesen Knopf und diess Video stumm oder nicht stumm zu schalten, oder drücke die HOCH oder RUNTER Knöpfe um die Lautstärke zu erhöhen oder abzusenken.Maximal Lautstärke.Videotranskript
     Beginn des Transkriptes, zum Ende springen.
     SPEAKER: In this video, I'm going to continue
     discussing some of the challenges of scaling TinyML,
     before we get into some potential solutions to deal with the issues.
     Specifically in this video, I'm going to focus very much on the challenges that
     emerge out of a single device.
     Contrast this with the previous video.
     In part one of the scaling issues, we talk broadly about the ecosystem.
     What are the trade offs that we see between a big computing
     system in the cloud context, versus an edge, versus an embedded?
     Here we want to go one level deeper.
     Once we understand that role, understand how we can potentially
     address the solutions.
     Now as a quick recap, if you remember, we
     have severe resource constraints, both on the hardware side
     as well as on the software side.
     On the software side, we don't have the rich and mature systems, like Windows
     and Linux, and all of that rich packaging and deployment services,
     like Docker Kubernetes, and so forth, in order
     to orchestrate the scalability aspects of things.
     Now in the embedded side, because we don't
     have that level of mature systems, it makes scalability really hard.
     And this gap between the cloud and embedded system
     is made worse by the fact that there's heavy fragmentation
     in the embedded ecosystem around machine learning,
     even when you consider one single device.
     Let me talk about this a bit more detail through this vertical stack
     that I'm showing you here.
     Now when you Zoom into this, a single embedded device that is actually
     running machine learning, you see this entire vertical stack, all of which
     has impact on what the ML artifact, at the end of the day, is.
     Starting from the bottom and working our way up,
     we'll see loads of specialization and fragmentation
     that leads to deployment challenges that limit our scaling ability.
     For instance, there are lots of different embedded devices out there.
     They come in many different shapes and sizes.
     More specifically, if you open up the hood
     and you look at the microcontroller units that are there,
     there is a rich variety of them out there.
     You have MCUs that span from being a Cortex-M class, M0 class
     microcontroller, all the way up to Cortex-M33 class microcontrollers.
     And in the future, there are going to be more.
     And each of these microcontrollers comes with different instruction set
     architectures, which means you cannot easily take code that's running on one
     device and be able to just plop it onto another device.
     But I'm pretty sure if you take the code that's running on my laptop
     and you run it on your laptop, it's going
     to run because it's x86 [INAUDIBLE].
     Now just to make matters even more worse,
     because of our interest in keeping these devices at the ultra-low power
     consumption level, we often end up making heavy specialization decisions
     as to whether we have an FPU or a floating point unit,
     whether we provide SIMD instructions, whether we have a DSP on board.
     Now all of this leads to an insane amount of fragmentation
     at even the lowest level of the hardware.
     Going one level up, you'll find that the underlying MCU, or accelerator, can
     be used in many different ways, depending on whether you have support
     for optimization kernels or not.
     For instance, here I'm showing you on the left, for arm microcontrollers,
     you may or may not have CMSIS kernel support.
     Real briefly, CMSIS-NN is a software library
     that is made up of a collection of efficient neural network kernels
     that are designed specifically to maximize the performance
     and minimize the footprint of neural networks
     on ARM Cortex-M class microcontrollers.
     The library comes with a whole load of different functions.
     It comes with specialized convolutional functions, activation functions,
     fully connected layer functions, pooling functions, [? software ?] functions,
     and even many other basic arithmetic functions that are highly tuned
     for the Cortex-M class microcontroller.
     So if you have it, then that's great.
     This is important, because if you-- whether you have it or not
     has an impact in terms of what the performance characteristics are
     or whether the code is even going to be portable or not across devices.
     Again, keep in mind, we're talking about scaling.
     So we're talking about having multiple devices here.
     And as I've alluded to during the model conversion stage of the MLOps pipe,
     there are a load of different inference frameworks out there.
     And while amazing folks are working on things like this onyx framework that
     tries to be the standard that allows interoperability across frameworks,
     I mentioned in that particular series of videos
     that we're still a long way from getting there.
     And even before you go into the deployment stage,
     there are loads of compiler tool chains that actually lead to these models
     that we're actually generating.
     And each tends to have its own little quirky things.
     So if you recall, for instance, we talked about issues like MMS operators
     that can actually lead to differences, depending
     on how the frameworks are actually implementing MMS operations.
     Going one level up, there are lots of different optimization algorithms
     like the ones we talked about, quantization, pruning,
     weight clustering, knowledge distillation, and so forth.
     So hopefully, you're now starting to appreciate
     why I've talked about all these things in the past
     as that complexity starts piling up as you go more
     and more towards the end of the MLOps.
     And that's where things start really getting hard for deployment.
     All these different pieces have to click in some capacity.
     Last but not least, when you're looking at neural network architectures,
     you can have different ways of actually building them,
     as we discussed previously with the neural architecture
     search and the AutoML pipeline.
     Of course, this design space can actually
     get much larger if you consider the facts
     that we can do feature engineering.
     We can do neural network hyperparameter tuning
     and also have auto optimizing compilers.
     So space is almost boundless down here.
     Bottom line, what I'm trying to get is that even if you just
     take a single device, there are so many different possibilities,
     because the stacks are so complex.
     They give so much flexibility.
     It almost creates us new headaches that we have to deal with.
     So this leads into this serious cross-product diversity
     problem, which really limits our ability to scale and efficiently deploying
     TinyML.
     So how do we cope with these things?
     What are the things that we have to worry about?
     Well, first and foremost, having specialized hardware
     at the lowest levels is important, because that's
     where you get energy efficiency from.
     But that, somehow, has to be balanced with the right use of compilers
     and optimization frameworks and tool chains,
     because you don't want to toss it out in order to simplify the deployment.
     That specialization is what makes things tick efficiently.
     So somehow, we have to support that part.
     We can't throw it out.
     Then when we talk about the compiler ecosystem
     on all these different libraries and so forth,
     well, that also leads to more fragmentation.
     However, yet again, in order to make the best utilization of the hardware
     that you have, you have to use these compilers,
     because that's where you get the execution
     efficiency when you look at it from a single device's output perspective.
     Ultimately, however, when you think, take a step back,
     and you look at the bigger picture with this stack,
     and then imagine it's replicated multiple times over,
     across a whole pool of different embedded devices out there,
     flexibility becomes a major nightmare for people like you and me,
     because switching from one hardware context going to another hardware
     context would effectively mean that we'd have
     to recompile the inference model for the target device.
     So given that the TinyML stack has so many levels of heterogeneity,
     from the hardware through to the compilers and above,
     the key challenge that we face is this notion of flexibility.
     And somehow, we have to address it.
     But have no fear.
     Having an appreciation and understanding these levels of heterogeneity
     is actually a really good thing, because remember,
     how despite the fact that we have rich hardware and software heterogeneity,
     folks went out and came up with the idea of this container.
     Well, they were able to do that because they had a strong understanding of all
     the factors that were actually impacting the ability to move goods
     through ecosystem efficiently.
     Similarly, with a strong understanding of all these issues, what we can do is
     try and think of innovative solutions, which ultimately leads us
     to TinyML as a service.
     Which we'll learn about next.
     And I'll see you then.
     Ende des Transkriptes, zum Anfang springen
     Downloads und Transskripte
     Video
     Videodatei herunterladen
     Transskripte
     Download SubRip (.srt) file
     Download Text (.txt) file
 
 \section{Challenges of Scaling TinyML Deployment}
 
 The TinyML computing stack is complex! As we just discussed there are:
 
 \begin{itemize}
  \item Lots of different model optimization methods
  \item Lots of specialized compilers for embedded devices
  \item Lots of inference frameworks for embedded deployment
  \item Lots of software kernels for unlocking the hardware’s performance
  \item Lots of specialized hardware and ML accelerators
 \end{itemize}


 Which of these causes of complexity do you think will have the greatest impact on the kinds of TinyML applications you are excited to deploy at scale? Which will not be an issue? How might you begin to approach solving some of this complexity?
 
\section{Scaling ML Deployment Needs Services}


\GRAPHICSC{1.0}{1.0}{TinyML/4 MLOps/Herunterladen.png} 

A figure showing that anything as a serivce encompasses: software-as-a-service, platform-as-a-service, infrastructure-as-a-service, and humans-as-a-service.

\subsection{Overview}

In this module, we will be focusing on a crucial part of model deployment known ‘as a service.’ This extends well beyond the Continuous Integration and Continuous Deployment requirements.

Often when you have to deploy models into production, you have to have the right infrastructure in place. One can choose to build all of the components themselves or try and lease them from others who are providing such services for a fee—leasing alleviates the customer to focus on their problem at hand, rather than build all the components from scratch with little to no experience just to deploy their solution. For instance, you can set up a web server with a push of a button today, thanks in part to all the “LAMP” (Linux, Apache, MySQL and PHP) infrastructure components that have all matured well and been integrated well with one another. This allows you to focus on the website you are designing, rather than figuring out how to bolt together all these different major software components together just to get your website hosted.

Similarly, when we want to deploy ML models into production at scale, we need the right underlying software infrastructure in place. Arguably, many of these components are still missing and if not they are limited in their capabilities. To understand the challenges and opportunities, we will be introducing Tiny Machine Learning as a Service (TinyMLaaS).

\subsection{Cloud Computing}

To talk about the anything-as-a-service (XaaS) paradigm, we first have to become familiar with the concept of cloud computing. When it comes to TinyMLaaS, there are challenges that emerge with scaling and deployment with respect to embedded systems via the cloud. So by understanding the different service models, we will be able to appreciate the challenges better.

Cloud computing refers to the use of externally hosted computational resources. Common platforms for this include the Google Cloud Platform, Microsoft Azure, and Amazon Web Services. Things that can be hosted on the cloud are applications (e.g., Dropbox, Slack), development environments (e.g., Amazon Sagemaker, Azure ML Studio), and even entire virtual machines (e.g., hosted websites, network devices). Cloud computing is an incredibly powerful idea, especially in the business world, as it reduces the need to have on-site computational resources. Essentially, this eliminates the need to spend capital on expensive servers and networking infrastructure, as they can instead be managed externally by a cloud provider. This essentially transfers capital expenses to operational expenses, making infrastructure management much cheaper. In the modern-day, most of our favorite websites are hosted by one of these cloud providers (e.g., Facebook, Netflix, LinkedIn, and Twitch are all hosted on Amazon Web Services).

\subsection{Anything-as-a-service}

\textbf{Anything-as-a-service (XaaS)} is a relatively new computing paradigm that refers to hosted cloud computing resources that are externally accessible to a customer. If you have ever used Google Docs, watched a video on Netflix, or hosted a website on Amazon Web Services, you have interacted with some form of XaaS. Although this concept may sound intimidating initially, the service essentially removes any complication to the user by providing them with the correct level of computational abstraction to perform the service they require. There are three common forms of XaaS that we will briefly describe.

\textbf{Software-as-a-service (SaaS)} is conceptually the simplest, and most common, of the as-a-service paradigm. In SaaS, the entire computing ecosystem is abstracted at the level of the application, with which the user is able to interact. The rest of the virtual machine is hosted externally through cloud resources, which are not visible to the user. This is why an internet connection is needed to use these services, which may only intermittently interact with the cloud resources (as in data storage and text editing applications), or maybe continuously streamed (as in video streaming applications). SaaS is great for many end-user applications, but provides relatively little flexibility, meaning it is not useful in developmental environments. This is where the other XaaS types come into play.

\textbf{Platform-as-a-service (PaaS)} refers to a set of resources that are contained within a virtual machine. For example, a suite of coding resources hosted on a virtual machine might be classed as PaaS. Naturally, this is harder to use than a single hosted application but offers more flexibility to the end-user, and is thus more useful for certain developmental applications (e.g., video editing or video game design).

\textbf{Infrastructure-as-a-service (IaaS)} is the lowest-level abstraction of XaaS, which is the level of the virtual machine. This is what most companies use to host their resources, as it offers maximal flexibility since the entire system (apart from the hardware itself) can be managed externally by the end-user. This is ideal for hosting websites, and even for an individual to host their own PaaS or SaaS applications on top of the infrastructure.

\textbf{Humans-as-a-service (HaaS)} is an emerging concept that is quite relevant to machine learning systems. In the vast majority of cases, we need humans to label the ground truth datasets for machine learning. These “human” services can be outsourced to the cloud through HaaS using existing platforms such as Amazon Mechanical Turk (AMT) and ClickWorker. Thanks to AMT we can readily tap into a wide range of people that would be willing to perform tasks that our computer systems cannot (yet) perform as efficiently as humans.

Of course, there are pros and cons to using the HaaS approach in practice that one should be aware of being considering it as a resource. For example, AMT offers a low rate of pay, and provides no incentives for providing high-quality responses, which may lead to an overall lower-quality of work. Additionally, services that require some form of specialization (such as AMT workers pretending to be a judge and deciding whether bail is granted in a particular legal case, testing the effectiveness of an algorithm-in-the-loop analysis) may not provide representative results of the actual implementation. Despite this, HaaS has already proven to be a highly useful resource for many datasets, such as labeling larges swatches of images, text, and audio data.

\subsection{Conclusion}

Why is XaaS important in the context of TinyML? Well, because we are now beginning to see applications that are described as machine-learning-as-a-service (MLaaS), which leverage cloud resources to perform machine learning tasks. For example, the Google Colaboratory, which we are all familiar with from this and previous courses, can be described as MLaaS. Naturally, this can be extended to TinyMLaaS in the context of TinyML, which depending on the context could mean several things. The remainder of this section will discuss the concepts of MLaaS and its extension to TinyMLaaS, as well as the unique challenges that come with them. 

\section{TinyMLaaS (Part 1): An Introduction}


    0:00 / 0:00Klicken Sie auf die Pfeiltaste nach oben, um das Geschwindigkeitsmenu aufzurufen. Regeln Sie dann mit den Hoch und Runter Pfeiltasten die Geschwindigkeit. Klicken Sie Enter, um die jeweilige Geschwindigkeit festzulegen.Geschwindigkeit1.0xKlicke auf diesen Knopf und diess Video stumm oder nicht stumm zu schalten, oder drücke die HOCH oder RUNTER Knöpfe um die Lautstärke zu erhöhen oder abzusenken.Maximal Lautstärke.Videotranskript
    Beginn des Transkriptes, zum Ende springen.
    In this section, I'm going to introduce you to an emerging concept called
    Tiny Machine Learning as a Service.
    TinyML as a Service, we aim to build a high-level abstraction
    of TinyML software that is as hardware and software agnostic as possible.
    If you recall, there's a lot of heterogeneity at the embedded system
    layer.
    We have a rich array of different hardware,
    we have many different kinds of compilation tools,
    and we have many different optimizations, all of which
    we can throw together.
    What this means is that we end up with a highly-specialized solution
    at the end of the day, even when we're looking at a single endpoint device.
    Now, when we talk about scaling this out, we often get into this pickle
    that we have too much fragmentation in the ecosystem
    because every one of these vertical stacks
    is going to look quite different when we scale things out.
    So that means that if I was to try and take some piece of code that's
    running on an embedded system and try and move it over to a different device,
    sort of how we talked about things in the context of dockers and containers,
    that's not going to be possible in the embedded ecosystem.
    The truth is that we don't want to expose this level of deep heterogeneity
    and fragmentation that's at the embedded ecosystem layer at the endpoint device
    all the way up to the cloud device, because that would complicate things
    at the cloud computing layer.
    So what this means is that we ideally want
    to be able to decouple the Cloud from the embedded ecosystem layer.
    The reason we want to do this is that we want the cloud computing
    system to be doing what it's already doing best
    and to keep doing it even better.
    In other words, we want our ML developers,
    who are working with big workstations and all that good stuff,
    to be able to focus on the experimentation
    and prototyping as in the ML development stage.
    And then they should enjoy all of the rich benefits of the software ecosystem
    for things like continuous integration and continuous testing and so forth.
    Now given that we want to keep that intact,
    then how can we decouple the Cloud from having
    to know all the details of the embedded ecosystem?
    Well, any time you want to get into any kind of decoupling in the computing
    environment, the answer is almost always create a level of indirection
    by introducing another abstraction layer.
    Abstractions are key, because it allows us
    to separate what it is that we're trying to do,
    from how it is that we actually go about trying to solve the problem.
    This is because we want the Cloud to be able to focus on just the model
    side of things, and we want to have a generic abstraction that
    allows us to transparently push those models to those embedded devices
    without exposing all the gory details of the Android ecosystem
    further up the food chain.
    And that is crucial.
    And abstractions have always worked well for us in the past.
    The reason abstractions work well is because when
    we introduce that abstraction layer, things
    that are communicating with things that are below the stack--
    they can be decoupled from what's on the top of the stack.
    In other words, the embedded ecosystem--
    that information no longer needs to be known at the cloud computing layer,
    because you can have some sort of custom format in the middle that
    allows this intermediate layer to actually just deal with the Cloud
    independently and have this intermediate layer
    also just deal with the low level details independently.
    That's beautiful in reality, and it's used very well in many different areas.
    Now I just want you to keep in mind that as we
    get into this concept of creating an abstraction layer, which
    is what TinyML as a Service is, I want you
    to keep in mind that TinyML as a Service is still in a nascent stage.
    It's still conceptual by many means.
    It was developed by folks at Ericsson, who
    sort of foresaw this problem coming down the pipe,
    as we want to scale TinyML deployments.
    And it's not yet fully backed by a real implementation out there yet.
    In other words, you can't quickly get your hands on something.
    The reason I'm mentioning this is that I don't
    want you to be disappointed tomorrow that, when you look up
    TinyML as a Service, you're not able to just
    quickly ready to use some software.
    But nonetheless, the reason I'm bringing this up
    is because I think it's got all the core fundamental building block concepts
    in place for prime implementation.
    So the way you should be interpreting what's coming next
    is really thinking about that these are some potential ways of solving
    the fragmentation problem in the embedded ecosystem
    so that you can more easily deploy models.
    Allow me to define TinyML as a Service.
    TinyML as a Service is a cloud or edge machine learning service.
    It's an intermediate layer whose job is to simplify
    the deployment of machine learning models
    into resource constrained devices.
    Now hopefully at this point, you understand the notion of something
    as a service, and if not, you probably skipped my reading
    and should go back and make sure you reread it.
    I tried to bring in something as a service,
    and I tried to show you things like SaaS, PaaS, IaaS,
    and so forth, because these are concepts that have helped us actually
    build upon each other in order to create some new services.
    And that's extremely important, because TinyML as a Service is, by and large,
    going to be building on some of those prior concepts.
    Similar to how dockers and so forth are actually
    solutions that have been built up on very mature software stack.
    So please go back and make sure you read that,
    because I need you to internalize that for this context.
    All right, so how does this really work?
    The input to this TinyML as a Service intermediate layer
    is really simply just going to be a model, and the output of it
    is going to be a binary image that fits the unique characteristics
    of the endpoint device.
    That's all there is.
    And what the TinyML as a Service is really doing
    is that it's hosting a wide set of different compilers
    with target different devices with various kinds of optimization
    capabilities.
    And it's the job of these compilers to actually convert a specific model
    into an appropriate format that's designed to target the endpoint device.
    Now in addition to providing just the intermediate abstraction
    layer that de-couples the Cloud from the embedded ecosystem of all the details,
    it's also the job of this layer to be able to figure out
    how to best customize that binary image for the target device.
    In other words, what kind of quantization, pruning, clustering,
    optimization-- various kinds of optimizations-- should be done?
    And that really is crucial, because that's
    what is really going to enable interoperability
    for various kinds of devices in this ecosystem.
    And that is how we'll be able to scale.
    Now let me give you a very high level overview of how TinyML as a Service
    would operate in four main steps.
    The components that you're seeing here are
    the Model Zoo that's up in the Cloud.
    Then you've got TinyML as a Service and you've got these endpoint devices
    at the bottom.
    Step one is that you want TinyML as a Service
    to figure out what the characteristics of the endpoint device are,
    so that it knows how to actually tailor a model for that endpoint device.
    So therefore it has to go and gather some information,
    like the CPU type, the amount of RAM or memory it has, the amount of flash
    size it has, and so forth.
    Because that information will come in quite handy in step two, where
    it's going to try and figure out what's the most suitable compiler that
    will actually be best for generating the target code
    that needs to be deployed on the end device.
    This could also mean that down here you, might do something fancy like NAS
    and apply various kinds of optimization, some of which
    we have talked about earlier.
    And then the next step is to really go ahead and build an OS image
    that actually maps to that particular target endpoint device,
    because different embedded system devices have different firm arrays
    and bootloaders that need to be worked with in a very specific way.
    And these are the details that we don't want
    to expose up to the higher levels of the cloud computing stack.
    And finally, we want to push the built image down into the endpoint device.
    Now, you can use things like a lightweight M2M protocol which is
    well-established for actually doing either a firmware over the air,
    a photo update, or a SOTA update--
    Software Over The Air update.
    There are lots of different ways of doing these updates.
    There are some well-established protocols like the M2M protocols
    that are industry standard, so if we build on that,
    then we can actually be building on existing technology
    which will enable us to have a sustainable solution around TinyML
    as a Service.
    So as you can see, the process is fairly simple when you break it down.
    So in summary, TinyML as a Service is a cloud or edge-based machine
    learning as a service that greatly simplifies the deployment of machine
    learning models into resource-constrained devices,
    and it develops this abstraction in order to ensure interoperability
    in the ecosystem.
    Now ideally, this would be the solution that
    will work regardless of what the platform capabilities are and so forth.
    Now going back to this diagram, though, now I
    hope that you're starting to appreciate the need for something
    like TinyML as a Service.
    I started off talking about the websites,
    and the reason is because the scaling of the websites
    to a billion different websites and it was enabled by the fact
    that we had the LAMP stack.
    Now, when it came to big machine learning models,
    the reason we were able to scale out is because we
    have mature infrastructure like dockers and containers
    and Kubernetes clusters in order to orchestrate these things.
    So when it comes to embedded devices, given
    that we have all these unique, quirky characteristics, it's important for us
    to understand how to tackle them systematically.
    And hopefully, what I've talked about here
    gives you a glimpse into what needs to be put in place.
    And in the next video, I'm going to peel the layers just a little bit further
    in order to help you understand what needs
    to be implemented in order to make TinyML as a Service click.
    And I'll see you there.
    Ende des Transkriptes, zum Anfang springen
    Downloads und Transskripte
    Video
    Videodatei herunterladen
    Transskripte
    Download SubRip (.srt) file
    Download Text (.txt) file

\section{TinyMLaaS (Part 2): Design Overview}

    0:00 / 0:00Klicken Sie auf die Pfeiltaste nach oben, um das Geschwindigkeitsmenu aufzurufen. Regeln Sie dann mit den Hoch und Runter Pfeiltasten die Geschwindigkeit. Klicken Sie Enter, um die jeweilige Geschwindigkeit festzulegen.Geschwindigkeit1.0xKlicke auf diesen Knopf und diess Video stumm oder nicht stumm zu schalten, oder drücke die HOCH oder RUNTER Knöpfe um die Lautstärke zu erhöhen oder abzusenken.Maximal Lautstärke.Videotranskript
    Beginn des Transkriptes, zum Ende springen.
    VIJAY JANAPA REDDI: Recall that TinyML as a service
    is a cloud or edge-based Machine Learning as a Service
    that attempts to greatly simplify the deployment of machine learning models
    into resource-constrained devices, and it does so
    by creating an abstraction that guarantees
    some sort of interoperability, which is crucial for us.
    Now, the three main components to Tiny Machine Learning
    as a Service, the compiler plugin, the orchestration protocol,
    and the inference module format.
    Let's dig into each one of these a little bit
    further just to have an appreciation for what they are going to be doing.
    The compiler plugin acts as a middleware layer
    between the TinyML as a Service front-end
    and the TinyML as a Service back-end.
    The front-end receives in response to requests from a wide variety
    of different linked devices.
    And these requests attempt to use a certain ML compiler
    from a library of different compilers that we're
    going to have inside the TinyML as a Service Engine.
    For instance, to compile down to Arduino or Spark
    or Texas Instrument or a Synthon device, you're
    going to have to have the right compiler.
    Now, one of the key things that this stage also does
    is it tries to figure out what are the customizations to perform for a given
    device's unique characteristics.
    Recall that we usually query the device first and then try and figure out
    what to do.
    That's why we can actually do those optimizations in this plugin interface.
    For instance, if the model has really tight resource constraints,
    you can apply a certain class of optimizations.
    Now, the next step is going to be around the orchestration protocol.
    Now, the orchestration protocol plays a key role
    in terms of the interactions between the end devices
    and the TinyML as a Service infrastructure.
    Now, we do this through some sort of well-defined APIs, of course.
    The functionality at large is really twofold.
    First, the orchestration protocol interacts with the end devices
    in order to gather information about their baseline characteristics, which
    I alluded to earlier.
    And it tries to understand what are some of the capabilities and constraints
    that the system is dealing with.
    This allows us to ensure that any request coming from any given device
    is properly managed by the TinyML as a Service back-end
    as in turn this will tailor the compiling process exactly
    according to the information that's transmitted by the orchestration
    component.
    Second, it also offers us the Firmware Over the Air or Software
    Over the Air updates that I mentioned in the previous video.
    You can use something like this lightweight M2M
    protocol for the orchestration protocol because those kinds of protocols
    are already well established, and they're standard industry features.
    Now, if you're not familiar with the lightweight M2M protocol,
    it's a standard interface that I'll give you
    a little bit of a reading in the additional resources
    that you can actually dig it up.
    Basically, you have a lightweight M2M server and a lightweight M2M client,
    and the job of the protocol is effectively
    to be able to orchestrate the communication between these two agents.
    Now, it's a way to effectively manage IoT devices,
    and it allows devices and systems from different vendors
    to coexist in a given IoT ecosystem where there tend to be differences.
    The last part is the inference module format itself.
    Now, standardization is key for the inference module stage
    because standardization allows us to have
    these different compilers and these different inference applications
    to be able to communicate interoperability in a way
    that you can actually run across different kinds of chipsets.
    Now, the standardization is exactly the thing
    that the compiler is going to output from the compiler phase.
    So we have to agree on what the standardization layer is
    going to look like because this is ultimately what the compiler plugin is
    going to be generating.
    Now, just as a quick recap, the goal of TinyML as a Service
    is to handle all the heterogeneity.
    And to do that, we have three core components
    that we'll have to put in place--
    the compiler layer whose job is to decouple
    what the front-end and the back-end are really servicing.
    We have the orchestration protocol in order
    to standardize how we actually communicate across different devices.
    And finally, we need to have a standard way
    to interface the output from the compiler
    itself so that it can easily be pushed to the end device.
    Now, if we get these three big building blocks together,
    then we'll really be able to make TinyML as a Service happen, which
    is crucial for being able to effectively scale TinyML deployments at large.
    In summary, I'm going to quickly recap how TinyML as a Service
    is going to work.
    We're going to start off by collecting the device information.
    We're going to find a model that we actually want to deploy from the model.
    We're going to squash the model.
    In other words, we're going to perform all kinds of optimizations
    that are needed.
    Then we'll build a runtime image, and we'll push that image out
    to the endpoint device through some sort of technology like the lightweight M2M
    protocols.
    And of course, you can continue to orchestrate the execution more of this
    during the prediction serving and the continuous monitoring
    aspects of the MLOps pipe, of course.
    That said, I hope you have an idea of some of the challenges
    and the opportunities that come about in trying to scale out TinyML at large.
    Ende des Transkriptes, zum Anfang springen
    Downloads und Transskripte
    Video
    Videodatei herunterladen
    Transskripte
    Download SubRip (.srt) file
    Download Text (.txt) file

\section{Summary of TinyMLaaS
    \subsection{About This Reading}
    
    We have spent the last few videos discussing the TinyML-as-a-service (TinyMLaaS) cloud paradigm. This is an emerging technology that has the potential to overcome many barriers that exist today in terms of managing interoperability of ML software across heterogeneous microcontroller units (MCUs; the unit in a resource-constrained IoT device that manages program execution). While still in its early stages of development, it is an example of the coupling of TinyML within the cloud ecosystem and will likely become a key technology used by TinyML engineers as IoT devices become more functional and ubiquitous. In this reading, we will recap the architecture of TinyMLaaS and some of its associated challenges.
    
    \subsection{TinyMLaaS Components}
    
    There are currently three main components to a TinyMLaaS ecosystem: (1) a compiler plugin interface, (2) an orchestration protocol, and (3) an inference module. 
    
    The TinyMLaaS runs as a piece of middleware in the cloud (or edge) and interfaces with both the cloud ecosystem (the TinyMLaaS back-end) and the embedded device (the TinyMLaaS front-end). The process is initially kickstarted when the embedded device sends a request to the TinyMLaaS client running on the cloud, communicating key information about its hardware constraints such as its available RAM, flash memory, CPU characteristics, and peripherals. The TinyMLaaS software then communicates this information to the cloud, which determines the appropriate model to extract from the model registry, and also the appropriate ML compiler and optimization that will correctly configure the model and its associated runtime to work on the embedded device. The compiler plugin interface is the component that creates the correct build for our specific embedded device by combining the appropriate model, compiler, and optimizations.
    
    \GRAPHICSC{1.0}{1.0}{TinyML/4 MLOps/5_7_12_image_1.png}   
    
    Graphical depiction of model repositories being downloaded into a TinyMLaaS build as well as ML compilers being applied as plugins to a TinyMLaaS build.
    
    After the build is configured, we need to put it on our device somehow. We also need a way to actually interface with the embedded device, otherwise we would not be able to receive the device information in the first place! This requires some kind of orchestration protocol, with some kind of client software running on the embedded device. Typically, the OMA Lightweight machine-to-machine (LwM2M) protocol is used for this purpose. This protocol has a LwM2M client running on the embedded device, and an LwM2M server running on the TinyMLaaS middleware. This protocol not only allows the device to send information about itself to the middleware, but allows us to send software/firmware updates over-the-air (SOTA/FOTA). This is how we are able to communicate our model build with the embedded device.

    \GRAPHICSC{1.0}{1.0}{TinyML/4 MLOps/5_7_12_image_2.png}   
    
    Graphical depiction of a TinyMLaaS build sending inference module results to a LwM@M server which operates over-the-air using SOTA/FOTA to an LwM2M client running on a device. The Device also sends device info back to the LwM2M server.
    
    Are we done? Not yet. While we have communicated the new build to our embedded device, we have not made the ML inference application available to the device. In order to do this, we also require an inference module, which interacts with the LwM2M client to essentially update the system. We can think of it kind of like an installation manager for our device, like when we try to run a new program on our laptop or desktop computer.
    
    \GRAPHICSC{1.0}{1.0}{TinyML/4 MLOps/5_7_12_image_3.png}   
    
    Graphical depiction of the Inference module communicating over LwM2M SOTA with a LwM2M client on a device. That device LwM2M client then updates the inference engine running on the device OS.
    
    \subsection{Challenges}
    
    This technology is quite powerful, as it removes a lot of the friction in setting up a model on a new device. However, as we have discussed, there are some challenges associated with TinyMLaaS. 
    
    Perhaps the most notable challenge to TinyMLaaS is security and privacy. One of the benefits of TinyML is that all of the data is kept on-device, with no communication to a centralized system. This approach minimizes the potential for man-in-the-middle attacks or for a bad actor to hack into our system. However, by setting up a communication protocol between these devices and the cloud, we reintroduce those concerns and stimulate the need for new lightweight security protocols to keep these devices secure. In the cybersecurity world, IoT security is one of the most pressing concerns of the field given their rapid expansion in recent years. Some other challenges include balancing the trade-off between ML performance and device resources, and managing distributed execution.

\section{Model Deployment Impact on MLOps}

    0:00 / 0:00Klicken Sie auf die Pfeiltaste nach oben, um das Geschwindigkeitsmenu aufzurufen. Regeln Sie dann mit den Hoch und Runter Pfeiltasten die Geschwindigkeit. Klicken Sie Enter, um die jeweilige Geschwindigkeit festzulegen.Geschwindigkeit1.0xKlicke auf diesen Knopf und diess Video stumm oder nicht stumm zu schalten, oder drücke die HOCH oder RUNTER Knöpfe um die Lautstärke zu erhöhen oder abzusenken.Maximal Lautstärke.Videotranskript
    Beginn des Transkriptes, zum Ende springen.
    VIJAY JANAPA REDDI: Congratulations on making it through the ML deployment
    stage of MLOps.
    Model deployment by and large concerns with packaging, testing,
    and deploying a model to a serving environment for online experimentation
    and production serving.
    In this kind of context, we have talked about things like shadow deployment,
    canary deployment, blue green deployment.
    some of these concepts that actually came up earlier on in the course,
    but nonetheless, they are just as applicable down here.
    In addition, one of the key things that I wanted you to learn about
    was that you need to have scalable mature platforms like the LAMP stack
    or docker and containers in order to be able to effectively bring
    new solutions to life.
    Now when it comes to TinyML, we have more problems
    there than we have solutions.
    And that's why I talked about a potential solution called TinyML
    as a service in order to kind of give you a little bit of breathing life.
    With that said, as usual, let's try and dig down and try and understand
    what's the impact of this stage on the rest of the MLOps pipe.
    One of the key things that I want to make sure I emphasize
    is that you can see down here just how far away the ML researchers who
    are doing experimentation and prototyping in the ML development stage
    are from the downstream deployment stage.
    And this is often why you tend to have issues in production,
    where the models that the researchers build
    do not end up working well in practice at the deployment stage.
    That said, it's not a problem that's completely infeasible to tackle
    if we understand these things.
    Which is why I'm glad you're taking this course.
    If you understand these kinds of inter-dependencies and relationships,
    then you would actually be able to kind of move this knowledge up ahead, and so
    that your ML researchers that are working on experimentation
    and prototyping are in fact cognizant about downstream effects and vise
    versa.
    Now another interesting thing that I hope you noticed
    is that when we're talking about continuous integration,
    continuous delivery, and production deployment,
    and so forth, we actually talked about this
    in the context of building an operationalization pipeline, which
    is actually crucial.
    And that that's the kind of pipeline that you're actually
    in a way reusing when it comes to the deployment stage as well.
    Of course, deployment brings about its own unique characteristics,
    but by and large, these are components that we use.
    Now again, this is a great example where for instance, as much
    as we're following this waterfall kind of a chart
    in order to kind of help break things down for you,
    you start seeing that it's very much an agile development cycle.
    Things are kind of tightly coupled together.
    So keep that in mind as we go through the course,
    because there are components that will come up in other stages,
    so you should be able to make those connections.
    Finally, I want to talk about how this actually
    relates down to continuous monitoring.
    Clearly, you can see that when you're actually deploying things,
    we talked a lot about the issues that we have in the context of TinyML
    around communication limitations, and bandwidth constraints, and so forth.
    Now all of that is actually going to come back into play
    when we talk about continuous monitoring where the inherent assumption is
    that we're able to monitor how the models are actually doing in the field.
    But some of the limitations that we talked about in the context of TinyML
    as a service are actually things that will come back
    up again in the monitoring stage, because somehow we
    need to have a standardized interface protocols and mechanisms in place
    in the TinyML context to be able to route all that data back
    up so that we can actually understand relationships of how the models are
    performing with respect to the data.
    What that said, as always, the folks that
    are going to be involved in this space are
    going to be the software developers in the DevOps team
    who's really responsible for taking a lot of the packaging tasks
    and building up those containers or lean and clean pieces of code that
    can actually be pushed out by the DevOps team into the field for production.
    That said, congratulations again on wrapping up the ML deployment stage.
    Ende des Transkriptes, zum Anfang springen
    Downloads und Transskripte
    Video
    Videodatei herunterladen
    Transskripte
    Download SubRip (.srt) file
    Download Text (.txt) file        

\section{Driving Mode Detection}

Provided By \href{https://streamanalyze.com/}{Stream Analyze}

\subsection{Introduction}
A manufacturer of mobile industrial short distance transportation vehicles wanted to implement edge AI solutions for detecting if the vehicle was carrying a load. The aim was to be able to monitor vehicle utilization across a huge fleet of vehicles globally to optimize routes. The result from the analysis was also to be shared with the customers-operators of the vehicles to enable joint development of optimized intelligent transportation solutions.

Vehicles produce a large amount of data on their CAN bus and even though all vehicles had mobile connectivity, sending such a large volume of data would quickly use up all of the available data bandwidth. Therefore an ML system is needed to classify the driving mode of the vehicle on-device which can then be sent to the cloud.

\subsection{Challenge}

The development of a driving mode analysis model is a complex challenge. To provide robust classification, a random forest classifier model was needed in conjunction with a Neural DriNetwork. The neural network acted as the feature extractor which then fed into the random forest classifier. The output of the random forest classifier was then adjusted using a probabilistic rule model, which mitigated un-probable driving mode transitions inferred by the random forest classifier (e.g. a state transition from ‘no load’ to ‘dropping off a load’).

Executing this model on edge devices was a challenge given due to the complexity of a cascade of models and the tight constraints of the platform. Further, as the converted model executes in the edge devices based on high volume CAN bus data streams, the system needed to be efficient.

\subsection{MLOps Solutions}

To address the challenges above, a complete solution was developed using SA Engine, an end-to-end real-time streaming analytics platform for edge devices and microcontrollers. The platform was chosen for this particular application as it supported the broad set of machine learning models needed for this application and provided support for MLOps tools which enabled the analysts to quickly test and deploy new versions of their models onto the vehicles.

In the lab, the neural network and random forest models were implemented using Python and then converted and integrated into the SA Engine pipeline through the SA Engine model conversion tools. Once converted, the models could be executed efficiently in the resource-constrained edge devices and could be automated through the SA Engine CI/CD pipeline.

The analysts used the generic CAN bus wrapper of SA Engine to build a compact and unambiguous CAN bus data model that extracted data streams from the CAN bus to the prediction model with minimal delay. This CAN bus data model was bundled and deployed together with the prediction model.

The result of the prediction model was then streamed using SA Engine’s communication protocol over 3G/4G to an SA Engine cloud instance for storage and further processing. Selected result streams were forwarded to other cloud-based infrastructure using SA Engine’s built-in interfaces for Kafka.

Looking Forward to the Future
The volume of data created by sensors on edge devices is predicted to grow at an unprecedented pace for the next decade and beyond. It is not possible to send all this data to the cloud due to technical and commercial limitations of wireless connectivity solutions. Furthermore, in this project, a combination of several techniques and model types was needed to produce reliable results.

\subsection{Additional Resources}

\begin{itemize}
  \item \href{https://streamanalyze.com/studio-welcome/}{Stream Analyze Studio}
  \item \href{https://docs.streamanalyze.com/#/docs/usermd/statistical_analysis_guide/docs}{Statistical analysis of large datasets}
\end{itemize}

\section{Case Study Forum}

Forum Questions

\begin{itemize}
  \item How do you think the data pressures discussed in this case study will impact the needs for edge processing, and in turn the technologies and methodologies needed to efficiently develop and scale edge solutions?
  \item How do you envision the future of developing real-world analysis and AI solutions? Could the selection between different classifier methods (statistical models, Neural Networks, or other ML models) be facilitated?
  \item What model deployment challenges were faced and how were they dealt with? Were there other MLOps considerations?
\end{itemize}


% 1.8: Prediction Serving

\section{Overview of Prediction Serving}

\GRAPHICSC{1.0}{1.0}{TinyML/4 MLOps/5-8-2_Prediction_Serving.png}  

    The MLOps pipeline. Going from ML Development to ML training to Continuous Training to Model Deployment to Prediction Serving to Continuous Monitoring. Along the way artifacts are produced at each step including: Code and Configurations, Training Pipelines, Registered Models, Serving Packages, and Serving Logs. All steps tie into the overarching themes of Data and Model Management.
    
    \subsection{Objective}
    
    While much of the attention in machine learning research has been focused on the process of training models (i.e., learning), there is a distinct set of issues associated with the process of serving and updating those models. In this section, we will look at the broader machine learning life-cycle and address the difficulties associated with giving predictions to customers. To this end, we will learn the following concepts and ideas in this section:
    
    \begin{itemize}
      \item  Different serving scenarios: Batch, Offline, Streaming, and Embedded
      \item Serving architectures: Blackbox vs. Whitebox methods
      \item Benchmarking embedded inference serving scenarios
    \end{itemize}
    
    \subsection{Motivation}
    
    Once we have our production environment ready, this is where the fun really begins. Let us imagine that our production environment is a restaurant. Our environment might have many different services that it can perform, in the same way, that a chef in a restaurant can (hopefully) cook more than one dish. Depending on what our customer orders, we may have to serve up different results. The idea of a “customer” and an “order” is slightly more subtle in the machine learning sense, since often the inferences performed by machine learning models are done on-the-fly or without knowledge from the customer. An example of this could be fraud prevention in a bank. Oftentimes, when you make a transaction, a machine learning model will decide whether the behavior from your account looks “normal” or “suspicious”. If you buy a coffee from a local store every day, the model will likely know this is typical behavior of you, so has no reason to be concerned. However, if your card is suddenly used to purchase a piano in Costa Rica, the model might view this more suspiciously (unless you mentioned this trip or purchase to the bank beforehand). Many of the models that are in production today are somewhat hidden from the customer, but the models are still serving up predictions.
    
    \subsection{Prediction Serving Pipeline}
    
    Prediction serving starts with the user, whereby a request is made, knowingly or unknowingly, to the service provider. This is typically done through a REST API, which will send a JSON query to the server. The server will then, based on the information in the query, output a prediction and respond to the user via the REST API. This procedure is common for singleton requests to provide online predictions in real time.
    
    In some situations, the model may require additional information that is not provided by the user. For a model that is meant to predict the likelihood that a customer will purchase a particular product, details about the customer, as well as the product, will likely be needed. The model will then have to take the product identifier and use this to interact with a product database to obtain relevant information about the product, which it will use in its prediction. Depending on how many external features might need to be queried to make a prediction, this can make the prediction serving pipeline quite complex.
    
    Interpretability is also important in prediction serving (especially in Europe due to the “Right to Explanation”) to provide rationale for why a particular prediction was provided. For example, why the price of a quote was the price provided, or why an individual was rejected from receiving a loan or a new credit card. These rationales are often saved alongside predictions in a serving log, so that the service provider can keep track of all of the predictions given to customers that it has served. In the remainder of this section, we will delve deeper into the nuances of prediction serving and how it is relevant to TinyML.
        
\section{Prediction Serving}

    0:00 / 0:00Klicken Sie auf die Pfeiltaste nach oben, um das Geschwindigkeitsmenu aufzurufen. Regeln Sie dann mit den Hoch und Runter Pfeiltasten die Geschwindigkeit. Klicken Sie Enter, um die jeweilige Geschwindigkeit festzulegen.Geschwindigkeit1.0xKlicke auf diesen Knopf und diess Video stumm oder nicht stumm zu schalten, oder drücke die HOCH oder RUNTER Knöpfe um die Lautstärke zu erhöhen oder abzusenken.Maximal Lautstärke.Videotranskript
    Beginn des Transkriptes, zum Ende springen.
    LARA SUZUKI: Hello.
    In our last lesson, we covered model deployment process and how to make
    machine learning pipeline securely available.
    Now it's time to move to the next stage of our MLOps framework, that is,
    of prediction serving.
    This stage is all about provisioning an online inference serving.
    The prediction serving stage of the MLOps process
    is concerned about serving the model that is
    deployed in production for inference.
    It includes tasks such as accepting prediction requests, serving data,
    and to serve responses, and also to provide prediction explanations.
    We can also log metrics related to system and model performance
    and metrics related to data sets.
    The MLOps capabilities employed in this stage
    are data set and feature repository and also model serving.
    After a model is deployed, the prediction service
    accept prediction requests-- the survey data-- and serve the responses.
    Online inference is near real time for high-frequency single [INAUDIBLE]
    requests, or many batches of requests.
    It uses interfaces such as rest or gRPC.
    The streaming inference in near-real time
    carries out predictions to an event processing pipeline.
    Offline batch inference for bulk data scoring
    is usually integrated with extract, transform, and load process,
    also known as ETL.
    It also serves embedded inference as part of embedded systems or edge
    devices.
    During prediction serving, functionalities
    such as looking at feature values that are related to requests
    is also employed.
    This lookup is performed in a feature store.
    The prediction serving should also provide explanations with insights
    into the rationale behind its predictions.
    There's also been logging of inference metrics,
    including latency, errors, amongst others, and prediction requests
    for continuous monitoring.
    Ende des Transkriptes, zum Anfang springen
    Downloads und Transskripte
    Video
    Videodatei herunterladen
    Transskripte
    Download SubRip (.srt) file
    Download Text (.txt) file        
   
\section{Prediction Serving Overview}

\subsection{Prediction Serving Scenarios
    
    \GRAPHICSC{1.0}{1.0}{TinyML/4 MLOps/serving.png}
    
    Prediction serving scenarios
    
    
    
    In the upcoming few videos, we will be stepping through the various model serving scenarios. We’ll learn about what each serving scenario is, what are some use cases of this scenario, what are the metrics that matter for such a scenario, as well as what are some pros and cons of the scenario. Briefly, here is a description of each scenario:
    
    \begin{description}
        \item [Online inference] generates machine learning predictions in real-time. It's sometimes called real-time or dynamic inference. 
        \item [Batch inference] is the process of making predictions from a set of data. Batch jobs are usually generated on a regular basis (e.g. hourly, daily). 
        \item [Streaming inference] is when you are working with data that is continuously flowing through, a good example of this is time-series data that
        \item [Embedded inference] is when you shrink wrap your model into an application. The application and the model are tightly coupled together for deployments in embedded devices like smartphones.        
    \end{description}

\section{Prediction Serving Scenarios: Batch}

    0:00 / 0:00Klicken Sie auf die Pfeiltaste nach oben, um das Geschwindigkeitsmenu aufzurufen. Regeln Sie dann mit den Hoch und Runter Pfeiltasten die Geschwindigkeit. Klicken Sie Enter, um die jeweilige Geschwindigkeit festzulegen.Geschwindigkeit1.0xKlicke auf diesen Knopf und diess Video stumm oder nicht stumm zu schalten, oder drücke die HOCH oder RUNTER Knöpfe um die Lautstärke zu erhöhen oder abzusenken.Maximal Lautstärke.Videotranskript
    Beginn des Transkriptes, zum Ende springen.
    VIJAY JANAPA REDDI: A machine learning model that is deployed in the field
    can be used in many different ways.
    And that's what prediction serving is about.
    And we'll learn all about prediction serving
    in the next coming couple of videos.
    Let's get started.
    First and foremost, there are four different types
    of serving models out there.
    The first is online inference, then you have batch inference,
    you have streaming inference, and then, finally, you have embedded inference.
    For each and every single one of these, I'm going to do five things.
    First, I'm going to talk about what it is.
    I'll talk about how it works, when is it actually useful.
    What are the crucial metrics that are associated with it.
    And finally, how do you determine when to use what?
    With that said, let's get started with batch inference.
    And by the way, this is the kind of thing
    that the ML engineer is most are going to be concerned with.
    You might be wondering, now, why is this not a devops sync?
    Because it's all about serving the predictions.
    Well, as you will see, while predictions serving does
    involve using some sort of operationalized pipeline
    that the ML devops team is going to be involved in.
    By and large, in the context of actually-- deployment
    from a batch inference sort of a standpoint,
    given its simplicity and its evaluation metrics,
    you will find that an ML engineer is, sort of, the key person in here.
    With that said, let's get started.
    First and foremost, batch inference is the process of making predictions
    from a large set of data.
    Now, batch jobs are usually generated on some sort
    of periodic cadence, an hourly, or daily, or weekly, or monthly,
    that kind of cadence.
    Now, this sort of inference is usually scalable.
    It uses some kind of technologies like MapReduce, or Hadoop,
    to Spark, those kinds of scalable frameworks.
    Because it's usually infinitely scalable.
    As in, given how much ever data you have, you can run a lot of inference
    on it.
    Let's take this example, here, of photo recognition.
    So let's say you have something like iPhoto or Google Photos.
    Now, as you may know, these software systems
    can actually recognize our faces.
    And in the future, when you're looking for pictures of some people,
    or a person that you know, you can quickly
    go in and click on one type of face.
    And then automatically you will find all the relevant faces in your album.
    Now-- which is really cool.
    Now the key question is how does that, sort of, happen?
    Well, batch inference is actually involved in that process.
    Another example really is something like Netflix or YouTube.
    Where, based on what you have seen in the past,
    you can get new content on things you like or, perhaps, things you dislike.
    Now, batch inference, by design, has several advantages.
    First and foremost, it's really that it's simple.
    That is, in fact, a major advantage.
    Let's say, for example, you have a whole set of pictures
    that you want to process for face identification.
    You would simply create a massive batch of all the images
    you have in your albums and shove them over to the machine learning model.
    And the model will, over time, just, kind of, learn to make predictions.
    And say, hey, OK, well, this is Vijay.
    OK, that's Dylan.
    OK, that's Brian and so forth.
    So it's all about, sort of, offline processing.
    Now, once you have, kind of, run a batch inference,
    you can save those details in a database and serve them later.
    So for instance, with the photos example,
    you can imagine how the data is really changing or not changing.
    Think about how often your photo album really changes.
    Mine doesn't change that fast.
    The face that's in the picture is the same face today, tomorrow,
    and the day after tomorrow.
    The file doesn't really change when you take an image.
    So that's why, in batch inference, one of the cool things that you can do
    is that you can actually run the batch inference once, and save the results,
    and simply serve them later.
    It's an extremely efficient way of doing things, instead of repeatedly doing
    inference.
    In other words, imagine that when you're searching
    for a certain person in your album, then you start running the inference.
    Well, imagine, if you end up going and looking for the same person, today,
    tomorrow, or the day after, well, wouldn't you
    be wasting a lot of energy doing the same process over and over,
    running a very costly inference model?
    So that would lead to a lot of wastage, in terms of the resources,
    in terms of energy, time, and just computational budgets that you have,
    right?
    So that's why in batch inference what you try and do
    is you try and save up the results.
    You run the inference once, and you save the results, and you serve them later.
    And the beautiful thing about this, from metric, sort of,
    standpoint, is that there are real time deadlines associated
    with batch processing.
    Now, latencies, loosely put down here, in the context of batch processing,
    can be in terms of hours or even days.
    So for instance, if you open up your iPhotos
    just after you have installed it, it won't recognize
    all the faces you have in your album.
    In a way, you first have to teach it what
    the name is that's associated with the face, so that it can actually learn.
    It would effectively learn that by working offline.
    Maybe it'll start up in the middle of the night, when it knows
    the computer is not going to be used.
    Or perhaps it'll just, kind of, run whenever your computer is free.
    That's a great example of batch inference.
    What that means, though, is that, like, you're
    not getting real time predictions.
    So that's a key thing to keep in mind.
    One of the things that's nice about batch inference
    is also the fact that it's actually relatively easier
    to implement than anything that is online,
    where you're actually interacting with it in real time.
    And this is a key feature to keep in mind.
    Because the system design becomes very simplified as a result of this offline,
    sort of, processing.
    And when it's offline processing, then the metric really
    is not so much about latency, as in how long
    it took to process a particular image.
    It really, just, kind of, becomes a throughput oriented, kind of, metric.
    As in, how much volume of data were you able to effectively churn through.
    That's what is really the key metric, when it comes to batch inference.
    Now, from an implementation perspective, as I've said, it's great.
    Because it does not factor in any real time latency constraints.
    In other words, it's not like you have to finish something
    in 10 milliseconds or 20 milliseconds.
    So therefore, your computing system, that
    has to support this sort of an inference mode, is really simple.
    There's no dynamic scheduling.
    There's no dynamic resource allocation.
    And all these complex things that will come up
    as we look into other types of inference scenarios.
    But this also means that the predictions can actually be served really fast.
    Because you're not doing online inference,
    which can be slow, in terms of the latency.
    So there's actually a benefit of doing the pre-computations
    and saving them in the database and then serving them out quickly.
    Now, there are plenty of large batch processing frameworks
    out there, like MapReduce, Spark, or Hadoop.
    And you can use any one of them.
    There are plenty of other ones as well.
    And it's easy to, kind of, implement this at scale.
    So if your data ends up growing bigger, and bigger, and bigger,
    it's not a problem.
    Because it's an offline process.
    And there are plenty of, like, scalable, kind of, technologies out there.
    You can easily work with them.
    It could also be that you don't need any fancy infrastructure.
    In fact, you might not even really need some, kind of,
    crazy MapReduce or Hadoop framework.
    It could very well just be a simple little cron
    job that, kind of, wakes up and populates the database periodically,
    every night for instance.
    Perhaps, that's how iPhoto works.
    I don't know.
    I'm just kind of guessing here.
    But you get the general idea as to what I'm talking about,
    in terms of the implementation simplicity.
    Now, one of the most important things you
    can benefit from with batch inference is that you can actually
    verify the predictions offline.
    Now, this is super key.
    Because you can never verify predictions online while
    you're trying to serve an end user in real time.
    There's just no way you can do that.
    Whereas, with batch inference, what you can actually do
    is, you run your inference, then you've got all these predictions
    that are sitting inside the database.
    Now, before you actually make them accessible to the end user,
    you can actually run some sort of pre-verification pass
    on it to make sure you're confident that all the predictions are actually good.
    What this also means, my friend, is that, because it's
    such a simplified environment, this means you can actually deploy things
    out and if something is not working, it's very easy to roll back out of it.
    So that's another beautiful thing about batch inference.
    All right, there's a whole bunch of cons around it, though.
    Now, one obvious example is a cold start penalty.
    In other words, it requires all the inputs ahead of time.
    Its batch processing, right?
    So let's say a new user signs up for something like Apple Photos.
    Well, if face recognition is done in the batch processing mode,
    every night for instance, then the user is not
    going to be able to get that personalized or tailored experience
    of having their pictures all automatically
    be recognized until Apple Photos runs whenever it's scheduled to run.
    Of course, there are some ways of, kind of, getting around this.
    So for instance, what you can do is, you can start off
    by giving the user some kind of recommendations based
    on someone else's, kind of, experience.
    Now, this is, kind of, usually, hard to do, specifically
    in the context of face recognition.
    Because you actually need to have the data.
    And you need to at least teach the system
    a handful of examples, like this is DJ, that's
    Dylan, that's Brian, and so forth.
    But nonetheless, if it's some kind of general geographical,
    kind of, pictures or something, it's possible to use some, kind of,
    recommendation based system tricks in order
    to get past this cold start penalty.
    But anyway, it is one of the issues.
    All right, the next thing that you have to really worry about is the fact that
    the predictions cannot leverage any real time characteristics.
    So for example, if I'm interacting with my application and maybe,
    perhaps, given a certain context or something,
    it could have done better predictions.
    But because this is an offline process, and there the predictions are just
    being served through the database in real time,
    you're not going to be able to dynamically adapt your predictions.
    Now that's a bit of a caveat.
    But that's something that, just, one has to, kind of, make a trade off on.
    Now, finally, because it's batch mode, you'll
    not be running the model all the time.
    This means you can only get your model predictions periodically.
    So once in a while, it's whatever it is.
    Maybe it's a nightly run, or maybe it's a weekly run,
    or maybe it's a monthly run, right?
    And again, this may be OK.
    But it really depends on the scenario that you're talking about.
    So for example, if you consider Netflix, does it really
    matter if, like, you get your predictions every hour?
    Maybe because you're not watching TV that much.
    Or it's like, as long as you get your predictions, perhaps,
    like, by the end of the day or something like that, things get updated.
    Well, then that's great.
    It really depends on the scenario as to when that, sort of,
    real time updates to your prediction database are needed.
    Now, finally, keep in mind that each scenario that we'll be stepping through
    has a critical metric associated with it.
    And that's really what I want to, kind of, summarize each of these scenarios
    what, right?
    So I've already mentioned this, but in the case of batch inference,
    throughput is the most important, sort of, a metric.
    It's an offline, kind of, a metric.
    But that is what you really care about.
    Because you want to be able to, just, kind of,
    crunch through a huge volume of data as quickly as possible.
    And you think about this, since the metric is actually going
    to influence how we build the systems.
    Now this is true whether it's a tiny mount system or not.
    Now, when we get into serving architectures,
    I'll, kind, of open up this, sort of, box a little bit further.
    But let's, first and foremost, step through all
    the other pending scenarios.
    Ende des Transkriptes, zum Anfang springen
    Downloads und Transskripte
    Video
    Videodatei herunterladen
    Transskripte
    Download SubRip (.srt) file
    Download Text (.txt) file

\section{Prediction Serving Scenarios: Online}

    0:00 / 0:00Klicken Sie auf die Pfeiltaste nach oben, um das Geschwindigkeitsmenu aufzurufen. Regeln Sie dann mit den Hoch und Runter Pfeiltasten die Geschwindigkeit. Klicken Sie Enter, um die jeweilige Geschwindigkeit festzulegen.Geschwindigkeit1.0xKlicke auf diesen Knopf und diess Video stumm oder nicht stumm zu schalten, oder drücke die HOCH oder RUNTER Knöpfe um die Lautstärke zu erhöhen oder abzusenken.Maximal Lautstärke.Videotranskript
    Beginn des Transkriptes, zum Ende springen.
    VIJAY JANAPA REDDI: In this video, we're moving on to another prediction
    serving scenario, which is the online mode.
    For online inference yet again, we'll be talking about what it is, how it works,
    when it's useful, and what are the key metrics.
    Now in this situation, you can see that I've reintroduced the DevOps team.
    And there's a valid reason for this.
    As you will see, the online scenario is far more
    complicated to actually implement, and it often
    requires much more infrastructure than the simple batch processing mode
    that we talked about previously.
    And this is exactly where your DevOps team kicks into action.
    That is in addition to the ML engineer, of course,
    who is always going to be handy when it comes to actual prediction server.
    All right, what is online inference?
    Online inference refers to the act of creating machine-learning predictions
    in real time in response to a request.
    Sometimes this is also known as predict on-demand.
    That's often referred to as real-time inference or dynamic inference as well.
    Now predictions that are generated through online inference
    can be generated by the end users at any time throughout the day any time.
    Now, typically these predictions are made during runtime based
    on some sort of observation of a data.
    Like on the right, for instance, I'm trying
    to demonstrate how requests are arriving at some sort of random rate
    at the inference engine on the server side.
    And you can see that it's some sort of a random arrival rate, which
    I'm trying to illustrate through this Poisson distribution kind of thing,
    which is usually how these things are modeled, if you're
    for instance just trying to do some initial testing
    phases of online deployments.
    But that said, the general point is that there's
    some sort of random arrival rate.
    And this is going to be crucial, because this is something that separates it out
    from the batch processing.
    Now remember, in the batch processing mode,
    the inference is just done as a massive bulk on everything together.
    Versus here you're seeing that each of those squares
    is going to be coming at some sort of interval,
    and you have to deal with that.
    Now online influence definitely broadens the scope
    of machine learning applications.
    Now instead of waiting for hours or days for having predictions,
    you would just be generating predictions on the fly,
    because the model is actually going to be triggered when the request actually
    comes in.
    Which is very different from the batch processing
    mode, where we just kind of do everything ahead of time
    and store the results in our database.
    Now just to kind of give you some examples
    of where online inference is used, well it's actually used everywhere.
    If you've ever used Uber or Lyft or any kind of food app or anything,
    you'll find that they all are relying on online inference.
    So for example, when you use an Uber app,
    you want to know when your food is actually
    going to get delivered to you, which requires a number of features
    to be combined in real time to generate a specific forecast about when
    that food is going to actually show up at your door.
    Now predictive maintenance is a more closely connected example for us
    in the TinyML context.
    Where for instance, you want to anticipate when a certain machine
    part might actually break.
    Whether it is the next couple of minutes or the next month
    or whatever that might be, right?
    Given some sort of sensor data that you're actually
    going to be looking at in real time.
    Now in general, online inference is more complex than the batch inference mode,
    primarily because you have a lot of additional goo
    that you have to put together in order to support those online requests that
    are actually coming in.
    That's one.
    Moreover, you also have to make sure that you
    able to meet certain critical latency constraints in the online phase.
    And that's what actually makes the online deployment much harder.
    So, for example, imagine if you're building
    a system that has a quote unquote "latency" of 24 hours.
    And I give you a single image that you have to label, something like that.
    If I have 24 hours, does it really matter
    how fast my entire infrastructure runs?
    Probably not, right.
    But think about the same thing where you have
    to do an inference in say something like 100 milliseconds, which
    is sort of like the key threshold for many online services,
    because online demands real-time interactive behavior for end users.
    Now during those 100 milliseconds, the system must retrieve the data,
    perform the ML model inference in real time,
    validate the output that comes out of the model,
    especially before it's presented to the end users,
    and then return the results through the network back to the end user.
    Which is kapow, mind blowing that you've got to get all this stuff done,
    all within 100 milliseconds.
    Doesn't matter where you are, right.
    So that's one of the reasons that you will
    find that it's not just the ML engineer who's involved in this process.
    The DevOps team also has to get involved.
    And that's why I was saying that you have to kind of understand
    the scenario in order to understand how to actually architect
    this sort of a system.
    I'll talk more about the details of the architectural design
    in a separate video that's coming up later on,
    but hopefully you're starting to get an idea
    of how this is different from batch inference
    and why they might be a lot more involved in here.
    And the key thing is that the entire end-to-end pipeline, right,
    I'm talking about more than just the machine learning model
    has to run within those 100 milliseconds.
    Which then means automatically that the engineers have to understand end-to-end
    pipeline, which includes the receiving, processing of all the features,
    deciding what kind of algorithm is going to be run--
    is it a big model or is it a small little model?
    Because it's going to affect your latencies.
    Then you've got to do some kind of post-processing more often than not as
    some kind of post-processing that you do that comes out
    of like the neural networks.
    Output goes into some kind of post-processing.
    And then of course, all of that has to get wrapped around web technologies,
    like REST APIs or something in order to be able to communicate.
    Now those things can actually degrade the performance,
    so a very important point here, my friend,
    is that when you're thinking of online inference,
    it's not just about the machine learning model.
    It's more than just ML.
    It's a lot of the goo that kind of holds everything together.
    So your engineer is going to have to be very
    cognizant about this kind of a thing.
    Now as I've alluded to this, in terms of metric,
    latency is the ultimate critical piece of this online inference nugget.
    I take, for example, the search engine, which
    uses a lot of machine learning technologies behind it, right.
    Which is trying to find something like a needle in a haystack
    based on your search query.
    Now I would think that you would get frustrated
    if you had to wait a whole minute just to get your results after you
    typed something into the Google box.
    But then again, Google is able to search through millions of results
    and very quickly get back to you, often within a fraction of a second, right.
    In milliseconds.
    It's mind blowing.
    Somehow, given that there's a latency constraint, I'll say 100 milliseconds--
    I'm just using 100 milliseconds as an example, it could be 200 milliseconds,
    250 milliseconds, just depends on the particular use case.
    But that is still a very tight threshold.
    Given that latency bound, you have to then handle
    some level of user requests, right.
    And that's where the metric ends up becoming something
    like a latency bandwidth throughput under some sort of like SLA guarantees,
    where you have to meet some number of queries per second, where
    the latency of an individual request does not exceed something like 100
    milliseconds or 200 milliseconds.
    You might get the perfect search result answer, but if it takes you one minute,
    no one cares about that result. So you got to kind of keep that in mind.
    All right, what are some of the pros and cons around here?
    One of the key benefits is that you can basically
    make the predictions on the fly in real time, which
    means that you can handle any kind of new contextual data
    that's coming at you.
    So you can actually also take in online features,
    which is really, really powerful.
    This is a great design for long-tail requests.
    In other words, if there are some requests to take a very long time,
    it's not that everybody's going to get penalized.
    Because you're handling things in real time,
    it's like you can actually serve a good number of requests with possibly 99.9%
    service level agreements being met.
    If you provide this service, clients will also
    profit from just paying for what they use.
    In fact, a lot of machine learning services today
    that are deployed in fact run this way for online inference
    because it benefits both the customer and the service provider.
    Now benefits a customer, because he or she is only paying for what they use.
    And it benefits the service provider, because it allows them
    to provision resources appropriately behind the scenes,
    so they don't have to overprovision resources just
    for you and me, who perhaps are not doing
    that many inferences continuously.
    But maybe there's another user who's actually
    continuously making a lot of requests.
    In that case, they can actually move more of the computing resources
    to support that particular request.
    But there are definitely cons.
    Now online inference typically has many more obstacles,
    as compared to batch inference, because as I have mentioned computationally
    the system is much more complicated because you
    have to meet these latency bounds, which makes it really, really hard.
    Throughput is almost typically easier to be
    met than actually a latency constraint.
    So in other words batch inference would actually
    be a lot easier to realize in practice than online inference.
    And this is again because simply, because you have
    to do so much more work, in real time.
    Extracting features, pre-processing, post-processing,
    validating the output, and all that for instance within 100 milliseconds.
    So now due to the latency limits, your MLOps engineer in your DevOps team
    has to be extremely careful to comprehend
    the end-to-end sort of a situation, given that its latency bandwidth
    throughput that matters.
    And as I mentioned, most of the time, it's beyond the model.
    There are other sources of overhead that you have to consider.
    Things like storage, network, like web service APIs and so forth.
    All of those start coming into play.
    And last, but not least, it may actually the model itself might actually end up
    limiting the capabilities that you have in terms of what you can actually
    do online, because if the model is too complex,
    then your latency is going to increase.
    So you get into this bit of a pickle situation.
    And then you got monitoring.
    This is something new that I have actually
    not touched on in the context of batch processing at all.
    Monitoring becomes quite key and important
    in the context of online inference.
    The reason is because of drift.
    Now, stop for a moment and try and recall what drift is.
    I introduced this concept in the continuous training stage,
    where I talked about model drift, I talked about concept drift,
    and then I talked about data skews and schemas changing.
    So some of those things, I'll again be coming back to
    in the continuous monitoring stage, which is actually
    coming up next after this MLOps stage.
    But because you're serving things online,
    you've got to make sure that the distribution is the input data,
    and the predictions that you're generating actually
    match from what you did in the training phase.
    And you have to kind of keep checking these distributions,
    because these distributions are different,
    then you need to re-trigger the model training.
    Because your new data that's coming into the system in the real world
    is obviously different than what your training data looks like.
    So your model is going to be mis-predicting.
    But you might actually not know this, unless you're continuously monitoring.
    And it's extremely key for online resources or online scenarios
    to actually be doing monitoring.
    Now here you can see, yet again, that the online inference scenario
    is different from the batch inference.
    And what you really care about here is the queries per second
    or some sort of latency bounded throughput.
    And that's the key metric that's actually
    going to drive much of your system infrastructure design.
    In the next video, I'll talk to you about streaming inference.
    Ende des Transkriptes, zum Anfang springen
    Downloads und Transskripte
    Video
    Videodatei herunterladen
    Transskripte
    Download SubRip (.srt) file
    Download Text (.txt) file

\section{Prediction Serving Scenarios: Streaming}

    0:00 / 0:00Klicken Sie auf die Pfeiltaste nach oben, um das Geschwindigkeitsmenu aufzurufen. Regeln Sie dann mit den Hoch und Runter Pfeiltasten die Geschwindigkeit. Klicken Sie Enter, um die jeweilige Geschwindigkeit festzulegen.Geschwindigkeit1.0xKlicke auf diesen Knopf und diess Video stumm oder nicht stumm zu schalten, oder drücke die HOCH oder RUNTER Knöpfe um die Lautstärke zu erhöhen oder abzusenken.Maximal Lautstärke.Videotranskript
    Beginn des Transkriptes, zum Ende springen.
    VIJAY JANAPA REDDI: We're into the third of our four scenarios.
    In this video, we're going to learn about streaming inference.
    It's one that is less commonly known than online and batch inference,
    but is still quite important and relevant,
    as I'll show you with some example use cases.
    And, yet again, I'm going to talk about what it is, how it works,
    when it's useful, and what are the key metrics.
    Streaming is also another type of on-demand scenario, which
    is akin basically to the online inference scenario
    that we saw earlier, except there are some crucial differences.
    Unlike the online inference scenario, streaming inference actually
    involves numerous different observations arriving
    in some sort of grouped or batched kind of form.
    And that's what I'm trying to show you here on the right side.
    As you can see, the square boxes are forming some sort of groups.
    And those groups are then going into the inference engine.
    Now, streaming inference is typically encountered
    in some sort of event processing pipeline where you're
    looking at a high volume of data, but the data
    has to have some sort of synchronous timestamp behavior or something
    to that effect.
    Now, streaming inference is also a little
    overused in terms of the terminology.
    Another meaning of streaming inference is often
    looking at data that is just continuously going
    by like a sensor data stream.
    So, in this particular video, I'm mostly focusing on the former
    where each query that comes into the inference engine
    has some n number of samples or observations
    that arrive at some sort of frequency.
    And the inference must be completed with some latency bound
    for that particular cohort, OK?
    I think the best example I can give you of this is really in an autonomous car.
    It's a classic illustration of streaming inference.
    Let's say you have a Tesla car.
    A Tesla primarily relies on camera feeds coming into it in order
    to enable its self-driving mode.
    Now, there's several different cameras that are all
    pointing in different directions.
    And all of that data must be processed in real time, much like online,
    and it has to happen concurrently.
    I hope you understand what I'm trying to say here.
    Let's say, for instance, that your car has to make a decision.
    It cannot simply look at the data that's coming in from one camera and make some
    decision and then move on to the next camera and then make another decision
    based on some data it sees there.
    It first has to build a global perspective of everything around it.
    And then, with that snapshot of the data from all the different feeds,
    it has to do some sort of inference decision-making.
    And that is the key difference between online inference and streaming
    inference.
    Another good example of dealing with numerous feeds coming at the same time
    into the inference engine is like when you're
    talking about smart cities, which often use CCTVs, for instance, where there's
    a lot of information to process.
    And, usually, you want all of that data to be handled
    in some sort of synchronous manner.
    Allow me to push this a little bit further because I actually
    find this quite exciting.
    Now, machine learning is becoming extremely important in application
    areas like scientific applications.
    Now, scientific experiments generate an insane amount of data.
    And when is machine learning good?
    Well, when there's lots of data.
    And scientific application domains produce an insane amount of data.
    So, on the right, you're seeing a photo of CERN's LHC.
    Now, this is the world's largest high-energy particle collider
    built by a small universe of 10,000 different people.
    But this gadget generates an insane amount of data.
    Let me show you how much data I'm talking about here.
    The chart on the right is actually from my research,
    and it depicts the amount of streaming data
    that is generated by some of these "FastML Science,"
    quote unquote, workloads.
    Now, FastML Science workloads are example workloads
    that you find in LHC that include things like sensor data compression, beam
    control, jet classification, basically, the area that's shown in the gray
    to the upper left of the figure.
    And, next to that gray box, I'm showing you
    labels called MLPerf Tiny and MLPerf Mobile.
    Now, these are more representative consumer-facing applications
    in the context of TinyML or the context of mobile smartphone use cases.
    And the data points that I'm showing you here
    are based on an industry standard benchmark
    that you'll learn about later in a couple of videos.
    And the cross hairs represent some approximate volume
    of data that they produce and the latencies that kind of correspond
    to that sort of application use case scenario.
    As you can see, there's a huge difference
    between what's on the left side with Fast Science or FastML
    for science workloads versus where these other cross hairs actually are.
    And the reason I'm showing you this is, with that much data
    being produced every single second in things
    like beam control or jet classification and so forth, you can't possibly
    store all that data to do something intelligent with it afterwards.
    You have to somehow do online machine-learning-driven intelligence
    to figure out what parts of the data that's actually coming out
    from this behemoth scientific instrument is actually interesting.
    And so machine learning there is really fascinating.
    And streaming inference plays an extremely critical role.
    All right, so I got a little carried away with that streaming inference
    example down there, but let's keep moving forward.
    Now, as you can see, as a result of having
    to deal with groups of data that's actually coming in
    from multiple sources, streaming inference
    is a bit of a mishmash of different kinds of components, both from online
    and a little bit from the batch settings.
    Now, as indicated on the right, let's say you have some sort of input stream
    that's coming in from some sort of CCTVs.
    Now, the data is arriving in some sort of synchronized methods.
    Or, in some sense, there's some layer that actually synchronizes it
    before the data comes to the inference engine.
    Now, you can imagine, for instance, you have two cameras,
    and both are transmitting data.
    And, as shown on the right, they have to somehow get buffered together.
    And that's that intermediate layer that I'm talking about.
    Now, as a systems engineer, for instance,
    you have the opportunity to kind of conceive or come up
    with smart and creative ways to figure out
    how you might be able to partition those groups of data into different groups
    so that you can actually concurrently feed them
    into parallel instances of your machine learning model, rather than having one
    big, behemoth sort of a machine learning model, which will just basically choke,
    given the amount of volume of data that it actually has to process.
    So here you're kind of looking at something quite interesting
    in streaming where you're trying to actually scale while you still
    have some sort of latency bound on you.
    So then the way you're scaling is by having replications of these machine
    learning models, possibly sitting in some sort of Docker containers,
    and you're intelligently kind of batching and doing
    sort of optimizations ahead of time in order to feed the engine.
    Obviously, at this point, it should become rather obvious to you
    that it's a rather complicated system to be able to engineer.
    And this is why I was telling you earlier,
    from a team-building perspective, we would need to have our ML engineers
    and we would need to have the DevOps team kind of working together
    with us on this.
    Now, from a metrics standpoint, you're mostly concerned
    with both latency and throughput.
    You're effectively aiming for some sort of latency threshold.
    Like, for example, when you have the Tesla car,
    it doesn't do you any good if you look at all the data that's coming in,
    but, in order to crunch through all of the data, it takes you,
    let's say, again one minute.
    Well, by then, you might have hit somebody.
    Or, well, really, the car would have hit somebody.
    Now, that's going to kill your use case in a jiffy.
    So you have some sort of latency bound.
    So maybe, yet again, maybe it's 100 milliseconds,
    or maybe it's 50 milliseconds, or maybe it's
    35 milliseconds because, in autonomous cars,
    for instance, you have to process all that data really quickly.
    And then you have to physically do something with the vehicle
    to either move the vehicle or whatnot.
    And that takes much more time.
    So, in a way, your computer has to be much faster
    than the cyber physical system itself.
    So that's an example where your latency bound really becomes critical,
    but, from a throughput, it's not really just raw throughput.
    The reason I say that is because you don't
    want to just chew through a lot of data that's coming in
    from just a single front camera, for instance, real quickly and claim,
    wow, look at my throughput.
    It's amazing.
    That's not going to work.
    You have to effectively get your throughput really
    from processing all of the cameras' inputs concurrently.
    Like that requires the ability to process n separate camera
    sources all concurrently.
    And that is the kind of throughput I'm talking about is like how many
    of those cameras.
    For instance, can you handle eight cameras?
    Can you handle 12?
    Can you handle 15, 30?
    How many can you support under a fixed latency constraint?
    That is super challenging to achieve.
    And that's what makes streaming really an interesting application scenario.
    All right, now, let's quickly look at some of the pros
    and cons of the streaming inference.
    One wonderful feature is that you can quickly scale it out
    by doing things like that grouping example
    that I showed you earlier with the CCTV kind of scenario.
    What you can also do is like streaming can be actually difficult to implement.
    However, it's also a really interesting application
    because you can look at multiple feeds, as in the Tesla car,
    or you can look at a single feed, like in the case of the LHC collider
    that I was talking about where, for example,
    you might just be interested in one single sensor input that's
    coming through, and it's streaming through really, really fast,
    and you have to handle that.
    I would say, it's a different form of streaming.
    It sounds like the same thing, but it's not.
    It's actually like online streaming kind of thing.
    Anyway, I'm trying to draw this distinction because I find
    that people often get confused with it.
    All right, from a cons sort of a perspective,
    hopefully, you can see why this is much more
    difficult because you have to have a lot of resources plugged into your backend.
    And that's going to make things really challenging for you.
    Now, furthermore, obviously, from a computational sort of a standpoint,
    you're going to need many more resources that you would actually have
    to support for streaming inference.
    And you have to architect things carefully both
    from a latency standpoint and a throughput standpoint.
    And, finally, just to kind of wrap this up,
    again, if we go back up to the metrics kind of viewpoint
    from a scenario standpoint, here, in streaming scenario,
    we're talking about latency-bounded throughput.
    And that latency-bounded throughput here is more about,
    for instance, how many different streams of inputs
    can you concurrently support in the context of a given latency threshold.
    All right, great.
    Next, we'll be moving on to embedded inference.
    Ende des Transkriptes, zum Anfang springen
    Downloads und Transskripte
    Video
    Videodatei herunterladen
    Transskripte
    Download SubRip (.srt) file
    Download Text (.txt) file

\section{Prediction Serving Scenarios: Embedded}

    0:00 / 0:00Klicken Sie auf die Pfeiltaste nach oben, um das Geschwindigkeitsmenu aufzurufen. Regeln Sie dann mit den Hoch und Runter Pfeiltasten die Geschwindigkeit. Klicken Sie Enter, um die jeweilige Geschwindigkeit festzulegen.Geschwindigkeit1.0xKlicke auf diesen Knopf und diess Video stumm oder nicht stumm zu schalten, oder drücke die HOCH oder RUNTER Knöpfe um die Lautstärke zu erhöhen oder abzusenken.Maximal Lautstärke.Videotranskript
    Beginn des Transkriptes, zum Ende springen.
    VIJAY JANAPA REDDI: [? At last, ?] we're at the embedded inference serving
    scenario stage.
    Yet, again I'm going to talk about what it is, how it works, when it's useful,
    and what are the metrics.
    All right, make a mental note on this particular slide.
    The names of those involved have actually changed.
    As you can see, I'm involving both the software engineer and the ML engineer
    in this.
    Previously, we used to have the DevOps team for streaming and online,
    but I'd like to focus more on the role that the software engineer and the ML
    engineer are going to play in this particular context.
    This is not to say that the DevOps team actually
    doesn't have anything to do in the context of embedded inference.
    That's actually not true.
    In fact, sometimes they have to deal with much more, because
    of the extreme heterogeneity that embedded devices have.
    But nonetheless, I want to focus on the ML engineer and the software engineer,
    because, from an implementation standpoint, it will be clearer to you
    what their role is.
    The key thing is that there--
    there's really no separate infrastructure required
    for this approach, as you'll see.
    The same application architecture, the release, the deployment procedure,
    basically can be reused for actual working with the model.
    And because we're using the same sort of technology stack as the model,
    the same engineers who are working on the main application
    can easily make modifications via the same stack.
    And this is kind of what leads to me wanting to focus more on the software
    engineer and the ML engineer as being the key people that
    are responsible for embedded inference.
    Now, embedded inference, like the streaming and online inference
    sort of a situation, is also an on demand scenario.
    Now, the emphasis is primarily on single requests that are coming in.
    The inference engine receives a single sample per request,
    then you wait for the inference to actually happen,
    and then you report the latency.
    So what you're really interested in is reporting that latency per inference.
    Now, let me contrast this real quickly with the online mode, which
    is kind of typical in this server sort of a situation.
    You might be getting confused here and saying, well,
    didn't you say the server is also getting
    these requests in some sort of arbitrary arrival rates,
    and then processing them?
    Yes, it's true.
    But when you're looking at a server, you're looking at volume of data
    that you can actually handle.
    So there are a lot of requests that are coming in,
    and you have to satisfy all of them under some sort of SLA agreement
    that you have.
    Versus in the situation here with embedded inference,
    you're interested in just the raw latency of a single request,
    versus in the server, it's more that you have a latency threshold,
    and then you're trying to see how many requests can you
    handle within that latency threshold.
    So hopefully that distinction is clear here.
    Now, let's take a look at where this embedded inference serving situation
    shows up.
    Well, your smartphone is a perfect example of this.
    It's a supercharged embedded system.
    So you take a picture, then you wait for some sort of sorcery to happen to it,
    and then out comes this beautiful picture that
    is way better than what the camera actually took in,
    and then you go back to whatever you're doing.
    In that kind of a situation, what do you really care about?
    What you really care about is the latency
    to process that individual frame that kind of came
    in from the camera to generating that picture on your display.
    That's what you really care about, and that's
    why it becomes a really latency-centric sort of an application.
    Any other TinyML application that we're talking about,
    things like [? keyboard ?] spotting, TinyML, and Ring doorbell
    to detect a person, and so forth are also basically
    under the same umbrella as here.
    How does it all fit together?
    Well, in the context of that embedded inference serving scenario,
    the application and the model are kind of package together.
    You can probably guess why this is the case, right?
    In the context of a smartphone app, for instance, rather than taking a picture
    and then sending it all the way up to the cloud,
    and then getting the result back, it's much better from a latency,
    from an energy efficiency, from a bandwidth, from a privacy argument,
    from a reliability standpoint.
    All of these things that we've talked about way ahead in the course--
    due to all of those reasons, it's way better to just have embedded inference.
    It kind of just makes more sense, because everything just gets better.
    Now, this scenario-- in one way, you could
    argue that it is a lot simpler than the streaming and the online ones,
    because you're there, you have to handle rather complex number of requests
    all coming in.
    So it's more than just a model, you have all this infrastructure,
    which ends up impacting how well that particular model serving is actually
    going to work.
    Here, you're really just kind of focused on raw latency.
    Now of course, this comes with its own set of problems.
    But at large, I would argue that this is a much more easier problem to solve.
    Now, let's quickly consider some of the advantages and disadvantages
    of embedded inference.
    The primary advantage is that you can do things
    on demand, which is fantastic, right?
    So remember, yet again, like anything in on demand,
    you might actually be able to use other online features that
    are supporting the inference request that you actually want to perform.
    So, for example--
    I'm just making this up.
    For example, you take a picture, and you want to quickly do some inference.
    Well, maybe based on the location, you might realize, oh,
    that's the Taj Mahal.
    I'm just using a crazy example here, but you kind of
    get the idea that if I'm in India, if I'm close to the Taj Mahal, then
    it's likely that that's what is actually interesting.
    Or maybe there's some of the additional features
    about that you might actually feed into the machine learning model
    in order to make it produce a much more interesting result.
    For example, if it's generating a caption of,
    what is this that I'm looking at?
    It might be able to generate a very rich sort of contextual response for you.
    That was actually a good one.
    I came up with that on the fly.
    The pros are also that you get all the benefits of lower [INAUDIBLE]
    bandwidth improvements, latency is much better.
    Energy efficiency, reliability, and privacy, which we already talked about.
    Now however, like anything else, you must weigh these benefits
    against all the drawbacks.
    Now, the disadvantages I've stated here are largely
    the same as what you would see in online inference,
    because it's somewhat comparable, in that you're
    making some bit of assumptions.
    Now, because you're running on an embedded device though,
    you're naturally going to be very computationally intensive, because you
    have very limited resources.
    And then you have latency sensitivity, which is yet another constraint
    that you have.
    Of course, monitoring is also going to be important, and very similar
    to the online case, because again, you had to worry about drift and so forth.
    And we'll talk more about online monitoring
    and the continuous monitoring stage next.
    Now, last but not least, because the model
    is being pushed into embedded devices, it's
    difficult to control the scalability or the flexibility of your implementation.
    Now, this is true because we're specifically, tightly
    coupling the model deployment and the model usage together.
    And this can create some serious nightmares down the road
    with maintenance overheads and technical depth that starts building up over time
    as you start scaling out.
    These are some of the real world challenges that Android faces,
    or Facebook, or now, Meta.
    [? That ?] these applications that are running
    on billions of devices, the developers really face these kinds of challenges.
    But it is the nature of the beast when we're dealing with embedded systems.
    So it's kind of just part of the job, and we just
    need to figure out how we deal with that efficiently.
    All right, now it's time to wrap this up.
    We have finally got embedded inference right at the bottom,
    and as you can take a close look down here at the slide,
    because it'll really help you understand what
    are some of the critical differences.
    Go all the way from the left of the slide, all the way to the right.
    Pause the video if you need to, and I want you to really internalize
    the key differences between these four different streaming scenarios
    that we actually have.
    And that said, obviously, in the context of embedded inference,
    latency is the key metric.
    Finally, they've covered all the different ways
    in which models can actually be served.
    The key thing you have got to realize here
    is that you might have an image classification model,
    or you might have an object detection model,
    or whatever-- pick your favorite model.
    But the meta point that you should take away from this serving scenario context
    is that machine learning models might actually be deployed and used
    in very different ways.
    And that is extremely important to keep in mind,
    because it's not only about the model, it's about the use case.
    And that use case drives how you actually
    engineer the serving architecture around that model.
    And that is why I'm actually talking to you about all this crazy stuff.
    All righty.
    With that said, we'll pick back up and [? do ?] like,
    what are some basic architectures or blueprints
    that we would actually put together in order
    to realize some of these different serving scenarios
    that we've talked about here.
    I'll see you there.
    Ende des Transkriptes, zum Anfang springen
    Downloads und Transskripte
    Video
    Videodatei herunterladen
    Transskripte
    Download SubRip (.srt) file
    Download Text (.txt) file
  
 \section{Prediction Serving Architectures}
 
     0:00 / 0:00Klicken Sie auf die Pfeiltaste nach oben, um das Geschwindigkeitsmenu aufzurufen. Regeln Sie dann mit den Hoch und Runter Pfeiltasten die Geschwindigkeit. Klicken Sie Enter, um die jeweilige Geschwindigkeit festzulegen.Geschwindigkeit1.0xKlicke auf diesen Knopf und diess Video stumm oder nicht stumm zu schalten, oder drücke die HOCH oder RUNTER Knöpfe um die Lautstärke zu erhöhen oder abzusenken.Maximal Lautstärke.Videotranskript
     Beginn des Transkriptes, zum Ende springen.
     VIJAY JANAPA REDDI: In this video, I'm going
     to talk to you about prediction serving architectures.
     Now, a prediction serving architecture is
     something that is a blueprint or a design
     that you would be applying in order to realize the prediction serving
     scenarios that we had talked about previously.
     Now, as a quick reminder, we have four different types
     of prediction serving scenarios.
     We have the batch one, we have the online, we have streaming,
     and we have the embedded.
     Now, in general, I'd like to say that there are
     three basic types of serving styles.
     Now, each has their own advantage and disadvantage,
     so I'm going to walk through each one of these.
     Let's start with the first one.
     This is all about deploying a serving instance as a container,
     or using those Docker concepts that I talked about previously.
     Sometimes this is commonly referred to as a black box-based design,
     but I'm going to refer to this as a container-based architecture.
     Now, containers are yet again amazingly powerful,
     because they allow us to shrink wrap the model and all the relevant code
     into a nice little Docker container instance that
     makes it easy for deployment.
     The model itself is kind of containerized and running
     on a separate system--
     or could be on the same system, but it's running in a dedicated container.
     The application that's actually wanting to interact with the model--
     for example, if I want to say, hey, is this picture a cat
     or is this picture of a dog--
     well, in that case, the application is going
     to interact not with the model directly, but rather, is actually
     going to interact with the serving system
     by invoking some sort of REST API that actually provides an interface.
     And the serving system actually makes the connection
     between the model and the request that's coming through using something
     like Remote Procedure Calls, or RPCs.
     This design has particular advantages, because it really
     decouples the system from how the model is actually deployed
     versus what the end user actually needs from a front-end interface.
     It makes methods and policies broadly easier to implement,
     because you might say, OK, I want the end
     user to be able to connect to a server.
     And how or what ends up happening once they connect to the server
     is actually a decoupled thing, because that does not
     involve the developer, who's dealing with that front-end REST API
     interfaces.
     And so in a very nice way, it allows you to distinguish
     what you do from how you go about actually doing it.
     And that's a remarkably powerful concept that's been used many times
     over building systems.
     And it allows for excellent scalability, and it
     allows you to have nice tolerance capabilities,
     given demand requests that come in.
     And one of the reasons you get that scalability inherently
     is because you're actually shrink wrapping your model inside a container.
     So if you want more instances, all you have to do is just simply spawn up
     more number of containers that you need.
     And as the request load increases, what you would be able to do
     is then be able to route traffic.
     Of course, this requires additional [? goo ?] inside
     in order to orchestrate all of this, but nonetheless, you
     can imagine how this might actually be a scalable way of doing things.
     However, there's a downside, which is that it's not really suitable if you're
     talking about ultra low latency kind of applications,
     because you have all this communication overhead,
     and you have these nice clean separations or abstractions
     between how you're interacting with a server, a front end,
     and not directly with the model.
     And it also has a lot of resource overheads
     that are not actually captured visually down here.
     The containers for example, are great.
     They provide resource isolation, and that
     allows you to quickly spin up more number of instances
     and have multitenancy if you need on a single system.
     Tenant is a buzzword to say that I can actually
     support multiple kinds of models that are all nicely isolated
     inside [INAUDIBLE] containers that are able to run collaboratively
     on a single host.
     Nevertheless, when you have a container, you
     have to realize that every container is going to have its own machine learning
     runtime sitting inside it.
     And it's going to have its own set of processes, and all the libraries,
     and all the things that go together to support that model.
     What that means is that now you can't share any resources across multiple
     of these containers, because every one of them
     is going to need its own memory footprint,
     is going to need its own compute cycles, and so forth.
     And so you cannot amortize that in any way through any kind of sharing
     because of that isolation guarantee that the containers [INAUDIBLE]..
     Now, furthermore, the RPC layer and the REST
     API-- they all had network communication overheads, which
     can be especially significant when you're talking about millisecond level
     prediction latencies.
     These are some of the cons.
     You get this beautiful architecture to easily scale and so forth,
     but because of the abstractions you create,
     you get into the latency-related issues.
     Let's talk about a different approach, which
     should look somewhat similar, because we just
     touched on this in the context of the embedded inference.
     Now, this is the place where, for instance, you import the models
     directly into the applications.
     Many of our applications today use some sort of ML model.
     If you're using your mail application or are you using your photos application,
     they all have machine learning models.
     And often, in that kind of a scenario, the ML model is already packaged in,
     and that is a classic example of [INAUDIBLE]..
     Now, this architecture can be used ideally
     both at the cloud side as well as more on the edge side,
     towards our interest of TinyML [INAUDIBLE]..
     Because it has to support a broad range of different use cases,
     be it on the cloud or on the edge side, this machine
     learning model that is deeply embedded inside
     must adhere to several types of key constraints
     that your DevOps, and your ML engineer, and your software developer
     all need to be aware about.
     In other words, you have to worry about the amount of memory it takes,
     the amount of computational cycles, the number of processors--
     processing units it actually needs to run physically on the [INAUDIBLE]..
     At the same time, despite those kinds of constraints,
     you actually have the benefits of having that containerization kind of a feel
     to it, because it's wrapped inside an application.
     And the best part about this, when you're
     talking about very high performance, is the fact
     that you don't have any remote procedure calls or REST APIs in this diagram,
     as you can see, which means that you immediately
     get the benefits of superior latency and system performance.
     So that means your application's going to be very responsive.
     Now, there's a third type of scenario, which is what's
     called white box kind of a method.
     And it gives you the most amount of flexibility.
     In this particular situation, it's really quite fascinating.
     This is a concept that was developed out of Microsoft
     and several other teams [INAUDIBLE] and so forth have actually
     of adopted the same kind of notion.
     And this is about very good flexibility and optimizations.
     And so what ends up happening here is that the models that we have
     are actually registered in with the runtime system,
     and this runtime doesn't think of them as much as models.
     It's actually interested in the operators that are actually inside all
     of the code that's in the models.
     And that's how it actually thinks about it.
     So it thinks of it as like, how can I schedule these different operators,
     different-- let's say a convolution as an operator,
     and those kinds of nuggets of building blocks
     is what it's actually thinking about.
     So the runtime knows about these kinds of characteristics
     that are there inside existing models.
     And what ends up happening is that the application basically
     can receive predictions by incorporating the runtime directly into their logic.
     That's one way of doing it.
     Similar to the way an embedded inference application
     scenario, where you actually plug the model directly in-- you
     can actually think of it that way.
     Or you can have this runtime system that's
     sitting separate, similar to what we talked
     about in the context of the container-based based method, where
     it's housed inside some sort of cloud framework.
     Or it could be a container by itself, and you're just connecting to it.
     So this gives you that flexibility, which is really nice,
     because it's the best of both worlds.
     And what's nice about this is this white box method--
     it's called the white box method because the runtime's actually
     able to look at what's inside the models,
     rather than treating the models as some kind of black box.
     It allows the runtime to do rather advanced optimizations,
     like operator reordering, which helps reduce latency.
     You can also reduce the memory consumption and the computation
     requirements by being able to reuse some of the computation that's actually
     happening.
     You can do this through very clever mechanisms, like hashing and so forth.
     And you can also use very creative scheduling of pipeline components.
     The various kinds of components can actually
     be pipelined to make it an optimal, or at least near-optimal allocation
     decisions when it's actually running on the machine.
     So this way, from a service operator standpoint,
     you actually get very high utilization, and you're not wasting
     any of the resources on your system.
     And so it's sort of the best of both worlds,
     from the container side and the embedded side,
     but it also brings in its own flavor of additional optimizations and design
     space that you can actually explore to make the system really efficient.
     All right, so in this video, we covered three different types of service
     architectures-- container-based approach,
     a direct import into the applications, and the white box approach.
     All right?
     So stop for a moment, and I want you to reflect on this,
     because it's going to be a quiz coming up on what are the key differences
     and what are the pros and cons between each one of these.
     I'll see you then.
     Ende des Transkriptes, zum Anfang springen
     Downloads und Transskripte
     Video
     Videodatei herunterladen
     Transskripte
     Download SubRip (.srt) file
     Download Text (.txt) file
 
\section{Embedded Inference Serving Benchmarks}

    0:00 / 0:00Klicken Sie auf die Pfeiltaste nach oben, um das Geschwindigkeitsmenu aufzurufen. Regeln Sie dann mit den Hoch und Runter Pfeiltasten die Geschwindigkeit. Klicken Sie Enter, um die jeweilige Geschwindigkeit festzulegen.Geschwindigkeit1.0xKlicke auf diesen Knopf und diess Video stumm oder nicht stumm zu schalten, oder drücke die HOCH oder RUNTER Knöpfe um die Lautstärke zu erhöhen oder abzusenken.Maximal Lautstärke.Videotranskript
    Beginn des Transkriptes, zum Ende springen.
    VIJAY JANAPA REDDI: All right, so far we have
    learned about two main things, which is about the prediction serving scenarios
    and the various serving architectures that we can put in place.
    But one of the key things that comes up over and over in TinyML with respect
    to serving is doing careful performance analysis before we actually
    put things into the serving type.
    So in this video, I'm going to talk about benchmarking
    embedded inference for serving.
    Now as you recall, there are wide and many different embedded systems
    out there.
    Now typically, if you're going to be deploying TinyML,
    in any kind of serving situation you're going
    to have to pick one particular system.
    But the key question becomes this-- how do you choose one device?
    Now you've seen this table before, and yes, these devices
    are rather different.
    But even if a particular road looked interesting to you, even with that,
    in reality, you have many different options.
    And now if you have taken my Deploying TinyML course,
    I've had you use Arduino Nano 33 BLE Sense.
    But have you thought about why I chose that particular device?
    Well, my team and I spend a lot of time thinking about it,
    and we found that to be the best pedagogical instrument that's
    out there that we can actually use in order to help you deploy TinyML models
    or learn how to deploy TinyML models.
    Now when it comes to serving time, as in the real world,
    we are to think about it the same way.
    We have to think about why is a particular device actually
    good for serving inference on it?
    And how do you go about finding that right device, given
    that there's so many different devices with so many different characteristics?
    And as a quick reminder, remember, in the context of embedded systems,
    there are many different challenges and those constraints can broadly
    be broken into the hardware side and the software side.
    And they often tend to have extremely tight requirements
    or tight constraints.
    So for example, we tend to have power usage, price, form factor,
    and compute performance constraints.
    Sometimes those constraints are actually requirements, and in other cases,
    they're just constraints, period.
    So somehow we have to navigate this space.
    So real quickly, if you remember, on the embedded systems side,
    the hardware has a lot of heterogeneity.
    We have many different kinds of microcontrollers.
    Even if we are not talking about microcontrollers,
    they tend to have dedicated DSP engines and so forth that help
    do a little bit of MR processing.
    So there's a rich array of them.
    How do you figure out which one of these particular devices
    has got the right computational capability for you?
    And of course, then they also tend to have
    rich number of constraints around the amount of memory,
    the flash, and the SRAM size and so forth, as well as power constraints.
    And on the software side, there are a lot of problems.
    We don't tend to have file systems, we don't
    tend to have basically library support, we
    don't tend to have operating system-level support,
    all of these things that we typically would take for granted.
    And things that, even if we take for granted,
    would run extremely efficiently on a large laptop or large desktop
    or workstation.
    So these factors start coming into some serious play when
    we're talking about considering how we're actually
    going to be serving inference on that endpoint device,
    because we have to pick a system that meets or can deliver
    the requirements that we actually have.
    And in reality, the ecosystem is even more complicated
    than just hardware and software.
    You've got many different kinds of models you can choose from.
    You've got different frameworks.
    You've got different libraries.
    This whole cross product gives me a headache every time
    I look at this chart.
    I generated this chart a long time back, and every time I look at it,
    I'm like, wow.
    It's only going to get more complicated.
    And it should give you a headache, too.
    If it doesn't, I'm going to be really concerned about this.
    Back on track.
    So this cross product means that doing apples-to-apples comparison
    is really challenging in the ecosystem.
    Because you have to be asking yourself, if there's a system X
    and there's a system Y and both of them claim,
    OK, I can do TinyML inference really well, OK, well that's great.
    And someone says, OK, I can do n number of inferences per second.
    OK, that's great.
    But are they running the exact same model?
    What tasks are they doing?
    What data set are they using?
    What's the batch size?
    What kind of quantization are they doing?
    What's the accuracy for the models?
    What software libraries are--
    you've got to consider all these different factors when you are actually
    going to that serving point.
    And that becomes extremely crucial.
    And that's where benchmarking becomes essential.
    The definition of a benchmark is that it's a standard or point of reference
    against which things may be compared or assessed,
    and it is crucial when you want to be able to intelligently navigate
    this rich array of different possibilities
    you have for picking a device.
    Now, benchmarking is used to compare solutions,
    make informed decisions on which device to pick.
    You can measure and track progress about how your device
    or how your models are performing, as long
    as you have some sort of rigorous benchmarking methodology.
    And over time, you can raise the bar, because you're measuring things
    and you're hoping to improve.
    Benchmarks are essential for being able to make improvements.
    But it requires a very strong and systematic methodology
    that is both fair and rigorous.
    In other words, it's got to work across a wide range of systems.
    And you also have to have a lot of community support and consensus
    around this, because you can't just cherry pick some sort of model
    and a task and a data set and accuracy, because that might be totally biased.
    So you have to have some sort of benchmark that
    actually meets the industry standard.
    And it needs to provide those clean standards
    for what are the most important use cases,
    and what are the workloads that are represented of those use cases?
    And how can you systematically compare heterogeneous hardware and software
    systems?
    It's got to be able to provide clarity and do
    some sort of complex characterization you're able to do off the system
    when you run those workloads.
    And you've got to be able to verify and reproduce results.
    And that's what a good benchmark would do,
    and that's, voila, where MLPerf comes in.
    MLPerf is an industry standard for benchmarking machine learning systems,
    and the reason I'm talking to you about this
    is that I want you to have a deep understanding of MLPerf,
    and at least know that it exists, because it
    is the de facto benchmark for machine learning systems, be it a TinyML device
    or be it a ginormous data center that has got a lot of GPUs and CPUs
    and so forth.
    They provide a wide range of different kinds of benchmarks,
    and I know quite a lot about this, in fact,
    because I actually helped create MLPerf.
    And that's one of the reasons I'm passionate about it
    and I want you to really understand it, because it
    provides a lot of value for you as a machine learning engineer or a systems
    person who's trying to figure out which is the right device for you.
    Let me quickly summarize, in six bullets, what exactly MLPerf does.
    It's an industry consortium and it enforces performance result
    replicability to ensure reliable results.
    And it enforces good best practices in order
    to improve the state of the art of ML.
    You want to encourage innovation in this space.
    And it allows you to accelerate machine learning
    in a fair and useful way for everybody.
    It also helps drive both the commercial entities and the research communities
    all together, because machine learning is very much a fast-moving space,
    so you want to bring the industry and academics and everybody
    together to be able to determine what are those workloads and so forth.
    And you've got to keep this benchmarking affordable,
    because everybody, be it a TinyML user or a big data center user,
    has got to be able to have easy access to a industry standard benchmark.
    All right, great.
    So now let's say we have this industry consortium that
    thinks about these kinds of workloads and benchmarks.
    Well, how do they systematically dissect this space?
    How does one go about sampling the workloads to actually benchmark?
    For example here, I'm showing you there's
    a wide range of different use cases.
    In the context of TinyML, it's like beta audio-based use
    case or some sort of sensor-based use case, where
    you're getting telemetry data.
    Or maybe it's a video sensor kind of a use case.
    There are many different use cases, so somehow, we
    have to distill this large space down.
    Not only from a use case, but also from all the different networks they have
    and all the different data sets that are associated with them.
    So somehow we have to have a method for doing this.
    And this is where industry consortiums often
    have good best practices that they bring,
    where they have a principled approach.
    For example, we need to, first and foremost, sit down and define
    what that benchmark task means.
    What's the definition of a benchmark?
    Is it just a machine learning model?
    Is it more than the machine learning model?
    All the goo, the pre-processing, the post-processing and so forth.
    Then, what is the task?
    Because there's so many different paths, so you can't benchmark all of them.
    Which are the most representative ones that will actually
    be useful for the industry to look at?
    And what does performance really mean here?
    Is it latency?
    Is a latency-bounded throughput?
    Which of those metrics that we talked about during scenarios
    are actually the relevant things here?
    Now you see why I thought about metrics there.
    How do folks actually use this benchmark?
    Like, to run on a rich array of different systems?
    So there has to be some standard way.
    What optimizations do you allow?
    What optimizations do you not allow?
    And how do you present or interpret the results,
    and what are the metrics that we need to have here?
    Is it just latency, or is it actually energy
    in the context of TinyML-based systems?
    So MLPerf does all the hard work of debating
    what these different use cases are and what the metrics are and so forth.
    And here, I'm showing you an example.
    For instance, for one instance of the benchmark for a certain version,
    there were four different tasks that they
    had picked that they believed were representative
    across a wide range of different TinyML use cases,
    going from keyword spotting, all the way to anomaly detection.
    Now, as you can see, there's a lot involved
    when you're talking about benchmarking.
    And this is where I'm hopeful you will get an appreciation for what
    the complexities are.
    So let's step through it.
    First and foremost, you got to understand, what problem are you
    trying to benchmark?
    OK, let's say in this particular case, it's
    anomaly detection, which is useful for predictive maintenance and TinyML use
    cases.
    Next, you've got to find a data set.
    Maybe something like TinyML or some other Hitachi or NASA-based data
    sets that are actually publicly available.
    You have to agree on that data set, because everyone's going to use it.
    Then comes a part about actually picking a network
    that you're actually interested in.
    Often, when you're talking about anomaly detection,
    autoencoders become quite useful.
    Then, they can actually work on different kinds of inputs.
    For example, they might be able to work on a spectrogram,
    or they might actually be able to work directly on time to use data.
    Then we have to go into the process of actually training
    this particular network.
    Now, if you remember, one of the key things that I told you
    is model provenance.
    In other words, you have to have the Training
    Code available, not just a model.
    Once you have the model, then you have to allow or figure out
    what kind of optimizations are actually permitted.
    Do you allow pruning, quantization, knowledge distillation?
    What kind of optimizations?
    Because sometimes, some optimizations change the structure of the network.
    If you remember, for instance, when we learned about doing pruning,
    that's going to change the network structure.
    Depending on how aggressive you are, that's
    going to change your network structure by a large amount or a small amount.
    So you have to constrain some of these things.
    Otherwise, you're comparing apples and oranges.
    Then finally, you have to provide some sort of reference implementation
    that can actually be executed on the system.
    And then, you actually do a benchmark run.
    So clearly, there are many different steps here,
    but if you can wrap all of this together in some sort
    of coherent and systematic way, then you've
    got a remarkably powerful tool that'll be able to allow you to compare system
    x versus system y.
    And then, you can make the right choice as to where you're actually
    going to be doing your serving.
    And of course, when it comes to the metrics,
    latency is undoubtedly going to be important in TinyML,
    but at the same time, you also got to worry about accuracy and energy.
    Because if a model that runs really fast--
    that's great, but if it's getting wrong accuracy predictions,
    well, that's not useful.
    And if it's consuming way too much energy--
    given the TinyML, we worry a lot about both latency and the battery
    consumption aspect--
    then we have to measure energy as well.
    Now energy measurement is also something that's remarkably hard to do,
    and there's an entire body of work just dedicated for doing this.
    Again, I encourage you to check out MLPerf for gaining
    access and more knowledge in the space.
    Having said all that, though, one thing that you
    have to be crucial about or careful about around benchmarking machine
    learning systems is that the ML ecosystem is continuously evolving.
    For example, every day there are new use cases coming out.
    Like here, I'm showing you the adidas GMR series,
    which is a smart shoe sole which can basically
    tell if you're kicking, passing, dribbling,
    and all those kinds of things.
    And it's really cool, and it's a smart shoe, basically.
    The second you put this little jagged kind of a piece into your shoe,
    it becomes a smart shoe.
    It's incredible.
    Now what's interesting here is that in this particular case,
    how does the system know, are you kicking
    or are you dribbling and so forth?
    Is that all happening in one model based on the sensor accelerometer gyroscope
    kind of data that's coming in?
    Or is there multiple models?
    Well, they might actually be running multiple models, which
    means that suddenly, you have to think about
    how can I actually support multiple models when I'm
    doing serving on an inputted device?
    Again, there are other interesting use cases in TinyML
    that are coming up around augmented reality,
    where, for instance, machine learning is being deployed
    to be able to tell, for instance, if your hand is actually in the scene
    and where your eyes are looking and so forth.
    And ML is showing that it can actually do
    a much better job than traditional algorithms that we had.
    And all of these have to be deployed on thin little frames like my glasses.
    That's the future.
    Now, these are exciting things, and these
    are the reasons why you want to have different kinds of application
    benchmarks there.
    So the point that I'm trying to get across
    is that the space is continuously evolving,
    and therefore, you need to have a rich class of benchmarks
    you can actually use to study the capabilities of your serving endpoint.
    In the future, I think we're going to see, not just a single model running.
    We're going to see pipelined DNNs, where things are executing back-to-back
    so it's more than one model running.
    You're going to see concurrent DNNs--
    in other words, things are happening simultaneously,
    like in the GMR shoe, which I explained with dribbling and kicking
    and so forth.
    Or you might actually have some kind of crazy combination between pipeline
    and concurrent models.
    So these benchmarks that come from MLPerf
    and so forth continue to stay up to date,
    and they are a great resource for us to understand how well our serving systems
    are actually doing.
    And it's extremely crucial to understand this,
    because you don't want to just go through all the effort of deploying
    the model and then not take a close look at how the serving is actually
    happening using the right set of metrics.
    And that's where institutions like MLPerf and so forth come
    in extremely handy.
    With all that said, hopefully you have an appreciation now
    as to why comparing systems is both challenging and yet crucial,
    and how entities like MLPerf. which create ML benchmarks,
    can actually be very useful for us in analyzing prediction
    serving performance.
    Ende des Transkriptes, zum Anfang springen
    Downloads und Transskripte
    Video
    Videodatei herunterladen
    Transskripte
    Download SubRip (.srt) file
    Download Text (.txt) file  
    
\section{Embedded Benchmarks: An Overview
    
    \subsection{Motivation}
    
    
    Benchmarking is a very important topic in computer science. Benchmarks are used to measure the relative performance of systems, programs, or other operations. This performance is measured through a variety of standardized tests and trials that measure, such as time taken to complete a task, or the power consumption of an operation.
    
    Benchmarks act as a figure of merit that help to determine the superiority of a particular approach, which can often dictate the path of future research. A good example of this is AlexNet, which debuted at the ImageNet Large Scale Visual Recognition Challenge in 2012. AlexNet was the prototypical deep neural network, the world’s first introduction to deep learning, and it’s statistical power was demonstrated clearly through benchmarking.
    
    The ImageNet challenge used classification accuracy on the ImageNet dataset as the performance benchmark to assess the model submissions of contestants. AlexNet achieved an error rate of 15.3\%. This error rate would be quite poor today, but at the time, the second best model had an error rate of 26.1\% - this is 10.8\% higher! Given that ImageNet consists of 1 million images, this means that it correctly classified approximately 100,000 more images than its closest competitor.
    
    \subsection{But What Constitutes A Benchmark?}
    
    A benchmark means different things to different communities. In the ML community, a benchmark refers to a dataset with some associated metrics attached to it (e.g., accuracy, model size) to measure model performance. To the hardware community, a benchmark is a set workload to run on hardware (e.g., a set of matrix calculations to perform) with an associated set of metrics monitored, typically MIPS/FLOPS (latency) or power consumption. The ML systems community dovetails the ML and hardware communities, which means that a benchmark must function as both.
    
    Benchmarks are important as they provide a means of standardization that allows incremental improvements to be clearly demonstrated, as in the case of AlexNet. If every paper was based on a different dataset, using a different set of metrics, it would be difficult to interpret the effectiveness of any given approach. By providing some standardization, we can (1) define a set of rules defining the task, as well as what optimizations are allowed and how they must be implemented (2) help to establish a set of metrics that are meaningful and easily interpreted for the given task, (3) require submitted implementations to be public to promote open-access and reproducibility, and (4) control for elements of randomness that may impact results. This standardization helps to combat some of the problems that are common across the research landscape, such as results that are reported but not public or reproducible, or solutions that are fine-tuned on a specific dataset but not generalizable.
    
    \subsection{Benchmark Challenges}
    
    Since the goals of two different systems (especially in ML) may be different (e.g., optimizing for latency vs. accuracy), different benchmarks are needed for different tasks. Even within the same broad task (e.g., computer vision), there may be applications that may not be exactly the same. For example, an object detection algorithm would require details for each item within the image, as opposed to just a classification value for the whole image as in image classification. Because of this, many benchmarks have been created to cover a variety of different tasks, utilizing different datasets, metrics, and rule sets for reporting purposes.
    
    Some challenges do exist with benchmarks. For example, some benchmarks may be “synthetic”, which do not represent real workloads but aim to measure performance relevant to real use cases. This can be problematic, since if the benchmark is not correctly aligned with its associated use case, improvements in the benchmark may not correspond to improved performance in real-world implementations. Naturally, this is also a potential issue if the benchmark is not correctly configured or contains some inherent biases (e.g., imbalanced data classes, specific image orientations) that may not be present in real-world implementations.
    
    \subsection{Examples}
    Many benchmarks exist today in both the ML and embedded systems communities. Common benchmarks for ML include: MLPerf inference, MLPerf Mobile, DAWNBench, and AIBench. Benchmarks within the embedded systems community have been around for longer than the ML ones and so have evolved over the years. Nowadays, some of the more common benchmarks include CoreMark (replacing Dhrystone), EEMBC’s ULPMark for ultra-low power systems (three variants: ULPMark-CoreMark, ULPMark-CoreProfile, ULPMark-PeripheralProfile) and Embench.
    
    While some other academic benchmarks exist, MLPerf Tiny is currently the only industry-standard TinyML benchmark, which we will go into in more detail in later sections.
    
\section{MLPerf Tiny}

\subsection{What is MLPerf Tiny}

    Now that we understand the importance of benchmarking for assessing performance, a natural question for us to ask is: what benchmarks exist for TinyML? This is where MLPerf comes in. MLPerf is a benchmark for ML performance, but has recently been extended to include MLPerf Tiny: a benchmark for measuring TinyML performance. Like any benchmark, it aims to provide fair comparisons between solutions by outlining a set of tasks and rules. MLPerf Tiny was created via a collaboration between numerous individuals from academia and industry who participated in a working group organized by MLCommons.org.
    
    \subsection{Use Cases}
    
    MLPerf Tiny currently supports four different ML application areas: keyword spotting, visual wake words, image classification, and anomaly detection. There are two submission windows per year, and two separate submission divisions that focus either on hardware or data-centric design. Each benchmark has a quality target, below which the submission will not be included.
    
    \begin{tabular}{cccc}
    Use Case	&    Dataset    (Input Size)&   Model   (TFLite Model Size)&    Quality Target   (Metric)\\
    eyword Spotting&	Speech Commands (49x10)	&DS-CNN (52.5 KB)&	90\% (Top-1)\\
    Visual Wake Words&	VWW Dataset (96x96)	&MobileNetV1 (325 KB)&	80\% (Top-1)\\
    Image Classification&	CIFAR10 (32x32)&	ResNet (96 KB)&	85\% (Top-1)\\
    Anomaly Detection&	ToyADMOS (5*128)&	FC-AutoEncoder (270 KB)&	85\% (AUC)\\
\end{tabular}
    
    \subsection{MLPerf design: flexibility and representativeness}
    
    MLPerf Tiny is more than a set of datasets like ImageNet that allow people to compare classification accuracies. One of the things that makes benchmarking TinyML devices more challenging than traditional models is additional functional requirements, such as latency, memory footprint, and power consumption. MLPerf Tiny must keep track of all of these, including model performance, in order to suitably compare performance across systems, which can exhibit heterogeneity at all levels of the system.
    
    \GRAPHICSC{1.0}{1.0}{TinyML/4 MLOps/target.png}   
    
    Overview of the TinyML Stack. Diversity is present at all levels of abstraction, making it challenging to enforce standardization for benchmarking purposes.
    
    For example, let’s say that an individual used the Visual Wake Words dataset to perform person detection. Assume that they create a model using a specialized hardware platform that has a very high throughput compared to other submissions for the person detection. If our benchmark was purely focused on throughput (i.e., inferences per second), then even if the benchmark achieved 0\% accuracy, it would appear to be the best model. 
    
    In reality, there is always some tradeoff between these parameters, and so MLPerf Tiny calculates accuracy, power consumption, and latency to allow these to be effectively compared. MLPerf Tiny also has two submission divisions, one which targets hardware (closed division) where only hardware parameters can be altered, and a second division (open division) focused on data-centric design, wherein the model architecture and training data can be altered.
    
        \GRAPHICSC{1.0}{1.0}{TinyML/4 MLOps/accelerator.png}   
    
    Overview of the modular design of MLPerf Tiny. Benchmarks can be submitted in two categories: open division and closed division. In open division, model architectures, training data, and other parameters are allowed to be altered. In closed division, only hardware changes can be made, in order to keep the task as standardized as possible for comparison between systems in terms of accuracy, power consumption, and latency (given in inferences per second). This means that the reference implementation's training data, training scripts, model architecture, and pruning is held constant in both divisions, while in the closed division the reference quantization scheme, compiler, inference framework, kernels, SoC, CPU, and accelerator are also held constant.
    
    \subsection{How to get involved}
    
    MLPerf Tiny is managed by MLCommons and is currently on v0.7. The project is still in its early stages and is likely to evolve beyond what we have discussed here. For the interested reader, we have provided links to the MLPerf TinyML code and paper, and instructions on how to run the inference benchmark and provide submissions.
    
    The code for MLPerf Tiny is available on \url{GitHub: https://github.com/mlcommons/tiny}
    
    The benchmark design is describe in the paper: \url{https://arxiv.org/pdf/2106.07597.pdf}
    
    If you are interested in running the benchmark, you should start by following the \href{https://github.com/mlcommons/tiny/tree/master/benchmark/reference_submissions/image_classification}{instructions on how to deploy the image classification reference submission}. Next you would use  EEMBCs \href{https://github.com/eembc/energyrunner}{EnergyRunner\textsuperscript{\texttrademark}} benchmark framework to connect to the system under test and run the benchmarks.        
        
\section{Prediction Serving Impact on MLOps}

    0:00 / 0:00Klicken Sie auf die Pfeiltaste nach oben, um das Geschwindigkeitsmenu aufzurufen. Regeln Sie dann mit den Hoch und Runter Pfeiltasten die Geschwindigkeit. Klicken Sie Enter, um die jeweilige Geschwindigkeit festzulegen.Geschwindigkeit1.0xKlicke auf diesen Knopf und diess Video stumm oder nicht stumm zu schalten, oder drücke die HOCH oder RUNTER Knöpfe um die Lautstärke zu erhöhen oder abzusenken.Maximal Lautstärke.Videotranskript
    Beginn des Transkriptes, zum Ende springen.
    VIJAY JANAPA REDDI: Congratulations on making it
    through the prediction serving portion of the MLOps pipeline.
    We covered a lot of different concepts in this.
    We learned about the differing serving methods that we have,
    like online inference.
    There was a batch inference.
    There was a streaming.
    Inference.
    There was an embedded inference.
    We talked about what they were.
    We talked about where are they applicable,
    what are some of the mechanisms that need to be put in place,
    and what are the metrics that are associated with them.
    Now, I also talked to you a little bit about benchmarking machine learning
    systems because when you're talking about serving,
    you need to be able to understand which system is going to be a good choice
    and which you should be doing serving.
    Is really machine learning system x the one that you should be picking,
    or is it system y.
    In order to answer that question, you have to be able to compare the systems.
    And therefore, you need industry standard benchmarks.
    So hopefully, you keep that in mind.
    But as always, we have to ask ourselves, what
    is the impact of the prediction serving states
    on the rest of the MLOps pipeline?
    Now, you have to realize that the way you do serving affects how you create
    or set up the match operationalisation stage.
    One of the first questions you'll need to answer
    when you're deciding how to deploy your ML models
    is whether to use batch inference, or online inference, or embedded
    inference, or streaming inference.
    This choice is obviously going to be driven heavily
    by who is using the inference system and how soon do they need the results.
    In other words, if the predictions don't need to be served immediately,
    then you can opt for a very simple batch inference method.
    And you can use that in the operational decision stage
    without too much headache.
    But if it turns out your predictions need
    to be served on an individual basis and within the time of a single request,
    then online inference is the way to go.
    Well, in that case, it's definitely going
    to affect how you do your CI/CD flows and how you
    do your production deployment methods.
    Either way, the bottom line that I'm trying to get to
    is that the sort of decisions that you make early up or down below here
    is actually going to have a ripple effect one way or the other.
    And this is key to keep in mind.
    Now, we also further discuss benchmarking, right.
    Now, this is something you might normally
    or should do during model conversion or model deployment.
    Where you do just kind of really depends how you organize your analysis pipe.
    But the key thing hopefully you're taking away really
    is that you need to have some of these core nuggets,
    like benchmarking, and performance evaluation,
    and so forth, somewhere in your MLOps pipe.
    And how you do that, well, that's your MLOps design at large.
    Finally, I want to close out by saying that what we do in prediction serving
    is extremely strongly tied to continuous monitoring that is coming up next.
    The reason is because when you're serving online requests,
    then you have to worry about drift.
    Well, in order to know if the distributions of that data that you're
    seeing and doing the online phase match to what your training data
    distributions were or not, you need to have the right infrastructure and bells
    and whistles put in place.
    That in itself as a whole kind of worms so we'll be getting into next.
    But hopefully, you're starting to see that there
    are lots of interdependencies here as to what happens within prediction
    serving in the rest of the pipe.
    Last but not least, I want to say that these
    are the kind of things that your ML engineer is really going
    to be dealing with on a daily basis.
    In addition, your software developer and your DevOps team
    is going to be quite deeply involved, depending
    on the particular type of inference serving that you're actually picking.
    That said, I hope you learned something, and I'll see you in the next video.
    Ende des Transkriptes, zum Anfang springen
    Downloads und Transskripte
    Video
    Videodatei herunterladen
    Transskripte
    Download SubRip (.srt) file
    Download Text (.txt) file        
        
% 1.9: Continuous Monitoring        
        
\section{Overview of Continuous Monitoring}

    Module Objectives
    
    \GRAPHICSC{1.0}{1.0}{TinyML/4 MLOps/5.9_Continuous_Monitoring.png}   
    
    The MLOps pipeline. Going from ML Development to ML training to Continuous Training to Model Deployment to Prediction Serving to Continuous Monitoring. Along the way artifacts are produced at each step including: Code and Configurations, Training Pipelines, Registered Models, Serving Packages, and Serving Logs. All steps tie into the overarching themes of Data and Model Management.
    
    \subsection{Objective}
    
    Once we have successfully deployed a model to our production environment, the next stage is to monitor its performance. This stage is often called continuous monitoring and is of great importance in MLOps. Primarily, continuous monitoring is done to validate model performance, often to ensure that it does not degrade over time. If the incoming data distribution gradually deviates over time, the model performance will gradually degrade due to a mismatch between the training and inference data distributions. In addition, there may be scenarios where other coupled services might experience errors or changes which result in poor predictions being outputted by our model. Continuous monitoring allows us to monitor our predictions and ensure that things are working as anticipated, and can raise alerts if any unexpected conditions occur. 
    
    To this end, in this section, we will learn the following concepts:
    
    \begin{itemize}
      \item Problem with Model Drift
        \begin{itemize}
          \item Concept drift
          \item Data drift
          \item Addressing drift
        \end{itemize}
      \item Challenges with Continuous Monitoring for TinyML
      \item Potential solution via Federated Machine Learning
    \end{itemize}
    
    \subsection{Motivation}
    
    As an example of continuous monitoring, an e-commerce model predicting a customer’s propensity to buy a product might be expected to degrade over time as consumer preferences change. In such a situation, retraining triggers might be added where a specific metric is monitored and, if model performance drops below a designated threshold, then the model is retrained using the most up-to-date data. Having retraining triggers may be preferable to continuous training in some scenarios, especially if models are large and computationally intensive to train. To achieve this, a continuous monitoring process consists of the following:
    
    \textbf{Logging.} Predictions served to the user are logged and designated “serving logs”. These predictions may be stored in their entirety or a subset of predictions may be stored if deemed suitable based on the context.
    
    \textbf{Periodic Checks.} The monitoring engine checks the serving logs periodically and computes statistics for the serving data, such as the distribution of served values.
    
    \textbf{Serving Comparison.} The schema generated by the monitoring engine is then compared to a “reference” schema which represents what we would expect from our prediction serving distribution. This stage allows the detection of distribution skews between the served data and the reference schema. In addition, if true labels are available for the serving data, these can be used to evaluate the effectiveness of the deployed model’s predictions.
    
    \textbf{Error/Anomaly Handling.} In the event that an anomaly or error occurs, or a threshold is passed indicating decaying performance of a model, flags can be raised. For example, a decayed model might send an email to the team responsible for monitoring models with relevant information, such that the team can assess the error and decide whether any interventions are necessary.
    
    \subsection{Model Decay}
    
    Typically, model decay is described in terms of concept drift and data drift.
    
    \begin{description}
      \item[Data drift] describes a growing skew between the data distribution used for training, and that currently exhibited in the newly served predictions. 
      \item[Concept drift] refers to an evolving relationship between the input parameters and the output of our model (i.e., the data distribution might be the same, but the fundamental mapping between the input data and our target variable has changed).
    \end{description}
    
    \subsection{Metrics}
    
    Different techniques can be used to evaluate the presence of concept drift vs. data drift. Data drift might be assessed using novelty or outlier detection, while concept drift might be assessed using feature attributional changes. Common metrics that are monitored by the continuous monitoring process are:
    
    \begin{description}
      \item[Resource utilization.] This refers to the usage of CPU, GPU, and memory resources.
      \item[Latency.] This refers to how fast predictions are served for streaming deployments and is a good indication of model service health. Generally, you would prefer there to be a small latency between a user request and the model’s response. Latency may even inadvertently impact the data distribution (you could imagine not purchasing an item because a model is running too slowly on a website and you lose interest).
      \item[Throughput.] This refers to how many predictions are served in a given time and is a key deployment metric. As throughput changes, this might influence various aspects of the serving system, including the data distributions.
      \item[Error Rates.] For example, the number of times the model predicted an individual would purchase an item, but was incorrect, and vice versa.
    \end{description}
    
    In the remainder of this section, we will discuss the above aspects of continuous monitoring in greater detail, and see how these tie in with the other aspects of MLOps that we have already discussed in previous sections.
    
    \subsection{Additional Resources}
    
    \begin{itemize}
      \item \href{https://docs.google.com/document/d/14uX2m9q7BUn_mgnM3h6if-s-r0MZrvDb-ZHNjgA1Uyo/edit}{Data Distribution Shifts and Monitoring} - Stanford CS 329S
    \end{itemize}
            
\section{Continuous Monitoring}

    0:00 / 0:00Klicken Sie auf die Pfeiltaste nach oben, um das Geschwindigkeitsmenu aufzurufen. Regeln Sie dann mit den Hoch und Runter Pfeiltasten die Geschwindigkeit. Klicken Sie Enter, um die jeweilige Geschwindigkeit festzulegen.Geschwindigkeit1.0xKlicke auf diesen Knopf und diess Video stumm oder nicht stumm zu schalten, oder drücke die HOCH oder RUNTER Knöpfe um die Lautstärke zu erhöhen oder abzusenken.Maximal Lautstärke.Videotranskript
    Beginn des Transkriptes, zum Ende springen.
    LARA SUZUKI: So far, we have covered MLOps aspects
    related to building machine learning pipelines, training, deploying,
    and serving them.
    Now it's time to move to the next stage of our MLOps framework,
    that is continuous monitoring.
    As we will see, this is the stage concerned
    with serving effective and efficient models in production.
    As described earlier, continuous monitoring
    concerns with monitoring the effectiveness and efficiency
    of a deployed machine learning model in production.
    Continuous monitoring is viewed to regularly and proactively
    monitor and assess the model performance, its effectiveness,
    for instance, whether there is any data drift in the schema or distribution
    of the data, and also concept drift.
    Regarding the deployment of a CI/CD pipeline,
    automated tests validate whether the module
    is to behaving as expected across all business KPIs and metrics.
    Once data scientists they deploy the model into the serving environment,
    they can start to log in predictions and explanations for each prediction.
    They can also set up monitoring routines to keep
    a check on what the model is witnessing in the production environment.
    They can also capture metrics regarding resources such as CPU, GPU, latency,
    QPS, error rate, amongst others.
    As we move towards a more sustainable use of technology,
    we can also capture metrics related to energy consumption and carbon
    footprint, compute, and also run time.
    In this final phase, the machine learning model
    is deployed for life prediction in batch, or online inference.
    When issues are detected, services can trigger full
    on processes such as human investigation or a retraining procedure.
    This could be something as simple as a notification to the model
    owner, or a batch word alert for business users.
    The output of this final phase is of course
    the live predictions, logs and records of the production inferences.
    We also have the alerts and triggers for monitoring to help teams to promptly
    react as needed.
    The core MLOps capabilities employed at this stage
    are dataset and feature repository, model monitoring,
    and machine learning metadata and artifact repository.
    The continuous monitoring process regularly and proactively assess
    the model performance.
    A typical continuous monitoring process looks like this.
    A sample of the request response payloads is captured.
    It loads the latest inference logs, generate a schema
    and computes the statistics for the serving data.
    Then it compares the generated schema to a reference schema
    to identify schema SKUs, and also compare the competitive statistics
    to baseline statistics to identify distribution skills.
    If the true [INAUDIBLE],, the ground truth, for the serving data
    is available the monitor engine uses it to evaluate
    the model predictive effectiveness in hindsight of the serving data.
    If anomalies are identified, or if the model performance is decaying,
    alerts can be issued to various channels, for instance email or chat,
    we can notify the owners of the model to examine the model
    or to trigger a new retraining cycle.
    As it can be noticed, there's a continuous evaluation
    of the model, data drift and concept drift,
    plus efficiency of resources such as CPU GPU, latency, and others
    as mentioned before.
    Ende des Transkriptes, zum Anfang springen
    Downloads und Transskripte
    Video
    Videodatei herunterladen
    Transskripte
    Download SubRip (.srt) file
    Download Text (.txt) file   
   
\section{Model Drift: The Big Picture}

    0:00 / 0:00Klicken Sie auf die Pfeiltaste nach oben, um das Geschwindigkeitsmenu aufzurufen. Regeln Sie dann mit den Hoch und Runter Pfeiltasten die Geschwindigkeit. Klicken Sie Enter, um die jeweilige Geschwindigkeit festzulegen.Geschwindigkeit1.0xKlicke auf diesen Knopf und diess Video stumm oder nicht stumm zu schalten, oder drücke die HOCH oder RUNTER Knöpfe um die Lautstärke zu erhöhen oder abzusenken.Maximal Lautstärke.Videotranskript
    Beginn des Transkriptes, zum Ende springen.
    In the previous video, Lara introduced us
    to the big picture of continuous monitoring.
    In this section, we're going to dive deeper into that big picture.
    And try and understand what are some of the critical building blocks.
    Let's get started with that.
    There's only one thing that is constant in life, and that is change.
    And that is very true when it comes to machine learning models as well.
    More formally, this is known as drift, in the machine learning ecosystem.
    To be even more specific, it's known as model drift.
    Modeled after model decay refers to the deterioration
    of a model's predictive power over time.
    All models deteriorate in some capacity or another.
    Low data quality, malfunctioning pipelines, or technical issues all
    leads to having poor prediction capabilities.
    For example, in this case, we're looking at accuracy prediction
    dropping over time in terms of months.
    We can observe that the model's prediction accuracy has gone down
    from 90% to around 55% in April.
    This is a totally hypothetical situation,
    but nonetheless, you can imagine that this would happen in the real life.
    And it is very critical to detect this sort of a drift
    in order to prevent inaccuracies in the model.
    Because failure to do so can actually have detrimental influence
    on your business outcomes.
    Because you're using ML to generate revenue for yourself.
    In the continuous monitoring stage, you need two sorts of drift detectors,
    fundamentally.
    You need something that's actually going to be watching out
    for the decline in accuracy.
    And another is to actually look for detecting drops in data consistency
    within the model.
    Now a drop in either one of these is really sort of critical,
    because it really ends up impacting how the model performs in real life.
    Broadly, you can bucket this into concept drift or data drift.
    In this particular video, I'm going to focus on the high order
    bits around drift, and in subsequent videos
    I'm actually going to go a little bit deeper
    into each one to give you the tangible pieces that
    actually fall into each of these.
    Let's revisit this example that I gave you during the continuous training
    stage of the MLOps pipeline.
    It's about how COVID has actually influenced our behavior over time.
    For example, it demonstrates that we tend to wash our hands more.
    It shows that we tend to isolate or distance ourselves more from people.
    We've stopped eating outside and started stockpiling food and water
    because we don't know what's going to happen next.
    Most notably, we all wear masks.
    So plainly, these are all extreme examples
    of deviations in human behavior from the norm, significant deviations.
    Now consider what would happen if a machine learning model was
    trained on some of our past behavior.
    Would it work well under these sorts of dramatically changed conditions?
    Well, that's where things get really interesting.
    Let's talk about some real issues that we
    have seen about how ML models have performed in the real world since COVID
    has come about.
    Because of our hastily kind of stockpiling food and water,
    the model's predictive capabilities in terms of store availability
    has declined from 93% to 61%, because our purchasing preferences are so
    radically different nowadays.
    Bankers have started to question credit algorithms,
    because it's credit algorithms who are trained when things were good,
    but a lot of people have lost their jobs.
    So now things are very different.
    So can you trust those machine learning algorithms
    to actually make good predictions?
    Furthermore, if you look at the stock market,
    there's been a lot of volatility in the market.
    Some of the funds are down by 21%, because their predictive capabilities
    are drastically different, because we're no longer doing stock trading
    the same way as we used to before.
    Because again, the market is behaving very differently.
    Of course, you know, now when you're sitting in front of a computer
    what are you doing?
    Most of the time you're working.
    In the past, used to probably be because you're
    FaceTiming with your friends or your family and having a great time.
    Now, we're all on Zoom, just kind of working most of the time.
    Take a look at me.
    I mean, for instance, all these videos are being recorded in my home office.
    I happened to start the TinyML course series right around when
    COVID launched, coincidentally.
    And I expected COVID be gone in like two to three months.
    But look, hey, it's been almost two years,
    and I'm still here standing in front of you recording.
    So the point is that our behaviors have changed dramatically.
    Getting back on point, last but not least, weather forecasts
    have become extremely less accurate.
    And the reason is because typically weather forecasts
    used to use data that used to come from commercial planes.
    But now, commercial planes have been canceled quite a bit.
    And so that means we're not getting access
    to all this rich information from that the planes used to collect.
    So as you can see, machine learning models under COVID
    have had a significant impact in terms of the predictive capabilities.
    So this is a very concrete example that hopefully you can really connect with.
    So coming back up to the MLOps pipeline, it's
    critical to understand how models' predictive capabilities change,
    and how that ends up affecting the rest of the pipeline.
    The ultimate thing that you really want to do
    is that you want to retrigger your training pipeline.
    Now it's extremely important to retrigger the training pipeline,
    because if you don't retrigger or you don't
    retrain the models from time to time, you're
    going to have poor prediction capability.
    So therefore, the continuous monitoring stage is an extremely crucial stage.
    Though it seems like it's right at the end.
    By no means is it any less important.
    In fact, I will emphasize that it's extremely important to get right.
    Now if I were to open up the continuous training stage
    and look at the building blocks in there,
    one of the first blocks that we actually pay attention
    is to the retriggering block.
    What retriggers it?
    Well, it's a continuous monitoring that sends a signal to the retraining
    block and then that's the one that actually goes ahead and fires up
    the rest of this complex pipeline, which we talked
    about in the context of these engines.
    We discussed how we can have an automated keyword data set
    generation pipeline that is able to adapt to languages and changes
    in the context of keyword spotting.
    We talked about how AutoML, neural architecture search,
    and several other techniques can be used during the continuous training process
    to actually generate highly optimized models that are highly tuned
    for the data that's coming in.
    Finally, we talked in the model evaluation
    step about different kinds of metrics that are
    associated with either use cases or just the training pipeline itself.
    Anyway, going back upwards on this image, to the training pipe,
    back up to the retraining stage and back up to the MLOps pipeline,
    you will see that drift detection is extremely key.
    This concept of drift detection, which the continuous monitoring
    pipe is going to do, is going to have a dramatic impact on that retraining
    stage.
    Now if I were to go back into the continuous monitoring pipe,
    then one of the first blocks you'll see is this alert.
    What actually triggers that alert that raises the retraining flag?
    Well, it really falls into three main stages.
    It's the concept drift, the data drift, and the continuous evaluation.
    And we'll be talking quite a bit about these three things
    in this particular module.
    Scaling further, you will see that there are many more components involved here
    in terms of model monitoring.
    So to that end, we need a large team of people.
    As you can see, there are numerous personas
    that I'm showing you that will play an important role
    in continuous monitoring.
    This is one of the first areas where I'm demonstrating
    that you need a lot of different domain expertise to be able to pull this off.
    So this should give you an idea of how important continuous monitoring is.
    With that said, let's dive into continuous monitoring's concept drift,
    and I'll see you there.
    Ende des Transkriptes, zum Anfang springen
    Downloads und Transskripte
    Video
    Videodatei herunterladen
    Transskripte
    Download SubRip (.srt) file
    Download Text (.txt) file

\section{Concept Drift}

    0:00 / 0:00Klicken Sie auf die Pfeiltaste nach oben, um das Geschwindigkeitsmenu aufzurufen. Regeln Sie dann mit den Hoch und Runter Pfeiltasten die Geschwindigkeit. Klicken Sie Enter, um die jeweilige Geschwindigkeit festzulegen.Geschwindigkeit1.0xKlicke auf diesen Knopf und diess Video stumm oder nicht stumm zu schalten, oder drücke die HOCH oder RUNTER Knöpfe um die Lautstärke zu erhöhen oder abzusenken.Maximal Lautstärke.Videotranskript
    Beginn des Transkriptes, zum Ende springen.
    VIJAY JANAPA REDDI: In this video, I'm going
    to introduce the notion of concept drift.
    Now, concept drift is a term that's used in machine learning and data science
    to describe changes in the relationships between the input and output
    data of a machine learning model.
    In other words, the way your model actually predicts the target variable
    that it's trying to predict is actually going
    to look different given a certain probability
    distribution of input samples that it's actually been trained on.
    Sometimes this is known as a covariate shift as well.
    Now, this is the kind of thing a data scientist is typically
    going to be really concerned with as it is his or her obligation to make sure
    that they can detect the concept drift.
    Now, as shown on this slide, it's an essential component
    of modern monitoring.
    There are many numerous forms of concept drifts that emerge.
    In practice, they can happen suddenly.
    They can happen incrementally.
    They can happen gradually.
    Or in fact, they could actually recur over time.
    And while not shown, another variant of this
    is simply that it could be a little brief burst of noise.
    But nonetheless, that is a certain type of drift you might actually encounter.
    What I want to do first now is actually go
    through each one of these at a high level and explain what they are,
    and then I'll give you very specific examples for each one.
    A sudden concept drift occurs when we transition from an old concept
    to a new concept all at once.
    Now, here I want to emphasize the old concept is
    one where we're seeing in green circles, and the new concept
    is one that we're actually looking at in the blue circles.
    And I'll be using this sort of a pattern to explain all the concepts.
    Now, gradual is when the change between the concept happens over time.
    The new concept occasionally pops up as we're
    seeing with the blue dots on top earlier in time.
    And then over time, it shifts into that.
    Another one is incremental.
    Incremental is somewhat close to gradual in that it's slowly shifting,
    but you're seeing a gradual probability distribution
    change about what your sampled characteristics really are.
    So that's a bit of a subtle sort of a difference
    between gradual and incremental.
    People sometimes tend to confuse them a little bit.
    And then finally, you have recurring or seasonal shifts.
    These are changes that recur over time from the initial observation you
    might make.
    So the same thing that you see with the green dots is coming up over and over.
    One example of a seasonal drift is really
    weather where because the weather has a certain pattern to it
    so it recurs over time.
    I've included all of them here in one big screenshot.
    So stop for a moment and make sure you really digest this
    because these are critical concepts that you just
    need to understand in terms of drift.
    Next, what I want to do is give very specific examples for each one of them.
    Let's take the sudden concept drift.
    You've already seen this.
    I previously mentioned that our behavioral patterns have changed
    dramatically with respect to COVID.
    We've got lockdowns.
    We have quarantines.
    So our behavior is completely ad-hoc all across the world.
    So no matter where you are on this planet,
    your behavior is definitely different.
    As a result, this is an excellent illustration for sudden concept drift.
    Your models which had been trained previously
    are not going to do so well under COVID situations.
    You're going to have to dynamically adapt to them.
    Gradual drift, again, is a process through which
    where you are occasionally seeing samples coming about.
    Now, a good example of gradual concept drift
    is how our smartphone behavior actually changed over time.
    I'm showing you made-up data down here.
    Nonetheless, this trend generally holds true, I'm sure.
    In early years, when we had feature phones,
    we would mostly be making audio calls, and we had little
    to do with the camera or the mobile internet.
    But as time went by, things changed.
    You can see between 2005 and 2010, we started
    to see more sophisticated phones coming about, truly smartphones coming about.
    For example, the iPhone came out in 2007 and radically transformed
    the way we use phones.
    For example, we started spending much more time actually doing
    texting on that beautiful glass device that they created.
    And when smartphones ultimately became the norm,
    we saw an increase in terms of our smartphone camera usage,
    and we also started using more of the mobile internet.
    So hopefully, you can see that this is a good illustration
    of gradual behavioral change.
    And I say gradual change here because there are always
    enthusiasts who sample some of these things ahead of time.
    So it's those little blue dots that I was showing you earlier.
    But over time, what we see is that that sort of becomes the norm.
    More and more people just start doing it.
    So that's why I'm using this as an example of gradual change.
    Now, shifting the focus back to the incremental drift, incremental drift,
    again, is where things are progressively changing, right?
    This, I think, is an excellent example of an incremental drift.
    It's to do with the discharge rate of lithium ion batteries.
    Did you know that your batteries, be it in your Tesla car
    or in your smartphone, they all behave differently over time?
    The amount of voltage that the battery puts out changes over time.
    And believe it or not, machine learning models
    are actually used in order to predict how the battery is aging
    or what the discharge curve is going to look like because the last thing you
    want is your device to be telling you have 90%
    charge when, in fact, it only has 20% charge.
    Or worse, if you're trying to do something that's extremely critical,
    but you only have 5% charge and the model totally mispredicts it and says,
    oh, you have 10%.
    So there's a big difference because when you're really running out of juice,
    you're going to be much more eager to try and finish up whatever task.
    So in this sense, it really matters, and people
    do a lot of user studies on this.
    Another reason why people actually use ML on these kinds of things
    is because no two batteries are actually alike.
    You use your phone differently than I do.
    And that means that the batteries age differently.
    They deteriorate differently, which means the discharge
    curves are going to look different.
    So online models are actually used in order
    to track how these batteries are discharging.
    However, you can basically see that the discharge rate is changing
    incrementally over time, and that's why I'm using this
    as an example for an incremental drift.
    Now, when it comes to this notion of recurring drift, well, recurring drift,
    unfortunately, I have to pick a rather sad example.
    I hate to be pessimistic, but a lot of people
    do suffer from recurring related issues.
    For example, people really suffer from seasonal depression,
    which is a medical issue, and it affects quite a lot of people.
    This is one of those times when people's behavior changes,
    and this happens regularly over the year.
    When holidays arrive, some of us feel sad, lonely, excluded, and so forth.
    And this is a real genuine thing, and it has a significant impact
    in terms of how people who are affected with this sort of depression actually
    behave.
    Now, consider that you're using machine learning in order
    to just provide them everyday services.
    Their behavior is going to be radically different.
    So this is an example where, for instance,
    you're seeing recurring drift, and you want to be able to handle it.
    Now why, once again, are we so concerned with these different types of drifts?
    We're talking about concept drifts now right now.
    And why are we so concerned about it?
    Well, it's really because if we're not careful about it,
    it can really bite us in our behinds down the road.
    So I'm trying to illustrate this notion of a concept drift
    here with the green and blue dots.
    And let's say that line that's cutting across is your concept.
    The concept is how do I separate the blue dots from the green dots,
    and your machine learning model has learned to do that.
    Now here on the right-hand side, you are looking at a graph
    where the data distribution actually looks different.
    Underlying distribution has changed so much so
    that if you were to figure out how to cut across the blue and the green dots,
    that line would actually look different, which
    means that it's a fundamentally different sort of a prediction you
    have to make.
    So this is a very nice illustrative way of understanding
    the notion of concept drift.
    And you have to be able to tackle this.
    This is where your model starts dropping its accuracy,
    and it could be because of sudden drifts,
    or it could be because of gradual rates, or it
    could be because of incremental drifts, or it
    could be because of recurring drifts.
    With that said, in the next video, I'm going
    to pick up talking about data drifts, or how do you
    sort of understand the data characteristics
    as the data might also shift.
    Ende des Transkriptes, zum Anfang springen
    Downloads und Transskripte
    Video
    Videodatei herunterladen
    Transskripte
    Download SubRip (.srt) file
    Download Text (.txt) file

\section{Data Drift}

    0:00 / 0:00Klicken Sie auf die Pfeiltaste nach oben, um das Geschwindigkeitsmenu aufzurufen. Regeln Sie dann mit den Hoch und Runter Pfeiltasten die Geschwindigkeit. Klicken Sie Enter, um die jeweilige Geschwindigkeit festzulegen.Geschwindigkeit1.0xKlicke auf diesen Knopf und diess Video stumm oder nicht stumm zu schalten, oder drücke die HOCH oder RUNTER Knöpfe um die Lautstärke zu erhöhen oder abzusenken.Maximal Lautstärke.Videotranskript
    Beginn des Transkriptes, zum Ende springen.
    VIJAY JANAPA REDDI: In this video, we're going to learn about data drift.
    Once again, this is in the domain of a data scientist.
    The change in statistical features of the target
    variable that the model is trying to predict,
    or that it's trying to forecast over time-- that is concept drift.
    Now, we're going to talk about data drift.
    In data drift, it's the input distribution.
    For example, the input feature vectors that we
    tend to receive will have different distribution
    characteristics inside them, and that is what data drift is really all about.
    For example, down here going back to this figure
    that I introduced previously in concept drift, on the extreme right
    you can see that the concept is still the same.
    In other words, that black line is still the same as the leftmost figure.
    However, the data, internally, has actually shifted.
    So this is a place, for example, your model
    might actually be predicting fine, but your input data distribution actually
    looks radically different.
    Now, there are two ways in which this can actually manifest.
    You might have a model that's predicting just fine for the original training
    data, but over time, there might be new customers
    and new users whose characteristics or features are actually quite different.
    And in that particular case, you will see that your concept might actually
    end up changing a bit.
    And so you have to understand both of these characteristics together.
    This type of drift affects many applications.
    Let's take, for example, Gboard.
    Gboard is a virtual keyboard for Android from Google,
    and it predicts or anticipates what you will
    be typing next based on all the words you've entered previously.
    In this instance of predicting the next word, for instance, your idea may drift
    or the concept might actually drift as your language evolves,
    or as you use new phrases.
    Or maybe, for instance, there are seasonal drifts.
    For example, you might be typing in a lot of "Happy New Year" or "Merry
    Christmas" and so forth during holidays--
    "Happy Hanukkah" and so forth.
    Now, this drift might naturally lead to making incorrect predictions,
    and it might deteriorate the application's quality over time.
    And that might be really annoying to you,
    because you're not able to use the keyboard well.
    Well, fortunately in this particular instance,
    it's actually quite easy to kind of backtrack things a little bit,
    because users often provide feedback by actually correcting what
    the Gboard is actually going to guess.
    So that becomes a very useful mechanism of doing continuous evaluation
    to actually get the data.
    More on continuous evaluation later on, though.
    All right.
    Now, let's say you've discovered that your data might actually be erratic,
    based on what you're finding.
    Your next question might be, what can I do about this?
    Well first and foremost, you need a mechanism in place
    to periodically monitor all of this data,
    in order to actually detect the drift.
    Now, this might be, yet again, in terms of the data drift
    or in terms of the concept process.
    So I'm going to touch a little bit using a couple of examples
    as to how you might want to go about monitoring data,
    because just because something looks different
    does not mean you have to retrigger your training pipeline.
    So, how do we actually, first and foremost, go and analyze the data?
    Is the key thing that you actually have to sit down and think about.
    Well, let's take a look at some data distribution
    and try and understand this problem.
    Here, I'm going to be showing you two sets of bars.
    One is what we call as a reference, which is what your model was originally
    trained on.
    And then the blue is what the current state of affairs are right now.
    So for example here, I'm showing you what
    happens when I'm looking at data distribution that's
    coming from different regions.
    For example, in region A and region B, you'll
    see that the current distribution characteristics
    are very similar to the reference distribution on which the model was
    originally trained.
    You're obviously going to see some level of difference
    and that's inevitable, because you've got a snapshot on which you're
    training on it.
    But the real world is a little noisy.
    So a little bit of delta is OK, as long as a general trends hold true.
    However, take a look at what's going on with regional samples with C and D.
    The real world of the current distribution has
    deviated quite dramatically from what the reference looks like.
    Now This should undoubtedly raise an eyebrow for your data scientist.
    Now, before you jump off the cliff and start retraining it
    and say, I've found a drift, and let me go
    retrain my entire pipeline to the new distributions that I've gotten, stop.
    Think about what actually might be causing the drift.
    Here's another example that I'm showing you where, for example, you're
    getting data that's coming in from sensors,
    and it turns out some set of sensors are producing wildly different data.
    Maybe they're just working just fine.
    If you're assuming that they're working fine
    and the data distribution looks like this,
    OK, well, you might just go ahead and retrigger the training pipe.
    And little did you know it turns out that there was actually
    some sort of fault in either the network connectivity that these sensors had--
    so maybe their data was actually not coming back properly--
    or maybe some of the sensors have actually become faulty.
    Or maybe someone replaced a sensor, so this new sensor is not
    producing data at the same rate.
    And so there are many reasons as to why your data quality might actually
    be bad.
    So before you just go jump off the cliff,
    trying to retrain a whole new model based on some distributional shifts
    that you're observing, you need to really double check your data quality.
    Now if you are 100% confident that your source of sensor data, for instance,
    is completely valid, then you would go on to the next step and say, OK, now
    what do I actually do about the fact that I actually have this difference?
    Well, let's begin by taking a close example at these data distributions
    here.
    We'll compare the reference again to whatever the real world production
    data looks like, and we'll try and interpret it.
    Taking a look at these input distributions as a starting point,
    we might want to go one level lower and say, let's try
    and understand not just the general distribution,
    but actually, let's look at the correlation between the features, which
    is really what is actually trying to drive your ML model.
    So if you look at it from that sort of a perspective,
    you might actually uncover some interesting examples.
    So here, for example, you're seeing that certain item has actually changed--
    for example, I'm just arbitrarily picking some sort of videos watched.
    And here, this might actually be a red flag for you
    that something is not quite right.
    This is because you're actually looking one level lower rather than just
    looking at the property distribution of some set of samples
    that you actually have.
    So the meta point is, really, the data scientist
    is going to specialize in a lot of this.
    Again, this is more of a tour of the core concepts.
    But a data scientist would typically go deeper
    into understanding the characteristics of the data,
    and then really try and understand, are the correlations
    that you are actually relying on to make predictions--
    are they fundamentally changing or are they not?
    If they are, then you've got a serious data-related problem on your hand.
    Let's see what happens if there is a genuine drift.
    So there are a whole bunch of different alternatives
    in terms of how you actually deal with it,
    and we'll learn about that in the next video.
    I'll see you there.
    Ende des Transkriptes, zum Anfang springen
    Downloads und Transskripte
    Video
    Videodatei herunterladen
    Transskripte
    Download SubRip (.srt) file
    Download Text (.txt) file

 \section{Forum}
 
 What example of data or concept drift have you seen lately in the real world? What applications that exist today do you think will face challenges with concept and data drift in the future?
 
 \section{Dealing With Drift}
 
     0:00 / 0:00Klicken Sie auf die Pfeiltaste nach oben, um das Geschwindigkeitsmenu aufzurufen. Regeln Sie dann mit den Hoch und Runter Pfeiltasten die Geschwindigkeit. Klicken Sie Enter, um die jeweilige Geschwindigkeit festzulegen.Geschwindigkeit1.0xKlicke auf diesen Knopf und diess Video stumm oder nicht stumm zu schalten, oder drücke die HOCH oder RUNTER Knöpfe um die Lautstärke zu erhöhen oder abzusenken.Maximal Lautstärke.Videotranskript
     Beginn des Transkriptes, zum Ende springen.
     VIJAY JANAPA REDDI: In this video, we're going to go over
     what to do when our model drifts.
     Now this is the kind of situation, for instance, your data engineer
     and your data scientist might actually want to work together,
     because they have to identify new sources of data
     and then try and comprehend what's inside
     that data in terms of the data quality.
     Now when we detect drift, we have a couple of options.
     We can do something as simple as, do nothing,
     or we might actually want to retrigger training pipelines,
     or we might want to rebuild our whole continuous training pipeline.
     There's a whole bunch of different situations you might find yourself in.
     In order to help navigate this space, I've
     listed out five key questions that I would ask myself as to when should
     I do something with respect to training or retraining.
     And we'll step through each one of these.
     The first question is about how rapidly is your model actually going
     to degrade in production.
     Model degradation can actually be fast or it can be slow.
     For example, in manufacturing processes where things often
     tend to be much more stable, the model degradation
     can actually take place over a year, for instance.
     This is a totally hypothetical situation.
     It really depends on what are you doing.
     For example, if you're monitoring some sort of real-time sensor
     that's in your manufacturing plant, in that case,
     you might have to train your model every week.
     For most of the regular processes in manufacturing,
     things don't change that fast.
     So in such a situation, you don't necessarily
     have to rush into doing something just immediately over time.
     So you'd look at the data and you might say, OK,
     what's a very small incremental change that I'm
     going to be able to have an impact on.
     So I'm going to choose to do nothing right now, because that is actually
     a good thing.
     Remember, just because you're going to push a model in production,
     does not mean things won't break.
     So there's some good reason as to why you might want to do nothing.
     However, if let's say, you're dealing with consumer applications or consumer
     use cases, well, they tend to change, like human behavior
     tends to change really fast.
     And it tends to be very costly too.
     So in that case, you might want to retrain
     your model very often in the context of a week or a month and so forth.
     Then you start getting into more nuanced questions around retraining.
     Well, when do I actually want to trigger this?
     What's the frequency with which we want to retrain?
     Because that really is a function of what's
     the frequency with which you actually receive data.
     You might find yourself in rather pickling situations
     where, for instance, you see the model starting to drift,
     but there's really nothing you can do.
     Your hands are tied.
     Why?
     Because you don't have the new data.
     You haven't got new labelers.
     Even though you might have the raw data, you
     haven't yet labeled the data, all of which takes time.
     And remember, when you do data engineering,
     you also got to make sure you just get the data that's been labeled,
     and you got to actually sanity check that the labels are correct and so
     forth.
     It's a very costly and tedious effort.
     It's not instantaneous.
     So you might need to think clearly about what you're going to do.
     There might not be a quick fix on here for you to quickly flip the switch
     and say, OK, well, let's start retraining.
     Because it's a function of whether you actually have data to help you or not.
     So instead of focusing your effort on whether to retrain the model,
     maybe you should then say, OK, well, how quickly
     can I actually get my data labelers to go and get the data.
     Instead of it being one month, is there anything
     I can do to shrink that time down to one week, for instance.
     Also remember, you don't always have to wait for all the data
     to come back in order to update the model.
     Instead, what you might be able to do is collect little bits of data.
     If you figure out what's causing your model to actually suffer,
     then you might identify there are certain key samples
     that you really need, some of which we actually
     talked about doing error analysis in the ML development stage.
     Which is, you might want to create buckets and identify
     which buckets are actually extremely critical
     in order to boost your model's accuracy and be
     able to focus on just collecting that data and rather than just arbitrary
     collecting a whole bunch of data.
     Because most of that data might not actually help you deal with the drift.
     So figure that out, and then pull that back in, and then
     say, OK, let me supplement that.
     In that way, you might actually be able to update the model,
     or you might actually have an ensemble of models.
     An ensemble means having more than one model doing prediction serving.
     And that can be quite useful.
     Because in some corner cases, you would have
     to route the traffic around and say, OK, well,
     these particular regions need this sort of a model to run,
     because this supplements the inaccuracies
     that the bigger model's seen, because it hasn't been really retrained.
     You might want to do this because retraining
     that bigger model might be very costly.
     Last but not least, you might also choose
     to want to use traditional rule-based systems instead
     of just machine learning-based systems.
     If you don't have all the data in order to train the machine learning model,
     which even in the most simplest cases might actually need a fair bit of data,
     you might then say, OK, well traditional rule-based, if, then, and so forth.
     If you use that sort of a traditional rule-based model,
     then you can actually take that and supplement your machine learning model
     with that kind of system alongside.
     These are things production companies actually do put into practice.
     And they work.
     Ultimately though, the thing that you're really struggling with
     is a trade-off of how quickly the model is going to decay
     and how aggressively it's going to decay across a various set of sample
     distributions, versus how quickly can we actually collect the data.
     And that's a key thing you have to figure out.
     And there's no silver bullet clear answer
     I can give you here, because it really depends on the situation.
     The third question is, how much data is really required for the model
     to be successful in the face of drift.
     Our ultimate goal given any neural network model
     is that you want to achieve a certain desired quality target.
     Well, that leads you into two situations.
     Now on the left-hand side, for instance, I'm
     showing you a model that's able to reach a certain target quality,
     but it might get saturated.
     In other words, even though you might have a lot of data,
     your network's capacity in terms of its ability to learn
     might be quite limited just because the architecture
     you have actually put in place.
     So in that case, you might have to go back and really rethink about, well,
     should I spend time getting more data, or should I actually
     be trying to push the network design a bit further so it actually
     has the capacity to learn.
     On the right-hand side perhaps, you have something
     like GPT-3 crazy that is arbitrarily big and can ingest huge volumes of data.
     And maybe in that case, it makes a lot of sense
     to continue throwing more data.
     So you have to ask yourself these kinds of questions in order
     to understand whether simply retraining a given network that you have
     in hand is going to make sense or not.
     Now how long should you keep old data?
     Just because you have data doesn't mean you ingest more data that's
     just arbitrarily coming in.
     What you really want is high quality data.
     Assuming the situation here, for example,
     the blue corresponds to some sort of old training data that you have.
     And then, you're using a model that's predicting
     based on new things that are coming in.
     So we're displaying a couple of buckets down here.
     And you find that during the first three rounds,
     for instance, the model performs really well.
     However, in the fourth bucket or time slot,
     you see that the model's accuracy goes down to 0.75.
     Now in that case, what you might want to do
     is, you might want to retrain the model after every two entries of new data
     that comes in.
     And then your model's accuracy goes up.
     But this leads into a variety of different situations,
     where it's not that easy to figure out.
     Because for example, sometimes when you drop data,
     it means that your model's accuracy will drop,
     because there are certain key characteristics in that data
     set that you're going to lose.
     So that means you have to keep all the data you have.
     You can't throw any arbitrary data out.
     You might be asking yourself, why would I want to throw some data out?
     Well, let's look at the situation for instance, when you throw some data out,
     it actually ends up boosting the data quality
     or boosting the prediction capability.
     Now a perfect situation where this might happen
     would be in the context of anomaly detection.
     Where, for example, if you train a model on a brand new engine
     that you just put in, well, the engine tends to wear.
     Wear is a very common thing that happens with any mechanical system.
     And when it wears, its behavior, as in its vibrations,
     its sounds, and so forth will change.
     And it'll settle into some sort of a cadence.
     And that cadence is where you really want the model to be trained on.
     Because that behavior is going to be different from the original brand
     new engine.
     And so, what that means is that if you're just collecting data,
     you actually want data that's more relevant.
     So you want the more steady state data rather than the brand new engine data.
     So in that kind of case, just having the raw lifetime
     data might actually not make any sense.
     So throwing away some of the old data might
     mean that your anomaly detector might be able to perform well
     on the more recent and relevant data.
     So that's an example.
     Now of course, there might be situations, for instance,
     like in the second row where you drop some data
     and really doesn't have an effect.
     But hey, in that case, it's OK maybe.
     If anything, it only ends up helping you having a shorter training time.
     So that can be a good thing.
     Now when should you retrain your model is the ultimate fourth question.
     You should really give this some sort of thought,
     because the answer is not so obvious.
     Simply because you have new data does not
     mean you should retrigger the model.
     You have to actually understand what the data quality is.
     Remember, training models in the continuous training pipeline
     using AutoML, NAS, and all those good things can actually be really costly.
     So if you bring in some new data, you might realize, for instance,
     that that data doesn't really have any rich attributes inside it.
     So in that case, it doesn't make any sense to retrain the data.
     That data already is effectively inside your original training data set.
     However, maybe the next bucket of data that you get,
     you find that it's got critical label data that was previously
     limited in your original set.
     In that particular case, upgrading the model
     might actually be worth all the effort of retraining it.
     But keep in mind, you also don't want to wait too long.
     Because as I previously mentioned, your continuous training pipe
     actually requires time to converge.
     When you make changes, there's no guarantee
     that your model is automatically going to converge.
     Something might throw your model off.
     At the end of the day, these things--
     I hate to say it-- are a little bit of a black box.
     So there's no guarantee you'll get conversions.
     You might run into all kinds of issues.
     And remember, you've got a lot of red tape.
     You've got to pass smoke testing.
     You've got to pass acceptance testing through the QA pipeline,
     and all of that stuff, which we talked about in the ML training stage.
     So there's a lot of red tape you have to get through after you train the model.
     So these things can really slow you down.
     So as you can see, there's a lot of cognitive thought
     process involved here.
     And we're just kind of scratching the surface.
     But hopefully, this is giving you an idea
     about the kind of issues that you need to make
     sure you're thinking about actively.
     Hopefully your data engineers and your data scientists
     are worrying about this.
     Now so far, I've only touched on the notion
     of either do nothing, in the manufacturing example, for example,
     because the decay is like, it takes a very long time.
     And you don't have to do everything on a weekly basis.
     Or you might want to retrain in the case of a consumer example.
     But there are other things that you might also have to do.
     For example, simply reusing the exact same network might not be enough.
     You might actually have to change the feature engineering that you actually
     have put in place.
     Because perhaps the data is indicating that the feature or the correlations
     between the features have changed.
     Which means now you actually have to spend much more time
     re-engineering the entire pipe.
     That continuous training pipe has actually changed.
     So that could be an invasive thing but needed sometimes.
     But remember, when you're pushing models out into production,
     just because you flip the switch doesn't guarantee that things
     will work in normal fashion.
     You need to have fallback mechanisms.
     You remember, we talked about counter redeployment, shatter deployments,
     blue green deployments, and so forth.
     These are things that we talked about in the production deployment stage.
     And it's actually quite key to have some sort of mechanism.
     Because you might push the new model, and something might actually break.
     So you need it have a mechanism to reverse things back.
     Finally, a crucial thing to keep in mind is just look at the data itself.
     Do very careful data analysis.
     I've kind of alluded to this earlier.
     In fact, in the model development stage, I
     talked about having different buckets for error analysis.
     Because sometimes when you do bucketing like that,
     and you prioritize certain types of errors that you want to fix better,
     that means you segment your data up, then you
     can actually end up with a much better model.
     Finally, as you can see, there's a lot of stuff that actually has to happen.
     So you need to have a lot of adaptive components inside your framework
     to be able to pull all this together.
     With all that said, hopefully you're starting
     to appreciate there are many critical moving pieces,
     and continuous monitoring is an extremely critical block
     in order to make sure that you have a successful AI project.
     Ende des Transkriptes, zum Anfang springen
     Downloads und Transskripte
     Video
     Videodatei herunterladen
     Transskripte
     Download SubRip (.srt) file
     Download Text (.txt) file
 
\section{Continuous Evaluation Challenges for TinyML}

    0:00 / 0:00Klicken Sie auf die Pfeiltaste nach oben, um das Geschwindigkeitsmenu aufzurufen. Regeln Sie dann mit den Hoch und Runter Pfeiltasten die Geschwindigkeit. Klicken Sie Enter, um die jeweilige Geschwindigkeit festzulegen.Geschwindigkeit1.0xKlicke auf diesen Knopf und diess Video stumm oder nicht stumm zu schalten, oder drücke die HOCH oder RUNTER Knöpfe um die Lautstärke zu erhöhen oder abzusenken.Maximal Lautstärke.Videotranskript
    Beginn des Transkriptes, zum Ende springen.
    VIJAY JANAPA REDDI: Welcome back.
    Hopefully you're convinced that modeling for model drift
    is extremely important and crucial.
    To that end, today I'm going to talk about continuous evaluation.
    Continuous evaluation is the process by which you effectively
    monitor the model's input and its predictions
    as it's going about doing them in the real world.
    And monitoring this data is extremely key,
    because that is usually how you detect drifts.
    Now, with traditional big ML systems, you
    don't tend to worry about continuous evaluation as much,
    as in you don't worry about the supporting infrastructure that
    is needed to realize continuous evaluation.
    You tend to assume that you're going to have network connectivity because it's
    plugged into a jack.
    It's got wall power, so you don't have to worry about battery and lifetime
    and so forth.
    However, in the context of TinyML, these things become extremely crucial
    in realizing continuous evaluation, because if you don't have continuous
    evaluation, you cannot realize whether you have model drift.
    OK?
    So you have to be very careful about this,
    and that's why, specifically, I'm focusing on continuous evaluation
    through the lens of TinyML challenges.
    And this is the kind of thing that your software engineer or your ML engineer
    is going to get involved, because there's a lot of infrastructure that
    has to be put in place in order to make sure you can actually
    monitor these tiny devices that are out in the wild.
    It now goes without saying that monitoring is critical, right?
    However, with TinyML, you have problems like, for instance, you
    don't know whether your device will have connectivity or not.
    You can't take it for granted that you will always have connectivity.
    You're also obviously super worried about battery life times and so forth.
    So the question that we have to answer is,
    how can we provide continuous monitoring in order
    to enable continuous training without transferring huge amount of data
    from TinyML devices in a very inefficient way?
    Well, the reason you have to think about this, again,
    goes back to this point about computation versus communication.
    As the graphic on the right shows, data movement
    is extremely expensive with respect to processing.
    This is just within the context of a single system,
    but when you scale it out across multiple systems
    and you're communicating data from one system to another,
    it's orders of magnitude much more costly.
    So you really have to worry about this, because TinyML devices
    don't have much energy to spare.
    So this brings us to the topic of data movement for continuous monitoring.
    If you want to monitor data, then you have
    to be cautious about how you communicate with the TinyML device
    or how the TinyML device communicates with the server.
    So here, I'm showing you bandwidth versus communication range.
    Now, you might be familiar with some of these technologies,
    like Wi-Fi and Bluetooth, which are widely used.
    However, you'll see that there's not a lot of range down here.
    Now, that can be quite limiting when it comes to TinyML deployments,
    though you have a lot of bandwidth down here.
    And bandwidth, often, you might not need a TinyML, really.
    You don't need very high bandwidth, because they're just
    filtering out and trying to figure out what's the most important data
    to actually communicate, for instance.
    Then, of course, you have cellular technologies like 4G and 5G,
    which are both good for range and bandwidth.
    However, there's a huge cost in terms of power consumption.
    And this is ideal for gadgets like smartphones, which can be connected,
    and transmitting a lot of data.
    However, they can be charged on a nightly basis, which is
    a convenience that we take for granted.
    And this is really where you need a new sort of technology that
    is going to be the backbone for continuous monitoring for TinyML
    devices, right?
    This is where low-power wide area networks,
    which were originally created for IoT devices, become extremely crucial.
    Now, I want to talk about this, because I want
    to make sure you're aware about this.
    Because your communication technology is going
    to really be your limiting factor when it
    comes to continuous evaluation and continuous monitoring.
    Therefore it's useful to have an appreciation for what is LPWAN,
    and specifically I'll talk about LORA.
    At a high level, here are some comparative factors
    across all three kinds of technologies.
    The green is good, as I'm showing you for LPWAN.
    What we can see here is that we have both a good amount of range--
    because TinyML devices can be put in remote places--
    cost is often low, battery requirements or energy requirements,
    especially in TinyML cases, is going to be a real limiting factor.
    We need to have devices that can actually communicate,
    but at the same time, despite the communication,
    be able to last for one year, three years, five years, 10 years,
    and so forth.
    These are real industry requirements.
    So we need to make sure we pick the right communication substrate
    in order to enable continuous evaluation.
    Now, I'm going to pick on LORA as an example--
    I'm talking about LORAWAN--
    but there are other kinds of technologies
    like SigFox and so forth that also follow
    similar properties which are all good.
    So I'm going to pick one, just to kind of give you some examples.
    Pause for a second on this video right here and take a look at this slide.
    It's got a lot of description about the rich capabilities
    that LORA offers as a technology, in terms of range, bandwidth, energy
    efficiency, and so forth.
    I want you to learn about this because I'm
    going to have a little quiz that specifically focuses on this slide.
    With that said, let's look at some example use cases.
    There are many of them.
    It can, for example, be used in vending machines to notify distributors
    when a product requires maintenance.
    Or, for example, when a certain item is running low,
    they can easily be alerted so they can actually predictively update
    the stock inside the vending machine.
    Pet owners can track their pets' movements,
    or people who are passionate about understanding migratory patterns
    can look over how birds, for instance, are moving over great distances, which
    is an excellent application of TinyML.
    Another example would be when household oil tanks are running low
    and oil companies can automatically receive alerts.
    Cargo containers and trucks, ships, and trains
    can be tracked for logistical purposes, which improves
    the efficiency of the overall system.
    All these are great examples of LORA, because of the low-power capabilities,
    and therefore, you find that it has a wide range of applications
    in industrial settings, agricultural, wildlife conservation, and smart cities
    and so forth.
    And the list goes on, as you can see down here.
    It's a very powerful piece of underlying communication technology.
    Now, the overall architecture is going to look something like this
    for our TinyML deployment with LORAWAN.
    Again, I'm talking about the big picture here,
    because I just want you to understand what's
    going to be needed in order to enable continuous monitoring
    in the field for TinyML.
    By and large, LORA has three core types of devices--
    Class A devices, Class B devices, and Class C devices.
    And they have different capabilities in terms of their communication
    capability, as well in terms of the battery or the energy span.
    Now, as you can see here, Class A device is the basic barebones endpoint device.
    That's the one that actually has the TinyML sort of attached to it.
    Now, Class A basically communicates and initiates
    the communication whenever it needs to send some sort of data up.
    And they're low power, and this is sort of a device that
    can be strapped on top of whatever you want to transmit back up.
    Now, Class B devices are more like beacons.
    So these endpoint devices typically talk to them-- beacons-- and these Class B
    devices have much more systematic or expected
    regular scheduled windows, in terms of communication
    with some sort of network.
    And they often have capabilities that sort of subsume
    what Class A devices provide.
    Class A are the smallest ones, and the Class B devices can relay that data up.
    And then finally, you have Class C devices.
    As you can see, Class C devices tend to have the lowest latency amongst
    all the different operating modes.
    But they tend to have fairly limited battery life,
    because they tend to consume a fair bit of energy,
    and often, Class C devices tend to be plugged in the wall.
    But by and large, the whole hierarchy will look like this.
    You have these endpoint devices--
    these Class A devices-- which are going to be your TinyML devices,
    and they really should be using something like LORA or LORAWAN
    to be able to communicate with endpoint gateways.
    Where a gateway, for instance, will be able to support something like 10,000
    Class A endpoint devices, and it can maybe
    transmit 10 messages per day per device.
    And so if you look at it like from a networking sort of a standpoint,
    you might be able to support 10,000 or 100,000 devices, or one million
    messages if you have 10 such gateways.
    Now, this infrastructure is something that already exists.
    So it's not so much that you are going to have to go and create these gateways
    and so forth.
    Nonetheless, you have to make sure that wherever
    you're putting in TinyML devices, if you want to do continuous monitoring,
    that they have the Class A sort of sensor communication device on top
    of them, so they can actually communicate with the endpoint--
    these gateways-- and then pass the data up to the servers.
    Again, I'm just touching things at a high level,
    because I want you to have an appreciation for how this thing all
    fits together.
    So let's take, for example, a smart device, and assume it's
    hooked up to some kind of smoke detector which
    is going to communicate and say, hey, if there's smoke,
    then the application needs to know at the other extreme.
    In such a situation, for example, you can
    see there's a heartbeat that's continuously
    going up from the smoke detector.
    It periodically sends the heartbeats up the network server
    to let it know that everything is working.
    The device is still operational and so forth.
    You'd need that just to make sure that you know your smoke alarm's actually
    working.
    And if you want to raise an alert, like the Class
    A device that's at the endpoint will send up the data.
    And you can see that in the yellow lines that I'm showing you here.
    All of this data is actually encrypted as it is going by,
    and that's another critical thing that you
    have to worry about for continuous monitoring,
    because you'll be sending up some of your model data-- your input
    and your prediction data-- all of that will be going up.
    So you don't want people to mistakenly be accessing all of that.
    And in certain situations-- for example, if the alert
    needs to be acknowledged-- in that case, the server
    can actually communicate back all the way to the device.
    So it's not that the device is only doing uplinks to the server.
    You can also actually have downlinks, which is useful.
    Moreover, what's actually important here is
    that you actually are able to time multiple of these devices--
    Class A devices-- together with some sort of gateway,
    and that's extremely critical.
    So once again, the point of all of this is
    to know that when you're talking about embedded inference,
    you've got computation versus communication.
    And when you're doing serving and you're trying
    to understand how serving is actually happening from a continuous monitoring
    standpoint, you need to have the right infrastructure in place.
    And that's really what I want you to take away.
    So make sure you do the readings that are coming up
    next, which tell a little bit about other kinds of communication mechanisms
    that are out there, so you have an appreciation for picking
    the right communication substrate, so you can actually continuously monitor
    the TinyML device in the field.
    Ende des Transkriptes, zum Anfang springen
    Downloads und Transskripte
    Video
    Videodatei herunterladen
    Transskripte
    Download SubRip (.srt) file
    Download Text (.txt) file 
  
\section{TinyML Communication Challenges & Technologies for Continuous Monitoring
    
    \GRAPHICSC{1.0}{1.0}{TinyML/4 MLOps/paste.png}   
    
    \subsection{About This Reading}
    
    Clearly, it is obvious that we have to monitor for model drift to ensure that the data distributions do not shift from what we expect when the ML model is serving embedded inference in the field. But to ensure that the data distribution has not changed, we need to monitor the TinyML devices that are in the field. This is where it becomes crucial for us to understand the pros and cons of different communication methodologies we have access to for TinyML devices. Continuously monitoring for data is very expensive because it drains the battery life extremely fast. In practice, nearly 50\% of TinyML device budget can sometimes be allocated just for communication purposes over the device’s projected lifetime. Given that we allocate so much energy resource for communication, it is crucial to understand the different technologies that exist, so that we can make informed decisions when we want to continuously monitor.
    
    \subsection{Overview of Communication Technologies}
    
    Over the past decade, we have seen massive improvements in the capabilities of wireless technologies. Twenty years ago, the prospect of millions of people transmitting photos, streaming videos online, and sending text messages over-the-air simultaneously would have been unthinkable. Nowadays, this situation is the norm, all enabled by increasingly sophisticated communication technologies. Given that intelligent IoT systems, like those leveraging TinyML, are network-enabled, it is good for us to get a general understanding of what communication technologies exist, how they work, and how they differ from one another.
    
    Communication technologies are generally compared using certain metrics. Most commonly these are (1) functionality, (2) data speed, and (3) operating range. Depending on our system requirements,  these metrics can help us to determine which communication technology is most appropriate for our given application. For example, utilizing a long range and high throughput communications protocol to send small amounts of data over a relatively short range (say, 10 meters) would be unnecessary, and likely consume other resources, like power, unnecessarily.
    
    \subsection{Peer-to-peer Communication}
    
    The first important type of wireless technology to discuss is peer-to-peer communication. This is a form of direct communication - for example, connecting wireless headphones to your smartphone, or using AirDrop to send files on Apple devices. The most prevalent protocol in this space is Bluetooth Classic. Bluetooth has a relatively short operating range, typically around 10 m, which also means that it is a low-power protocol in comparison to other technologies. More recently, we have seen Bluetooth Low-Energy (BLE) and Zigbee become popular alternatives, which are even more power efficient technologies than Bluetooth Classic. One of the disadvantages of these technologies is their relatively low data speeds. Thus, Bluetooth is advantageous for short-range communications with low data transmission rates.
    
    The next most prevalent peer-to-peer technology is WiFi Direct. This is not quite the same as WiFi, since it does not require an access point, but uses the same frequency (2.4 or 5 GHz) for communication between devices and has comparable speed and throughput. The major advantage of WiFi Direct over Bluetooth is that data speeds are around 100x faster, but this comes at the cost of increased power consumption.
    
    Another peer-to-peer communication technology that is becoming increasingly commonplace is Near Field Communication (NFC). This is the technology that is being used for contactless payment systems - every time you buy a coffee by tapping your credit card against the screen instead of inserting the card, or paying for something with ApplePay, a smartwatch, or similar item. NFC communicates by using an electromagnetic field generated by two coupled coils, instead of emitting radio waves like in other forms of wireless technologies. The main benefit of NFC is that it can be entirely passive, requiring zero power. This is why it can be utilized by credit cards even though we never need to charge our credit cards. Wireless charging mats also utilize NFC to charge your devices. One of the caveats of NFC is that it only works over extremely short ranges, on the order of a few centimeters, and so is not suitable for even most short-range applications.
    
    \subsection{Mesh Technologies}
    
    Moving on from peer-to-peer communications, which are one-to-one networks, we now look at mesh networks, which are many-to-many networks. There are two subsets of mesh technologies that we will look at. Firstly, technologies that offer low-power, short-range, and low data rates, and secondly, those offering low-power, long range communications.
    
    \subsection{Low Power + Short Range Mesh Technologies}
    
    There are four that fall into this range: BLE, Zigbee, Z-Wave, and 6LoWPAN. These are ideal for battery-powered devices that need to transmit information over a relatively short range.
    
    One of the major benefits of mesh technologies is that they allow communication through other nodes in the network, which essentially act like middlemen. Instead of two nodes, A and B, needing to be directly connected in a peer-to-peer network to communicate with each other, in a mesh network, these can communicate between a third node, C. Thus, with a large enough network, data could still be sent over relatively long distances using a short-range and low-power communication technology. 
    
    \textbf{BLE} is the most common of these communication technologies, and is what was used on the Arduino Nano BLE Sense used in our prior courses. This allows data speeds up to 1 Mbps (Bluetooth Classic provides 2-3 Mbps) and mesh sizes for up to 32,767 devices. The operating distance for BLE is somewhere between 50-100 m, and so this could allow for meshes that span many miles if configured properly. Unless you have a good reason to use a different communication technology, BLE is often the best choice. It is the most widely supported of the technologies, easy to set up and configure, and consumes very little power.
    
    \textbf{Zigbee} offers similar capabilities to BLE and works at the same frequency (2.4 GHz). However, one thing that makes it stand out is that it can support mesh sizes up to a staggering 65,000 devices. Zigbee is typically used for home automation technologies, and is thus widely supported in this application area for things like smart lighting, thermostats, energy monitoring, security systems, doorbells, etc.
    
    \textbf{Z-Wave} is a competing technology for Zigbee and BLE. The main thing that stands out about Z-Wave is its lower frequency of 908 MHz for communication. A lower frequency means longer operating range, and also since it does not use the saturated 2.4 GHz band that many devices, such as microwaves, tend to use, there is significantly less interference. The disadvantage of this technology is the reduced data transmission rate as a result of the lower communication frequency. Also, the maximum supported mesh size of Z-Wave is only 232 devices, which is significantly less than BLE or Zigbee, but probably still sufficient for most applications.
    
    \textbf{6LoWPAN} is a competitor for Zigbee in the space of home automation, and is differentiated by its use of an IP-based network similar to WiFi.
    
    \subsection{Low power + long distance technologies}
    
    Many important IoT applications require low-power and long-distance communications, and the same is likely to be true for TinyML. Often, the type of network used for this application is called a Low Power Wide Area Network (LPWAN). This type of network is desirable in situations where data is being collected over a wide area and needs to be transmitted to a location for uploading to the cloud. The most promising wireless technologies for this application are LoRa, NB-IoT, and LTE-M.
    
    \textbf{LoRa/LoRaWAN} stands for long range wireless area network, and allows for communication for distances up to 6 miles while still consuming relatively little power. This wireless technology was developed by Semtech in 2012 and has become popular in recent years since the rise of the IoT device. LoRa uses sub-GHz range frequencies for transmission, meaning it has relatively low data rates but a long operating range and reduced interference. Since LoRa is not a cellular technology, it cannot communicate with mobile networks, making it less complex, cheaper, and simpler to set up. A LoRa gateway device will be required at the end of the network if the data is to be submitted to the cloud for data storage and processing.
    
    \textbf{NB-IoT} is another popular technology in this space, and is a cellular technology. This makes it more complex to set up, more expensive, and also more power-consuming. However, it provides direct access to the internet and also higher quality cellular connections. Currently, NB-IoT is only supported in Europe and not the U.S., and is only capable of transmitting very small amounts of data.
    
    \textbf{LTE-M} covers the applications that are not suitable for LoRa or NB-IoT, being specifically designed for long-distance cellular access that also require high data rates. It is significantly less power-consuming than the standard LTE used in smartphones and is also cheaper to implement due to the reduced bandwidth.
    
    \subsection{Which Should I Use?}
    
    The decision of which communication protocol is suitable for a given project is dependent on the context of the problem. Will you be using large numbers of devices that will be far apart and require cellular access? Or perhaps only a few devices that will each be placed in a single room in your house? Once you understand your specific use case, you can estimate what parameters you might need and then select the most suitable wireless technology based on those estimates.  
        
\section{On-device Training: Limitations and Opportunities}


\subsection{About This Reading}

    Machine learning models are great when they work, but once deployed their performance can often degrade over time due to the phenomenon of both concept drift and data drift. One of the ways to tackle these issues and ensure models continue to remain high-performing is to update them with new data. To do this, we have to continuously monitor the ML model deployments, which in our case can be costly as communication is costly in terms of energy consumption. An alternative way to keep the models updated is through on-device training.
    
    \subsection{What Is On-Device Training?}
    
    On-device training refers to the training of a machine learning algorithm on the device where it will be deployed. That is, the device was not trained on an external device (like a laptop or desktop computer) and exported to the deployment device, which is what we have done in all of our applications thus far. There are additional challenges that come with on-device training that make it far from simple to implement, especially on the resource-constrained devices we work with in the field of TinyML.
    
    Smartphones and other devices have been able to perform on-device training for a while. However, when it comes to resource-constrained devices like microcontrollers, this is considerably more challenging.
    
    \subsection{Tiny In-Device Training Challenges}
    
    Firstly, in order to train an algorithm, we need data to be able to achieve this. Unfortunately, this also means we need a place to store that data, which can be difficult if only a few kB of memory is available for data storage. This would ultimately limit how much data could be stored and would dictate how often the algorithm would need updating.
    
    Secondly, training the algorithm would also be difficult due to the lack of numerical precision available on most resource-constrained devices. Propagating gradients in the computational graph of a neural network requires high precision because these values are usually quite small. The 8-bit precision of a microcontroller is not really sufficient for this procedure and will result in vanishing or exploding gradients due to the limited numerical precision available.
    
    Another important issue is the ability to perform training itself. Training requires additional features in the runtime that are not necessarily required during inference. By incorporating training characteristics, this will add bloat to the library required to generate the runtime, limiting the size of the network that can be deployed alongside the runtime, while still fitting into the same limited memory. Furthermore, the peak memory loads are often higher during training which might also limit the size and performance of the deployed model.
    
    Clearly, on-device training is not currently feasible on some of the more resource-constrained devices. However, this is still an active area of research, and the hope is that one day training will be able to be performed on all manner of devices, regardless of resource constraints. For the time being, we will have to rely on some of the more resource-rich devices that contain sufficient memory, compute power, and numerical precision to perform on-device training.        
        
\section{Continuous Monitoring with Federated ML}

    0:00 / 0:00Klicken Sie auf die Pfeiltaste nach oben, um das Geschwindigkeitsmenu aufzurufen. Regeln Sie dann mit den Hoch und Runter Pfeiltasten die Geschwindigkeit. Klicken Sie Enter, um die jeweilige Geschwindigkeit festzulegen.Geschwindigkeit1.0xKlicke auf diesen Knopf und diess Video stumm oder nicht stumm zu schalten, oder drücke die HOCH oder RUNTER Knöpfe um die Lautstärke zu erhöhen oder abzusenken.Maximal Lautstärke.Videotranskript
    Beginn des Transkriptes, zum Ende springen.
    VIJAY JANAPA REDDI: At this point, I expect
    you to have a good grasp on continuous monitoring, specifically,
    through the lens of continuous evaluation.
    I hope you got to appreciate the communication can become
    a major problem or a major bottleneck for being
    able to do continuous monitoring with TinyML devices.
    Now, in this video, I'm going to talk a little bit about federated machine
    learning, which is a very hot topic, especially in the context of TinyML.
    That's one of the reasons I want to touch on it.
    I also want to touch on it, because it does help us address some
    of the pitfalls around communication.
    Interestingly, it brings about two things.
    One is it kind of helps us with the communication pitfall,
    but also helps us with a new angle when we're
    doing continuous monitoring, which is around data privacy.
    Let's go back to the example of the GBoard.
    If you recall, we previously stated that we
    have to maintain the model that's on GBoard updated.
    Because users might actually be subject it to drift.
    Users, for example, may use new words.
    They might use new languages and so forth.
    And as a result, the model on device needs to be constantly updated.
    Now, in the context of traditional continuous monitoring,
    one could imagine that the contents of GBoard experiences,
    the words and so forth that are being typed,
    are being sent up to the server where the server gathers all this data--
    be it the messages and be it their large number of participants
    that we have across the globe for applications like GBoard--
    and train a global model that can then be distributed back
    to everyone periodically.
    Now, normally, this isn't an issue if you're
    talking about wall-powered devices or cell
    phones that can be quickly recharged.
    And if you wanted to, you could potentially
    undertake continuous monitoring almost every day if it was possible.
    Given that there are billions of Android devices out there,
    it's quite possible to actually realize this at a much higher frequency then
    we would otherwise be able to do if we had only a handful of devices.
    Well, all that said, though, we still have some concerns.
    And then one of the key concerns is around privacy.
    This is a new element that I'm introducing
    in the context of continuous monitoring.
    It's not about just having data that you can send back up to the server.
    It's also about keeping that data secure.
    So, for example, here, you want to monitor and make sure
    that you're not incorrectly transmitting some sensitive data.
    Sensitive data could be like private messages
    that you exchange with your loved ones or some sort of credit card information
    that you might be typing into your GBoard application.
    You want to ideally minimize and prevent gathering some sort of critical data.
    And you want to dispose some of that critical data on a regular basis
    rather than sending that up.
    You also want to employ things like encryption mechanisms
    to safeguard the data that is actually being sent up
    in order to train better models.
    You might want to also make sure you've got the right consent if you are taking
    some amount of data from the user.
    So you've got to think about all these issues.
    In the context of manufacturing, for instance, if you have a sensor,
    sensors can hold very rich data about how well you're doing.
    And that might become a key thing that your competitor might actually
    want to take a look at, because that will tell them
    how to quickly catch up with you.
    So privacy becomes an extremely important part
    when we're talking about continuous monitoring.
    So what can we do?
    Well, one of the key problems that generally we
    have when we're trying to do continuous monitoring and so forth is you're
    looking at raw data that you want to bring up to the server.
    Well, that's where, really, most of the problems come about.
    And that's where federated machine learning
    becomes a really nice alternative.
    What federated machine learning does is that rather
    than having a centralized approach, it takes a decentralized approach
    to model training.
    And in doing so, it maintains privacy inherently.
    In short, in classical machine learning, what we're doing is
    we're bringing all the data up to the code,
    which is federated learning is all about bringing
    the code to the data at the endpoint.
    And what ends up happening is that you end up working and creating
    local model updates.
    And only the updates, not the raw data, are actually periodically sent back up
    to the server.
    So that would be like the gradient changes
    that were made to the networks-- those would be sent out to the server.
    And amazingly, that is enough for the server
    to be able to construct a new model from time to time.
    And that model can then be periodically sent back to the endpoint devices.
    And this way, you are effectively able to keep data locally at the endpoint
    while also facilitating some of the model updates and periodic fashion
    so you're getting the best of both worlds.
    So why is this very useful?
    Well, the nice thing is the fact that it allows for local models
    to be continually trained.
    Because in a way, you're having continuous monitoring
    and continuous training at the endpoint.
    And this way in a way, the endpoint device
    might be able to deal with real-time drifts that might be coming about.
    And at the same time, the model periodically gets all the updates.
    So therefore, your big models can actually
    get updated from time to time as well.
    Another thing is that because you're able to deal with updates locally,
    this means that your latencies are much faster
    to deal with those necessary updates.
    So your continuous monitoring loop is effectively
    running much, much faster here.
    So it's like having two loops-- one that is
    running tightly near the device and one that does a bigger update.
    So this is another way of effectively dealing with drift.
    And, of course, this obviously minimizes the overheads
    in the cloud size for the infrastructure.
    But I would argue that this is more of a subtle point.
    Perhaps, a bigger point is, really, the fact that we're
    able to preserve data privacy.
    Now, there is no free lunch with federated machine learning.
    There's some key things that you just have to be aware about.
    The first is the fact that you might have potentially unbalanced training
    samples.
    Now, the data that is received from different users,
    it's coming about as a result of their local data sets,
    which might not be typical of the general distribution
    that you would normally observe.
    Because in federated learning, you don't get to see all of the data
    at a single snapshot.
    Rather, you're getting little tidbits of data coming up.
    So therefore, you might end up with an unbalanced training data set.
    Another thing to keep in mind is that federated learning relies largely
    on the fact that you have a large number of devices.
    So because deep learning typically needs a lot of data,
    your federated learning inherently is going
    to need a large number of endpoint devices.
    And the amount of data that's generated by all these endpoint devices
    will obviously vary quite significantly.
    So it's quite different from having a central repository where
    you can be guaranteed about the amount of data you're getting.
    Here, you're really subject to how many devices you have.
    In the context of GBoard kind of applications
    where you have billions of users, because there are billions of Android
    users, well, in that case, it's great because you
    have a large number of clients.
    Do you have enough clients for your TinyML devices?
    That's a question only you can answer.
    Third bullet is that you might sometimes have unreliable network connections,
    which can actually cause problems.
    Because you need to be able to get those gradient updates from time to time
    to be able to do a big model update, which is crucial, as you know,
    for model drift.
    To address some of these issues which are actually known,
    the system tends to come up with a federated learning plan
    where the idea is that it tries to capture what the task requirements are
    in the left-hand side from the development side of the world.
    And then it tries to transform that computation graph
    into something that is effective when it's actually
    deployed in the field taking into account resource
    constraints, the scale of your deployment
    as in the number of devices that are needed for being effective at actually
    doing federated learning.
    So similar to how MLOps streamlines the whole process of ML
    creation and development and deployment, federated
    learning with the federated plan is trying
    to have some sort of operational kind of a setup that
    is in place that deals with the development
    side along with the deployment side.
    I'm only mentioning this, because these are
    things that add another level of complexity
    if you're thinking of adopting federated learning into your MLOps pipe.
    But there are concerns.
    Mainly, it falls primarily into two aspects--
    the system issues and the privacy concerns
    associated with federated learning.
    There are still some pitfalls.
    On the system side, you, obviously, have to deal with the computation
    and communication issues.
    On the computation side, we are saying that we're
    going to try and do some local on-device training which
    means you've got to have the right computational capabilities.
    If you want to do that, you've got to make sure you pick the right device.
    And then, of course, apply all the things
    that we learn in model conversion like quantization, pruning,
    [? sparcification, ?] weight clustering and all those kinds
    of cool optimizations.
    On the communication side, you still got to make
    sure you rely on something reliable and something that is actually
    energy aware.
    Because you don't want to, again, use up too much energy communicating.
    So you still have to understand those nuggets that I've talked about earlier.
    Those become the foundation of building blocks
    if you want to use this sort of a high-level level training method.
    There are some privacy problems that are starting to become obvious even
    with federated learning.
    The gradients, as it turns out, end up leaking information
    about the user's training data, which can become
    a big issue for federated learning.
    Now, these model inference attacks, as they are known,
    are a type of attacks where you can recover the client's original training
    data from the incremental changes that are going back up to the server-- so
    yet another thing to keep in mind.
    But it's not a showstopper.
    Because there are ways like doing secure aggregation
    and injecting noise into the actual data that
    allows you to hide some of these issues.
    But these are solutions that are still being heavily-explored.
    Nonetheless, these are things you have to keep in mind.
    Therefore my point is in order to maintain privacy,
    federated machine learning must be carefully designed.
    And when you're designing the systems with it,
    you've got to make sure you've got the nice, right equipment that supports it.
    So this might be an additional level of complexity for MLOps.
    But it might also open up new doors for how you deploy things.
    Now, as you can see, federated learning is still
    a relatively nascent field from the way I'm laying things out to you.
    So this is why, here, I'm highlighting that.
    It might be a ML researcher and the DevOps team
    that I actually get involved in here in order to try and help you
    deploy out federated learning.
    With that said, I'm excited to tell you that we're
    close to wrapping up this section.
    I'll see you in understanding the impact of continuous monitoring
    on the MLOps pipeline.
    Ende des Transkriptes, zum Anfang springen
    Downloads und Transskripte
    Video
    Videodatei herunterladen
    Transskripte
    Download SubRip (.srt) file
    Download Text (.txt) file        
   
\section{Federated Learning GBoard}


    \subsection{Overview}
    
    \textbf{}ederated learning} is a relatively new technology allowing the machine learning models deployed on edge devices to be updated using locally available data. Incredibly, this can be achieved without communicating the local data to a centralized server, making it a privacy-preserving form of distributed machine learning. In this reading, we will review how federated learning works, and how it is used by Google to update GBoard, the phone keyboard for their Android mobile operating system.
    
    \subsection{Recap of Federated Learning}
    
    Instead of working with individual data samples, federated learning works with client devices. That is, a randomly selected sample used in traditional machine learning will instead be cast as a randomly selected client containing N training data samples. Then, instead of sending the data that is locally stored, the weights and biases of the local model can be sent instead. Let’s walk through the procedure step-by-step.
    
    First, the global model, stored in a centralized repository, is broadcast to all devices (or those that meet certain criteria, such as operating system or country of origin). Then, the model is locally trained on newly acquired data over time - essentially, continuous training. However, since there will likely not be enough information on a single device to optimize performance, federated learning allows us to crowdsource data from other users’ devices to help improve the performance of our model. Our local model’s weights and biases are then communicated to the centralized server, which aggregates these localized weights with those from other users’ devices (such as by averaging their values). The updated model can then be rebroadcast to devices, and will hopefully demonstrate superior performance to the prior model.
    
    \subsection{Pros \& Cons of Federated Learning}
    
    There are pros and cons to the federated learning approach. Federated learning is inherently privacy-preserving since no raw user data is transferred. Furthermore, it exhibits less communication overhead, since weights and biases are generally smaller in size than the corresponding data, which would also then take up space in the centralized server. However,  the performance of a federated model is typically lower compared to one built on the full dataset (which can be a larger model which is subsequently compressed using model distillation to fit on the end-user device), which might make it unsuitable for some applications.
    
    The aggregation procedure is one of the most important aspects of federated learning. There are two common aggregation procedures: fedSGD and fedAVG. When using fedAVG, the procedure is as described above, where our model is trained multiple times using local data, and the weights and biases of our network are updated. These updated weights and biases are then sent to the centralized server and then averaged across available client devices. In fedSGD, the model weights are not necessarily updated, but the gradients for updating are determined based on local data. These gradients are then sent to the centralized server, which are then averaged and used for updating. Using fedSGD provides superior performance since it guarantees convergence, but has a higher communication overhead since the gradients must be sent at every iteration. In comparison, fedAVG has a fairly low communication cost but cannot guarantee convergence.
    
    As you may be thinking, these aggregation procedures are quite naive implementations. This area is still an active research area, and so new methods are being developed continuously. For the interested reader, some promising methods worth looking into are federated matched averaging (FedMA) and probabilistic federated neural matching (PFNM).
    
    \subsection{Google’s GBoard for Android and iOS devices}
    
    A great example of federated learning on edge devices is Google’s GBoard. GBoard is the keyboard used in a number of Android and iOS devices. The system works by recording what the user writes (i.e., the context), what recommendations were shown, and whether the user clicked the suggestion.
    
\GRAPHICSC{1.0}{1.0}{TinyML/4 MLOps/Gboard.gif}

    A gif showing that if you click the google “G” on an Android phone you can search for things like “umami burger” which will bring up the address that you can text to a friend. This sort of interaction is saved to the cloud to help predict what you want to type and search in the future.
    
    Since the application needed to provide training updates across millions of devices, Google developed their own federated learning algorithm to minimize communication overhead. Updating was also only done on devices that were charging and connected to WiFi. They also developed a secure aggregation protocol to further improve privacy, ensuring that no user’s average update can be inspected before averaging takes place.
    
    \subsection{Back To TinyML}
    
    While GBoard does not explicitly use TinyML, it is a good example of how federated learning can be applied to TinyML applications. Most resource-constrained devices that fall under the purview of TinyML do not have sufficient resources to train a model from scratch but do have access to small amounts of localized data. Federated learning provides an alternative route wherein local TinyML models could be updated by sending gradient updates to a centralized server and aggregated.
    
    However, there are some additional challenges associated with the implementation of federated learning in TinyML. Most notably, some devices may have insufficient resources to model binary representations which allow gradient updates to be performed accurately. Despite these challenges, this is likely to be an important area of focus given its potential utility,  and is thus well worth knowing about!        
        
\section{Continuous Monitoring Impact on MLOps}

    0:00 / 0:00Klicken Sie auf die Pfeiltaste nach oben, um das Geschwindigkeitsmenu aufzurufen. Regeln Sie dann mit den Hoch und Runter Pfeiltasten die Geschwindigkeit. Klicken Sie Enter, um die jeweilige Geschwindigkeit festzulegen.Geschwindigkeit1.0xKlicke auf diesen Knopf und diess Video stumm oder nicht stumm zu schalten, oder drücke die HOCH oder RUNTER Knöpfe um die Lautstärke zu erhöhen oder abzusenken.Maximal Lautstärke.Videotranskript
    Beginn des Transkriptes, zum Ende springen.
    VIJAY JANAPA REDDI: Congratulations on completing the MLOps Pipeline
    around continuous monitoring.
    I know it's probably a lot for you to process, but don't give up,
    we're almost there.
    We're getting to the finish line.
    In this section, I stressed on three core aspects of continuous monitoring.
    And I stressed on the notion of concept drift,
    then I stressed on the notion of dealing with data drift,
    and analysis that you might want to perform before you decide
    to send up the alert for retraining.
    And finally, I also talked about continuous evaluation,
    and especially in the context of continuous evaluation,
    I emphasized some of the communication problems
    we'll see when we come to TinyML, and how
    we need to be cognizant about that underlying framework that
    is needed to be put in place in order to be able to actually have
    continuous monitoring capabilities.
    I also briefly touched on federated machine
    learning as a potential solution.
    One of the reasons I did that is because more often than not, a lot of people
    are often asking me about, OK, is federated machine
    learning going to be something that is actually useful, and so forth.
    But I wanted to kind of frame some of its opportunities and challenges
    in the context of MLOps pipe.
    So, that said, you learned quite a bit.
    I hope you're happy about that.
    But of course, this brings us to the ultimate question,
    what is the impact of continuous monitoring on the rest of the MLOps
    pipeline?
    It's the million dollar question that I love to ask you every single time.
    All right, let me pick a couple of interesting samples down here.
    I hope, at this point, you see how continuous monitoring very tightly
    ties to continuous training.
    The model monitors that we have in the continuous monitoring stage
    are critical, because they're going to set the retriggers off
    in the continuous training.
    And even when they do set off, we have to be careful as
    to when we want to monitor.
    So we talked a little bit about making sure we do
    want to retrain, because sometimes when we retrain,
    we might get into issues like conversions and so forth.
    So it's not obvious, always, that just because you have drifts,
    you immediately retrain.
    Last but not least, I also hope you saw how federated machine learning is
    actually something we might want to do.
    Now, normally, one would think that I would
    talk about federated machine learning in the continuous training stage,
    because it's part of the training pipe.
    But the reason I kind of held it off to talk
    about it in the context of continuous monitoring
    is because I hope you see how continuous monitoring is
    a little bit of a mixture around communication versus computation,
    and data privacy related aspects.
    So that's why I wanted to kind of hold it out until the end.
    But hopefully, that kind of fits the bigger picture in your head.
    Also, keep in mind that you must have a lot of this infrastructure in place
    at the training operationalization part, even though we're
    talking about it down here.
    Hopefully, it should be obvious that, oh,
    continuous monitoring means I'm going to retrigger those training pipes over,
    and over, and over, so the CI/CD flows that you put in place need
    to be sure that they're actually ready in order
    to handle these sorts of continuous retrainings that are going to happen.
    Now, to push even further back, I hope you
    can see how your ML researcher might need to be cognizant about the fact
    that they need to collect drift information.
    These are topics that are rarely ever discussed
    when people are talking about ML development,
    because they often tend to focus on the model architectures and all
    that good stuff.
    But hopefully, you're starting to appreciate just how much
    is involved in this kind of process.
    The model that you have come up with right at the beginning
    needs to be continuously monitored in the field,
    and you had to close this massive loop.
    So I really hope you have a deep appreciation
    for the large number of people that are going to be helping you out.
    Now, when I say a large number of people,
    I don't mean to say that you need an army of tens or hundreds of people
    for this.
    You could very well be a startup with just eight or 10 people,
    but if you've got the right skill set in there,
    then you wouldn't have to worry about it, as long as you know that you've
    got to have all these nuggets.
    But if you don't, then you might find yourself in a pickle.
    And that's why I wanted to focus on the personas that are all involved in here.
    And as you can see, your ML engineer, your ML researcher,
    your data scientist, your data engineer, your software engineer,
    your DevOps team, all of them come into play
    when it comes down to continuous monitoring as part of the MLOps pipe.
    It's extremely critical that you have an appreciation for your colleagues
    on the team for this.
    Now, hang in there, because we're almost to the end of this course.
    And I'll see you in the next video.
    Ende des Transkriptes, zum Anfang springen
    Downloads und Transskripte
    Video
    Videodatei herunterladen
    Transskripte
    Download SubRip (.srt) file
    Download Text (.txt) file        
        
\section{Forum on the Privacy vs Performance Trade Off}        

In this forum we would like you to read and explore a real world case study of Federated Learning from Google. You can find the article here: \url{https://ai.googleblog.com/2017/04/federated-learning-collaborative.html}

\begin{itemize}
  \item How do you feel about this feature? Do you think it is useful? Exciting? 
  \item What are some other products you have that you wish used federated learning to improve themselves? 
  \item Do you think it is a privacy violation to share your model updates to the cloud? Do you think this solves the privacy issue by not sending your raw data and just the model updates?
\end{itemize}
        
\section{Optional: Colab on Federated Learning}        

Now that you’ve explored the theory behind Federated Learning it's time to see it in action through code. In this next Colab you are going to explore training of the MNIST dataset using Federated Learning. To begin, click the link below:

\url{https://colab.research.google.com/github/tinyMLx/colabs/blob/master/5_9_9_FederatedLearning.ipynb}

        
% 1.10: Data & Model Management
        
\section{Data and Model Management}

    0:00 / 0:00Klicken Sie auf die Pfeiltaste nach oben, um das Geschwindigkeitsmenu aufzurufen. Regeln Sie dann mit den Hoch und Runter Pfeiltasten die Geschwindigkeit. Klicken Sie Enter, um die jeweilige Geschwindigkeit festzulegen.Geschwindigkeit1.0xKlicke auf diesen Knopf und diess Video stumm oder nicht stumm zu schalten, oder drücke die HOCH oder RUNTER Knöpfe um die Lautstärke zu erhöhen oder abzusenken.Maximal Lautstärke.Videotranskript
    Beginn des Transkriptes, zum Ende springen.
    LARA SUZUKI: In this lesson, we will cover the last step of the MLOps
    framework for machine-learning systems.
    We will navigate you through how to support auditability, traceability,
    and compliance of your machine-learning systems.
    Data and model management processes are a central, cross-cutting function
    for governing ML artifacts to support auditability, traceability,
    and compliance.
    Data model management can also promote shareability, reusability,
    and discoverability of ML assets.
    The core MLOps capabilities employed here
    are dataset and feature repository, model registry,
    machine-learning metadata, and artifact repository.
    Data model management provides a unified system
    to manage datasets and ML features.
    It creates, maintains, and reuses high-quality data
    and allows us to maintain the same dataset entities from which we use
    in the machine-learning environment.
    We also benefit from the feature management system,
    where data scientists can discover and re-use available feature sets
    for their entities.
    We can also establish a central definition of features
    and avoid training-serving skew by using the feature repository
    as a data source for experimentation, continuous training,
    and online serving.
    We can also serve up-to-date feature values from the feature repository
    and provide a way of defining and sharing new entities and features.
    In data management system, the dataset is
    used for a particular machine-learning task or use case.
    This system maintains a script for creating datasets and splits,
    a single dataset definition, and realization
    within the team to using various implementations and hyperparameters.
    We also have metadata and annotation that
    can be useful for team members who are collaborating
    on the same dataset and tasks.
    We can also provide a reproducibility and lineage and tracking of datasets.
    In a model management system, the data that has been collected and used
    for model training and evaluation is accurate, unbiased,
    and used appropriately without any privacy concern or privacy violations.
    The models are evaluated and validated against effectiveness quality measures
    and fairness indicators.
    That is for us to ensure that they are fit for deployment in the production
    environment.
    The models are interpretable, and their outcomes are explainable, if needed.
    The performance of the deployed models are
    monitored using continuous evaluation, and the model performance metrics
    are also tracked and reported.
    Potential issues in training or prediction
    serving are traceable, debuggable, and reproducible.
    Then a machine-learning metadata tracking system
    is integrated with different MLOps processes.
    Artifacts produced by the other processes
    are usually automatically stored in a machine-learning artifact
    repository, along with information about the process executions.
    It also tracks experimentation parameters and pipeline configuration
    for reproducibility and for tracing lineage.
    It serves for data scientists to search, discover, and export
    existing ML models and artifacts.
    In a model governance system, there's registering, reviewing, validating,
    and approving models for deployment that can be automated, semi-automated,
    or fully manual.
    There's also the reporting on the performance of the deployed models.
    The information [INAUDIBLE] in the ML metadata and model registry
    can be used to store in model registry.
    We can add or update model properties and track model versions and property
    changes.
    Evaluate, we can compare a new challenger model to the champion model
    by looking not only at the evaluation metrics, such as accuracy, precision,
    recall with specificity, and so on, but also at business KPIs.
    Additionally, model owners need to be able to understand and explain
    the model predictions, check, review, request changes, and approve the model
    to help to control for risks, such as business, financial, legal, security,
    privacy, reputational, and ethical concerns.
    Release, we can invoke the model deployment process to go live.
    Report, aggregate, visualize, and highlight model performance metrics
    that are collected for continuous evaluation process.
    These ensure the quality of the modeling production.
    Important topics here are explainability,
    a clear view of the lineage and accountability,
    review decisions in accordance with an organization's ethical and legal
    responsibilities.
    Now, let's take a final look at the entire end-to-end MLOps process.
    Here, we can see the MLOps processes we have covered in this course.
    It includes the eight processes we have covered,
    and note that the processes in gray are the ones that are iterative processes.
    In this final diagram, we illustrate the MLOps
    processes defined on the top of Google Cloud Vertex AI platform.
    All services and capabilities to ensure a complete end-to-end MLOps process
    is supported in Vertex AI.
    Ende des Transkriptes, zum Anfang springen
    Downloads und Transskripte
    Video
    Videodatei herunterladen
    Transskripte
    Download SubRip (.srt) file
    Download Text (.txt) file        
  
\section{verview of Data and Model Management}

\GRAPHICSC{1.0}{1.0}{TinyML/4 MLOps/1000.png}
    
    The MLOps pipeline. Going from ML Development to ML training to Continuous Training to Model Deployment to Prediction Serving to Continuous Monitoring. Along the way artifacts are produced at each step including: Code and Configurations, Training Pipelines, Registered Models, Serving Packages, and Serving Logs. All steps tie into the overarching themes of Data and Model Management.
    
   \subsection{Objective}
    So far, we have discussed the operations of our system in terms of training, deployment, prediction serving, and continuous monitoring. However, we have, until now, left out the most important aspect of this system: the data and models used within our system! And that is where model management, data management and model governance all come into play. In this section, we will briefly understand these three different concepts and discuss their relevance.

\GRAPHICSC{1.0}{1.0}{TinyML/4 MLOps/pasted image 100.png}
    
    A figure shows how Datasets and Feature Management interact with other steps of the MLOps pipeline through model training, monitoring, and serving enginers, as well as through experimentation environments. It also shows that Datasets and Feature Management is made up of many different discrete technologies and steps including inputs of discovery, point-in-time queries, create-update-annotate, and dataset and feaure repositories; both online and batch serving; and both online and batch ingestion through ETL and streaming systems as well as update statistics and metrics steps.
    
    \subsection{Dataset Management}
    
    In dataset management, features describe fundamental entities related to our application, such as products, customers, and locations. While individual features can be used in a large number of datasets, an individual dataset is typically used for a specific model or application. For example, a recommendation system might require information about both customers and products, while a propensity model might only require information about a customer. To produce the datasets used for training, some transformations or operations will typically be performed on the features included within the company’s databases. These sets of operations might be described in scripts, which will need to be saved and monitored for future use in some kind of repository. These scripts might have to be versioned depending on whether they are for training, testing, or deployment. We will also need to know how the data was split, along with any other information about the data in order to ensure we can reproduce the results in the future. All of these activities fall under the purview of dataset management.
    
   \subsection{Model Management}
   
    Keeping track of models becomes increasingly difficult when they are all simultaneously being constructed, tested, and deployed. To prevent models from being misplaced or the wrong models tested or deployed, it is important to have a strong model management system in place. Model management involves tracking metadata related to the model, such as what training data and hyperparameters were used to train the model, as well as model governance. These aspects help to ensure that (1) data being used to train a model is accurate, unbiased, and does not violate licensing or privacy requirements, (2) models are evaluated using metrics, such as those related to quality control or fairness indicators, (3) model outcomes are interpretable and explainable (necessary for GDPR), and (4) that issues during model training or deployment are reproducible, traceable, and can be debugged.
    
    For model metadata and artifact tracking, artifacts related to the machine learning model are captured and stored in a repository. Metadata that might be captured might include pipeline run  ID, start and end time of training, run status, and environment configuration. Artifacts that might be captured include evaluation metrics, hyperparameters, and schemas. Note that model metadata refers to aspects of the model runtime, whereas artifacts describe aspects of the  model itself. Keeping track of artifacts and model metadata is helpful for data scientists and ML engineers to ensure reproducibility and traceability since the parameters used during experimentation as well as the configuration environment are known. These aspects also make searching through large numbers of models for a specific experiment easier.
    
    \GRAPHICSC{1.0}{1.0}{TinyML/4 MLOps/pasted image 101.png}
    
    A figure showing how the ML metadata and artifact repository can hold all sorts of metadata including: hyperparameters, evaluation metrics, statistics, schemas, models, profiling legs, among others. These items come from ML development, continuous training, model deployment, prediction serving, and continuous monitoring, and need to be actioned through discovery, adds/updates, analysis, visualization, and exports.
    
    \textbf{Model governance} dictates what happens to a model during its development. This entails registering (i.e., instantiating within the model registry), reviewing, validating, and approving a model for deployment. The process of model governance is highly context-dependent, and will vary between application, organization, and regulatory environment. Some implementations may require more checks and balances than others, and the same implementation may require changes to the governance model over the course of its lifetime due to changes in the regulatory environment as well as other factors. The process of model governance can itself be manual, semi-automated, or fully automated.
    
        \GRAPHICSC{1.0}{1.0}{TinyML/4 MLOps/pasted image 102.png}
        
    A figure showing how Model Governance can be operationalized. This is through trained models being registered and manifested into a model registry. This passes to a store step for recording of model properties and for official tracking. This passes to an evaluate step to compare new models to champion models and to explain and validate the challenger model. This masses to a check step which interacts wiht the ML metadata & artifact repository to review and approve the model, and then send it to a model serving engine. This passes to a relase step to either deploy or roll back the updates and finally to a report step to visualize and aggregate performance.
    
    Automated model governance typically relies on information from the model metadata repository, as well as the model registry. This allows the governance model to perform various tasks. The first task is to be able to add and update model properties, as well as to track model versions (such as the version used in training compared to that used for deployment). This makes it easy to keep track of models for future reference, for traceability and repeatability purposes. A new “challenger” model can also be incorporated into the model governance system, which can then compare it to the available model and test it across various evaluation metrics, as well as business key performance indicators to discern whether deploying the new model would yield superior results compared to the existing model.
    
    The governance model can then also check the model for non-functional requirements, such as fairness indicators, privacy concerns, and regulatory compliance (e.g., ensuring no code was used that is distributed under a software license containing a share-alike attribution policy). The governance model can also control how deployment is implemented (e.g., whether a blue-green or canary deployment is used for a particular model) and how traffic is split between the deployments (if a deployment requiring multi-way traffic is designated). Finally, the governance model details how the continuous monitoring protocol works, including how and which metrics are aggregated and reported.
    
    \subsection{Additional Resources}
    
    \begin{itemize}
      \item \href{https://neptune.ai/blog/machine-learning-model-management}{Challenges with Model Management}
      \item \href{https://thodrek.github.io/cs839_sp20/}{Stanford Course on Data Management}
    \end{itemize}
  
% 1.11: Responsible AI: Transparency & Sustainability

\section{Responsible AI Overview}

\subsection{What to Expect in This Course}

    As the famous saying goes, with great power comes great responsibility. While the superior performance and broad applicability of ML and TinyML have the potential to empower individuals and generate positive social impact, this can only be achieved if the developers use the technology responsibly. Only by grounding our work in responsible AI throughout the development and deployment of a project, can our scaled deployments of ML and TinyML applications go on to enact positive change in the world. 
    
    Responsible AI is an all-encompassing term that means a variety of different things, such as holding developers accountable, providing model transparency while maintaining individual privacy, and facilitating sustainable development and social impact to the best of our ability. While for some applications, some of these factors may not be as applicable, keeping these factors in mind is a good first step. Some criteria for social impact and sustainable development can be found in the United Nations 17 Sustainable Development Goals.
    
    Responsible AI has been born out of the growing concern amongst the computer science and data science community in recent years that the use of large amounts of user data meant that these fields of expertise might now have profound impacts beyond their specific applications and analysis, but were virtually unrestricted in ethical terms. In fact, it is only in the past few years that ethics has become a part of computer science and data science curriculums at universities. This, coupled with the increasingly profound environmental impact of ML, and the growing use of “black box” models, means that people are becoming acutely aware of the potential dangers and negative impacts of ML. As such, moving forward, designing applications responsibly will help to maximize the possibility that our work results in a positive change in the world.
    
    To this end, in this module we will survey a few key concepts of responsible AI that you need to consider when scaling machine learning deployments:
    
    \begin{itemize}
      \item Sustainable AI at Scale
      \item ransparency in ML Models
      \item The Opportunities for Social Impact
    \end{itemize}

    To learn more about these topics, I’ve asked my good friend and colleague Susan Kennedy, who is a philosopher, to come and share a few thoughts with us. Susan is an expert in ethics, applied ethics, and social/political philosophy, with a particular focus on emerging technologies. She will help us understand what sustainability and transparency mean in terms of ML. 
    

    
    Responsible AI is a huge field, there could be an entire course on that itself. So the objective here is not to give you a complete overview of Responsible AI, rather we want to get you thinking about it as part of MLOps.        
        
\section{Sustainability of TinyML}

    0:00 / 0:00Klicken Sie auf die Pfeiltaste nach oben, um das Geschwindigkeitsmenu aufzurufen. Regeln Sie dann mit den Hoch und Runter Pfeiltasten die Geschwindigkeit. Klicken Sie Enter, um die jeweilige Geschwindigkeit festzulegen.Geschwindigkeit1.0xKlicke auf diesen Knopf und diess Video stumm oder nicht stumm zu schalten, oder drücke die HOCH oder RUNTER Knöpfe um die Lautstärke zu erhöhen oder abzusenken.Maximal Lautstärke.Videotranskript
    Beginn des Transkriptes, zum Ende springen.
    SUSAN KENNEDY: One ethical concern that becomes more salient
    as we move towards scaling out TinyML to hundreds, thousands,
    or even millions of devices is the environmental impact.
    So to backtrack for just a moment, in course one,
    Vijay offered an overview of some of the basics of microcontrollers
    and how these compare to microprocessors.
    With one important difference being the MCUs
    consume significantly less power than a microprocessor.
    A typical MCU draws power in the order of milliwatts or microwatts,
    which is 1,000 times less power consumption than a microprocessor.
    And it's because of this energy efficiency that TinyML devices can
    run on battery power for weeks, months, and even years,
    all while running ML applications on the edge.
    But microcontrollers in themselves are not new.
    There's an estimated 250 billion microcontrollers in use today.
    Now with TinyML, it's expected that these devices
    will become even more pervasive.
    So while we might not ordinarily think about computing resources
    when we're working with a singular device,
    when we're faced with the possibility of scaling out TinyML,
    an important consideration to keep in mind
    is that the production and maintenance of a large quantity of MCUs
    can lead to substantial carbon emissions.
    In this video, we'll be addressing the sustainability of TinyML,
    which is focused on sustainable data sources, power
    supplies, and infrastructures as a way of measuring and reducing the carbon
    footprint from manufacturing hardware to training and running an ML application.
    When we're addressing the environmental impact of ML practices
    in terms of its carbon footprint, there's two different ways
    to categorize the costs.
    The first kind of environmental cost is what
    typically comes to mind when people think of carbon emissions.
    It's the operational or recurring costs associated with product use.
    For example, the energy consumption required for training and running
    inference on an ML device.
    But there's a second kind of environmental cost
    that's often overlooked, and that's the capital costs of ML.
    These are one-time activities like the manufacturing and production of MCUs
    the supply chain underlying this, as well as end-of-life processing.
    For example, consider having thousands of TinyML devices
    that run on battery power.
    End-of-life processing could refer then to the cost
    of properly disposing of these batteries so that they don't end up in landfills.
    Interestingly, recent research has shown that the carbon footprint from ML
    operations is diminishing, thanks to advancements
    in system energy efficiency and the increasing use of renewable energy.
    However, capital costs remain high, and more efforts
    to create sustainable supply chain activities are needed.
    And while quantifying the environmental impact of ML
    in terms of operational and capital costs can be helpful,
    it's not immediately obvious how to factor these costs into one's decision.
    And that's why it's important to reflect on what the ultimate goal is.
    In other words, what does sustainability mean, and why is it important.
    Well, sustainable development, as defined by the World Commission
    on Environment and Development in their 1987 report,
    is "Development that meets the needs of the present
    without compromising the ability of future generations
    to meet their own needs."
    So at the heart of sustainable development
    lies a tension between innovation and the equitable distribution
    of resources across society from one generation to the next.
    And determining what an equitable distribution of resources
    amounts to in the context of sustainable ML is a matter of ethics.
    Broadly speaking, equity refers to a fair distribution
    of burdens, benefits, and resources.
    So unlike a principle of equality, equity
    doesn't mean that everyone must share exactly the same benefits and burdens,
    as it's possible that a fair distribution entails
    some individuals getting a larger share and others a smaller one.
    For example, imagine people of different heights
    attempting to watch a ball game over a fence,
    and we have a limited number of blocks that we
    can stack to make step stools for them.
    If everyone got the same amount, it's not
    guaranteed that everyone would be able to see over the fence.
    But it might be possible to find an equitable
    distribution of these blocks in a way where
    everyone could see over the fence.
    So if you look at this image on the slide,
    you can see each individual shares the same eye level
    but with a different share of the blocks.
    This is what equity is all about.
    But sustainable development calls for an equitable distribution
    of resources in two different ways.
    The first is in terms of intergenerational justice,
    where resources are shared across and within a society.
    For instance, making sure that people have access
    to hardware and educational resources to build their own TinyML devices,
    and that the benefits of this technology are widely
    enjoyed, and the burdens do not disproportionately
    impact one community.
    But sustainable development also calls for intergenerational justice,
    where resources are shared between the present generation and future ones.
    For example, if the raw materials for producing MCUs
    are depleted by the present generation, making
    it impossible for future generations to meet their own needs,
    this wouldn't be considered sustainable.
    And such a big part of sustainability is considering what
    more obligations we might have to people who don't even exist yet.
    But as an old Greek proverb states, a society
    grows great when old men plant trees in whose shade
    they know they will never sit.
    In light of pressing environmental concerns facing the ML community,
    industry leaders have made some really encouraging progress.
    For instance, these are a couple of examples of the sustainability pledges
    that companies have made.
    Google has been carbon-neutral since 2007,
    meaning its carbon emissions are offset.
    And it aims to be carbon-free by 2030.
    Amazon aims for 100% renewable energy by 2025 and carbon-neutral by 2040.
    And Microsoft aims to be carbon-negative by 2030,
    where it removes more carbon than it produces,
    and furthermore, to remove historical carbon emissions by 2050.
    It's hard not to feel inspired by these steps forward to more sustainable AI.
    And if you're wondering how you can make sustainable development
    a part of your TinyML journey, the good news
    is that there's helpful tools and information already
    available and free to use, and more is on the way.
    For example, one great tool you can check out
    is Microsoft's sustainability calculator, which provides transparency
    into your carbon footprint associated with Dynamics 365 and Azure cloud
    usage.
    For example, in the bottom left-hand corner
    of the dashboard pictured on the slide, you
    can see an example of decreasing carbon emissions between 2017 and 2020.
    And the graph immediately to the right of this
    shows that carbon emissions have decreased
    even though cloud usage has increased.
    By tracking actual and avoided emissions over time,
    this information makes it much easier to make decisions that
    support more sustainable development.
    And while there's no one-size-fits-all strategy for sustainability,
    factoring the environmental impact and decisions about which cloud platform
    to use, which manufacturer to purchase hardware from,
    and so on, can all make a difference, as a seemingly negligible difference when
    building a single TinyML device can really add up when you look to start
    scaling out your product.
    Ende des Transkriptes, zum Anfang springen
    Downloads und Transskripte
    Video
    Videodatei herunterladen
    Transskripte
    Download SubRip (.srt) file
    Download Text (.txt) file        
    
    
\section{AI Sustainable}
    
    
    The field of artificial intelligence (AI) and machine learning (ML) is growing rapidly and with it the demand for computing power. This is great news for businesses and consumers alike, as it means more powerful and efficient applications are becoming available all the time.
    
    However, this increased demand for computing power comes at an environmental cost. Computing power requires electrical energy, and this energy must be obtained through some means. This energy could come from a variety of sources: wind, solar, geothermal, wave, fossil fuels, and other sources. For both consumers and businesses, this is an upstream process whose specifics are often largely unknown, much like the plumbing or heating systems in an individual’s home.
    
    The idea of sustainable energy focuses on the use of energy sources that are, in principle, infinitely renewable (i.e., they can be “sustained” in perpetuity). Because of the nuance involved in energy production, sustainability is not a binary characteristic. For example, the sun has a finite lifetime, albeit very long, which effectively makes it a renewable energy source from a human perspective. On the other hand, coal, oil, and natural gas all exist in finite quantities, although these quantities might differ substantially. Ensuring that the energy used in an AI or ML application is sustainable energy is difficult but is becoming an increasingly prominent issue.
    
    This issue is most important for data centers, which utilize vast numbers of GPUs and other devices in large storage facilities. These types of computers generate significant heat and consume enormous amounts of energy to run their powerful processors; this means that the energy used to power AI and ML is quickly becoming a major contributor to climate change. For ​​instance, in the Continuous Training stage, we discussed how training a single AI model for natural language processing can emit as much carbon as five cars!
    
    So what can be done about this? For the upstream process of energy production, we can help to promote energy governance from a sustainable perspective and also hold companies accountable for their choices of energy. In terms of the models themselves, we aim for two goals: efficiency and minimalism.
    
    Efficiency can be achieved by developing more efficient computing technologies. The best way to characterize efficiency is using a metric like performance per Watt, which can be measured on most architectures using benchmarks such as LINPACK. For a full data center, power usage effectiveness is often used instead as a more holistic metric. Current architectures are already very efficient, but nowhere near the limit of computational power efficiency, known as the Landauer limit. Additionally, the current onus is on the development of more powerful architectures, not power-efficient ones. As a result, our demand for computing using scale-out resources is putting increasing pressure on energy resources, further exacerbating existing environmental issues.
    
    Minimalism can be achieved by optimizing our algorithms so that they require less processing power. This can entail various things, such as reusing models (i.e., transfer learning) so that redundant computations are not performed, and using appropriately sized models, datasets, and architectures for the application, along with the optimal compilation flags and perhaps parallelization. As you might have already guessed, minimalism is really where TinyML comes to life.
    
    One more issue of sustainability we have not touched upon yet is our devices themselves. Every device that we use had to be constructed before we acquired or purchased it. This includes other finite resources such as rare metals, which should not be overlooked from a sustainability perspective. These materials have to be mined, transported, processed, and also deconstructed once their useful life has expired. All of these processes require energy and other additional resources. Thus, in the context of minimalism, we must strive to limit device consumption. Given the huge increase in demand for internet-connected microcontroller devices, this may be the most important challenge to be cognizant of in terms of the sustainability of TinyML.
    
    In the end, the most sustainable option of all is to limit energy usage overall to a minimum. It is up to us as consumers and developers to make sure that AI and ML remain sustainable. We must work together to find ways to reduce energy consumption without sacrificing performance or usability. With a little effort, we can create a future where artificial intelligence and machine learning are both efficient and environmentally friendly.    

\section{Model Cards for Transparency}

    0:00 / 0:00Klicken Sie auf die Pfeiltaste nach oben, um das Geschwindigkeitsmenu aufzurufen. Regeln Sie dann mit den Hoch und Runter Pfeiltasten die Geschwindigkeit. Klicken Sie Enter, um die jeweilige Geschwindigkeit festzulegen.Geschwindigkeit1.0xKlicke auf diesen Knopf und diess Video stumm oder nicht stumm zu schalten, oder drücke die HOCH oder RUNTER Knöpfe um die Lautstärke zu erhöhen oder abzusenken.Maximal Lautstärke.Videotranskript
    Beginn des Transkriptes, zum Ende springen.
    SUSAN KENNEDY: ML models can be powerful tools
    to transform the world we live in, but in order
    to make responsible decisions about how and when to use them,
    we need to rely on accurate and accessible information.
    But unlike the nutrition labels that we find on our food packaging,
    this information doesn't typically accompany ML models
    in such an organized and transparent manner.
    But what if it could?
    If you remember in Course 2, we talked about transparent reporting
    of information about data sets in the form of data sheets or data nutrition
    labels.
    Closely related to these documentation procedures
    are model cards for model reporting, which
    was an idea proposed in a Google research paper published in 2019.
    Model cards are short documents meant to accompany trained machine learning
    models, and provide detailed information about how the model performs,
    how it was evaluated, the context in which
    the model is intended to be deployed, and other relevant information.
    By doing so, model cards are thought to be
    a great step towards increasing transparency of a model's capabilities
    and limitations, facilitating communication
    and the exchange of information between model builders, product
    developers, and other relevant stakeholders,
    fostering accountable ML practices by providing model information required
    for effective public oversight and accountability,
    and democratizing ML technology by providing information to product users,
    so that they can make informed decisions about how to use them.
    At the heart of this proposal for model cards
    is a concern over the lack of transparency of ML practices.
    That's why the contents of model cards are designed in such a way as
    to offer thorough documentation and insight into a model.
    For example, a model card begins with model details
    section, with basic information like the person or organization who
    developed it, when it was created, as well as
    information about parameters, features, and fairness constraints.
    The intended use section captures what the intended uses and users are,
    as well as a warning label of sorts about use cases and users
    that fall outside of scope.
    The factors section offers information about how the environment might
    affect model performance, like whether background noise affects wake word
    detection.
    And it also covers what instruments were used to capture the input to the model,
    as a choice of recording instruments can impact the audio and video input,
    and subsequently, how the model performs.
    And while the relevant metrics will depend on the model being used,
    this section is an opportunity to explain
    what particular measures are being used for model performance,
    and why they were chosen.
    The evaluation data covers what data sets were used to evaluate the model,
    why they were chosen, and how the data was pre-processed.
    And similarly, training data includes, at the very least, basic details
    about distributions over various factors in the training data set.
    Quantitative analyses gives us more details about model performance,
    but broken down, for example, by demographics or other factors.
    The ethical considerations section is an opportunity
    to identify the various ethical challenges that cropped up
    along the way in development, and share how you took these into account,
    either respect to data, the model's potential impact on people, risks
    and harms, and more.
    The last section, caveats and recommendations,
    is an opportunity to share additional information not covered
    in the previous sections, which can be really
    helpful in directing those hoping to do further work with the model.
    By increasing transparency for machine learning models,
    this fosters better communication between all stakeholders.
    For example, model cards have obvious benefits
    for ML practitioners, organizations, and impacted individuals,
    but they can also be beneficial for policymakers, journalists, and advocacy
    groups and more.
    And it's important to keep this diversity of potential stakeholders
    in mind when you think about how to communicate certain information,
    as you want to make sure that the language is accessible to experts
    and non-experts alike, and that you're making efforts
    to anticipate the questions that people from different backgrounds might have.
    While this might seem like a lot to take in,
    the good news is that Google offers a model card
    toolkit to help you with the process of compiling and visualizing
    this information.
    You can see an example of a model card for object detection here.
    It offers written descriptions, like how many different object classes the model
    can detect, as well as visual representations
    for things like model performance, where a PR curve
    graph shows how model performance varies across two evaluation data
    sets-- in this case, the Open Images validation set,
    and the Google internal test set.
    With excellent tools and resources like this at your fingertips,
    it'll be much easier to increase your model's transparency,
    and reap all the downstream benefits that come with it.
    Ende des Transkriptes, zum Anfang springen
    Downloads und Transskripte
    Video
    Videodatei herunterladen
    Transskripte
    Download SubRip (.srt) file
    Download Text (.txt) file
  
\section{TinyML for Social Impact}

    0:00 / 0:00Klicken Sie auf die Pfeiltaste nach oben, um das Geschwindigkeitsmenu aufzurufen. Regeln Sie dann mit den Hoch und Runter Pfeiltasten die Geschwindigkeit. Klicken Sie Enter, um die jeweilige Geschwindigkeit festzulegen.Geschwindigkeit1.0xKlicke auf diesen Knopf und diess Video stumm oder nicht stumm zu schalten, oder drücke die HOCH oder RUNTER Knöpfe um die Lautstärke zu erhöhen oder abzusenken.Maximal Lautstärke.Videotranskript
    Beginn des Transkriptes, zum Ende springen.
    SUSAN KENNEDY: When ML projects are developed to prevent, mitigate,
    or even resolve problems that adversely affect
    human life or the well-being of the natural world,
    they fall in the category of AI for social impact.
    With the power of TinyML, there are tons of opportunities
    to address some of the world's most pressing challenges.
    For example, in course one, you might remember
    learning about how TinyML can be used for predictive maintenance
    and industry, to help farmers identify and treat plant diseases
    and for wildlife conservation by using sensors to mitigate
    conflicts between elephants and humans.
    The list goes on.
    Now, earlier in course four, you learned about benchmarking
    and how this contextualizes your performance.
    For example, some key performance indicators
    might be model accuracy, inference speed, and energy consumption.
    And this can all be really useful information
    to help you make improvements or adjustments.
    But when it comes to social impact projects,
    success isn't measured by benchmarks so much as by whether your model actually
    has--
    you guessed it-- a positive social impact.
    And knowing whether we've accomplished this goal or not
    can be difficult to measure.
    It involves going out into the field and conducting complex studies.
    For example, you might need to conduct field studies
    to understand the problem more clearly.
    And then test your product over time to see if this intervention actually
    works.
    Sometimes, this can take years.
    And so feedback, loops, and iteration can be slow as a result. Moreover,
    it can be difficult to measure impact or effectiveness of an intervention.
    For instance, if we want to know whether the solution we deployed
    was effective in preventing some outcome,
    it's impossible to simultaneously know what
    would have happened if we hadn't intervened-- otherwise known
    as the counterfactual.
    In short, evaluating whether your model creates
    the positive social impact you developed it for can be a challenging task.
    And similar to how social impact projects differ
    in how they measure success, they also differ with respect
    to what ultimately counts as success.
    So let's suppose you conduct a small field study
    and determine that your product's effective.
    True social impact requires a bit more than this.
    Not only do you need an effective solution,
    but you need the solution to be integrated into practice.
    And unfortunately, effectiveness does not guarantee uptake.
    So that's why it's crucial to come up with a plan for proper implementation
    of your solution.
    First, you might want to consider determinants of implementation.
    In other words, what features of the context impact uptake?
    For instance, does the area you plan to deploy in
    have weak Wi-Fi connectivity or none at all?
    Do product users understand how to incorporate this into their workflow?
    And will it be difficult for batteries to be
    replaced when they run out of power?
    Once you understand what kinds of things might prevent people
    from using your solution, then you can focus
    on implementation strategies, which refers to tactics that
    can be employed to increase uptake.
    So let's suppose we discover there's weak Wi-Fi connectivity.
    Can you develop your device in a way that
    still performs in this environment?
    And if battery life poses a problem, can you
    incorporate a self-sustained source of energy like solar power?
    By this point, you're probably noticing that when we aim for social impact
    as the outcome of a project, it'll influence
    the way we go about every stage of the ML workflow from design to deployment.
    And in general, social impact projects do have different requirements
    throughout the end-to-end ML workflow.
    Understanding both the technical and social elements at play
    is crucial and typically requires extensive stakeholder participation,
    collaboration, and feedback throughout each stage.
    But the good news is that after four courses of TinyML
    you'll be equipped with the knowledge and skills to do this.
    Now, when it comes to ethics of AI for social impact,
    beneficence is one of the foremost principles at play.
    Beneficence means acting for the benefit of others
    by helping them further their important and legitimate interests
    and by preventing or removing possible harms.
    And while beneficence may be a necessary condition for social impact,
    it's insufficient, in part, because whatever
    the positive impact the project might have
    could be offset by the creation or amplification of other risks and harms.
    So just because we aim to do good doesn't mean we can't inadvertently
    make a bad problem worse.
    Well, let's suppose that we're successful in acting
    for the benefit of others.
    One ethical question that might arise still
    is, how those benefits should be allocated?
    For instance, you could aim to maximize the total expected benefits
    across individuals.
    Or alternatively, you could give priority to the worse-off individuals,
    perhaps, by making sure your solution benefits underserved populations
    even if that means sacrificing some of the total expected benefits.
    Which one you should do will depend on the details of the case.
    But recognizing that beneficence may be a little more complicated than it first
    appears and spending some time thinking through different options
    of how to allocate benefits is a great first step in the right direction.
    Now, it's worth noting that ML for social impact
    presents both challenges and opportunities.
    One challenge is the lack of sufficient data.
    But you can revisit course two to learn about data collection
    and tackle this problem yourself.
    An opportunity is the need for low-cost hardware.
    And this makes TinyML a perfect candidate.
    Similarly, while connectivity issues might pose a challenge,
    this means there are unique opportunities
    for ML on the edge to offer a solution.
    And lastly, as energy efficiency is becoming increasingly important,
    the low power of TinyML offers a huge advantage
    compared to other alternatives.
    If there is a problem you're passionate about solving,
    chances are that you can find a way to harness the power of TinyML
    to create real social impact.
    Ende des Transkriptes, zum Anfang springen
    Downloads und Transskripte
    Video
    Videodatei herunterladen
    Transskripte
    Download SubRip (.srt) file
    Download Text (.txt) file  
 
%1.12: Summary

 
\section{Course Summary}

    Congratulations on completing the course! You are now equipped with all of the tools necessary to understand and approach the scaled deployment of TinyML applications with MLOps! To see just how far we have come, let’s review the key things we learned throughout this course.
    
    This course focused on MLOps in the context of TinyML. Because MLOps is becoming an important aspect of applied machine learning in industry, analogous to DevOps for the software engineering world, we had to differentiate between MLOps for BigML and TinyML. We defined BigML in terms of hardware, software, and network capabilities, such as having access to memory off-chip, using traditional software workflows, and utilizing cloud infrastructure. In contrast, TinyML involves resource-constrained devices that operate at the edge of the cloud, and the tiny memory footprint precludes the use of traditional software workflows, stimulating the need for novel lightweight frameworks such as TensorFlow Lite Micro.
    
    These differences between BigML and TinyML presented novel challenges in the implementation of TinyMLOps. For example, continuous training becomes largely untenable due to (1) the inability to store data on-device and (2) the lack of numeric precision (e.g., lack of a floating point unit as part of the processor architecture) required to propagate gradient terms during the training process. The limited communication also presents additional challenges, since communication processes are power-intensive, making it costly to stream data to a centralized storage unit. Because of these limitations, novel techniques needed to be developed to combat these issues, such as the use of federated learning across a network of devices. Further challenges were identified such as scalability and interoperability, which led to novel methods such as TinyMLaaS.
    
    To this end, throughout the course we learned both how traditional MLOps works for BigML applications, and how these MLOps approaches are translated into the world of TinyML. Based on the challenges we identified above, we presented the current state-of-the-art solutions that TinyML engineers use to approach these problems, covering the gamut from training operationalization and model development through to prediction serving and continuous monitoring. Since TinyML is a proto-engineering field, it is inevitable that these tools and techniques will change over time as the field evolves and advances. However, knowing the current approaches and their corresponding challenges will make it easier to rationalize and understand further advancements in the field - who knows, maybe you will be the one producing those advancements!
    
    Whether you took this course to develop new skills for your career, to enrich your resume, or just for personal development, finishing this course is no small feat. Whatever the case, we hope that this will just be the start of your journey deploying TinyML applications at scale and that this course will act as a stepping stone for many future advancements in TinyML. We wish you success in your future careers, and remember - the future of ML is tiny and bright.  
    
\section{Key Concepts of MLOps}

    0:00 / 0:00Klicken Sie auf die Pfeiltaste nach oben, um das Geschwindigkeitsmenu aufzurufen. Regeln Sie dann mit den Hoch und Runter Pfeiltasten die Geschwindigkeit. Klicken Sie Enter, um die jeweilige Geschwindigkeit festzulegen.Geschwindigkeit1.0xKlicke auf diesen Knopf und diess Video stumm oder nicht stumm zu schalten, oder drücke die HOCH oder RUNTER Knöpfe um die Lautstärke zu erhöhen oder abzusenken.Maximal Lautstärke.Videotranskript
    Beginn des Transkriptes, zum Ende springen.
    VIJAY JANAPA REDDI: Congratulations on making it to the end.
    In this video, I'm going to quickly recap
    some of the key concepts for scaling TinyML using MLOps.
    Now I hope you recall this figure from earlier on, which
    comes from the hidden technical depth in machine learning systems
    paper by Sculley and his team at Google, where
    they showed that the machine learning code is a small part of the bigger
    picture.
    More specifically, I hope you'll recall this quote,
    "Developing and deploying machine learning systems
    is relatively fast and cheap, but maintaining them over time
    is difficult and expensive."
    And to that end, I hope you see how MLOps
    is a set of processes that are aimed at reliably and efficiently designing,
    deploying, and maintaining machine learning models in production.
    In this course, we covered a lot of ground,
    starting from how we develop models in the ML development stage, all the way
    down to how we actually do continuous monitoring when the models are
    in the field serving inferences.
    To that end, we landed on answering these five
    key questions throughout the course--
    first, we learned about what is MLOps, and how is it similar or different
    from DevOps and AIOps?
    Then we understood why it is important to have systematic processes like MLOps
    in place so that we can scale our deployments.
    Third was what of all the different stages of MLOps,
    as you can see on this slide.
    Fourth is how are all the various stages interrelated with each other.
    For instance, we talked about the various feedback loops
    that we see across the various stages, as in how one MLOps stage can actually
    influence another, with examples such as how continuous monitoring can actually
    re-trigger continuous training and so forth.
    Last but not least, within each of these stages, we uniquely dissected
    the challenges that stand out with respect to scaling TinyML.
    For instance, we discussed challenges with respect
    to continuous integration and continuous development and deployment.
    We talked about how TinyML is dealing with embedded physical devices, which
    leads to a lot of challenges in terms of being able to robustly test them.
    And we talked about how we don't have containers and things like dockers
    and ought to be able to shrink wrap the software so we can easily push them out
    into deployment.
    I also talked about some of the challenges
    around continuous monitoring for instance, where we need to have
    or we need to employ intelligent hardware and software technologies.
    For instance, we talked about how we need
    to use communication technologies that are low power, for instance, like LoRa.
    And then we talked about how we might be able to use technologies
    like federated learning in order to minimize data communication while also
    getting benefits of privacy and security.
    Now last but not least, we saw how all these various components of MLOps
    clicked together like a giant jigsaw puzzle.
    And they're all continuously working with one another,
    so there's a lot that's actually going on.
    And to pull all of this off, we saw how different people
    with different expertise are required, ranging from your ML engineer
    to your business analyst who actually helps you start with a clear problem
    definition.
    I hope now you see why I was saying that you
    needed to have a T-shaped mindset as you go through the course.
    You need to have depth in an area, but you also need the breadth
    in order to understand how to take an ML project into life
    and make it a successful project.
    It's one thing to deploy models on a single device,
    but it's a whole different thing to scale it to many devices
    and serve many different users.
    With all that said, where do we go next?
    Well, I have a couple of quick thoughts that I'd
    like to share with you in the upcoming video.
    And hopefully you take a moment to reflect on all the exciting things
    that you've picked up throughout the entire course.
    Ende des Transkriptes, zum Anfang springen
    Downloads und Transskripte
    Video
    Videodatei herunterladen
    Transskripte
    Download SubRip (.srt) file
    Download Text (.txt) file

\section{What's Next}

    0:00 / 0:00Klicken Sie auf die Pfeiltaste nach oben, um das Geschwindigkeitsmenu aufzurufen. Regeln Sie dann mit den Hoch und Runter Pfeiltasten die Geschwindigkeit. Klicken Sie Enter, um die jeweilige Geschwindigkeit festzulegen.Geschwindigkeit1.0xKlicke auf diesen Knopf und diess Video stumm oder nicht stumm zu schalten, oder drücke die HOCH oder RUNTER Knöpfe um die Lautstärke zu erhöhen oder abzusenken.Maximal Lautstärke.Videotranskript
    Beginn des Transkriptes, zum Ende springen.
    VIJAY JANAPA REDDI: Now that you have all these new skills
    to build meaningful data sets, design, and deploy models
    and operationalize them through MLOps, the question is, what's next?
    How can we turn our MLOps pipeline into a new product,
    something that can provide value?
    But to be able to answer that question, though, first and foremost,
    we need to understand how products provide value
    and what is value from a business standpoint?
    And the reason I want to talk about this is more often than not,
    I meet a lot of people who are very excited about this new and emerging
    machine learning area called TinyML.
    And everything looks like a nail to them.
    They've got this hammer.
    And they think he can solve the whole world.
    But that's not true.
    So I want to help you understand this through a process called
    Jobs To Be Done, which was created, in part, by a Harvard Business School
    Professor named Clayton Christensen.
    Now, the idea behind this framework is that the value of a product
    is a job that it accomplishes for us.
    For example, a customer might go into a shop
    and buy a skateboard, all the parts that go with the skateboard.
    But that's not what they're really buying.
    What they're really buying is the experience of skateboarding
    and the social and emotional factors that go along with that.
    So valuable products are a means to an end.
    Now, we'll explore this in the context of a TinyML through a case study.
    As you might recall, TinyML has a wide range
    of applications across many different industries.
    Now, let's imagine, for instance, that we have developed a novel image
    classification technique that's super efficient and beats state-of-the-art.
    At first glance, this might seem like it's a great opportunity for us
    to go out and sell the technique to everyone who does anything
    with image classification.
    But, in fact, it's far more valuable to identify a specific industry
    where you can solve a specific, high-value problem for a single use
    case.
    Let's dive into this by taking space systems as an example.
    The primary goal is space systems broadly is
    to inform critical decisions on Earth.
    And many of these systems tend to do this through imagery.
    Right now, if you look across the planet,
    there are about 1,500 satellites that are actually orbiting
    that are geospatial imaging satellites.
    And by 2028, we're expected to have 50,000 of these geospatial imaging
    satellites, which means it's a growing market.
    And there are hundreds of industries that
    are just starting to learn how to utilize satellite imagery, which means
    the customer demand continues to grow.
    And if you stop for a moment and you look at these geospatial imaging
    satellites, you'll find that they tend to operate with very low power.
    They tend to have limited compute capability.
    They tend to have latency constraints and bandwidth constraints.
    They tend to have privacy concerns.
    All of these are starting to sound very much like TinyML systems.
    Because they share many of the same exact challenges or constraints.
    So there's clearly an opportunity here to apply our TinyML solution
    to these geospatial imaging satellites.
    But the question you should be asking yourself
    is what progress or job are the geospatial imaging satellites trying
    to do that they would want to hire our image classification
    technique to solve their problem?
    How can we connect our innovation to a direct benefit to the customer
    by increasing, for instance, their profits?
    Now, the best way to figure this out is to understand how their product
    provides value to their customers.
    Let's take a hypothetical company called SpaceCo
    whose job is to provide pictures of a particular point on Earth
    in as little time as possible.
    In other words, they provide value to the customer
    but providing useful information fast.
    Now, when a customer order comes in, SpaceCo
    directs one of their satellites to take an image of a specific location
    as soon as possible.
    They do this by tasking one of the satellites
    to go to a specific location.
    And then take an image ASAP.
    And then the satellite sends the image back to the ground station.
    Now, one of the key challenges down here is the latency problem.
    We don't know until minutes later what the image actually looks like.
    The customer might expect to receive an image like this
    and will only pay for a high-quality image.
    However, it turns out that 70% of the time, we
    get back images that look more like this,
    which means SpaceCo doesn't actually provide value to its customer
    and doesn't make a profit on this image.
    Cloud coverage is a major issue for these sorts of geospatial imaging
    satellites.
    But what if we could build a satellite that could autonomously
    wait to snap the photo.
    We could use our new highly-efficient image classification algorithm
    to solve this problem.
    We could design a TinyML model to recognize
    that there are clouds in the image while in orbit
    and just waiting around that target.
    And then when the cloud clears, we can actually take a clear picture.
    Great.
    So now we've figured out that they really want is a nice, clean picture.
    And our immediate reaction here might be to jump the gun and say,
    let's license them a visual keyword, no-cloud recognition model.
    And then they can deploy it however they want.
    Great, right?
    Wrong.
    We also need to help them solve the MLOps problem.
    Recall that SpaceCo wants this.
    But there are several things to think about.
    They're not an expert on this technique that you and I have dealt with.
    It's also potentially complex and costly to hire
    someone who can maintain and modify the model as needed over time.
    So the real job here is to think about how can we
    reduce the complexity that we're adding into their system
    and yet provide value that we can achieve by providing them a tool that
    can easily integrate into their existing product
    that you and I manage and improve for them, aka, MLOps solution?
    Now that you've figured out what the customer really wants,
    now we know what the real job is.
    We have to have a full-stack solution for this particular problem.
    And that's where you realize that because they
    don't have the in-house chops to build and deploy models,
    we can employ our knowledge that you've gained through this course
    and then hire the right experts from your network to create a team that
    can actually solve this MLOps problem.
    Now, you can go out, and you can build a company whose mission
    is to build and deploy TinyML models and workflow
    management solutions to democratize satellite imagery
    and bring about world peace.
    And that's a simplified example how to transfer a technological innovation
    into a product that provides real value to a customer.
    In summary, it's useful to approach products and commercialization
    through the lens of the Jobs To Be Done framework.
    The framework highlights that successful products unlock a capability
    or remove a barrier for a customer without adding additional complexity
    and issues to them.
    And so with what you've learned about the MLOps pipeline,
    I hope you see how you can manage that complexity for them.
    And in doing so, provide new value to the customer.
    So with all that said, I hope you're feeling
    inspired to go out and build something new that's going to change the world.
    Congratulations once again on finishing the MLOps course.
    Ende des Transkriptes, zum Anfang springen
    Downloads und Transskripte
    Video
    Videodatei herunterladen
    Transskripte
    Download SubRip (.srt) file
    Download Text (.txt) file

\section{What are you excited about doing and learning next?}

Now that you’ve completed MLOps for Scaling TinyML, what are you excited about doing with your newfound knowledge? How are you going to apply these concepts to your work and personal projects? What would you be excited about learning next?

\section{Closing Remarks}

    0:00 / 0:00Klicken Sie auf die Pfeiltaste nach oben, um das Geschwindigkeitsmenu aufzurufen. Regeln Sie dann mit den Hoch und Runter Pfeiltasten die Geschwindigkeit. Klicken Sie Enter, um die jeweilige Geschwindigkeit festzulegen.Geschwindigkeit1.0xKlicke auf diesen Knopf und diess Video stumm oder nicht stumm zu schalten, oder drücke die HOCH oder RUNTER Knöpfe um die Lautstärke zu erhöhen oder abzusenken.Maximal Lautstärke.Videotranskript
    Beginn des Transkriptes, zum Ende springen.
    VIJAY JANAPA REDDI: I want to wrap up this course
    with a couple of closing remarks.
    We started this course with the mind-blowing fact that over 87%
    of data science or machine learning projects never see the light of day.
    And I honestly believe that this is because most people can't
    see the forest for the trees.
    They tend to get so caught up or excited about machine learning training
    or model optimizations and so forth, that they tend to miss the big picture.
    And that's why in this course, I chose intentionally
    to paint the broad strokes across MLOps.
    It's important that we understand what it is, why it's helpful,
    and how it can help you take your project into real life.
    Now look, there is no one course that's going
    to get all concept of MLOps picture perfect right.
    So don't let this course be the end of your MLOps journey.
    Rather, think of this is the beginning.
    Go out, go learn more, go take more courses,
    and put your knowledge into practice.
    And do something that's going to change the world.
    With that said, thank you for hanging out with me throughout this course.
    It's been a pleasure.
    I look forward to seeing all the amazing things you're going
    to be doing out there in the world.
    Ende des Transkriptes, zum Anfang springen
    Downloads und Transskripte
    Video
    Videodatei herunterladen
    Transskripte
    Download SubRip (.srt) file
    Download Text (.txt) file
        
  